% METADATA DEFINITIONS (Hidden from PDF output)














\documentclass[10pt, a5paper, twoside]{article}


    \raggedbottom

% \flushbottom

\usepackage{fontspec}
\usepackage{polyglossia}
\usepackage{geometry}
\usepackage{fancyhdr}
\usepackage{tcolorbox}
\usepackage{ifthen}
\usepackage{tikz}
\usetikzlibrary{shadows}
\usepackage{xcolor}
\usepackage{microtype}
\usepackage{graphicx}
\usepackage{float}
\usepackage[object = vectorian]{pgfornament}
\usepackage{tocloft}
\usepackage{xltabular}
\usepackage{booktabs}
\usepackage{array}

% --- Table Alignment Control ---
\newcommand{\tabledefaultalign}{\raggedright}
\newcolumntype{K}{>{\tabledefaultalign\arraybackslash}X}
\newcommand{\settablealign}[1]{%
    \ifthenelse{\equal{#1}{Right}}{%
        \renewcommand{\tabledefaultalign}{\raggedleft}%
    }{%
        \renewcommand{\tabledefaultalign}{\raggedright}%
    }%
}


% --- Robust Purnaviram Line Break Control (LuaTeX) ---
\usepackage{luacode}
\begin{luacode*}
function fix_devanagari_punctuation_breaks(head)
  local GLYPH = node.id("glyph")
  local GLUE  = node.id("glue")
  local KERN  = node.id("kern")
  local PENALTY = node.id("penalty")

  for n in node.traverse_id(GLYPH, head) do
    -- 0x0964 = । (Purnaviram), 0x0965 = ॥ (Deerghapurnaviram)
    if n.char == 0x0964 or n.char == 0x0965 or n.char == '?' then
      
      -- 1. Prohibit break BEFORE punctuation
      local curr = n.prev
      while curr and (curr.id == GLUE or curr.id == KERN or curr.id == PENALTY) do
        if curr.id == GLUE or curr.id == KERN then
            local p = node.new(PENALTY)
            p.penalty = 10000 -- Block break
            head = node.insert_before(head, curr, p)
        elseif curr.id == PENALTY then
            curr.penalty = 10000 -- Override existing penalty
        end
        curr = curr.prev
      end
      
      -- 2. Encourage break AFTER punctuation
      local nxt = n.next
      if nxt and nxt.id == GLUE then
         local p = node.new(PENALTY)
         p.penalty = -500 -- Preferred break point
         head = node.insert_before(head, nxt, p)
      end
    end
  end
  return head
end
luatexbase.add_to_callback("pre_linebreak_filter", fix_devanagari_punctuation_breaks, "Fix Devanagari Breaks")
\end{luacode*}

% Support settings for the above
\sloppy
\emergencystretch=5cm
\hyphenpenalty=10000
\exhyphenpenalty=10000
\linepenalty=1000

% Layout settings for better line breaking
\sloppy
\emergencystretch=5cm


\geometry{
    top=1in,
    bottom=1in,
    left=0.5in,
    right=0.5in,
    headheight=15pt,   % Ensures fancyhdr has enough room for the text
    headsep=0.3cm
}

% --- spacing configuration ---
\linespread{1.0}
\setlength{\parskip}{0pt plus 3pt minus 1.5pt}
\setlength{\parindent}{20pt}

% --- Devanagari Page Numbering Macros ---
\makeatletter
\def\devanagaridigits#1{\expandafter\@devanagaridigits\number#1@}
\def\@devanagaridigits#1{%
  \ifx#1@\else
    \ifcase#1 ०\or १\or २\or ३\or ४\or ५\or ६\or ७\or ८\or ९\fi
    \expandafter\@devanagaridigits
  \fi
}

\def\devanagarialph#1{\expandafter\@devanagarialph\number#1}
\def\@devanagarialph#1{%
  \ifcase#1\or अ\or आ\or इ\or ई\or उ\or ऊ\or ऋ\or ॠ\or ऌ\or ॡ\or ए\or ऐ\or ओ\or औ\or क\or ख\or ग\or घ\or ङ\or च\or छ\or ज\or झ\or ञ\or ट\or ठ\or ड\or ढ\or ण\or त\or थ\or द\or ध\or न\or प\or फ\or ब\or भ\or म\or य\or र\or ल\or व\or श\or ष\or स\or ह\or ळ\or क्ष\or ज्ञ\else\@ctrerr\fi
}
\makeatother

\newcommand{\ornamentalpage}{%
    %
        \smash{%
        \begin{tikzpicture}[baseline=(page.base)]
            \node[color=theme, inner sep=0pt, yshift=0pt] (ornl) {\pgfornament[width=1.5cm]{65}};
            \node[anchor=west, xshift=5pt, inner sep=0pt] (page) at (ornl.east) {\thepage};
            \node[anchor=west, color=theme, xshift=5pt, inner sep=0pt, yshift=0pt] (ornr) at (page.east) {\pgfornament[width=1.5cm]{65}};
        \end{tikzpicture}%
        }%
    %
}

\pagestyle{fancy}
\fancyhf{} 


    \renewcommand{\thesection}{\devanagaridigits{\value{section}}}
    \renewcommand{\thesubsection}{\thesection.\devanagaridigits{\value{subsection}}}
    \renewcommand{\thesubsubsection}{\thesubsection.\devanagaridigits{\value{subsubsection}}}


% -- Rules --

    \renewcommand{\headrulewidth}{0.4pt}



    \renewcommand{\footrulewidth}{0.4pt}


% -- Header/Footer Positioning --

    % Mirrored (Outer/Inner logic)
    % Header Left -> Left Side on Odd (Outer), Right Side on Even (Outer)
    
    % Header Right -> Right Side on Odd (Outer), Left Side on Even (Outer)
    
    % Header Center
    
    
    % Footer Left
    
    % Footer Right
    
    % Footer Center
    


% --- Theme Color Definition ---
\definecolor{theme}{named}{black}

% --- Vishesh Vishaya Suchi (Special Table of Contents) ---
\newcommand{\visheshcontents}[5]{%
    \cleardoublepage
    \begingroup
    
    % 1. Use specialkhanda for title (handles marks, TOC entry, and fancy style)
    \specialkhanda{#1}%
    
    % 2. Suppress standard ToC title and force fancy style
    \renewcommand{\contentsname}{}% 
    \renewcommand{\cftaftertoctitle}{\thispagestyle{fancy}}% 
    \tocloftpagestyle{fancy}%
    
    % 3. Style Mahakhanda (Part)
    \cftpagenumbersoff{part}% Hide page numbers for Mahakhanda
    \renewcommand{\cftpartfont}{\hfill\raisebox{#5}{\mbox{\pgfornament[width=#4]{#2}}}\quad\Large\bfseries\color{theme}}% Left Orn + Title
    \renewcommand{\cftpartafterpnum}{\quad\raisebox{#5}{\mbox{\pgfornament[width=#4]{#3}}}\hfill\null\par\vspace{1em}}% Right Orn + Gap Below
    \setlength{\cftbeforepartskip}{2.5em}%
    \setlength{\cftpartnumwidth}{0pt}% 

    % 4. Style Khanda (Section)
    \renewcommand{\cftsecfont}{\large\bfseries}%
    \renewcommand{\cftsecleader}{\cftdotfill{\cftdotsep}}%
    \setlength{\cftbeforesecskip}{0.5em}%

    \tableofcontents
    \endgroup
    \pagebreak
}

% --- Custom Segments (Khanda) ---
\newcommand{\mahakhanda}[1]{
    \cleardoublepage
    \thispagestyle{empty}
    \vspace*{\fill}
    \begin{center}
        \textbf{\Huge \color{theme} #1}
    \end{center}
    \vspace*{\fill}
    \addcontentsline{toc}{part}{#1}
    \setcounter{section}{0}
    \pagebreak
}

\newcommand{\specialmahakhanda}[1]{
    % \cleardoublepage
    \clearpage
    \thispagestyle{empty}
    \nopagebreak
    \vspace*{\fill}
    \begin{center}
    \begin{tikzpicture}
        \node[inner sep=20pt, fill=theme!10, rounded corners=10pt, drop shadow] (title) {
            \parbox{0.8\textwidth}{\centering \textbf{\Huge \color{theme} #1}}
        };
    \end{tikzpicture}
    \end{center}
    \vspace*{\fill}
    \addcontentsline{toc}{part}{#1}
    \setcounter{section}{0}
    \pagebreak
}

\newcommand{\specialmahakhandatwo}[2]{
    \cleardoublepage
    \thispagestyle{empty}
    \vspace*{\fill}
    \begin{center}
    \begin{tikzpicture}
        \node[inner sep=20pt, fill=theme!10, rounded corners=10pt, drop shadow] (title) {
            \parbox{0.8\textwidth}{\centering \textbf{\Huge \color{theme} #1} \par \vspace{0.6cm} \textbf{\Large \color{theme} #2}}
        };
    \end{tikzpicture}
    \end{center}
    \vspace*{\fill}
    \addcontentsline{toc}{part}{#1}
    \setcounter{section}{0}
    \pagebreak
}

\newcommand{\khanda}[1]{
    \markboth{#1}{}
    \vspace{0.8cm}
    \noindent\makebox[\textwidth][c]{{\huge \color{theme} \textbf{#1}}} \par
    \addcontentsline{toc}{section}{#1}
    \vspace{0.4cm}
}

\newcommand{\upakhanda}[1]{
    \vspace{0.6cm}
    \noindent\makebox[\textwidth][c]{{\Large \color{theme} \textbf{#1}}} \par
    \addcontentsline{toc}{subsection}{#1}
    \vspace{0.3cm}
}

\newcommand{\upaupakhanda}[1]{
    \vspace{0.4cm}
    \noindent\makebox[\textwidth][l]{{\Large \color{theme} \textbf{#1}}} \par
    \addcontentsline{toc}{subsubsection}{#1}
    \vspace{0.2cm}
}

% --- Custom Paragraph ---
\newcommand{\customparagraph}[1]{
    % \penalty 10000 
    % \vspace{0.2cm}
    % \vspace{0.2cm plus 0.1cm minus 0.05cm}
    \noindent \textbf{#1} \par
    \penalty 10000 % Prohibits a page break right here
    % \vspace{0.1cm}
    \vspace{0.1cm plus 0.05cm minus 0.02cm}
    \penalty 10000 % Prohibits a page break right here
}

% --- Alignment and Spacing Macros ---
\newcommand{\vspacer}[1]{\vspace{#1 cm}}

% --- Shloka Environment ---
\newcounter{shlokaline}
\newenvironment{shloka}[1][alternating]{
    \par\vspace{0.3cm plus 0.3cm minus 0.1cm}
    % \begingroup
    \noindent\begin{minipage}{\textwidth} % <-- Replaced \begingroup
    \parindent=0pt
    \linespread{1} \selectfont
    \large
    \def\shlokaalign{#1}
    \setcounter{shlokaline}{0}
    \ifthenelse{\equal{\shlokaalign}{alternating}}{
        \let\oldpar\par
        \def\par{\oldpar\stepcounter{shlokaline}}
        \everypar{\ifodd\value{shlokaline}\raggedleft\else\raggedright\fi}
    }{
        \ifthenelse{\equal{\shlokaalign}{left}}{\raggedright}{}
        \ifthenelse{\equal{\shlokaalign}{right}}{\raggedleft}{}
        \ifthenelse{\equal{\shlokaalign}{center}}{\centering}{}
    }
}{
    % \par\endgroup
    \goodbreak
    \par\end{minipage} % <-- Replaced \endgroup
    \nopagebreak
    % \vspace{-0.5cm}
}

% --- Vishesh Shloka Environment (with Ornament) ---
\newenvironment{visheshshloka}[1][alternating]{
    \begin{shloka}[#1]
}{
    \end{shloka}
    % \vspace{-0.8cm}
    % \begin{center}
    %     \color{theme}\pgfornament[width=0.2\textwidth]{88}
    % \end{center}
    % \vspace{-0.3cm}
}

% --- Vishesh Shloka Two Environment (with Custom Ornament) ---
\newenvironment{visheshshlokatwo}[2][alternating]{
    \def\visheshornament{#2}
    \begin{shloka}[#1]
}{
    \end{shloka}
    \vspace{-0.5cm}
    \begin{center}
        \color{theme}\pgfornament[width=0.2\textwidth]{\visheshornament}
    \end{center}
    \vspace{0.1cm}
}

% --- Image Macros ---
% --- Vishesh Butta (Repeated Ornaments) ---
\newcommand{\visheshbutta}[5]{%
    \par\vspace{0.4cm}
    \noindent\begin{minipage}{\textwidth}
        \ifthenelse{\equal{#1}{Center}}{\centering}{}%
        \ifthenelse{\equal{#1}{Left}}{\raggedright}{}%
        \ifthenelse{\equal{#1}{Right}}{\raggedleft}{}%
        \foreach \n in {1,...,#3}{%
            \pgfornament[width=#4]{#2}\hspace{#5}%
        }%
    \end{minipage}\par\vspace{0.4cm}
}

\newcommand{\customimage}[3]{
    \begin{figure}[H]
        \ifthenelse{\equal{#2}{Center}}{\centering}{}
        \ifthenelse{\equal{#2}{Left}}{\raggedright}{}
        \ifthenelse{\equal{#2}{Right}}{\raggedleft}{}
        \includegraphics[width=#3\textwidth]{#1}
    \end{figure}
}

% --- Special segments with TikZ gradient ---
\newcommand{\specialkhanda}[1]{
    \markboth{#1}{}
    \vspace{1.2cm}
    \noindent\makebox[\textwidth][c]{{\huge \color{theme} \textbf{#1}}} \par \nointerlineskip
    \addcontentsline{toc}{section}{#1}
    \noindent\begin{tikzpicture}[remember picture, overlay]
        
            \shade[left color=white, right color=theme!80!white] (0,-0.4) rectangle (0.5\textwidth, -0.2);
            \shade[left color=theme!80!white, right color=white] (0.5\textwidth,-0.4) rectangle (\textwidth, -0.2);
        
    \end{tikzpicture}
    \vspace{0.8cm}
}

\newcommand{\specialupakhanda}[1]{
    \vspace{0.8cm}
    \noindent\makebox[\textwidth][c]{{\Large \color{theme} \textbf{#1}}} \par \nointerlineskip
    \addcontentsline{toc}{subsection}{#1}
    \noindent\begin{tikzpicture}[remember picture, overlay]
        
            \shade[left color=white, right color=theme!60!white] (0,-0.3) rectangle (0.5\textwidth, -0.15);
            \shade[left color=theme!60!white, right color=white] (0.5\textwidth,-0.3) rectangle (\textwidth, -0.15);
        
    \end{tikzpicture}
    \vspace{0.6cm}
}

\newcommand{\specialupaupakhanda}[1]{
    \vspace{0.4cm}
    \noindent\makebox[\textwidth][l]{{\Large \color{theme} \textbf{#1}}} \par \nointerlineskip
    \noindent\begin{tikzpicture}[remember picture, overlay]
        
            \shade[left color=theme!40!white, right color=white] (0,-0.2) rectangle (\textwidth, -0.1);
        
    \end{tikzpicture}
    \vspace{0.4cm}
    \addcontentsline{toc}{subsubsection}{#1}
}

% --- Font Configuration ---

    \setmainfont[Script=Devanagari, AutoFakeBold=1.8]{Tiro Devanagari Sanskrit}


\begin{document}
\sloppypar
\visheshcontents{विषय सूची}{70}{70}{1cm}{0pt}

\specialmahakhanda{\textbf{प्रथम-प्रकाशः}}

\specialkhanda{\textbf{शास्त्र-परिचयः}}

\customparagraph{\textbf{विषयप्रवेशः}}

`शिष्यते उपदिश्यते वा अनेन' इति व्युत्पत्त्या \textbf{`शास्'} धातोः \textbf{`ष्ट्रन्'} प्रत्यये कृते \textbf{`शास्त्रम्'} इति शब्दो निष्पद्यते । अस्य सामान्यार्थः— आज्ञा, आदेशः, नियमश्चेति भवति । विश्ववाङ्मये `शास्त्र'शब्दस्य प्रयोगः प्राचीनकालादेव प्रवृत्तो दृश्यते । पौरस्त्यवाङ्मये विविधविधागतविशेषणैः सह शास्त्रशब्दः प्रयुज्यते; यथा— धर्मशास्त्रम्, न्यायशास्त्रम्, राजनीतिशास्त्रम्, काव्यशास्त्रमित्यादीनि । व्युत्पत्तिलभ्यार्थानुसारेणैव यत्र धर्म आदिश्यते नियम्यते वा, तत् \textbf{धर्मशास्त्रम्}; तथैव न्याय-राजनीति-समाज-अर्थ-काव्यादीनां यत्र नियम-आदेशादीनामुल्लेखो भवति, तत्र तत्तद्विषयकं शास्त्रं कथ्यते । शास्त्रीयग्रन्थेषु पञ्चाङ्गम् (अधिकरणम्) अनिवार्यं वर्तते । पञ्चाङ्गैर्विना कोऽपि ग्रन्थः शास्त्रपदभाग् न भवति; यथा उक्तं मम्मटाचार्येण—

\begin{shloka}[center]

\textit{\textbf{विषयो विशयश्चैव पूर्वपक्षस्ततोत्तरम् ।}}

\textit{\textbf{निर्णयश्चेति पञ्चाङ्गं शास्त्रेऽधिकरणं स्मृतम् ॥ (मम्मटः १९८३ : २)}}

\end{shloka}

अत्र तु विषयप्रसङ्गेन काव्यशास्त्रस्य विषयो विविच्यते । कवेः कर्म काव्यं भवति । तच्च निश्चितं प्रयोजनमुद्दिश्यैव विरच्यते, अत एतादृशं काव्यं नियमबद्धं भवितुमर्हति । काव्यं सुगठितम्, सौन्दर्ययुक्तम्, व्यवस्थितञ्च कर्तुं शास्त्रीयपक्षस्यावश्यकता भवति । शास्त्ररचनां सर्वज्ञा विद्वांस एव कुर्वन्ति । वाण्या मार्गद्वयं दृश्यते— एका प्रज्ञा, अपरा च प्रतिभा । तत्र प्रज्ञया शास्त्रस्य, प्रतिभया च काव्यस्य प्रणयनं भवति । तदुच्यते—

\begin{shloka}[center]

\textit{\textbf{द्वे वर्त्मनी गिरां देव्याः शास्त्रं च कविकर्म च ।}}

\textit{\textbf{प्रज्ञोपज्ञं तयोराद्यं प्रतिभोद्भवमन्तिमम् ॥}}

\end{shloka}

शास्त्रस्य प्रणयनं प्रज्ञया भवति, प्रज्ञा च अभ्यासात् समुपजायते । तस्माद् अविदुषा शास्त्रकारेण प्रणीतं शास्त्रं प्रयोजनहीनं भवति । अत एव अग्निपुराणकारः समुल्लिखति—

\begin{shloka}[center]

\textit{\textbf{व्युत्पत्तिर्दुर्लभा तत्र विवेकस्तत्र दुर्लभः ।}}

\textit{\textbf{सर्वं शास्त्रमविद्वद्भिर्मृग्यमाणं न सिध्यति ॥}}

\end{shloka}

काव्यशास्त्रे काव्यनिर्माणार्थम् आवश्यकतत्त्वानां साङ्गोपाङ्गं विवेचनं क्रियते । कमनीयकाव्यस्य कृते कानि कानि आवश्यकानि तत्त्वानि सन्ति ? रस-ध्वनि-अलङ्कार-गुण-रीत्यादीनां काव्यतत्त्वानां स्वरूपं महत्त्वञ्च किम् ? तेषु कस्य प्राधान्यं वर्तते ? कस्मिन् स्थाने केन प्रकारेण तेषामङ्गीकारो विधीयते ? इत्यादि जिज्ञासाशमनमेव काव्यशास्त्रस्य मुख्यमुद्देश्यमपेक्षते । येन विना कवेः कर्मरूपं काव्यं सर्वस्वीकारयोग्यं विलक्षणानन्दप्रदञ्च न भवति, एतादृशं काव्यस्य स्वरूपादिनियामकं शास्त्रीयविषयमध्येतुमुन्मुखानां जिज्ञासुजनानां सहाय्यं भवेदिति विचार्य काव्यशास्त्रीयविवेचनमत्र विधीयते ।

\specialupakhanda{\textbf{शब्दप्रयोगवैविध्यम्}}

संस्कृतवाङ्मये काव्यशास्त्रस्य कृते स्वकीयविचारस्थापकैराचार्यैः विविधानि नामानि प्रयुक्तानि सन्ति  । तेषु प्राचीनतमो विद्यते \textbf{काव्यालङ्कारः} । काव्यस्य शास्त्रीयविषयस्य विश्लेषणकर्तृभिराचार्यैः प्राचीनकाले स्वग्रन्थस्य नाम एव काव्यालङ्कारोऽकार्षुः । भामहरुद्रटवामनादयः स्वग्रन्थस्य नाम तादृशमेव कृतवन्तः । तेषां मतेऽलङ्कारः केवलं काव्यशोभाविषयक एव न, अपि तु  सर्वेषां काव्यतत्त्वानामपरपर्याय इति ज्ञायते । साहित्यशास्त्रीयपरम्पराया विकासक्रमे प्रारम्भिकाचार्याः काव्यशास्त्रस्यार्थेऽलङ्कारशब्दस्य प्रयोगमकार्षुः । प्रतापरुद्रियेण ‘यद्यपि रसालङ्काराद्यनेकविषयकमिदं शास्त्रं तथापि छत्रीन्यायेनालङ्कारशास्त्रमित्युच्यते’ इति कथयित्वा \textbf{अलङ्कारशास्त्र}मिति प्रयोगाभिप्रायः प्रकटीकृतः । पश्चाज्जातः भोजदेवः स्वकीये सरस्वतीकण्ठाभरणाख्ये लक्षणग्रन्थेऽलङ्कारशास्त्रस्य नाम \textbf{काव्यशास्त्र}मकार्षीत्, यथा—

\begin{shloka}[center]

\textbf{काव्यं शास्त्रेतिहासौ वा काव्यशास्त्रं तथैव च ।}

\textbf{काव्येतिहासशास्त्रेतिहासस्तदपि षड् विधम् ।। (सरस्वतीकण्ठाभरणम्)}

\end{shloka}

काव्यालङ्कारवत्  काव्यशास्त्रवच्च  \textbf{साहित्यशास्त्र}स्य प्रयोगोऽपि प्राचीनकालादेव अजायत । काव्यशास्त्रार्थे शब्दोऽयमर्वाचीने कालेऽतीव प्रख्यातिमगात् । चतुर्दशशताब्द्यां जातः विश्वनाथः स्वकीयलक्षणशास्त्रस्य नाम साहित्यदर्पण इत्यकार्षीत् । एकादशशतकस्य आचार्यो रुय्यकः साहित्यमीमांसा इति नामकं ग्रन्थं विरचय्य काव्यशास्त्रे साहित्यशब्दं प्रयुक्तवान् । साहित्यशब्दस्याभिप्रायमभिव्यञ्जयन् आचार्यः कुन्तकोऽपि वक्रोक्तिजीवितेऽघोषयत्—

\begin{shloka}[center]

\textbf{साहित्यमनयोः शोभाशालिताप्रतिकाप्यसौ ।}

\textbf{अन्यूनमतिरिक्तत्वमनोहारिण्यवस्थितिः।।}

\end{shloka}

काव्यमीमांसा इति लक्षणग्रन्थस्य प्रणेता राजशेखरोऽमुं  ‘पञ्चमी साहित्यविद्येति यायावरीयः’ इति कथयित्वा साहित्यविद्येति नाम्ना स्वीकरोति । एतादृशानि विविधानि नामानि संस्कृतसाहित्ये प्राप्यन्ते । तेषु च प्राचीनतमो विद्यते \textbf{क्रियाकल्प} इति । अस्य निर्देशो वात्स्यायनस्य कामशास्त्रे परिगणितासु चतुःषष्टिकलासु समुपलभ्यते । इदं नाम काव्यक्रियाकल्पस्य सङ्‌क्षिप्तमाभाति । तथा च वाल्मीकीयरामायणस्योत्तरकाण्डस्य चतुर्नवतितमे सर्गे क्रियाकल्पशब्दस्य  काव्यविद्-शब्दस्य च प्रयोगः प्राप्यते । तेषु काव्यविद्-शब्दः पुरुषार्थे प्रयुज्यते । क्रियाकल्पशब्देन तु काव्यसौन्दर्यपरीक्षणसमर्थपुरुषो ज्ञायते । यथा— \textbf{`क्रियाकल्पविदश्चेति तथा काव्यविदो जनान् ।'} आभिश्चर्चाभिः लक्षणशास्त्रस्य कृते काव्यालङ्कारः, अलङ्कारशास्त्रम्, काव्यशास्त्रम्, साहित्यशास्त्रम्, साहित्यविद्या, क्रियाकल्पश्चेत्यादयः शब्दभेदाः समुपलभ्यन्ते । एतेषु काव्यशास्त्रमलङ्कारशास्त्रञ्चे’ति द्वे नामनी  पौरस्त्यशास्त्रेतिहासे प्रसिद्धे स्तः ।

\customparagraph{\textbf{कलाशास्त्रम्}}

संस्कृतवाङ्‌मये कलाशास्त्रमिति नाम्नापि शास्त्रीयचर्चा विहिता । तत्र कविकर्मभूतं काव्यमेव न, अपि तु  विदुषां गभीरज्ञानार्जिताः सर्वा रचनाप्रकारा विश्लिष्यन्ते । कलेति  नाम्ना ताः सर्वा रचनाः  सङ्‌गृह्यन्ते । विश्ववाङ्‌मये प्राचीनकालादेव कलाया महत्त्वं सर्वत्रानुभूयते । अधुना कला मानवसमुदायस्य कृतेऽपरिहार्यो विषयः सञ्जातः । पाश्चात्त्यवाङ्‌मये व्यापकार्थे कलातत्त्वस्य चिन्तनं विहितं दृश्यते । तत्र सामान्यतया मानवश्रमेण स्वजीवनसञ्चालनार्थं कृताः सर्वाः कृतयः सौन्दर्यस्याधारभूताः कला इति कथ्यन्ते । पौरस्त्यवाङ्‌मयेऽपि चतुःषष्टिसङ्ख्यायुतेषु कलाप्रभेदेषु प्रकारान्तरेणापि मानवस्य सर्वाणि कर्माणि सङ्गृहीतानि विद्यन्ते ।

‘कल सङ्ख्यानम्’ इति सङ्ख्यानार्थकात् कलधातोः ‘कला’ शब्दो निष्पद्यते । ‘चक्षिङ् व्यक्तायां वाचि’ इति धातोः ‘ख्या’ आदेशात् निर्मितस्य समुपसर्गपूर्वकस्य सङ्ख्यानशब्दस्यार्थः सङ्कलनम्, गणनम्, स्पष्टाभिव्यक्तिश्च भवति । तात्पर्यार्थत्वेन मननम्, चिन्तनञ्च बोध्यते । कलाकृतिरित्याशये कलाशब्दस्य प्रयोगः ‘कं लाति’ इति व्युत्पत्त्यनुसारम् आनन्दप्रदायकेऽर्थे च भवति । कल्यते अस्यामिति भावप्रक्रिययापि मानवीयक्रियया उत्पन्ना इत्यर्थे कलाशब्दो निष्पद्यते । अतश्चिन्तनमननादिज्ञानसामर्थ्येन स्पष्टरूपेण च प्रकटीभूता निर्मितिरेव कलाशब्देन गृह्यते । पाश्चात्त्यवाङ्‌मये मानवश्रमनिर्मितानि वास्तुमूर्तिचित्रसङ्गीतकाव्यादीनि सर्वाणि वस्तूनि कलारूपेण  गृह्यन्ते । परं पौरस्त्यकाव्यशास्त्रे तु साहित्यस्य, सङ्गीतस्य, कलायाश्च पृथक् चर्चा प्राप्यते । संस्कृतविद्वांसः काव्यं साहित्यं वा कलातिरिक्तं तत्त्वं स्वीकुर्वन्ति । राजशेखरोऽपि कलाया भिन्नमेव साहित्यमुल्लिखति परं साहित्यस्य कृते कला जीवनरूपा इति स्वीकरोति। शुक्रनीतौ  कलाया महत्त्वं प्रतिपादयन्नुल्लिख्यते यत्-

\begin{shloka}[center]

\textbf{‘यद्यत्स्याद् वाचिकं सम्यक् कर्म विद्याभिसंज्ञकम् ।}

\textbf{शक्तो मूकोऽपि यत्कर्तुं कलासंज्ञं तु तत्स्मृतम् ।।’}

\end{shloka}

इत्यादिरूपेण मूकस्यापि कृते कर्मत्वेन कला परिभाषिता विद्यते । शुक्रनीतेश्चतुःषष्टिसङ्ख्यकाः कलाः, प्रबन्धकोषस्य द्विसप्ततिसङ्ख्यकाः कलाः, क्षेमेन्द्रस्य अष्टोत्तरद्विशतमिताः कलाश्चेति सर्वाः कला मतमिदं दृढीकुर्वन्ति । अधुना तु सर्जकेन श्रमसिर्जिता काऽपि  अभिव्यक्तिः ‘कला’ इति स्वीक्रियते । भौतिकं सांस्कृतिकञ्चेति प्रयोजनद्वयमाधारीकृत्य कलाया उपयोगिकला, ललितकला चेति प्रकारद्वयं  प्रतिपाद्यते । अधुना पाश्चात्त्यसमालोचकेन हेगेलमहोदयेन कृतस्य कलाभेदस्य सर्वोपयोगितया चर्चा लभ्यते । अत्र पौरस्त्यवाङ्‌मये चतुःषष्टिसङ्ख्यासु यासां कलानां प्रसिद्धिरस्ति, तासामेवाऽऽधुनिकी व्याख्या अधुना प्रचलिता इति लक्ष्यते ।

कलाशास्त्रकाव्यशास्त्रयोः समधिकं पार्थक्यं नैव दृश्यते । पाश्चात्त्यपरम्परायां काव्यं कलायामेव सङ्‌गृह्यते  । तत्र कलायाः पञ्चसु भेदेषु काव्यमपि गण्यते । परं पौरस्त्ये कलाकाव्ययोः पृथग् विवेचनं क्रियते । काव्यमभिलक्ष्य कृता लक्षणग्रन्थाः काव्यशास्त्रम्  अलङ्कारशास्त्रमिति वा नाम्ना विकसिता विद्यन्ते । कलाशास्त्रमिति नाम्ना केवलं वास्तुमूर्तिप्रभृतिहस्तनिर्मिता कला एव समीक्षिता प्राप्यते। अतः काव्यशास्त्रकलाशास्त्रयोः सामीप्यसम्बन्धो विद्यते ।

\specialupakhanda{\textbf{पौरस्त्यवाङ्मये काव्यशास्त्रम्}}

प्रतिभावन्तः कवयः काव्यनिर्माणाय प्रयतन्ते । तत्र अज्ञातपूर्वतया एव उपमादीनामलङ्काराणाम् आत्मप्रकाशो भवति । वैदिकसाहित्यं केवलं पौरत्यवाङ्‌मयस्यैव न, अपि तु विश्ववाङ्‌मयस्य कृते प्राचीनततमं मन्यते ।

\customparagraph{\textbf{वैदिकसाहित्ये काव्यशास्त्रम्}}

विश्ववाङ्मयस्य प्राचीनतमे आर्षग्रन्थे ऋग्वेदेऽपि अज्ञातपूर्वतया एवालङ्काराणामुपस्थितिर्लभ्यते । तत्र उपमा, रूपकम्, अतिशयोक्तिः, अनुप्रास इत्यादीनामलङ्काराणां सौन्दर्यं प्रकाशते । नाट्यशास्त्रकारो भरत सृष्टिकर्तुर्ब्रह्मणोऽनुरोधात् साधारणजनानां कृते सारल्येन भावग्रहणं भवेदिति विचार्य चतुर्वेदसारं सङ्‌गृह्य पञ्चमवेदाख्यं नाट्यशास्त्रं रचितमिति नाट्यशास्त्रे समुल्लिखति यत्-  	 \textbf{तस्मात्सृजापरं वेदं पञ्चमं सार्ववर्णिकम्।}

\begin{shloka}[center]

\textbf{एवं सङ्कल्प्य भगवान्सर्ववेदाननुस्मरन् ।}

\textbf{नाट्यवेदं ततश्चक्रे चतुर्वेदाङ्गसम्भवम्।।}

\textbf{जग्राह पाठ्यमृग्वेदात् सामभ्यो गीतमेव च ।}

\textbf{यजुर्वेदादभिनयान् रसानथर्वणादपि ।। (नाट्यशास्त्रम्)}

\end{shloka}

ऋग्वेदस्य संवादसूक्तेषु काव्यतत्त्वानां मूलं दृग्गोचरीभवति । तेषु च अगस्त्यलोपामुद्रासंवादः, यमयमीसंवादः, पुरुरवा-उर्वशीसंवाद इत्यादीनि प्रामुख्यं भजन्ते । तेषु रसध्वन्यलङ्कारादीनां काव्यतत्त्वानां सौन्दर्यं विलसत्येव । ऋग्वेदः पद्यकाव्यस्य, यजुर्वेदो गद्यकाव्यस्य, सामवेदो गीतिकाव्यस्य च प्राथमिकं स्रोतो  मन्यते । वेदभाष्येष्वपि आचार्यैः काव्यशास्त्रीयप्रयोगाणां स्पष्टतया सङ्केतो विहितः । वेदानां तात्पर्यं ज्ञातुं संरचितानि वेदाङ्गान्यपि क्वचित् काव्यशास्त्रीयतत्त्वानि निरूपितानि । निरुक्ते उपमावाचकानां द्वादशानां पदानां निर्वचनं लभ्यते। इत्यादिभिरुदाहरणैरलङ्कारशास्त्रस्य चर्चा वैदिककालादेव सञ्जाता इति ज्ञायते ।

\customparagraph{\textbf{पुराणे काव्यशास्त्रम्}}

पौरस्त्यकाव्यविकासपरम्परायां पौराणिकग्रन्था आधाररूपा विद्यन्ते ।  काव्यरचनायां महर्षिव्यासप्रणीतानि पुराणानि उपजीव्यग्रन्थरूपेण स्वीक्रियन्ते । इत्येव न कतिचिद् ग्रन्थेषु काव्यशास्त्रीयतत्त्वान्यपि वर्णितानि विद्यन्ते । समयस्य निश्चयाभावात् तत्र वर्णितानि काव्यशास्त्रीयतत्त्वानि पश्चवर्त्तिनामाचार्याणां कृते आधारभूतानि पूर्ववर्त्तिनां काव्यशास्त्रिणां समीक्षारूपाणि वा इति सारल्येन वक्तुं न शक्यते तथापि तेषां काव्यशास्त्रीयविवेचनं महत्त्वमावहत्येव ।

काव्यशास्त्रीयविषयवर्णितेषु पौराणिकग्रन्थेषु प्राथम्येन \textbf{अग्निपुराण}स्य नाम समायाति । तत्र ३३७ तमाद् अध्यायात् ३४७ तममध्यायं यावद् अग्निदेवः काश्यपं (वशिष्ठं) प्रति काव्यशास्त्रस्योपदेशो विहितवान् । तत्र काव्यस्वरूपम्, रसः, गुणः, दोषः, रीतिः, अलङ्कार इत्यादीनां  काव्यशास्त्रीयतत्त्वानां विस्तारपूर्विका चर्चा प्राप्यते। एतदेव नापितु नाट्यतत्त्वविषयेऽपि तत्र साङ्गोपाङ्गं समीक्षणं विद्यते ।

\textbf{विष्णुधर्मोत्तरपुराणे}ऽपि प्रायशः सर्वेषां काव्यतत्त्वानां चर्चा लभ्यते । अस्य तृतीये खण्डे मार्कण्डेयवज्रयोः संवादेऽपि काव्यशास्त्रीयविषयाः समुल्लिखिता विद्यन्ते । तत्र विष्णुधर्मोत्तरपुराणकारेण रूपकरसयोः सन्दर्भे नाट्यशास्त्रस्यानुसरणं कृतं विद्यते । अन्यत्रापि केषुचित् स्थलेषु प्राग्जातानां काव्यशास्त्रिणामनुकरणम्, कुत्रचिद् नवीनमतस्थापनमपि प्राप्यते । तृतीयखण्डस्य द्वितीयाध्यायादारभ्य शास्त्रीयविषया विवेचिता विद्यन्ते । तेषु चित्रसूत्रम्, नृत्तशास्त्रम्, गीतम्, शब्दशास्त्रम्, अलङ्कारः, महाकाव्यम्, प्रहेलिका, द्वादशरूपकम्, आहार्याभिनयः, सामान्याभिनयः, रसः, एकोनपञ्चाशद् भावाश्चोल्लिखिताः सन्ति ।

पुराणानां रचनासमयानिश्चितत्वाद् नाट्यशास्त्रं नाम ग्रन्थ एव पौरस्त्यकाव्यशास्त्रस्य प्रथमः काव्यशास्त्रीयोपलब्धिः । भरतादारभ्य आधुनिककालपर्यन्तम् अलङ्कारशास्त्रस्य विकाशमयी धारा निरन्तरं प्रवहति । एतस्मिन् समयावधौ काव्यशास्त्रस्य विवेचनाय नैकाः काव्यशास्त्रीयग्रन्थाः प्रकाशिताः । तत्र सर्वादौ भरतमुनेर्नाट्यशास्त्रं विद्यते । नाट्यशास्त्रे मुख्यतया रसनिरूपणम् नाट्यशास्त्रीयतत्त्वानां निरूपणञ्च प्राप्यते । तदनु प्राचीनाः काव्यशास्त्रिणः काव्यविषये बहूननि मतानि प्रदत्तवन्तः । काव्ये किमात्मस्थानीयं तत्त्वं भवति ? इति विषये काव्यशास्त्रिणां मतसाम्यं मतवैषम्यं च दृश्यते । भामहादयोऽलङ्कारशास्त्रिणोऽलङ्कार एव काव्यस्य प्राणभूतं तत्त्वमिति प्रतिपादितवन्तः । ध्वनिकारोऽऽनन्दवर्धनः स्वकीये ध्वन्यालोकग्रन्थे रसध्वनिरेव काव्यस्यात्मस्थानीयं तत्त्वमिति नवीनमतमुपस्थापितवान् । सोऽत्र काव्ये ध्वनिविरोधिनां मतान्युपस्थाप्य खण्डितवान् । आनन्दवर्धनस्यैव मार्गानुगामी मम्मटाचार्यः काव्यप्रकाशग्रन्थं विरचय्य ध्वनिमार्गमेव स्वीकृतवान् । कतिपय स्थानेषु विमतिं प्रदर्श्य विश्वनाथेन साहित्यदर्पणाख्यो लाक्षणिकग्रन्थोपस्थापितः । एतेषां पूर्वाचार्याणां काव्यशास्त्रीयमतेषु भिन्नमतमुपस्थापयन् पण्डितराजजगन्नाथेन रसगङ्गाधराख्यः प्रौढशास्त्रीयग्रन्थो रचितः । एते प्रमुखाः पौरस्त्यकाव्यशास्त्रिणः । एतेषां मध्येऽपि प्रौढविचारप्रतिपादकाः काव्यशास्त्रिणः सञ्जाताः । तेषु काव्यालङ्कारस्य प्रणेतुः रुद्रटस्यालङ्कारवर्गीकरणम्, काव्यादर्शस्य रचयितुर्दण्डिनो भामहमार्गानुसरणम्, काव्यालङ्कारसूत्रवृत्तेः प्रणेतुर्वामनस्य रीतिप्राधान्यम्, अलङ्कारसर्वस्वकारस्य  रुय्यकस्यालङ्काराणां वैज्ञानिकं वर्गीकरणम्, वक्रोक्तिजीवीतस्य रचयितुः कुन्तकस्य वक्रोक्तिसम्प्रदायप्रतिपादनञ्च पौरस्त्यकाव्यशास्त्रीयेतिहासस्य कृतेऽसाधारणमवदानं मन्यते । एतेषां योगदानेन पौरस्त्यकाव्यशास्त्रस्य गरिमा गुरुतरा सञ्जाता ।

\specialupakhanda{\textbf{पाश्चात्त्यवाङ्मये काव्यशास्त्रम्}}

पौरस्त्यवत् पाश्चात्त्यवाङ्मयेऽपि काव्यशास्त्रविषये विस्तृता गभीराश्च समीक्षाः समुपलभ्यन्ते । तत्र प्रारम्भिकसमये दार्शनिकप्रवरौ `प्लेटो' (Plato) तथा `अरस्तु' (Aristotle) इति नामानौ मूर्धन्यौ विचारकौ काव्यस्य शास्त्रीयपक्षस्य विवेचनं प्रारभताम् । ततश्च परवर्तिनि काले `होरेस' (Horace), `डायोनिसियस' (Dionysius), `लोन्जाइनस' (Longinus) प्रभृतयः काव्यशास्त्रिणः काव्यशास्त्रीयपरम्परां समवर्धयन् । पश्चात् प्राग्जातैर्विद्वद्भिर्बीजरूपेण समुप्तस्य काव्यशास्त्रीयविचारस्य प्रवर्धनार्थं दाँते (Dante), ब्वायलो (Boileau), लेसिङ (Lessing), वर्डस्वर्थ (Wordsworth), कलरिज (Coleridge), म्याथ्यू आर्नोल्ड (Matthew Arnold), क्रोचे (Croce), आइ. ए. रिचार्ड्स (I.A. Richards), टी. एस्. एलियट (T.S. Eliot) प्रभृतयो बहवो विख्याताः समीक्षकाः प्रादुर्बभूवुः, यैः पाश्चात्त्यसमीक्षापद्धतिरुत्कर्षं नीता ।

प्लेटोमहोदयः सत्यतत्त्वमेव परमं लक्ष्यं मत्वा नाट्यकाव्यं जगतोऽनुकृतिरूपत्वं (Mimesis) स्वीकरोति । सत्यनैतिकतत्त्वयोः सुदृढः समर्थकोऽयं दार्शनिकः काव्यकलां `सत्याद् द्विवारं दूरस्थां' (Twice removed from reality) अनुकृतेरप्यनुकरणात्मिकाञ्चेति मत्वा आदर्शगणराज्यात् (Republic) काव्यकलाया बहिष्कारं घोषयति । परम् अरस्तुमहोदयः प्लेटोमतं परिष्कुर्वन् अनुकरणस्य नवीनं सृजनात्मकं रूपं प्रस्तौति । स कथयति यत्- कोऽपि कविः केवलं प्रतिकृतिकारमेव नापितु `निर्माता' (Maker) अपि मन्यते। अनेन स्वकीये `पोएटिक्स' (Poetics) इति कालजयिनि ग्रन्थे काव्यस्य विस्तृता शास्त्रीयसमीक्षा कृता, यत्र त्रासदी (Tragedy) इत्यस्याः स्वरूपं निरूपितम् । अयमेव पाश्चात्त्यकाव्यशास्त्रे सुप्रसिद्धस्य `विरेचनसिद्धान्तस्य' (Theory of Catharsis) तथा `अनुकरणसिद्धान्तस्य' च मुख्यः प्रतिष्ठापको वर्तते । एतदतिरिक्तं रोमनविद्वान् होरेसः स्वकीये `आर्स पोएटिका' (Ars Poetica) इति ग्रन्थे काव्यस्य औचित्यं लोककल्याणञ्च प्रतिपादितवान् । अपरे च `लोन्जाइनस' महोदयेन स्वकीये `पेरि इप्सुस्' (On the Sublime) इति ग्रन्थे `उदात्ततत्त्वसिद्धान्तः' (Theory of Sublimity) प्रवर्त्तितः, येन काव्यस्य भावपक्षः सुदृढो जातः ।

पाश्चात्त्यकाव्यशास्त्रस्य परवर्तिनि युगे `वर्डसवर्थ'महोदयस्य काव्यभाषासिद्धान्तः, `कलरिज'महोदयस्य `कल्पनासिद्धान्तः' (Theory of Imagination) चेति द्वावेव सिद्धान्तौ स्वच्छन्दतावादस्य (Romanticism) सुदृढौ आधारस्तम्भौ सञ्जातौ । ततश्च लेसिङ-कान्ट-हेगेल-प्रभृतीनां जर्मनविदुषां विचारैः प्रभावित इटालीदेशीयः `क्रोचे' इति समालोचकः `अभिव्यञ्जनावादस्य' (Theory of Expressionism) प्रवर्तनं कृतवान् । एतस्य मते वस्तुनो द्रव्य-रूपात्मकौ द्वौ भागौ भवतः, तयोश्च रूपमेव सौन्दर्यस्याधारतत्त्वं भवति, इदमेवाभिव्यञ्जनायाः पर्यायभूतम् । अन्तःकरणदर्शनं प्राधान्येन स्वीकुर्वाणः क्रोचे ज्ञानस्य द्वौ भेदौ मनुते— अन्तःप्रेरणा (Intuition) विचारश्च । तत्र सहजानुभूतिः, अभिव्यञ्जना, कला च पर्यायवाचिनः शब्दाः सन्ति । कल्पनाजन्यं प्रातिभज्ञानं तथा तर्कजन्यं व्यावहारिकज्ञानञ्च भवति । एतयोर्मध्ये प्रातिभज्ञानस्यैव अभिव्यञ्जना भवति, सैव च कला वर्तते । आधुनिकसंस्कृतसमीक्षकाः क्रोचेमहोदयस्य अभिव्यञ्जनासिद्धान्तम् आचार्यकुन्तकस्य `वक्रोक्तिसिद्धान्तेन' सह बहुधा तोलनं कुर्वन्ति, तस्मात् प्रभावितं वा कथयन्ति ।

आधुनिकयुगे `आइ. ए. रिचार्ड्स'महोदयस्य काव्यसिद्धान्तो मनोवैज्ञानिकः, गभीरो व्यापकश्च मन्यते । अस्य मते काव्यं केवलं यशसे मनोरञ्जनाय वा न भवति, अपितु मानवे आवेगानां सन्तुलनम्, भावनाया उद्भावनं परिष्कारश्च काव्यस्य मुख्यं प्रयोजनम् अस्ति । अस्य `सम्प्रेषणसिद्धान्तः' (Theory of Communication) काव्यशास्त्रे वैशिष्ट्यं भजते । अयं विरेचनसिद्धान्तस्यापि नूतनां मनोवैज्ञानिकीं व्याख्यां कृतवान् । अस्य विचारोऽस्ति यत्- वयं त्रासस्य (Horror) कारणेन पलायनपथे गच्छामः, करुणया (Pity) च आकृष्टा भवामः । अनयोर्द्वयोर्विरोधीरूपाणाम् आवेगानां समन्वयेन हृदये विरेचनजन्यो मानसिकशान्तिरूप आनन्दः समुत्पद्यते । अपरतः क्रोचे-गेटे-रवीन्द्रनाथठाकुरप्रभृतयो विद्वांसः साहित्यम् `आत्माभिव्यक्तेः' (Self-expression) रूपं स्वीकुर्वन्ति, परं `टी. एस्. एलियट'महोदयस्तु एतद्विपरीतं `वैयक्तिकतावादं' (Theory of Impersonality) प्रतिष्ठापयति । अस्य मते कवेर्व्यक्तित्वं तस्य कृतिश्च सर्वथा पृथगेव भवतः; अर्थाद् व्यक्तिगतो भावः काव्यगतो भावश्च भिन्न एव। काव्यनिर्माणकाले कविः स्वकीयेभ्यः व्यक्तिगतरागोद्वेगादिभ्यः तटस्थो भवति, अतः कला आत्माभिव्यक्तिर्नहि, अपितु आत्माभिव्यक्तेः पलायनमस्ति (Escape from personality) । एलियटमहोदयेन उपस्थापितौ `परम्परासिद्धान्तः' (Theory of Tradition) तथा `वस्तुनिष्ठप्रतिरूपकसिद्धान्तः' (Objective Correlative) चेति द्वावेव सिद्धान्तौ आधुनिकसमीक्षापद्धतौ महदुपकारकौ सिद्धौ स्तः । एवम्प्रकारेण कलाप्रधानपाश्चात्त्यकाव्य- शास्त्रीयेतिहासः समृद्धतरो ज्ञायते ।

सामान्यतया देश-काल-परिधिवशात् संज्ञाभेदे सत्यपि पौरस्त्य-काव्यशास्त्रस्य पाश्चात्त्य-काव्यशास्त्रस्य च मूललक्ष्यं समानमेव वर्तते । यत्र पौरस्त्याचार्या रस-ध्वनि-वक्रोक्त्यादिभिः सिद्धान्तैः काव्यस्यान्तरात्मानं भावपक्षञ्च परिचिन्वन्ति, तत्रैव पाश्चात्त्याः समीक्षकाः अनुकरण-विरेचन-कल्पना-वैयक्तिकतादिभिर्दृष्टिभिः कलाया मनोवैज्ञानिकं व्यावहारिकञ्च स्वरूपं विश्लेषयन्ति । भाषाया आवरणं भिन्नं स्यात्, परं काव्यसौन्दर्यस्य रसस्वादस्य च अनुभूतिः सार्वभौमिका सार्वजनीना च भवति । अतः साहित्यस्य व्यापकसिद्धान्तानां सम्यगवबोधाय, तुलनात्मकसमीक्षापद्धतेः विकासाय च पौरस्त्यकाव्यशास्त्रीयविषयाध्ययनक्रमे पाश्चात्त्यकाव्यशास्त्रस्य अध्ययनं नितरामावश्यकं प्रतिभाति।

\customparagraph{\textbf{पौरस्त्यपाश्चात्त्ययोः काव्यशास्त्रीयविषययोस्तुलना}}

पौरस्त्य-पाश्चात्त्ययोः काव्यचिन्तनं विलोक्यते चेद् उभयोश्चिन्तनेषु मौलिकं वैलक्षण्यं दार्शनिकं गम्भीरत्वञ्च वैशिष्ट्यं परिलक्ष्यते । पाश्चात्त्यजगति काव्यं `कला' (Art) इत्यस्यैव अङ्गरूपेण स्वीकृत्य तस्य विवेचनं भौतिक-मनोवैज्ञानिक-सामाजिकदृष्ट्या कृतमस्ति, परं पौरस्त्यपरम्परायां काव्य-कलयोः सूक्ष्मस्तात्त्विकश्च भेदो विहितः । पौरस्त्यविदुषां मते काव्यस्य सम्बन्ध अलौकिक-रसचेतनया सह वर्तते, अतः काव्यजन्य आनन्दः `ब्रह्मानन्दसहोदरः' इति गीयते । कलानां (शिल्प-सङ्गीत-चित्रकलादीनां) सम्बन्धस्तु मुख्यतया व्यावहारिक-सौन्दर्येण सह ऐहिकानन्देन च भवति । पाश्चात्त्यचिन्तने कलाया व्यापकतत्त्वे एव काव्यं समाविशति, यदा तु पौरस्त्यचिन्तने काव्यं कलाभेदेभ्य उच्चतरम्, पुरुषार्थचतुष्टयप्रदञ्च विशिष्टं स्थानं भजते । गुणालङ्कारादीनां बाह्यसौन्दर्यतत्त्वानां विवेचने उभयत्र साम्यता परिलक्ष्यते, तथापि काव्याङ्गानां सूक्ष्मातिसूक्ष्मवर्गीकरणदृष्ट्या पौरस्त्यकाव्यशास्त्रं समधिकं सुव्यवस्थितं वैज्ञानिकञ्च दृश्यते ।

पौरस्त्यसाहित्यशास्त्रे आचार्यैः काव्यस्य स्वरूपम्, हेतुः, प्रयोजनम्, शब्दशक्तयः, दोष-गुणालङ्काराणां वर्गीकरणञ्च येन प्रकारेण साङ्गोपाङ्गं विश्लेषितम्, न तथा पाश्चात्त्यकाव्यशास्त्रे प्राचीनकाले समुपलभ्यते । पाश्चात्त्यानां मते भाव-विचार-कल्पना-शैलीप्रभृतयो विषया एव काव्ये प्रधानत्वेन अपेक्ष्यन्ते, परं पौरस्त्यानां मते तु रस-ध्वनि-अलङ्कार-रीति-गुण-औचित्य-वक्रोक्तिरिति षट्‌सम्प्रदायानां माध्यमेन काव्यस्य बाह्यावरणम् अन्तरात्मा चेति द्वावेव सम्यक् निरूपितौ स्तः । पाश्चात्त्यजगतोऽनुकरण-विरेचन-उदात्तता-प्रभृतयः सिद्धान्ताः काव्यस्य बाह्यप्रभावं प्रक्रियां वा बोधयन्ति, यदा तु पौरस्त्यानां रसध्वन्यादयः सिद्धान्ताः सहृदयस्य अन्तश्चेतनां तरलीकुर्वन्ति ।

ऐतिहासिकविकासक्रमदृष्ट्या पौरस्त्यकाव्यशास्त्रस्य उत्तरकालस्तु किञ्चित् सङ्कोचकाल इव  प्रतीयते । पण्डितराजजगन्नाथानन्तरं पौरस्त्यपरम्परायां मौलिकसिद्धान्तप्रतिपादका नवीनपथप्रदर्शका आचार्याः प्रायो नोत्पन्नाः । केवलं प्राचीनसिद्धान्तानां व्याख्या-टीकासम्बद्धग्रन्था एव प्रतिपादिताः । एतद्विपरीतं पाश्चात्त्यजगति उत्तरआधुनिककाले नूतना विचारवाहिनः, समालोचनापद्धतयश्च (यथा— संरचनावादः, उत्तरसंरचनावादः, मनोविश्लेषणवादः, मार्क्सवादः इत्यादयः) निरन्तरं प्रवहन्त्यः सन्ति, येन वर्तमानकाले पाश्चात्त्यसमीक्षापद्धतिर्बहुमुखी गतिशीला च दृश्यते तथापि प्राचीनकालादारभ्य पौरस्त्यकाव्यशास्त्रेण या उपलब्धय आविष्कृताः, ता विश्ववाङ्मयस्य विचारबीजप्रदानदृष्ट्या पाश्चात्त्यम् अतिरिच्यैव वर्तन्ते । पौरस्त्यकाव्यशास्त्रस्य गम्भीरता, रस-मीमांसा, ध्वनिसिद्धान्तस्य दार्शनिकपृष्ठभूमिश्च विश्वसाहित्ये अद्यापि अद्वितीया एव अस्ति ।

\specialupakhanda{\textbf{कालविभागः}}

कवेः कर्मभूतस्य काव्यस्य शास्त्रीयपक्षस्य चर्चा वैदिककालादेव समुपलभ्यते । तदानीं वेदपुराणादिग्रन्थेषु काव्यशास्त्रीयतत्त्वानां सङ्केताः प्राप्यन्ते, तथापि तस्मिन् युगे काव्यशास्त्रस्य स्वरूपं सुस्पष्टं नासीत् । अग्निपुराणे विद्यमानं काव्यशास्त्रीयविवेचनप्रकारं विलोक्य कतिपये समालोचकाः अस्य ग्रन्थस्य वर्तमानं रूपं परवर्तिकालोनं मन्यन्ते, अतः पौराणिककालस्य ग्रन्थत्वेन अस्य प्राचीनता न सर्वैः समाद्रियते । तस्मात् काव्यशास्त्रीयविषयस्य प्रामाणिको व्यवस्थितश्च प्रारम्भः नाट्यशास्त्रप्रणेतुर्भरतमुनेः समयादेव स्वीक्रियते । नाट्यशास्त्रकर्तुर्भरतमुनेः समयः विक्रमपूर्वद्वितीयशताब्दीमारभ्य विक्रमपश्चात्-द्वितीयशताब्दीं यावत् निश्चीयते, अत एव तस्मात् कालात् काव्यशास्त्रस्य सैद्धान्तिकस्वरूपस्य बीजारोपणं सञ्जातमिति मन्तुं शक्यते ।

एतादृशान् विविधान् ऐतिहासिकतथ्यान् विचिन्त्य विद्वांसः काव्यशास्त्रविकासस्य समयं मुख्यतश्चतुर्षु भागेषु विभजन्ति । तदत्र उपस्थाप्यन्ते यथा-

\begin{itemize}
\item \textbf{प्रारम्भकालः (अज्ञातकालात्- वि.सं. ६०० यावत्) :} वैदिकादिसङ्केतकालादारभ्य भामहपूर्वकालं यावत् ।


\item \textbf{रचनाकालः (वि.सं. ६०० - ८०० यावत्) :} आचार्यभामहादारभ्य आनन्दवर्धनाचार्यपूर्वकालं यावत् ।


\item \textbf{निर्णयात्मककालः (वि.सं. ८०० - १००० यावत्) :} ध्वनिकाराद् आनन्दवर्धनाचार्यात् आरभ्य आचार्यमम्मटपर्यन्तम् ।


\item \textbf{व्याख्याकालः (वि.सं. १००० - १७०० यावत्) :} मम्मटसमयादारभ्य जगन्ननाथपर्यन्तम् ।


\end{itemize}
वैक्रमाब्दसप्तदशशततः (वि.सं. १७००) अनन्तरं संस्कृतकाव्यशास्त्रस्य इतिहासे नवीनसिद्धान्तप्रतिपादका मौलिकग्रन्थाः प्रायो दुर्लभा एव, अतः समालोचकैरयं समयः `क्षयकालः' इत्युद्घोषितः । इममेव तथ्यं पुरस्कृत्य समालोचकैः पृथक् कालगणना संयोजिता अस्ति, यथा-

\begin{itemize}
\item \textbf{ह्रासकालः (वि.सं. १७०० तमाद् वैक्रमाब्दाद् अद्यावधि) :} पण्डितराजजगन्नाथसमयादारभ्य अधुनापर्यन्तम् ।


\end{itemize}
एतादृशे सार्वत्रिके कालविभागे सति, ध्वनिवादिनः केवलं ध्वनेः विकासाधारेण पृथक् कालविभागं कुर्वन्ति । तेषां मतानुसारेण ध्वनिसिद्धान्त एव काव्यशास्त्रस्य प्रधानभूतो वर्तते, अतस्ते निम्नानुसारं कालविभाजनं प्रस्तुवन्ति, यथा-

\begin{itemize}
\item \textbf{पूर्वध्वनिकालः (अज्ञातकालात् वि.सं. ८०० यावत्) :} आनन्दवर्धनकालं यावत् ध्वनितत्त्वस्य प्रस्फुटनकालः ।


\item \textbf{ध्वनिकालः (वि.सं. ८०० ततः १००० यावत्) :} आनन्दवर्धनकालात् मम्मटकालेऽन्तर्भूतः सिद्धान्तप्रतिष्ठाकालः ।


\item \textbf{पश्चाद्ध्वनिकालः (वि.सं. १००० ततः १७०० यावत्) :} मम्मटकालात् पण्डितराज-जगन्नाथकालं यावत् व्याख्याकालः ।


\end{itemize}
अस्य कालविभागस्य प्रवर्त्तकाः केवलं ध्वनिमेव काव्यात्मत्वेन अङ्गीकुर्वन्ति, अतो मतमिदं काव्यशास्त्रस्य व्यापकतादृष्ट्या एकपक्षीयत्वाद् न सर्वमान्यं प्रतीयते । तस्माद् युक्तायुक्तविवेचनानन्तरं प्रायः सर्वैः अङ्गीकृतत्वात् प्रथममतस्याधारेणैव काव्यशास्त्रस्य व्यवस्थितं कालविभागं स्वीकृत्य, तत्तत्कालिकानां आचार्याणाम् अवदानदृष्ट्या अलङ्कारशास्त्रस्य सङ्क्षिप्तं विवेचनम् अत्र विधीयते ।

\specialmahakhanda{\textbf{द्वितीय-प्रकाशः}}

\specialkhanda{\textbf{प्रारम्भकालः}}

\begin{center}

\textbf{(अज्ञातकालात् — वि.सं. ६०० तमवर्षपर्यन्तम्)}

\end{center}

कालोऽयं पौरस्त्यकाव्यशास्त्रस्य बीजरूपो मन्यते । अत्र वैदिकग्रन्थेषु पुराणग्रन्थेषु च काव्यशास्त्रसम्बन्धिनो बहवः शब्दाः प्रयुक्ता दृश्यन्ते । अस्मिन् काले शिलाली, कृशाश्व, वृद्धभरतः, तण्डुः, नन्दिकेश्वरः, मतङ्गः, भरतः, भरतकोहल, वत्तिलः, नखकुट्टः, अश्मकुट्ट, विष्णुगुप्तः, काश्यपः, वररुचिः, ब्रह्मदत्तः, नन्दिस्वामी, मेधावीरुद्रः, भट्टी, श्रीमत्स्थविरः’ इत्यादयः काव्यशास्त्रमर्मज्ञा विद्वांसः सञ्जाताः । एतेषु मेधावीरुद्रस्यावदानमुल्लेखनीयं प्रतिभाति, तथापि स्पष्टं शास्त्रीयविवेचनं तु भरतमुनेरेव प्राप्यते । अतो भरतमुनिरेवास्य कालस्य प्रमुखाचार्यो मन्यते । भरतमुनिना प्रणीतं नाट्यशास्त्रं लौकिकसाहित्यस्य कृते महता उपकारेण स्वीक्रियते । ग्रन्थेऽस्मिन् भरतमुनिः काव्यशास्त्रीयतत्त्वानां विशदं विवेचनं सुबोध्यतया प्रस्तौति  ।

\specialupakhanda{\textbf{भरतमुनिः}}

संस्कृतसाहित्ये शास्त्रीयपरम्परायाः प्रारम्भकर्ता भरतमुनिरस्ति । पौरस्त्यपौराणिकेतिहासे पञ्च भरतनामका जनाः प्रसिद्धाः सन्ति । तेषु नाट्यशास्त्रकर्ता भरतमुनिः पृथग् आसीत् । कालिदासानुसारं नाट्यशास्त्रकारो भरतमुनिर्देवलोकेऽपि नाटकानि प्रदर्शितवान् । तस्योल्लेखो विक्रमोर्वशीयनाटके प्राप्यते यत्-

\begin{shloka}[center]

\textbf{मुनिना भरतेन यः प्रयोगो भवतीष्वष्टरसाश्रयप्रयुक्तः ।}

\textbf{ललिताभिनयं तमद्य भर्ता मरुतां द्रष्टुमना सलोकपालः ॥ २/१८}

\end{shloka}

भरतमुनिना नाट्यशास्त्रे केवलं चत्वार एवालङ्कारा उल्लिखिताः । पश्चाज्जातानामलङ्कारशास्त्रिणामलङ्कारसङ्‌ख्या त्रिंशदधिका प्राप्यते । अलङ्कारविवेचनदृष्ट्या कालिदासस्योल्लेखनेन चास्य समयः कालिदासप्राग् विक्रमपूर्वद्वितीयशताब्द एव मन्यते । अनेन सर्वसाधारणजनस्य कृते पञ्चमवेदरूपं नाट्यवेदार्थं नाट्यशास्त्रं चकार’ इति अस्यैव नाट्यशास्त्राद् ज्ञायते । इदं सत्यमेवास्ति यत् नाट्यशास्त्रस्याधारभूमौ लौकिकसाहित्यशास्त्रस्य सुगठितो विकासः सञ्जातः । नाट्यशास्त्रस्य मुख्यमुद्देश्यं नाट्यविवेचनमेवास्ति । किन्तु तत्र छन्दःशास्त्रस्य, अलङ्कारशास्त्रस्य, सङ्गीतशास्त्रस्य चापि विवेचनं समुपलभ्यते । नाट्यशास्त्रस्य षट्‌त्रिंशद् अध्यायेषु ये विषया विवेचिताः सन्ति, तेषां नामोल्लेखनमत्र क्रियते । तत्र \textbf{प्रथमेऽध्याये} नाट्यशास्त्रस्योत्पत्तिः, \textbf{द्वितीयेऽध्याये} नाट्यशालाया निरूपणम्,  \textbf{तृतीयेऽध्याये} रङ्गदेवतासमर्चा, \textbf{चतुर्थेऽध्याये} ताण्डवविषयकाणाम्  अष्टोत्तरशतानां करणानां द्वात्रिंशदङ्गहाराणां च वर्णनम्, \textbf{पञ्चमेऽध्याये} पूर्वरङ्गस्य विस्तृतं विधानम्, \textbf{षष्ठेऽध्याये} रसस्य विस्तृतं विवेचनम्, \textbf{सप्तमेऽध्याये} भावानां विस्तृतं विवेचनम्, \textbf{अष्टमेऽध्याये} उपाङ्गैरभिनयवर्णनम्, \textbf{नवमेऽध्याये} हस्ताभिनयस्य वर्णनम्, \textbf{दशमेऽध्याये} शरीराभिनयस्य वर्णनम्, \textbf{एकादशेऽध्याये} पारिविधानम्, \textbf{द्वादशेऽध्याये} मण्डलविधानम्, \textbf{त्रयोदशेऽध्याये} रसानुकूलगते प्रचारवर्णनम्, \textbf{चतुर्दशेऽध्याये} प्रवृत्तधर्मस्य व्यञ्जनावर्णनम्, \textbf{पञ्चदशेऽध्याये} छन्दसां विभागवर्णनम्, \textbf{षोडशेऽध्याये} वृत्तानां सोदाहरणं लक्षणम्, \textbf{सप्तदशेऽध्याये} वागभिनयस्य लक्षणालङ्काराणां काव्यदोषाणां काव्यगुणानाञ्च वर्णनम्, \textbf{अष्टादशेऽध्याये} भाषाभेदाभिनयप्रयोगयोः वर्णनम्, \textbf{एकोनविंशेऽध्याये} नाटकीयानां पञ्चसन्धिनां सन्ध्यङ्गानाञ्च वर्णनम्, \textbf{द्वाविंशेऽध्याये} चतसृणां वृत्तीनां विधानम्, \textbf{त्रयोविंशेऽध्याये} आहार्याभिनयवर्णनम्, \textbf{चतुर्विंशेऽध्याये} सामान्याभिनयवर्णनम्, \textbf{पञ्चविंशेऽध्याये} बाह्यानामुपचाराणां वर्णनम्, \textbf{षड्विंशेऽध्याये} चित्राभिनयस्य वर्णनम्, \textbf{सप्तविंशेऽध्याये} सिद्धिव्यञ्जनानां निर्देशः, \textbf{अष्टाविंशे} सङ्गीतशास्त्रस्य वर्णनम् तथा आतोद्यवर्णनञ्च, \textbf{एकोनत्रिंशेऽध्याये} ततातोद्यवर्णनम्, \textbf{त्रिंशेऽध्याये} सुषिरातोद्यविधानवर्णनम्, \textbf{एकत्रिंशेऽध्याये} तालवर्णनम्, \textbf{द्वात्रिंशेऽध्याये} ध्रुवाविधानम्, \textbf{त्रयस्त्रिंशेऽध्याये} वाद्यस्य विस्तृतं वर्णनम्, \textbf{चतुर्त्रिंशेऽध्याये} प्रकृतिविचारवर्णनम् \textbf{पञ्चत्रिंशेऽध्याये} भूमिकारचनायाः वर्णनम्, \textbf{षट्त्रिंशे चाध्याये} नाट्यस्य भूतलेऽवतरणस्य विस्तृतं विवेचनं कृतं वर्तते ।

\customparagraph{\textbf{काव्यशास्त्रीयं योगदानम्}}

भरतमुनिना नाटकमुद्दिश्य कृतं शास्त्रीयविवेचनं काव्यशास्त्रस्य कृते विकासपूर्वाधारार्थं मार्गदर्शकं जातम् । पश्चवर्त्तिभिराचार्यैः एतानेव नाट्यशास्त्रीयान् विषयानङ्गीकृत्य सपरिष्कारं  प्रसिद्धाः शास्त्रीयग्रन्थाः रचिताः । तेषु रसविवेचनप्रसङ्गे प्रस्तुतं रससूत्रं तु पौरस्त्यकाव्यशास्त्रस्य प्राणभूतमस्ति । अस्य रससूत्रमेव रसप्रतीतिविषये प्रामाणिकाधारोऽस्ति । अस्मिन्नेवाधारे विद्वांसः स्वकीयैः प्रौढैस्तर्कशक्तिभिरुत्पत्तिरनुमिति–भुक्तिरभिव्यक्तिश्चेति प्रसिद्धा रसनिष्पत्तिविषयकाश्चत्वारः सिद्धान्ताः प्रतिपादिताः । ते च रससम्प्रदाये उत्तरोत्तरेणोत्कृष्टाः मान्याश्च सन्ति । आचार्यधनञ्जयस्य दशरूपकम् इत्यस्योपजीव्यमपि नाट्यशास्त्रमेव । तथा च नाटकछन्दोरसगुणादीनां विवेचने विद्वांसो नाट्यशास्त्रमाधारीकृत्य स्वकीयान् विचारान् आविष्कुर्वन्ति । एतादृशैः प्रमाणैर्वक्तुं शक्यते यत्— भरतस्य नाट्यशास्त्रमेव संस्कृतकाव्यशास्त्रस्य मेरुदण्डत्वेन सुशोभते ।

\customparagraph{\textbf{दृश्यकाव्यविवेचनम्}}

पौरस्त्यवाङ्मयान्तर्गतस्य दृश्यकाव्यस्यैव न, अपितु संस्कृतकाव्यशास्त्रस्यैव प्राथमिकी प्राप्तिर्भरतमुनेर्नाट्यशास्त्रमस्ति । ग्रन्थोऽयं बीजत्वेन स्वीकृत्य पाश्चात्त्यसमालोचकैरपि नवीनविचाराः प्रतिपादिता दृश्यन्ते । नाट्यशास्त्रे नाटकमुद्दिश्यैव अन्ये शास्त्रीयसिद्धान्ता निरूपिताः  सन्ति । तत्र रसालङ्कारगुणादीनां काव्यातत्त्वानामपि नाट्यार्थमेव विवेचनं कृतम् । प्रजापतेरुपदेशेन  वैदिकज्ञानार्जनासमर्थानां जनसाधारणानां कृते पञ्चमवेदमहत्त्वाधायकस्य  नाट्यशास्त्रस्य रचनार्थं भरतमुनिः प्रेरितो जातः’ इति स एव नाट्यशास्त्रे उल्लिखति । तत्र नाट्यकाव्ये प्रयोगयोग्यानां सर्वेषां तत्त्वानां विवेचनं प्राप्यते । मुनिरेव तत्रोल्लिखति यत्–

\begin{shloka}[center]

\textbf{न तज्ज्ञानं न तच्छिल्पं न सा विद्या न सा कला ।}

\textbf{न स योगो न तत्कर्म यन्नाट्येऽस्मिन्न दृश्यते ।। १/११६}

\end{shloka}

तेषु रसविषयको मुनेर्विचारस्तु पौरस्त्यकाव्यशास्त्रस्यैव कृते महत्तमोपलब्धिर्मन्यते । पश्चाज्जातानाम् आचार्याणां कृतेऽस्य शास्त्रीयविचारः पथप्रदर्शकरूपेण स्मरणीयो जातः । धनञ्जयस्य दशरूपकाख्यो नाट्यसिद्धान्तग्रन्थः, विश्वनाथस्य साहित्यदर्पणस्य नाट्यविवेचनात्मकः षष्ठपरिच्छेदश्च भरतमुनेरेव विचाराधारे निर्मितो जातः । अतः नाट्यकाव्यस्य विषये भरतमुनेर्योगदानमविस्मरणीयमस्तीति मन्तुं युज्यते ।

\customparagraph{\textbf{अलङ्कारविचारः}}

आचार्यो भरतमुनिरलङ्कारस्य विषये सामान्यरूपेण उल्लिखति । स नाट्यकाव्यार्थसौन्दर्यप्रवर्धकान् चतुरोऽलङ्कारान्नेव प्रस्तौति । यथा—

\begin{shloka}[center]

उपमा रूपकं चैव यमकं दीपकं तथा ।

\end{shloka}

\customparagraph{\textbf{गुणग्रहणम्}}

गुणस्य विषये भरतमुनेर्विचारः प्राचीनमतानुयायिनां कृते आधारोऽस्ति । अस्यैव मताधारे दण्डीवामनादयः स्वकीयानि मतानि प्रस्तुवन्ति । एतस्य गुणसम्बन्धिविचारोऽस्ति यत्—

\begin{shloka}[center]

\textbf{श्लेषप्रसादसमतासमाधिमाधुर्यमोजपदसौकुमार्यम् ।}

\textbf{अर्थस्य च व्यक्तिरुदारता च कान्तिश्च काव्यार्थगुणा दशैते ।।}

\end{shloka}

एते भरतस्य दशगुणा यथास्यात्तथा प्राचीनमतानुयायिनो वामनादयः स्वीकुर्वन्ति ।

\customparagraph{\textbf{रसनिरूपणम्}}

रसविषये काव्यशास्त्रीयसमाजेऽस्य नाट्यशास्त्रे प्रस्तुतस्य रससूत्रस्य सर्वाधिकं महत्त्वमस्ति । एतस्यैवाधारे विभिन्नाः सिद्धान्ताः विकसिताः । सर्वे रसवादिनो रसनिष्पत्तिविषये रससूत्रमेवाधारत्वेन स्वीकुर्वन्ति । काव्यशास्त्रीयेतिहासे सर्वाधिकप्रसिद्धस्य रससम्प्रदायस्य संस्थापकरूपेण अस्यैव नाम गृह्यते ।

\customparagraph{\textbf{रससम्प्रदायसंस्थापकः}}

पौरस्त्यकाव्यशास्त्रीयेतिहासे रससम्प्रदायस्य सर्वोत्कृष्टं स्थानमस्ति । नाट्यशास्त्रप्रणेता भरतमुनी रसविषयेऽऽधारस्तम्भरूपं सूत्रात्मकं मतं प्रदर्श्य रससम्प्रदायस्य ‘संस्थापकः’ इति सर्वोत्कृष्टविशेषणेन विभूषितो जातः। अस्य रसनिष्पत्तिविषयकं सूत्रं विद्यते यत्-  `विभावानुभावव्यभिचारिसंयोगाद्  रसनिष्पत्तिः’ इति । अस्यैव रससूत्रस्य व्याख्यानक्रमे श्रीशङ्कुक–भट्टलोल्लट–भट्टनायक– अभिनवगुप्ताचार्याणामुत्पत्ति–अनुमिति–भुक्ति–अभिव्यक्तिश्चेति चत्वारो रसनिष्पत्तिविषयकाः सिद्धान्ता आविर्भूताः । एतेषु अभिव्यक्तिवादस्य प्रणेता अभिनवगुप्तस्य मतमेव समीचीनमुत्तरवर्त्तिभिराचार्यैरनुमोदितञ्च ज्ञायते । अस्यैव रससूत्राधारे उत्तरवर्त्तिभी  रसवादिभिः  समालोचकै रसस्वरूपनिष्पत्तिविषये खण्डनमण्डनादिकं विनैव स्वीकृतत्वाद् भरतमुनिं रससम्प्रदायस्य संस्थापकरूपेण आद्रियन्ते ।

आचार्यभरतमुनेर्नाट्यशास्त्रं केवलं नाट्यकलाया मार्गदर्शकमेव नापितु समग्रस्य संस्कृतकाव्यशास्त्रस्य, सङ्गीतशास्त्रस्य, छन्दःशास्त्रस्य च अक्षया आधारशिला वर्तते । मुनिना प्रणीतं ‘विभावानुभावव्यभिचारिसंयोगाद् रसनिष्पत्तिः’ इति सर्वमान्यं रससूत्रं पौरस्त्यकाव्यशास्त्रस्य प्राणभूतं तत्त्वमस्ति, यमधिकृत्य भट्टलोल्लट-श्रीशङ्कुक-भट्टनायक-अभिनवगुप्तादिभिः प्रवरैराचार्यैः स्वकीयप्रौढतर्कैर्रसनिष्पत्तेर्व्यापकरूपेण व्याख्या कृता । भरतमुनिना निरूपिताश्चत्वारोऽलङ्काराः, दश गुणाश्च परवर्त्तिनां भामह-दण्डि-वामनाद्याचार्याणां कृते नवीनकाव्यसिद्धान्तप्रतिपादनस्य मुख्योपजीव्यमभूवन् । “न तज्ज्ञानं न तच्छिल्पम्” इति मुनेः उद्घोषानुसारं नाट्यशास्त्रं मानवजीवनस्य विविधज्ञानविज्ञानस्य कलात्मकं सङ्गमस्थलम् अस्ति । धनञ्जयकृत-दशरूपकादारभ्य विश्वनाथकृत-साहित्यदर्पणपर्यन्तं प्रायः सर्वेऽपि दृश्यकाव्यविवेचका आचार्या भरतमुनेरेव सिद्धान्तान् आदृतवन्तः । अतः पौरस्त्यसाहित्यपरम्परायाम् आदिगुरुत्वेन प्रतिष्ठितस्य भरतमुनेर्नाट्यशास्त्रमिति महान् ग्रन्थो संस्कृतवाङ्मयस्य मेरुदण्डत्वेन सुशोभते । यस्य आलोकं विना पौरस्त्यसमालोचनापद्धतेरितिहासः सर्वथा अपूर्ण एव । पश्चाद्वर्तिनः सर्वेऽपि काव्यशास्त्रीयसम्प्रदायाः प्रत्यक्षं परोक्षं वा नाट्यशास्त्रस्य ऋणाद् मुक्ता भवितुं न शक्नुवन्तीति अवितथं वाक्यम् ।

\specialupakhanda{\textbf{मेधाविरुद्रः}}

नैकैः काव्यशास्त्रिभिः स्वग्रन्थस्योद्धरणेषु मेधाविरुद्रस्य नामोल्लेखः कृतः । विद्वानयं जन्मान्ध आसीदिति राजशेखरः स्वग्रन्थस्य प्रतिभायाः प्रभावनिरूपणप्रसङ्गे उल्लिखति यत्—‘प्रत्यक्षप्रभावतः पुनरपश्यतोऽपि प्रत्यक्ष इव यतो मेधाविरुद्रकुमारदासादयो जात्यन्धाः श्रूयन्ते ।’ नमिसाधुकथनानुसारेणापि मेधाविरुद्रः अलङ्कारशास्त्रस्य प्रणेता आसीत्, यथा– काव्यालङ्कारटीकायां ‘ननु दण्डीमेधाविरुद्रभामहकृतानि सन्त्येवालङ्कारशास्त्राणि ।’ एतस्यानुसारेणैव मेधाविरुद्रः ‘नामाख्यातोपसर्गनिपाताः’ इति चतुरः शब्दप्रकारान् स्वीकरोति । आचार्यो भामहोऽपि स्वग्रन्थे उपमादोषप्रकरणे ‘सप्त मेधाविनोदिताः’ इति उल्लिखति । एभिः प्रमाणैः स्पष्टो भवति यत्— मेधाविरुद्रः भरतभामहयोर्मध्ये एको महान् आचार्य आसीत् । अस्य प्रामाणिकग्रन्थस्तु अधुना नोपलभ्यते, परं पश्चवर्त्तिनामाचार्याणामुल्लेखेन काव्यशास्त्रेऽस्य महद्  योगदानं दृश्यते ।

\specialupakhanda{\textbf{प्रारम्भकाल-उपसंहारः}}

पौरस्त्यकाव्यशास्त्रपरम्परायाः प्रारम्भकाल उत्तरवर्तिनामलङ्कारशास्त्रिणां कृते आधारत्वेन स्वीक्रियते । अज्ञातकाले सञ्जातानां साङ्केतिकतत्त्वानां सङ्‌ग्रहेण विधिवत् समारम्भ एव प्रारम्भिककाले भवति । पौरस्त्यालङ्कारशास्त्रस्य प्रारम्भकाल यदा सञ्जातः, तस्मात् कालात् प्राक् वैदिकपौराणिककालयोः प्रणीतेषु ग्रन्थेषु शास्त्रीयतत्त्वस्य सङ्केतादिकं लभ्यते । प्रारम्भकालस्य समयावधिरेव अज्ञातकालाद् भामहस्योदयं यावद् निर्धारितो विद्यते । अज्ञातकालेऽपि बहवः काव्यशास्त्रमर्मज्ञाः सञ्जाता इति  भरतभामहप्रभृतयः काव्यशास्त्रिणः ससम्मानं तेषां नामानि उल्लिखन्ति । सैद्धान्तिकरूपेण प्रारम्भकालस्य प्रवर्धकस्तु भरतमुनिरेव मन्यते । यद्यप्यनेन नाट्यकाव्यमुद्दिश्य शास्त्रीयो ग्नन्थो रचितस्तथापि तत्र सर्वेषां काव्यभेदानां बीजतत्त्वानि समुल्लिखितानि विद्यन्ते । पश्चाज्जातेषु कालेषु रसालङ्कारादयो ये सम्प्रदाया विकसिताः,  तेषां सर्वेषां बीजं तु नाट्यशास्त्रमेव । भरतस्य रसविषयकं विवेचनं तु पश्चवर्त्तिनां कृते पाठ्यविषयमिवाचरति । तत्रैवाश्रिता विविधाः सिद्धान्तवादा विकसिताः । अलङ्कारगुणादिविषयेऽपि उत्तरवर्त्तिनो भरतमुपजीव्यत्वेन स्वीकुर्वन्ति । अस्मिन्नेव काले मतङ्गमेधाविरुद्रकाश्यपप्रभृतयः शास्त्रीयविषयमर्मज्ञाः सञ्जाता इति यत्र तत्र तेषामवदानोल्लेखेन ज्ञायते परं स्पष्टतया तेषां ग्रन्थादिकं नोपलभ्यते । सामान्यतया पौरस्त्यालङ्कारशास्त्रस्याधारभूमित्वेन प्रारम्भकाल इति मन्यते ।

\specialmahakhanda{\textbf{तृतीय-प्रकाशः}}

\specialupakhanda{\textbf{रचनात्मकः कालः}}

\begin{shloka}[center]

\textbf{(भामहादारभ्य आनन्दवर्धनाचार्यपूर्वकालः)}

\end{shloka}

\begin{shloka}[center]

\textbf{(६०० तमात् ८०० तमपर्यन्तम्)}

\end{shloka}

काव्यशास्त्रीयकालविभागस्य समयोऽयं संस्कृतकाव्यशास्त्रस्य सैद्धान्तिकविकासदृष्ट्या सातिशयो महत्त्वपूर्णोऽस्ति ।  भामहादारभ्य आनन्दवर्धनपर्यन्तमर्थात् ६०० वैक्रमाब्दादारभ्य ८०० वैक्रमाब्दं यावत् समयस्यास्य प्रसरो विद्यते । अस्मिन्नेव काले चतुर्ण्णां सम्प्रदायानां विकासो  मौलिकग्रन्थानां  प्रणयनञ्चाभवत् । रसरीतिध्वन्यलङ्कारसम्प्रदायानां स्थापनयैव विश्ववाङ्मये संस्कृतकाव्यशास्त्रस्य महत्ता प्रवर्धिता । उद्भटरुद्रटादयोऽलङ्कारसम्प्रदायस्य, दण्डीवामनादयो रीतिसम्प्रदायस्य, भट्टलोल्लटभट्टनायकशङ्कुकादयो रससम्प्रदायस्याचार्याः सञ्जाताः, तथा च ध्वनिसम्प्रदायस्य प्रवर्तक आनन्दवर्धनाचार्योऽपि अस्मिन्नेव समये सञ्जातः । एतेषां सम्प्रदायानां मौलिकग्रन्थनिर्मातारस्तेषां समर्थकाश्च विद्वांसोऽस्मिन् समये समधिकाः सञ्जाताः । तेषां विदुषां परिचयोऽवदानञ्चात्र कालक्रमेण विविच्यते ।

\specialupakhanda{\textbf{भामहः}}

पौरस्त्यकाव्यशास्त्रेतिहासे भामहः प्रारम्भिक आचार्योऽस्ति । प्राग् जातेन भरतमुनिना दृश्यकाव्यमेवोद्दिश्य सिद्धान्ताः प्रतिपादिताः  । परं भामहेन काव्यस्यालङ्कारादीन् सौन्दर्यपक्षान् अवलम्ब्य विस्तृतं विवेचनं कृतम् । अत एवानेन स्वकीयशास्त्रीयग्रन्थस्य नाम `काव्यालङ्कारः' इति कृतम् । आनन्दवर्धनस्य ध्वन्यालोके पूर्ववर्णितार्थोऽपि रसादिपरिग्रहे नवत्वमायाती’ति प्रसङ्गे प्राक् उदाहृतं पद्यं भामहस्य द्वितीयञ्च वाणभट्टस्यास्ति । अतः वाणभट्टात् पूर्ववर्ती भामहो वर्तते । भामहः काव्यालङ्कारस्य पञ्चमपरिच्छेदे न्यायनिर्णयावसरे बौद्धदार्शनिकानां सिद्धान्तेन साकं स्वकीयं प्रगाढं परिचयं प्रकटयामास । अत्र भामहकृतं प्रत्यक्षप्रमाणलक्षणं दिङ्नागमतेन सह साम्यं भजति । दिङ्नागः प्रत्यक्षकल्पनापोदम्’ इति प्रत्यक्षलक्षणं करोति । अत्र कल्पनाशब्देन जातेर्नाम्नश्च कल्पनाऽभिधीयते । धर्मकीर्तिर् लक्षणेऽस्मिन् भ्रान्तपदं योजयित्वा भ्रान्तिरहितं कर्तुं यतते स्म । किन्तु भामहोऽस्य लक्षणेन सह परिचितो नास्ति । अतो भामहस्य समयो दिङ्नागादनन्तरं धर्मकीर्तेश्च प्राक् स्वीक्रियते । दिङ्नागस्य कालः पञ्चमशताब्द्यां तथा धर्मकीर्तेः समयः सप्तमशताब्द्याः पूर्वार्धः स्वीक्रियते । अनयोर्मध्ये षष्ठशताब्दी एव भामहस्य समयो निश्चीयते ।

\customparagraph{\textbf{काव्यशास्त्रीयं योगदानम्}}

अधुना भामहस्य काव्यालङ्काराभिधान एक एव शास्त्रीयो ग्रन्थ उपलभ्यते । भामहालङ्कार इति नाम्नापि प्रसिद्धो ग्रन्थोऽयम् अलङ्कारशास्त्रस्य सम्पूर्णतत्त्वविवेचनसम्बद्धा प्राचीनतमा उपलब्धिर्मन्यते । अत्र षट् परिच्छेदाः सन्ति । एषु ३९९ सङ्‌ख्यामिता कारिका विद्यन्ते । परन्तु भामहोल्लेखानुसारेण काव्यालङ्कारस्य श्लोकसङ्ख्याः चतुश्शतम् (४००) ज्ञायते ।                                                 (काव्यालङ्कारः, भूमिका, पृ. २) यथा—

\begin{shloka}[center]

\textbf{षष्ठ्या शरीरं निर्णीतं षट्षष्ठ्या त्वनलङ्कृतिः ।}

\textbf{पञ्चशता दोषदृष्टिः सप्तत्या न्यायनिर्णयः ।।}

\textbf{षष्ठ्या शब्दस्य शुद्धिः स्यादित्येवं वस्तुपञ्चकम् ।}

\textbf{उक्तः षड्भिः परिच्छेदैर् भामहेन क्रमेण वः ।। (भामहालङ्कारः)}

\end{shloka}

\customparagraph{\textbf{काव्यलक्षणम्}}

अस्य काव्यलक्षणं प्रारम्भिकं सङ्क्षिप्तञ्च विद्यते, तद्यथा— शब्दार्थौ सहितौ काव्यमिति । शब्दार्थौ च वक्राभिधेयौ । भामहस्य इदमेव काव्यलक्षणमुत्तरवर्त्तिनामाचार्याणां कृते उपजीव्यं जातम्। प्रायः सर्वेषां काव्यलक्षणे पृष्ठभूमिरूपेण इदमेव काव्यलक्षणं दृश्यते । अत्र विशेषणादिप्रयोगद्वारा स्पष्टं परिष्कृतं तु तैर्विहितमेव । अतो भामहस्य काव्यलक्षणमुत्तरवर्त्तिनां कृते बीजरूपं मन्यते ।

\customparagraph{\textbf{काव्यप्रयोजनम्}}

काव्यप्रयोजनविषयेऽपि भामहस्य मतमुत्तराचार्याणां कृते प्रेरकं दृश्यते । अस्य मतं  यथा स्यात्तथैव  स्वीकृत्य कतिपयैराचार्यैः स्वकीयमैक्यं प्रस्तुतम्, कतिपयैर्भामहस्य प्रयोजनमेव परिष्कृत्य प्रस्तुतम्  । एतेनापि अस्य काव्यप्रयोजनस्य महत्त्वं प्रदर्शितम् । भामहस्य काव्यप्रयोजनं यथा—

\begin{shloka}[center]

\textbf{धर्मार्थकाममोक्षेषु वैचक्षण्यं कलासु च ।}

\textbf{करोति कीर्तिं प्रीतिं च साधुकाव्यनिबन्धनम् ।। (भामहालङ्कारः)}

\end{shloka}

\customparagraph{\textbf{काव्यहेतुः}}

भामहः काव्यनिर्माणे कारणरूपेण विभिन्नविषयाणाम् उद्धरणं करोति । एतस्यैव मताधारे पश्चात् काव्यशास्त्रीया आचार्याः केवला प्रतिभा वा प्रतिभा, व्युत्पत्तिः, अभ्यासश्चेति हेतुत्रयसंवलितं काव्यकारणं स्वीकुर्वन्ति । अनेनापि प्राथमिकं कारणं तु प्रतिभा एव स्वीकृता । यथा—

\begin{shloka}[center]

\textbf{गुरूपदेशादध्येतुं शास्त्रं जडधियोऽप्यलम् ।}

\textbf{काव्यं तु जायते जातु कस्यचित् प्रतिभावतः ।। (भामहालङ्कारः)}

\end{shloka}

शास्त्रज्ञानार्थमपि कवित्वस्य (प्रतिभायाः) आवश्यकता भवतीति कथयति यत्—

\begin{shloka}[center]

\textbf{अधनस्येव दातृत्वं क्लीवस्येवास्त्रकौशलम् ।}

\textbf{अज्ञस्येव प्रगल्भत्वम् अकवेः शास्त्रवेदनम् ।।}

\end{shloka}

काव्यकरणशीलानां  पुरुषाणां कृते  शब्दः, छन्दः, कोशः, इतिहासप्रभृतीनां तत्त्वानामपि अध्ययनम् (नैपुण्यम्) आवश्यकं भवतीति कथयति यत्—

\begin{shloka}[center]

\textbf{शब्दश्छन्दोऽभिधानार्था इतिहासाश्रया कथा ।}

\textbf{लोकयुक्तिकलाश्चेति मन्तव्याः काव्ययैर्वशी ।।}

\end{shloka}

काव्यस्य कृते काव्यशास्त्रज्ञविदुषामुपासनं करणीयम्, अर्थात् तेभ्यः काव्यशास्त्रीयज्ञानं ग्रहणीयं भवतीति कथयति यत्—

\begin{shloka}[center]

\textbf{शब्दाभिधेये विज्ञाय कृत्वा तद्विदुपासनाम् ।}

\textbf{विलोक्यान्यनिबन्धाश्च कार्यकाव्यक्रियादरः ।।}

\end{shloka}

अकरणात् मन्दकरणं श्रेयः’ इति लोकन्यायम् अत्र नाङ्गीकरणीयम्, निन्दनीयकाव्यरचनया तु मृत्युवत् हीनत्वं लभ्यते’ इति–अवबोधनार्थं कथयति यत्—

\begin{shloka}[center]

\textbf{सर्वथा पदमप्येकं न निगद्यमवद्यवत् ।}

\textbf{विलक्ष्मणा हि काव्येन दुस्सुतेनेव निन्द्यते ।।}

\textbf{नाकवित्वमधर्माय व्याधये दण्डनाय वा ।}

\textbf{कुकवित्वं पुनः साक्षान्मृतिमाहुर्मनीषिणः ।।}

\end{shloka}

\customparagraph{\textbf{काव्यभेदाः}}

काव्यभेदविषये भामहो विभिन्नाधारे स्वमतं प्रस्तौति । तेषु  रचनाशैल्याधारे—  गद्यं पद्यं च द्विविधम्, भाषाधारे— संस्कृतम्, प्राकृतम्, अपभ्रंशश्चेति त्रिविधम्, कथावस्तुगताधारे— देवादिचरितयुतम्, उत्पाद्यवस्तुवर्णितम्, कलाश्रयम्, शास्त्राश्रयञ्चेति चतुर्विधम्, तथा च विधागताधारे— सर्गबन्धम्, अभिनेयार्थम्, आख्यायिका, कथा, अनिबद्धञ्चेति पञ्चविधम् । एतेषां सर्वेषां भेदानां सस्वरूपमुल्लेखनं  कृतमाचार्येण । तद्यथा काव्यालङ्कारे—

\begin{shloka}[center]

\textbf{शब्दार्थौ सहितौ काव्यं गद्यं पद्यं च तद् द्विधा ।}

\textbf{संस्कृतं प्राकृतं चान्यदपभ्रंश इति त्रिधा ।।}

\textbf{वृत्तदेवादिचरितशंसि चोत्पाद्यवस्तु च ।}

\textbf{कलाशास्त्राश्रयञ्चेति चतुर्धा भिद्यते पुनः ।।}

\textbf{सर्गबन्धोऽभिनेयार्थं तथैवाख्यायिकाकथे ।}

\textbf{अनिबद्धञ्च काव्यादि तत्पुनः पञ्चधोच्यते ।। (भामहालङ्कारः)}

\end{shloka}

अनेन प्रकारेण भामहः काव्यभेदविषये सविस्तारं चर्चां विदधाति । पश्चाज्जातैराचायैरेतस्यैव काव्यभेदविषयको विचारः परिष्कारसहितः स्वीकृतः ।

\customparagraph{\textbf{महाकाव्यस्वरूपम्}}

भामहस्य महाकाव्यलक्षणं सर्वप्राचीनमस्ति । महाकाव्यस्य स्वरूपविषये शून्यावस्थायामनेन यन्  महाकाव्यलक्षणं निर्मितं तत् संस्कृतकाव्यशास्त्रस्यैव कृते विशिष्टा प्राप्तिर्मन्यते । अनेन महाकाव्यस्वरूपं तु सङ्क्षिप्तमेव प्रस्तुतम्, तथापि तदुत्तरवर्त्तिनामाचार्याणां कृते पथप्रदर्शकं जातम् । तद्यथा भामहालङ्कारे—

\begin{shloka}[center]

\textbf{सर्गबन्धो महाकाव्यं महताञ्च महच्च यत् ।}

\end{shloka}

\begin{shloka}[center]

\textbf{अग्राम्यशब्दमर्थ्यञ्च सालङ्कारं सदाश्रयम् ।।}

\textbf{मन्त्रदूतप्रयाणाजिनायकाभ्युदयैश्च यत् ।}

\textbf{पञ्चभिर्सन्धिभिर्युक्तं नातिव्यख्येयमृद्धिमत् ।।}

\textbf{चतुर्वर्गाभिधानेऽपि भूयसार्थोपदेशकृत् ।}

\textbf{युक्तं लोकस्वभावेन रसैश्च सकलैः पृथक् ।।}

\textbf{नायकं प्रागुपन्यस्य वंशवीर्यश्रुतादिभिः ।}

\textbf{न तस्यैव बधं ब्रूयादन्योत्कर्षाभिधित्सया ।।}

\textbf{यदि काव्यशरीरस्य न स व्यापितयेष्यते ।}

\textbf{न चाभ्युदयभाक्तस्य मुधादौ ग्रहणं स्तवे ।। (भामहालङ्कारः)}

\end{shloka}

सङ्क्षेपेऽपि लक्षणेऽस्मिन् महाकाव्यनिर्माणस्य कृते आवश्यकाः प्रायशः सर्वे विषयाः समाविष्टाः सन्ति । एतस्यैवाधारे दण्डीविश्वनाथादयः परिष्कारप्रवर्धनपूर्वकं महाकाव्यलक्षणं प्रस्तुवन्ति । एतेनापि कारणेन अस्य महाकाव्यलक्षणस्य महत्त्वम् अन्यूनमिति ज्ञातुं शक्यते ।

\customparagraph{\textbf{गुणविवेचनम्}}

गुणस्य विषये भामहः प्राचीनोऽपि नवीनमतप्रवर्तको मन्यते । अयमाचार्यो भरतस्थापितं दशगुणवादं परिष्कृत्य त्रिगुणवादं स्थापयति । भामहो भरतोल्लिखितानां दशगुणानां खण्डनमण्डनादिकार्यं विनैव सामान्यतया त्रिगुणमतं प्रस्तौति, तथाप्यस्य त्रिगुणात्मिका धारणा प्रायशः सर्वैरुत्तरवर्त्तिरभिराचार्यैरनुमोदिता । अस्य गुणविचारो यथा—

\begin{shloka}[center]

\textbf{माधुर्यमभिवाञ्छन्तः प्रसादं च सुमेधसः।}

\textbf{समासवन्ति भूयांसि न पदानि प्रयुञ्जते ।।}

\textbf{केचिदोजोऽभिदित्सन्तः समस्यन्ति बहून्यपि ।}

\textbf{यथा मन्दारकुसुमरेणुपिञ्जरितालका ।।}

\textbf{श्रव्यं नातिसमस्तार्थं काव्यं मधुरमिष्यते ।}

\textbf{आविद्वदङ्गनाबालप्रतीतार्थं प्रसादवत् ।। (भामहालङ्कारः)}

\end{shloka}

\customparagraph{\textbf{औचित्यतत्त्वसङ्केतः}}

सम्प्रदायपरम्परायाश्चरमसम्प्रदायरूपेण विकसितस्य औचित्यसम्प्रदायस्य कृते भामहः प्रेरकं विचारं प्रदर्शयति । स औचित्यपूर्वकं स्थापयेच् चेद् दोषोऽपि गुणत्वं भजतीति उल्लिखति स्वग्रन्थे काव्यालङ्कारे—

\begin{shloka}[center]

\textbf{सन्निवेशविशेषात्तु दुरुक्तमपि शोभते ।}

\textbf{नीलं पलाशमाबद्धमन्तराले स्रजमिव ।।}

\textbf{किञ्चिदाश्रयसौन्दर्यात् धत्ते शोभामसाध्वपि ।}

\textbf{कान्ताविलोचनन्यस्तं मलिमसमिवाञ्जनम् ।। (भामहालङ्कारः)}

\end{shloka}

पश्चात् औचित्यसम्प्रदायस्थापने क्षेमेन्द्रस्य कृतेऽस्य विचारः सम्प्रदायरूपगृहनिर्माणस्य शिलान्यासवत् सञ्जातः’ इति मन्यते ।

\customparagraph{\textbf{अलङ्कारसम्प्रदायसंस्थापकः}}

काव्यस्य सौन्दर्यतत्त्वरूपेण व्यवस्थितस्यालङ्कारसम्प्रदायस्य प्रवर्तको भामहोऽस्ति । यद्यपि भामहात् प्रागपि अलङ्कारस्य विषये किञ्चित् सङ्केतः प्राप्यते, तथापि सैद्धान्तिकरूपेण केनापि विवेचनं न कृतम् । अनेन तु समहत्त्वम् अलङ्कारस्वरूपं  प्रदर्शितम् । वक्रोक्तिरेव अलङ्कारस्यात्मा भवति । वक्रोक्तिमयेनालङ्कारेण विना काव्यं कमनीयं न भवतीति कथयन् भामहोऽलङ्कारतत्त्वं काव्यस्य कृते महत्त्वपूर्णं स्वीकरोति, यथा—

\begin{shloka}[center]

\textbf{सैषा सर्वत्र वक्रोक्तिरनयार्थो विभाव्यते ।}

\textbf{यत्नोऽस्यां कविना कार्यः कोऽलङ्कारोऽनया विना ।।}

\textbf{न कान्तमपि निर्भूषं विभाति वनिताननम् । (भामहालङ्कारः)}

\end{shloka}

अनया रीत्या अलङ्कारविवेचनं कुर्वन् भामहः ३९ सङ्ख्यकान् अलङ्कारान् विविच्य अलङ्कारसम्प्रदायस्य मेरुदण्डवत् प्रसंशितोऽस्ति ।

आचार्यभामहः पौरस्त्यकाव्यशास्त्रस्य इतिहासे अद्वितीयः काव्यशास्त्रीयबीजप्रकाशको वर्तते । नाट्यशास्त्रस्य दृश्यकाव्याश्रितात् परिधेः काव्यशास्त्रं बहिर्निष्कास्य, तस्मै स्वतन्त्रविद्यारूपेण प्रतिष्ठाप्रदानस्य श्रेयो भामहायैव गच्छति । तेन प्रणीतः `काव्यालङ्कारः' उत्तरवर्त्तिनां सर्वेषामपि आचार्याणां कृते सिद्धान्तबीजप्रदाता उपजीव्यग्रन्थश्च सञ्जातः । तस्य `शब्दार्थौ सहितौ काव्यम्' इति कालजयि काव्यलक्षणम्, भरतमुनेर्दशगुणानां स्थाने त्रिगुणवादस्य नवीनमार्गस्य स्थापना, औचित्यतत्त्वस्य स्पष्टसङ्केतबीजम्, तथा च `सैषा सर्वत्र वक्रोक्तिः' इति उद्घोषेण अलङ्कारस्य व्यापकबीजतत्त्वप्रतिपादनं परवर्तिनां दण्डि-वामन-आनन्दवर्धन-मम्मटाद्याचार्याणां मार्गदर्शकं सिद्धम् अभवत् । `न कान्तमपि निर्भूषं विभाति वनिताननम्' इति ब्रुवाणोऽयं मनीषिप्रवरो विश्ववाङ्मये काव्यसौन्दर्यमीमांसाया आदिप्रवर्तकत्वेन अद्यापि सहृदयैः सादरं स्मर्यते, तस्य च एतद् अनुपमम् अवदानं साहित्यशास्त्रस्य इतिहासे अहरहः शाश्वतम् अक्षुण्णञ्च रूपमाश्रित्य अवस्थितमस्ति ।

\specialupakhanda{\textbf{आचार्यो दण्डी}}

पौरस्त्यालङ्कारशास्त्रस्य प्रवर्तको भामहो वर्तते । तस्य प्रवर्धकत्वेन दण्डी सञ्जातः । अस्य समयविषये विदुषां मतैक्यं नास्ति । पूर्वापरप्रसङ्गमवलोक्यैवास्य समयविषयेऽनुमातुं शक्यते । राष्ट्रकुटनरेशस्य नपत्तुङ्गस्य रचना कन्नडभाषामयी राजमार्गाभिधाना वर्तते । ग्रन्थेऽस्मिन् अनेके श्लोकाः काव्यादर्शे अक्षरशः समुपलभ्यन्ते । नपत्तुङ्गस्य कालो नवमशताब्दस्य पूर्वार्धो वर्तते । वामनप्रणीतस्य काव्यालङ्कारस्यानुशीलनेन ज्ञायते यद् — वामनो दण्डिनः पूर्णतया परिचितोऽस्ति । दण्डी केवलं द्वयोरेव रीत्योः वर्णनं करोति । किन्तु वामनः तयोरन्तरालवर्त्तिन्याः पाञ्चाल्या अपि निर्देशं करोति । अनेन निस्सरति यद् — दण्डिनः समयोऽष्टमशताब्दादनन्तरं नहि भवितुमर्हति । दण्डिनः ‘अरत्नालोकसंहार्यमवार्यं सूर्यरश्मिभिः—’ इत्यादिरचनायां कादम्बर्याः शुकनाशोपदेशस्य यौवनदोषाणां वर्णनस्य छाया प्राप्यते । तथा च रत्नभित्तिषु सङ्‌क्रान्तै—’ इत्यादिसु माघस्य रचनाप्रभावोऽपि दृश्यते । अस्य रचनायां भर्तृहरेर्वाक्यपदीयस्योऽपि प्रभावः प्राप्यते । तेन काव्यादर्शस्य द्वितीये परिच्छेदे कर्मणो निवर्त्यविकार्यप्राप्यादिभेदास्त्रयो  वाक्यपदीयात् गृहीताः सन्तीति समालोचकः डा.के.पि.पाठकमहोदयः समुल्लिखति । एतैः प्रमाणबाहुल्यैः दण्डिनः समयः सप्तमशताब्दस्य पूर्वार्धः’ इति अनुमीयते । अस्य मातुर्  नाम `गौरी', पितुश्च `वीरथः' आसीत् । जन्मस्थानञ्चास्य ‘काञ्चिनगरम्’ आसीदिति अस्यैव अवन्तीसुन्दरीकथायां प्राप्यते ।

\customparagraph{\textbf{काव्यशास्त्रीयं योगदानम्}}

दण्डिनः कृतयः काव्यशास्त्रेतरेषु विषयेष्वपि प्राप्यन्ते ।  दशकुमारचरितम्, अवन्तीसुन्दरीकथा चेति द्वौ प्रसिद्धौ कथाग्रन्थौ स्तः । दशकुमारचरितस्य गद्यवैशिष्ट्यं विलोक्यैव समालोचकाः ‘दण्डिनः पदलालित्यम्’ इति प्रशंसन्ति । अस्य काव्यशास्त्रीयं योगदानमपि वैशिष्ट्यपूर्णमस्ति ।

दण्डिनः काव्यादर्शाख्यः शास्त्रीयग्रन्थ उपलभ्यते । तत्र त्रयः परिच्छेदाः सन्ति । तेषु षष्ट्युत्तराणि  षट्शतानि  पद्यानि  विद्यन्ते । प्रथमे परिच्छेदे– काव्यपरिभाषायाः, काव्यभेदस्य, महाकाव्यस्य, गद्यकाव्यस्य, कथाख्यायिकयोः, रीतिगुणयोः काव्यकारणस्य च विवेचनं प्राप्यते । द्वितीये परिच्छेदे– अलङ्कारलक्षणम्, पञ्चत्रिंशत्सङ्ख्यकानामलङ्काराणां सोदाहरणं निरुपणञ्च विद्यते । तृतीये परिच्छेदे– यमकस्य पञ्चदशोपरित्रिंशद्भेदयुतं विस्तृतं वर्णनं लभ्यते । तथा च तत्र चित्रबन्धकाव्यस्य गोमुत्रिकासर्वतोभद्रवर्णनियमादिकस्य, षोडशप्रकाराणां प्रहेलिकानाम्, दशप्रकाराणां दोषाणाञ्च वर्णनं समुपलभ्यते । आचार्यो दण्डी रसध्वन्यादिविषये स्पष्टतया किमपि नोल्लिखति, अतोऽयमलङ्कारसम्प्रदायस्यानुवर्त्तीति ज्ञायते । अस्य कवित्वशक्तिरपि विलक्षणा आसीत् । स्वकीये लक्षणग्रन्थेऽपि उदाहरणरूपेण स्वरचिताः श्लोका एव प्रयुक्ताः सन्ति ।

विश्लेषणानन्तरं कथयितुं शक्यते यद्— आचार्यो दण्डी संस्कृतवाङ्मयस्य बहुमुखी प्रतिभा आसीत्, विलक्षणप्रतिभासम्पन्नो दण्डी काव्यपक्षं  शास्त्रीयपक्षञ्च प्रवर्धयितुं कालेऽस्मिन् महद् योगदानमकरोत् ।

\customparagraph{\textbf{काव्यलक्षणम्}}

काव्यलक्षणविषये दण्डिनाऽपि पूर्वेषाम् आचार्याणां मतमेव सविशेषणम् उल्लिखितम् । इष्टार्थेन अभिन्ना पदावली (शब्दः) एव काव्यशरीरं भवतीति काव्यादर्शे प्रस्तौति यत्—

\begin{shloka}[center]

\textbf{शरीरं तावदीष्टार्थः व्यवच्छिन्ना पदावली । (काव्यादर्शः)}

\end{shloka}

\customparagraph{\textbf{काव्यभेदाः}}

आचार्यो दण्डी काव्यभेदविषये पृथङ् मतमुपस्थापयति । स तु काव्यस्य वर्णनशैलीगताधारे छन्दरहितम्, छन्दसहितम्, अनयोर्मिश्रणञ्चेति त्रिविधान् काव्यभेदान् स्वीकरोति, तद्यथा काव्यादर्शे—

\begin{shloka}[center]

\textbf{गद्यं पद्यं च मिश्रं तत् त्रिविधैव व्यवस्थितम् ।}

\end{shloka}

\customparagraph{\textbf{अलङ्कारविचारः}}

काव्यादर्शकारोऽलङ्कारविषये ध्वनिवादिनो  मतमनुसरति, तथापि काव्यस्य सुदीर्घकालपर्यन्तं स्थायिरूपेण प्रशंसनीयं त्वलङ्कृतं काव्यमेव भवति, अतः काव्येऽलङ्कारस्याङ्गीकरणम् अवश्यमेव कर्तव्यमिति कथयति काव्यादर्शे—

\begin{shloka}[center]

\textbf{काव्यशोभाकरान् धर्मानलङ्कारान् प्रचक्षते ।}

\textbf{ते चाद्यापि विकल्प्यन्ते कस्तान् कार्त्स्येन वक्ष्यति ।।}

\textbf{काव्यं कल्पान्तरस्थायी जायते सदलङ्कृतिः ।  (काव्यादर्शः)}

\end{shloka}

\customparagraph{\textbf{रीतिस्थापनम्}}

आचार्यो दण्डी अन्येषां पूर्वाचार्याणां रीतिभेदानां विषये अस्पष्टं सूक्ष्मान्तरञ्च कथयित्वा स्पष्टभेदयुतौ द्वावेव रीतिभेदौ प्रस्तौति, काव्यादर्शे यथा—

\begin{shloka}[center]

\textbf{अस्त्यनेको गिरां मार्गः सूक्ष्मभेदः परस्परम् ।}

\textbf{तत्र वैदर्भिगौडियौ वर्ण्येते प्रस्फुटान्तरौ ।।}

\end{shloka}

\customparagraph{\textbf{गुणसमीक्षा}}

गुणविषये अयमाचार्यो भरतस्यानुकरणं करोति । नामसङ्ख्ये च भरतस्य मतं स्वीकृत्य वैदर्भमार्गस्य प्राणरूपेण गुणविषये मतमिदं प्रस्तौति ।

\begin{shloka}[center]

\textbf{श्लेषः प्रसादः समता माधुर्यः सुकुमारता ।}

\textbf{अर्थव्यक्तिरुदारत्वमोजः कान्तिसमाधयः ।।}

\textbf{इति वैदर्भमार्गस्य प्राणा दशगुणाः स्मृताः ।  (काव्यादर्शः)}

\end{shloka}

\customparagraph{\textbf{महाकाव्यस्वरूपम्}}

महाकाव्यविषये दण्डिनो विचारो व्यापको विद्यते । भामहस्य महाकाव्यलक्षणस्य परिष्कृतं विस्ताररूपं चास्य महाकाव्यलक्षणमस्ति । प्रबन्धकाव्यस्य कृते आवश्यकाः प्रायशः सर्वे विषयाः प्रदर्शिताः सन्ति । अत्र काव्यादर्शे वर्णितं महाकाव्यस्वरूपम् उल्लिख्यते यत्—

\begin{shloka}[center]

\textbf{सर्गबन्धो महाकाव्यमुच्यते तस्य लक्षणम् ।}

\textbf{आशीर्नमस्क्रियावस्तुनिर्देशो वापि तन्मुखम् ।।}

\textbf{इतिहासकथोद्भूतमितरद्वा सदाश्रयम् ।}

\textbf{चतुर्वर्गफलायत्तं चतुरोदात्तनायकम् ।।}

\textbf{नगरार्णवशैलर्तु चन्द्रार्कोदयवर्णनैः ।}

\textbf{उद्यानसलिलक्रीडा मधुपानरतोत्सवैः ।।}

\textbf{विप्रलम्भैर्विवाहैश्च कुमारोदयवर्णनैः ।}

\textbf{मन्त्रदूतप्रयाणाजिनायकाभ्युदयैरपि ।।}

\textbf{अलङ्कृतमसङ्क्षिप्तं रसभावनिरन्तरम् ।}

\textbf{सर्गैरनतिविस्तीर्णैः श्रव्यवृतैः सुसन्धिभिः ।।}

\textbf{सर्वत्र भिन्नवृत्तान्तैरुपेतं लोकरञ्जकम् ।}

\textbf{काव्यं कल्पान्तरस्थायि जायते सदलङ्कृती ।।  (काव्यादर्शः)}

\end{shloka}

पश्चात् विश्वनाथादयो महाकाव्यलक्षणकर्तारो विद्वांस इदमेव महाकाव्यलक्षणं परिष्कारपूर्वकं प्रस्तुवन्ति ।

आचार्यो दण्डी केवलं काव्यशास्त्री एव न, अपितु सुकविरप्यासीत् । तस्य दशकुमाराख्यं गद्यकाव्यं पदलालित्यविषये सहृदयैरभिनन्दितम् । अवन्तिसुन्दरीकथाऽपि मनोज्ञं गद्यकाव्यमास्ते । स्वकीये लाक्षणिकग्रन्थे काव्यादर्शेऽपि दण्डी उदाहरणत्वेन प्रायशः श्लोकपद्यान् स्वरचितान्नेवोल्लिखति । दण्डिनः सुकवित्वमवलोक्यैव समालोचकैरभिहितं यत्–

\begin{shloka}[center]

\textbf{जाते जगति वाल्मीकौ कविरित्यभिधाऽभवत् ।}

\textbf{कवी इति ततो व्यासे कवयस्त्वयि दण्डिनि ।।}

\end{shloka}

आचार्यो दण्डी पौरस्त्य-साहित्यशास्त्रस्य इतिहासपरम्परायां कश्चन अद्वितीयो युगप्रवर्तको मनीषी वर्तते । स केवलं सिद्धान्तप्रतिपादको लाक्षणशास्त्री एव नासीत्, अपितु `पदलालित्य'गुणेन समलङ्कृतः सुकविरपि आसीत् । काव्यादर्शे प्रयुक्तास्तस्यैव स्वरचिताः श्लोका अस्य तथ्यस्य निदर्शनाः सन्ति । दशकुमारचरितादिषु विलसितं गद्यसौन्दर्यम् अद्यापि सहृदयानां चेतांसि आल्हादयन्ति । भामहप्रवर्तिताम् अलङ्कारपरम्परां सुदृढां कुर्वता दण्डिना काव्यादर्शे यद् दशगुणानां प्राणात्मत्वस्थापनम्, महाकाव्यस्य च यद् व्यापकं वैज्ञानिकञ्च स्वरूपं निरूपितम्, तत् परवर्तिनां मम्मट-विश्वनाथाद्याचार्याणां कृते सर्वथैव आधारशिला सञ्जातम् । अत एव वाल्मीकि-वेदव्यासाभ्यां सह तस्य तुलनां कुर्वती `कवयस्त्वयि दण्डिनि' इति प्राक्तनी उक्तिस्तस्य सार्वभौमं महत्त्वम्, बहुमुखीं प्रतिभाम्, काव्यशास्त्रीयपरम्परायां तस्य अविस्मरणीयं योगदानञ्च पूर्णतः प्रमाणीकरोति ।

\specialupakhanda{\textbf{आचार्यो भट्टोद्भटः}}

काव्यशास्त्रीयविभाजनस्य रचनात्मककालस्य भट्टोद्भटोऽपि विशिष्टा प्रतिभाऽस्ति । अस्य समयविषये राजतरङ्गिणीकारः कल्हणो महाराज्ञो जयापीडस्य सभापण्डित आसीदुद्भट इति उल्लिखति राजतरङ्गिण्याम्—

\begin{shloka}[center]

\textbf{विद्वान् दीनारलक्षेण प्रत्यहं कृतवेतन ।}

\textbf{भट्टोऽभूदुद्भटस्तस्य भूमिभर्तुः सभापतिः ।।}

\end{shloka}

जयापीडस्य राज्यसमयः ८३६ तमात् वैक्रमाब्दात् ८७० तमं वैक्रमाब्दं यावत् स्वीक्रियते । आचार्य आनन्दवर्धनोऽपि भट्टोद्भटस्य नाम असकृत् संस्मार । परम् आचार्योद्भट आनन्दवर्धनस्य वा ध्वनिसिद्धान्तस्य विषये कुत्रापि चर्चा नाकरोत् । अतः भट्टोद्भटस्य समयः अष्टमशतकस्य पूर्वार्ध एवेति निश्चीयते ।

\customparagraph{\textbf{काव्यशास्त्रीयं योगदानम्}}

आनन्दवर्धनाप्पयदीक्षितरुय्यकजगन्नाथादयोऽस्य योगदानस्य चर्चां सम्मानेन कुर्वन्ति । अधुनाऽस्य त्रयो ग्रन्थाः स्वीक्रियन्ते-  भामहविवरणम्, कुमारसम्भवकाव्यम्, अलङ्कारसारसङ्ग्रहश्च । यद्यपि अस्यालङ्कारसारसङ्ग्रहाभिधो ग्रन्थ एव प्राप्यते । अनेनैव अस्य वैदुष्यं सूचितं भवति । शास्त्रीयग्रन्थेऽस्मिन् केवलमलङ्कारतत्त्वमेव विवेचितं विद्यते । तत्र षड्‍वर्गाः सन्ति । तेषु एकोनाशीतिसङ्ख्यकाभिः कारिकाभिरेकचत्वारिंशदलङ्काराणां लक्षणं कृत्वा तेषामुदाहरणार्थं स्वकीयकुमारसम्भवकाव्यादेव श्लोकानुद्धृताञ्चकार  भट्टोद्भटः । वर्गानुक्रमेण अस्य लक्षणग्रन्थेऽलङ्कारस्योपस्थापनक्रमोऽत्र प्रदर्श्यते, यथा- **प्रथमे वर्गे—**पुनरुक्तवदाभासस्य, त्रिविधानुप्रासस्य, परुषोपनागरिकाग्रम्याकोमलाख्यानां चतसृणां वृत्तीनाम्, अलङ्कारेषु उपमायाः, रूपकस्य, दीपकस्य, प्रतिवस्तूपमायाश्च विवेचनमास्ते । \textbf{द्वितीये वर्गे—} आक्षेपार्थान्तरन्यास- व्यतिरेकविभावनातिशयोक्ति-समासोक्त्यादीनामलङ्काराणां सोदाहरणं वर्णनं समुपलभ्यते । **तृतीये वर्गे—**यथासङ्ख्योत्प्रेक्षास्वभावोक्तिनामलङ्काराणां विश्लेषणं प्राप्यते । \textbf{चतुर्थे वर्गे—} प्रेयरसवदुर्जस्विपर्यायोक्तिसमाहितोदात्तश्लिष्टादयोऽलङ्काराः वर्णिताः । **पञ्चमे वर्गे—**अपह्नुति- विशेषोक्तिविरोधतुल्ययोगिता-प्रस्तुतप्रशंसाव्याजस्तुतिनिदर्शनासहोक्त्युपमेयोपमापरिवृत्तिसङ्करादीनामलङ्काराणां वर्णनं कृतमाचार्येण । **षष्ठे वर्गे—**अनन्वयससन्देहसंसृष्टिभाविक-काव्यलिङ्ग-दृष्टान्ता- नामलङ्काराणां विवेचनं समुपलभ्यते ।

अनेन प्रकारेण भट्टोद्भटोऽलङ्काराणां सविस्तारं विवेचनमेव न, अपितु काव्यशास्त्रीयविषये स्थाने स्थाने स्वकीयं मौलिकविचारमपि प्रस्तौति । अनेनाभिधाव्यापारगुणार्थविषये तथा च उपमाया भेदविषये स्वकीयं मौलिकं मतं प्रस्तुतम् । मतमिदमुत्तरवर्त्तिनामाचार्याणां कृतेऽपि मार्गनिर्देशकः सञ्जातः । अतो भट्टोद्भटोऽलङ्काराणां विस्तृतविवेचनकर्ता, मौलिकविचारप्रस्तावकः, विदुषां वरेण्यः सर्वैः समालोचकैः समादृतश्चासीदित्यत्र न कस्यापि विप्रतिपत्तिः ।

\specialupakhanda{\textbf{आचार्यो वामनः}}

पौरस्त्यकाव्यशास्त्रे वामनस्य स्थानं वैशिष्ट्यपूर्णमस्ति । अयं भट्टोद्भटस्य प्रतिद्वन्द्वी आसीत् । राज्ञो जयापीडस्य राजसभायां इमौ द्वावेव सभापण्डितौ आस्ताम्, तथापि अन्योन्यविषये द्वावेव मौनिनौ स्तः । परमुद्भटः काव्यस्यालङ्कारपक्षे एव स्वकीयं वैदुष्यं प्रदर्शयति । वामनस्तु सर्वेषु शास्त्रीयविषयेषु गभीरम्, मौलिकञ्च विचारं प्रस्तौति । अतः काव्यशास्त्रस्य सर्वपक्षीयविवेचनस्य प्रारम्भकर्तुर्भामहस्य परवर्त्तिषु लक्षणशास्त्रिषु वामनस्य स्थानं स्मरणीयमस्ति । राजतरङ्गिण्याः प्रणेता कल्हणोऽपि वामनाचार्यं  जयापीडस्य मन्त्रिष्वन्यतमं  स्वीकरोति । जयापीडस्य समयः ७७९ ईसवीतः ८१३ इसवीं यावत् मन्यते । तथा च राजशेखरेण वामनस्य रीतिसम्प्रदायस्योल्लेखः कृतः । आनन्दवर्धनेनापि ध्वन्यालोकस्य तृतीये उद्योते सङ्घटनाप्रसङ्गे वामनस्य मतमुल्लिखितम् । ध्वन्यालोकलोचनकारस्याभिनवगुप्तस्य विवेचनेनापि राजशेखरानन्दवर्धनाभ्यां प्रागेव वामनस्य समयो मन्तुं युज्यते ।  काव्यालङ्कारसूत्रवृत्तौ उत्तररामचरितस्य ‘इयं गेहे लक्ष्मी—’ इत्यादिश्लोकस्योल्लेखनेनापि भवभूतेः समयः (७००—७५०) पश्चादेव अस्य समयो मन्यते। एतैः प्रमाणैः सप्तमशताब्द्या उत्तरार्धात् अष्टमशताब्द्याः पूर्वार्धं यावत् वामनाचार्यस्य समयोऽनुमीयते । अस्य स्थानविषये तु उपरिवर्णितसमयप्रमाणेनैव कश्मिर एवेति स्थिरीभवति ।

\customparagraph{\textbf{काव्यशास्त्रीयं योगदानम्}}

आचार्यो वामनो रीतिसम्प्रदायस्य प्रवर्तकोऽस्ति । सूत्रशैल्यां रचिता ‘काव्यालङ्कारसूत्रवृत्तिः’ संस्कृतवाङ्मये शास्त्रीयपक्षस्य महत्त्वाधायको ग्रन्थो मन्यते । ग्रन्थेऽस्मिन् पञ्चाधिकरणानि सन्ति । तत्र च द्वादशसङ्ख्यका अध्यायाः, एकविंशत्युत्तरत्रिशतानि सूत्राणि सन्ति । तेषु प्रथमाधिकरणे— काव्यप्रयोजनस्याधिकारिणश्च वर्णनमास्ते । अत्रैव ग्रन्थकारो रीतेः काव्यात्मत्वं प्रतिपाद्य रीतेस्त्रयाणां भेदानां, काव्यस्यानेकेषां प्रकरणानाञ्च वर्णनं करोति । द्वितीयेऽधिकरणे—पददोषाणां वाक्यदोषाणां वाक्यार्थदोषाणाञ्च विवेचनं प्राप्यते । तृतीयं तु गुणविवेचनाख्यमधिकरणम् । तत्र गुणालङ्कारयोः पार्थक्यं विविच्य दशानां शब्दगुणानां दशानामर्थगुणानाञ्च विश्लेषणं वर्तते । चतुर्थेऽलङ्कारिकाभिधेऽधिकरणेऽलङ्काराणां विस्तारेण वर्णनं प्राप्यते । पञ्चमे प्रायोगिकाधिकरणे सन्दिग्धानां शब्दानाम्, शब्दशुद्धेश्च समीक्षा कृता काव्यशास्त्रिणा ।

रीतेर्काव्यात्मत्वप्रतिपादनमस्य महनीयं वैशिष्ट्यं विद्यते । रीतेः सिद्धान्तस्तु प्राचीनतमो विद्यते। परं तस्य काव्यात्मत्वनिर्धारणं तु वामनेन एव कृतम् । अयं प्राचीनैः स्थापितयो रीतिर्भेदयोर्मध्ये तृतीयायाः पाञ्चाल्याश्चापि प्रकारत्वेन स्वरूपं प्रतिपादयामास । प्राक् गुणालङ्कारौ काव्यस्य बाह्यशोभाधायकौ मन्यते स्म । अनेन तु गुणालङ्कारयोर्भेदं कृत्वा ‘विशेषो गुणात्मा’ कथयित्वा गुणतत्त्वं काव्यात्मस्थानीभूतत्वेन स्वीकृतम् । वामनात् प्राक् दश एव गुणसङ्ख्या आसीत् । तत्र वामनेन शब्दार्थरूपेण पृथक्कृत्वा विंशतिर्गुणाः  प्रतिपादिताः । अस्य मते उपमा एव मुख्योऽलङ्कारः। अन्ये तु उपमाया एव प्रपञ्चभूताः सन्ति । अन्यैः आचार्यैः रसतत्त्वं काव्यस्य बाह्यसाधनरूपं स्वीकृतम् । परं वामनस्तु रससत्तायाः स्वीकर्ता ज्ञायते ।

एतादृशैर्मौलिकविचारप्रवाहैर्वामनस्योपस्थितिः पौरस्त्यवाङ्मयस्य कृते महत्तमा उपकृतिर्दृश्यते। पौरस्त्यकाव्यशास्त्रे प्रसिद्धस्य रीतिसम्प्रदायस्य प्रवर्तकोऽयं काव्यशास्त्रीयेतिहास- रसिकानां कृतेऽविस्मरणीयो विद्यते ।

\customparagraph{\textbf{काव्यलक्षणम्}}

वामनो रीतिवादी आचार्योऽस्ति । अस्य मते काव्यम् अलङ्कारेणैव ग्रहणीयं भवति । काव्यशरीरं तु गुणालङ्कारसमन्वितं शब्दार्थमयं भवति । एतादृशस्य काव्यस्य आत्मा तु गुणसमन्विता विशिष्टपदरचनास्वरूपा रीतिरेव । अनेन प्रकारेणायं काव्यस्वरूपे रीतितत्त्वमेव प्राधान्येन स्वीकरोति, तद्यथा काव्यालङ्कारसूत्रवृत्तौ—

\textbf{काव्यं ग्राह्यमलङ्कारात्, काव्यशब्दोऽयं गुणालङ्कारयोः शब्दार्थयोर्वर्तते । रीतिरात्मा काव्यस्य, विशिष्टपदरचना रीतिः, विशेषो गुणात्मा’} इति । \textbf{(काव्यालङ्कारसूत्रवृत्तिः)}

\customparagraph{\textbf{काव्यहेतुः}}

आचार्यो वामनः काव्यकारणविषये विस्तारपूर्विकां चर्चां करोति । स्वकीये काव्यालङ्कारसूत्रे ‘काव्याङ्ग’नाम्ना काव्यहेतुरूपान् विषयान् प्रस्तौति । तत्र प्राग् लोकः, विद्या, प्रकीर्णञ्च त्रिविधमुल्लिखति । एतेषां व्याख्यायां स्पष्टरूपेण काव्यकारणमुपस्थापयति, तद्यथा—

लोकः, विद्या, प्रकीर्णञ्च काव्याङ्गानि । लोकवृत्तज्ञानम् एव ‘लोकः’ इत्यस्याशयः । शब्दस्मृत्यभिधानकोशाच्छन्दोविच्छित्तिकलाकामशास्त्रदण्डनीतिपूर्वा विद्या, अर्थाद् व्याकरणशास्त्रम्, शब्दकोशः, छन्दःशास्त्रम्, कलाशास्त्रम्, कामशास्त्रम्, दण्डनीतिः, इत्यादिषु ज्ञानार्जनम्, पूर्वा इत्यनेन गणितविद्यादिपरिग्रहः । लक्ष्यज्ञत्वमभियोगो वृद्धसेवाऽवेक्षणं प्रतिभानमवधानञ्च प्रकीर्णम्, अर्थात् काव्यपरिचयः (लक्ष्यज्ञत्वम्) काव्यबन्धोद्यमः (अभियोगः) काव्योपदेशगुरुशुश्रूषणम् (वृद्धसेवा) पदाधानोद्धरणम् (अवेक्षणम्) कवित्ववीजम् (प्रतिभानम्) चित्तैकाग्र्यम् (अवधानम्) इति प्रकीर्णस्याशयः ।

एतया लोकविद्याप्रकीर्णशब्दानां व्याख्यया प्रतिभाव्युत्पत्त्यभ्याससंवलितं काव्यकारणमेव सुस्पष्टं  भवति ।

\customparagraph{\textbf{काव्यप्रयोजनम्}}

वामनो यशसः , आनन्दस्य च  प्राप्तये काव्यं भवतीति स्वीकरोति । अर्थात् सत् काव्यात् दृष्टः (प्रत्यक्ष) रूपेण प्रीतिः तथा अदृष्टरूपेण कीर्तिः प्राप्यते । तद्यथा—

\begin{visheshshloka}[center]

\textbf{काव्यं सद् दृष्टाऽदृष्टार्थं प्रीतिकीर्तिहेतुत्वात् । (काव्यालङ्कारसूत्रवृत्तिः)}

\end{visheshshloka}

तथा च वृत्तिभागे—

\begin{shloka}[center]

\textbf{प्रतिष्ठां काव्यबन्धस्य यशसः सरणिं विदुः ।}

\end{shloka}

\begin{shloka}[center]

\textbf{अकीर्तिवर्त्तिनीं त्वेवं कुकवित्वविडम्बनम् ।।}

\textbf{कीर्तिस्वर्गफलामाहुराससारं विपश्चितः ।}

\textbf{अकीर्तिं तु निरालोकनरकोद्देशदूतिकाम् ।।}

\textbf{तस्मात् कीर्तिमुपादातुमकीर्तिं च निवर्हितुम् ।}

\textbf{काव्यालङ्कारसूत्रार्थः प्रसाद्यः कविपुङ्गवैः ।।}

\end{shloka}

\customparagraph{\textbf{अलङ्कारः}}

आचार्योऽयम् अलङ्कारतत्त्वं सौन्दर्यप्रयोजकं स्वीकरोति । सौन्दर्योत्पादका अलङ्कारा एव काव्यस्य ग्राह्यत्वे प्रयोजका भवन्ति । तस्मात् काव्येऽलङ्काराणां स्थानं महत्त्वपूर्णमस्तीति कथयति यत् काव्यालङ्कारसूत्रवृत्तौ—

\textbf{सौन्दर्यमलङ्कारः । काव्यं ग्राह्यमलङ्कारात् । गुणप्रतिपादितायाः काव्याशोभाया अतिशयो हेतवस्त्वलङ्काराः ।} \textbf{(काव्यालङ्कारसूत्रवृत्तिः)}

\customparagraph{\textbf{गुणः / रीतिः}}

वामनो रीतिसम्प्रदायस्याचार्योऽस्ति । अयं काव्यशरीरस्यात्मरूपां रीतिं स्वीकरोति । अस्य मते काव्यस्य कृते पदानां समायोजनमेव मुख्यो विषयोऽस्ति । एतेन विना काव्यं स्वकीयप्रयोजनसिद्धये सामर्थ्ययुतं न भवति । गुणात्मतया व्यवस्थिता रीतिरेव काव्यस्यात्मा भवितुं शक्नोति, तद्यथा—

रीतिरात्मा काव्यस्य । विशिष्टा पदरचना रीतिः । विशेषो गुणात्मा । काव्यशोभायाः कर्तारो धर्मा गुणाः ।

\customparagraph{\textbf{रीतिसम्प्रदायस्य संस्थापकः}}

काव्ये वर्णपदादीनां काव्यानुकूलव्यवस्थापनमेव रीतिर्मन्यते । काव्ये रीतितत्त्वं प्राधान्यतया मन्यमानानां विदुषां विचारसङ्ग्रह एव रीतिसम्प्रदायोऽस्ति । अस्य सर्वप्रथमः प्रतिपादको वामनोऽस्ति। अत एवाचार्योऽयं रीतिसम्प्रदायस्य संस्थापकत्वेन सम्मानितोऽस्ति । यद्यपि भामहादयो विद्वांसोऽपि रीतेर्विवेचनं कृतवन्तः, परं वामनः काव्यालङ्कारसूत्रवृत्ताख्ये स्वग्रन्थे रीतेर्महत्त्वपूर्णं स्थानं प्रादात् । एतेन रीतिरात्मा काव्यस्य’ इत्युल्लिख्य रीतितत्त्वं काव्यात्मस्थानीभूतं स्वीकृतम् । यद्यपि सर्वेषां कृते स्वीकारयोग्यं तु मतमिदं नाभूत्, तथापि काव्यशास्त्रे एको वैचारिकवैशिष्ट्यालङ्कृतः सैद्धान्तिकसम्प्रदायस्तु सञ्जात एव ।

पौरस्त्यकाव्यशास्त्रस्य इतिहासे वामनाचार्यो युगप्रवर्तकः, मौलिकविचारकः, विशिष्टसिद्धान्तकारश्च वर्तते। काश्मीरनरेशस्य जयापीडस्य सभापण्डितोऽयम् अष्टमशताब्द्यां सूत्रशैल्या ‘काव्यालङ्कारसूत्रवृत्तिः’ विरच्य साहित्यशास्त्रस्य सर्वपक्षीयविवेचनस्य सुदृढं मार्गं प्रदर्शितवान्। “रीतिरात्मा काव्यस्य” इत्युद्‌घोषयन् अयं प्रसिद्धस्य रीतिसम्प्रदायस्य संस्थापकत्वेन राजते । ‘विशिष्टा पदरचना रीतिः’ तथा ‘विशेषो गुणात्मा’ इति स्वीकृत्य अयं काव्यस्यान्तरिकं बाह्यञ्च पक्षं सम्यक् संस्थापितवान्। अनेन ‘काव्यशोभायाः कर्तारो धर्मा गुणाः’ इत्युल्लिख्य गुणतत्त्वं काव्यात्मस्थानीयत्वेन अङ्गीकृतम्, प्राचीनदशगुणानां शब्दार्थभेदेन विंशतिधा विस्तारः कृतः, अलङ्काराश्च ‘सौन्दर्यमलङ्कारः’ इति रूपेण काव्यशोभातिशयहेतवः प्रतिपादिताः, उपमा च सर्वालङ्कारमूलत्वेन अङ्गीकृता। काव्यलक्षणे गुणालङ्कारसमन्वितं शब्दार्थमयं स्वरूपं प्रदर्श्य, काव्यहेतुत्वेन लोक-विद्या-प्रकीर्णानाम् (प्रतिभा-व्युत्पत्ति-अभ्यासानाम्) तथा काव्यप्रयोजनत्वेन च प्रीति-कीर्त्योः (आनन्द-यशसोः) सप्रमाणं निरूपणं विहितम्। अतः परवर्त्तिषु लक्षणशास्त्रिषु आचार्यवामनस्य एतद् व्यापकम्, मौलिकसैद्धान्तिकञ्च विवेचनं संस्कृतसाहित्यशास्त्रस्य इतिहासे सर्वदा अनुकरणीयम् अविस्मरणीयञ्च विराजते।

\specialupakhanda{\textbf{आचार्यो रुद्रटः}}

अलङ्कारस्य प्रौढविवेचनया अयं काव्यशास्त्री भामहदण्डीवामनानन्तरं सञ्जातः’ इति मन्तुं युज्यते । धनिकमम्मटप्रतिहारेन्दुराजादिका विद्वांसः स्वग्रन्थेषु रुद्रटमतमुल्लिखन्ति । राजशेखरः काव्यमीमांसायां ‘काकुवक्रोक्तिर्नाम शब्दालङ्कारोऽयमिति रुद्रटः’ इत्युल्लिख्य स्वपूर्ववर्त्तीति प्रमाणितवान् । राजशेखरस्य समय ईसवीय ९२० तमे वर्षे स्वीक्रियते । ध्वन्यालोकलोचनकारेण रुद्रटस्य भावलक्षणं सोदाहरणमुल्लिखितम् । मम्मटोऽपि अस्य हेतुव्यतिरेकसमुच्चयाद्यलङ्कारमतं खण्डयति । 	अतो रुद्रटस्य समयो नवमशताब्दस्य पूर्वार्धो भवितुमर्हति । अस्य पितुर्नाम भट्टवावुकः आसीत् । शतानन्दोपनाम्ना विभूषितोऽयं रुद्रटेति नाम्ना प्रसिद्धो जातः । प्रायः काव्यालङ्कारग्रन्थाः कश्मीरे निर्मिताः प्राप्यन्ते । तथा च नाम्ना अपि अयं कश्मीरवासीति अनुमातुं शक्यते ।

\customparagraph{\textbf{काव्यशास्त्रीयं योगदानम्}}

अलङ्कारवर्गीकरणं रुद्रटस्य महनीयं योगदानं दृश्यते । वास्तव्यौपम्यातिशयश्लेषाख्येषु चतुर्भागेषु कृतमलङ्काराणां वैज्ञानिकं वर्गीकरणं काव्यशास्त्रीयेतिहासे प्राथमिक प्रयासोऽस्ति । तत्र वास्तव्यमूलकाः २३, औपम्यमूलकाः २१, अतिशयमूलकाः १२, श्लेषः १ समेत्य सप्तोपरिपञ्चाशदलङ्काराः (५७) प्रतिपादिताः । तेषु एकत्रिंशदलङ्काराः (३१) स्वयमेवाविष्कृताः ।

अस्य काव्यशास्त्रीयग्रन्थः काव्यादर्शोऽस्ति । आर्याछन्दसि प्रणीते ग्रन्थेऽस्मिन् षट्दशसङ्ख्यकेषु अध्यायेषु अष्टचत्वारिंशदोपरिसप्तशतसङ्ख्यकानि पद्यानि निर्मितानि । तेषु प्रथमे—काव्यप्रयोजनम्, काव्यहेतुः, कविमहिमा चोल्लिखिताः सन्ति  । द्वितीये—काव्यलक्षणम्, शब्दभेदः, वृत्तिर्वक्रोक्तिरनुप्रासश्च प्रदर्शिताः । तृतीये, चतुर्थे, पञ्चमे च यमकश्लेषचित्रालङ्काराः सविस्तारेण वर्णिताः । षष्ठे—काव्यदोषः, सप्तमे—अर्थस्य लक्षणम्, वाचकशब्दस्य भेदाः, अर्थालङ्काराणां वर्गीकरणम् तथा वास्तवमूलकालङ्काराश्च सोदाहरणं विवेचिताः । अष्टमे नवमे च—औपम्यातिशयमूलकालङ्काराः प्रस्तुताः सन्ति । दशमे—अर्थश्लेषस्य दशभेदाः, एकादशे—अर्थदोषाश्च वर्णिताः । द्वादशात् पञ्चदशाध्यायपर्यन्तं रसनिरूपणम्, षट्दशेऽध्याये—महाकाव्यमहाकथाख्यायिकालघुकथादयः काव्यभेदाः प्रदर्शिताः ।

अयमलङ्कारशास्त्री प्रथमोऽस्ति यः रसस्य महत्त्वं स्वीकरोति । अनेन रसस्य विषये प्राचीनानामलङ्कारशास्त्रिणामनुकरणमकृत्वैव स्वतन्त्ररूपेण विवेचितः । ‘तस्मात् कर्तव्यं यत्नेन महीयसा रसैर्युक्तम्’ इति कथयित्वा काव्ये रसस्यानिवार्यता सङ्केतिता । तथा च प्राचीनानां नवरसाः, तेषामुपरि दशमरसे प्रेयानपि स्वीकरोति रुद्रटः । अनेन नायकनायिकयोर्भेदविषये कृतं विवेचनं परवर्त्तिनामाचार्याणां प्रेरकमभवत् । भावस्य विश्लेषणप्रसङ्गे उदाहृतं पद्यद्वयं ध्वनिवादिनो गुणीभूतव्यङ्ग्यस्य तथा ध्वनिकाव्यस्योदाहरणं मन्यन्ते,  तद्यथा भावालङ्कारस्योदाहरणम् —

\begin{shloka}[center]

\textbf{ग्रामतरुणं तरुण्या नववञ्जुलमञ्जरीसनाथकरम् ।}

\textbf{पश्यन्त्या भवति मुहुर्नितरां मलिना मुखच्छाया ।।}

\end{shloka}

पद्यमिदं ध्वनिवादिनो मते गुणीभूतव्यङ्गस्योदाहरणम् । तथा च भावस्येव द्वितीये उदाहरणे ध्वनिवादिनो विद्वांसः ध्वनेरुदाहरणं स्वीकुर्वन्ति । तद्यथा—

\begin{shloka}[center]

\textbf{एकाकिनी यदबला तरुणिर् यथाह–}

\textbf{मस्मिन् गृहे गृहपतिश्च गतो विदेशम् ।।}

\textbf{किं याचसे यदिह वासमियं वराकी ।}

\textbf{श्वश्रूः ममान्धबधिरा ननु मूढः पान्थः ।।}

\end{shloka}

अतो ध्वनिवादिकाव्यशास्त्रिणो रुद्रटाचार्यो ध्वनितत्त्वेन सह परिचितः आसीदिति स्वीकुर्वन्ति ।

\customparagraph{\textbf{काव्यलक्षणम्}}

रुद्रटः ‘ननु शब्दार्थौ काव्यम्’ इति कथयित्वा काव्यलक्षणे पूर्ववर्त्तिनां पक्षपाती दृश्यते । स्वाग्रजस्य भामहस्य मतं यथा स्यात्तथा स्थापनेन विना काव्यस्वरूपविषयेऽस्य किमपि नावीन्यं नास्ति । स्वकीये काव्यालङ्कारे काव्यस्वरूपविषये प्रस्तौति यत्—

\begin{shloka}[center]

\textbf{ननु शब्दार्थौ काव्यं शब्दस्तत्रार्थवान्–अनेकविधः ।}

\textbf{वर्णानां समुदायः—————।’ (रुद्रटालङ्कारः)}

\end{shloka}

\customparagraph{\textbf{काव्यभेदाः}}

काव्यभेदविषये रुद्रटः प्राक् उत्पाद्य—अनुत्पाद्यभेदेन द्विधा विभजति । तत्र चायं स्वरूपभेदेन महत् लघु चेति द्विविधमुल्लिखति । तद्यथा—

\begin{shloka}[center]

\textbf{सन्ति द्विधा प्रबन्धाः काव्यकथाख्यायिकादयः काव्ये ।}

\textbf{उत्पाद्यानुत्पाद्ये महल्लघुत्वेन भूयोऽपि ।। (रुद्रटालङ्कारः)}

\end{shloka}

\customparagraph{\textbf{काव्यप्रयोजनम्}}

यशः, धनम्, असामान्यसुखम्, अभीष्टकामना च काव्यप्रयोजनं भवति । भयानकरोगान्मुक्तिः, चतुर्वर्गफलप्राप्तिश्च काव्यादेव जायते । एषु प्रारम्भिकाणि  पञ्च प्रयोजनानि कवेरेव, अन्ये तु कविसहृदययोः कृते काव्याल्लभ्यन्ते, तद्यथा काव्यालङ्कारे यथा—

\begin{shloka}[center]

\textbf{(कविगत)   	ज्वलदुज्ज्वलवाक्प्रसरः सरसं कुर्वन्महाकविः काव्यम् ।}

\textbf{स्फुटमाकल्पमनल्पं प्रतनोति यशः परस्यापि ।।}

\textbf{(पाठकगत)    ननु काव्येन क्रियते सरसानामवगमश्चतुर्वर्गे ।}

\textbf{लघु मृदु च नीरसेऽभ्यस्ते हि त्रस्यन्ति शास्त्रेभ्यः ।।( रुद्रटालङ्कारः)}

\end{shloka}

\customparagraph{\textbf{काव्यहेतुः}}

अस्य मतं काव्यहेतुविषये प्रौढमस्ति । शक्तिर्निपुणताभ्यासश्च समुदयरूपेण काव्यनिर्माणे कारणमस्ति । तस्य कृते कर्ता जनः दोषरहितः तथा आलङ्कारिको भवितुमर्हति, तद्यथा काव्यालङ्कारे—

\begin{shloka}[center]

\textbf{अधिगतसकलज्ञेय सुकवेः सुजनस्य सन्निधौ नियतम् ।}

\textbf{नक्तं दिनमभ्यसेदभियुक्तः शक्तिमान् काव्यम् ।। (रुद्रटालङ्कारः)}

\end{shloka}

\customparagraph{\textbf{अलङ्कारविवेचनम्}}

अस्य काव्यालङ्काराख्ये ग्रन्थे पञ्चशब्दालङ्काराणां सप्तचत्वारिंशद् (४७) अर्थालङ्काराणाञ्च विवेचनं प्राप्यते । पूर्ववर्त्तिनस्तु २६ अर्थालङ्काराणामेव विवेचनमकुर्वन् । अनेन समुच्चय–अनुमान–परिवृत्ति–परिकर–परिसङ्ख्या–कारणमाला–अन्योन्य–सार–एकावली–इत्यादयः २१ अलङ्काराः स्वयमेवाविष्कृताः । एते नूतनालङ्काराः प्रायशः परवर्त्तिनामाचार्याणां कृतेऽपि मान्या जाताः । अतो नवीनालङ्कारप्रतिपादनमस्य महत्त्वपूर्णं योगदानं मन्यते । अस्यालङ्कारवर्गीकरणं तु सर्वथा नावीन्यं ज्ञायते । रुद्रटोऽलङ्काराणां विभाजकतत्त्वरूपेण वास्तव्यौपम्यातिशयश्लेषान् अङ्गीकृत्य तेनैव नाम्ना चतुर्षु भागेषु विभाजितवान् । तेषु कस्यापि वस्तुनः स्वरूपवर्णनरूपं वास्तवमूलकं भवति । तत्र सहोक्तिसमुच्चययथासङ्ख्यपरिसङ्ख्याव्यतिरेककारणमालादयः २३ अलङ्काराः समाविशन्ति । पदार्थस्य सम्पूर्णताबोधनार्थं तत्सदृशपदार्थस्य वर्णनम् औपम्यमूलके भवति । तत्र उपमा–उत्प्रेक्षा–रूपक–अपह्नुति–आदयः २१ अलङ्काराः प्रविशन्ति । कस्यापि अर्थस्य वा धर्मस्य नियमः स्वकीयप्रसिद्धिबाधस्य कारणेन लोकातिक्रान्तत्वं प्राप्नोति तदा प्रधान्येन अतिशयमूलकं भवति । तत्र विभावना–तद्गुण–विरोध–असङ्गति–विषमादयः १२ अलङ्काराः आगच्छन्ति । श्लेषवर्गे तु एक एव, तद्यथा काव्यालङ्कारे—

\begin{shloka}[center]

\textbf{अर्थस्यालङ्काराः वास्तव्यौपम्यमतिशयश्लेषाः ।}

\textbf{एषामेव विशेषा अन्ये तु भवन्ति निशेषाः ।। (रुद्रटालङ्कारः)}

\end{shloka}

\customparagraph{\textbf{रसविवेचनम्}}

रसविवेचनप्रसङ्गे रुद्रटः कथयति यत्– ‘काव्याद् रसिकाः शीघ्रतया कोमलतापूर्वकं चतुर्वर्गफलप्राप्त्यर्थं प्रवृत्ता भवन्ति । ते नीरसकाव्यात् भयमाप्नुवन्ति । अतः प्रयत्नेन रसपेशलं काव्यं करणीयम् । रसस्याभावे शास्त्रवत् काव्येऽपि उद्वेगमुत्पद्यते’ इति । वचनेनानेन आचार्योऽयं रसतत्त्वपरिचितः समर्थकश्चासीदिति वक्तुं युज्यते । भरतपश्चादयमेव रसविवेचकोऽस्ति । परमनेन विभावानुभावव्यभिचारिभावविषये किमपि नोल्लिखितम् ।

आलङ्कारिकोऽपि रुद्रटो रसविषये महत्त्वपूर्णं मतं प्रस्तौति । यदि काव्ये रसो न भवेच्चेत् पाठकानां कृते शास्त्रवत् काव्यमुद्वेजनजनकमेव स्यात् । अतः सत्काव्यस्य कृते  रसनिबन्धनमत्यावश्यकं भवतीति कथयति यत्—

\begin{shloka}[center]

\textbf{तस्मात्तत् कर्तव्यं महीयसा रसैर्युक्तम् ।}

\textbf{उद्वेजनमेतेषां शास्त्रवदेवान्यथा हि स्यात् ।। (रुद्रटालङ्कारः)}

\end{shloka}

एतादृशेन रसविचारेण कतिपय समालोचका रुद्रट रसवादी आचार्यः’ इत्यपि वदन्ति ।

\customparagraph{\textbf{रससङ्ख्या}}

रससङ्ख्याविषयकं रुद्रटमतं नावीन्यं भजति । एषः प्राचीनानां मतमेव न स्वीकरोति, अपितु स्वकीयं भिन्नविचारमपि प्रस्तौति । अस्य मते अष्टौ एव रसा न, अपितु दश रसा भवन्ति । भरतस्वीकृतेषु अष्टरसेषु तेनैव नाट्येऽस्वीकृतः शान्तो रसोऽपि युक्तिसहितं स्वीकरोति रुद्रटः । तथा च स्वप्रतिभाबलात् ‘प्रेयान्’ इत्याख्यं दशमं रसं प्रस्तौति, तद्यथा—

\begin{shloka}[center]

\textbf{शृङ्गारवीरकरुणा बीभत्सभयानकाद्भुताः हास्यः ।}

\textbf{रौद्रः शान्तः प्रेयानिति मन्तव्या रसाः सर्वे ।।}

\end{shloka}

शान्तरसविषये युक्तिप्रदर्शनं यथा—

\begin{shloka}[center]

\textbf{रसनाद्रसत्वमेषां मधुरादीनामिवोक्तमाचार्यैः ।}

\textbf{निर्वेदादिष्वपि तन्निकाममस्तीति तेऽपि रसाः ।। (रुद्रटालङ्कारः)}

\end{shloka}

एभिर्विवेचनैर्रुद्रटो न केवलमलङ्कारपक्षपाती रसवाद्यप्यस्तीति स्वीकर्तुं शक्यते । काव्यलक्षण–काव्यहेतु–काव्यप्रयोजनालङ्कारविस्तारादिप्रदर्शनपूर्वकं रसमहत्त्वप्रकाशनञ्च कृत्वा काव्यशास्त्रिणानेन संस्कृतकाव्यशास्त्रेतिहासोऽलङ्कृत इति सत्यमेव ।

आचार्यरुद्रटः संस्कृतकाव्यशास्त्रीयपरम्परायां समादरणीयः काव्यशास्त्री वर्तते। भामह-दण्डि- वामनानां परवर्ती अयं काश्मीरवासी विद्वान् स्वकीयया मौलिक-वैज्ञानिकदृष्ट्या काव्यशास्त्रं नूतनोत्कर्षं निनाय । अस्य ग्रन्थः काव्यालङ्कारः साहित्यचिन्तनस्य नूतनं मार्गं प्रदर्शयति। अस्य नवीनमलङ्कारवर्गीकरणं वैज्ञानिकमभिलक्ष्यते । अनेनालङ्काराणां वास्तव्य-औपम्य-अतिशय- श्लेषभेदेन चतुर्धा वैज्ञानिकं वर्गीकरणं कृत्वा, एकत्रिंशत् (३१) नवीनालङ्काराणाम् आविष्कारेण स्वकीयमसाधारणं शास्त्रपाण्डित्यं प्रदर्शितम्। यद्यपि एषः प्रधानतया अलङ्कारशास्त्रित्वेन प्रसिद्धः, तथापि काव्ये रसस्यानिवार्यतां प्रदर्श्य \textit{‘तस्मात् कर्तव्यं यत्नेन महीयसा रसैर्युक्तम्’} इत्युल्लेखयन् अनेन रसतत्त्वस्य महत्ता प्रतिष्ठापिता । प्राचीनान् रसान् अतिरिच्य `प्रेयान्' इति दशमरसस्य परिकल्पना च कृता। राजशेखर-मम्मट-आनन्दवर्धनादिभिर्विद्वद्भिः स्वग्रन्थेषु रुद्रटमतस्य सविशेषमुल्लेखः कृतः, येन तस्य ऐतिहासिकं महत्त्वं स्वतः सिद्ध्यति। काव्यलक्षण-हेतु-प्रयोजनानां निरूपणपूर्वकं रसालङ्कारादीनां यत् प्रौढविवेचनं रुद्रटेन कृतम्, तेन संस्कृतकाव्यशास्त्रस्य इतिहासः सर्वथा समृद्धोऽभवत्। अतः काव्यशास्त्रपरम्परायाम् आचार्यरुद्रटस्य योगदानं सर्वदा चिरस्मरणीयं प्रतीयते।

\specialupakhanda{\textbf{रचनात्मककाल-उपसंहारः}}

संस्कृतकाव्यशास्त्रस्य इतिहासे अयं रचनात्मककालः सैद्धान्तिकक्रान्तेः सुवर्णमयी आधारशिला वर्तते। अस्मिन् काले काव्यशास्त्रिभी रस-अलङ्कार-गुण-रीति-प्रभृतीनि विभिन्नानि काव्यतत्त्वानि सूक्ष्मेक्षिकादृष्ट्या विवेचितानि, यद्यपि काव्यस्यात्मतत्त्वविषये आचार्याणां मध्ये महत्यो विप्रतिपत्तयो आसन् । अस्मिन् युगे भामहप्रभृतयो विद्वांसो अलङ्कारमेव काव्यस्य आत्मतत्त्वं मन्यन्ते स्म। अलङ्कारप्रेम्णा अस्य प्राधान्यं परिपोषयन्तः काव्यशास्त्रिणो विद्वांसः स्वप्रणीतानां लक्षणग्रन्थानां नामसु अपि — **'**काव्यालङ्कारः', `काव्यालङ्कारसारसङ्ग्रहः', `काव्यालङ्कारसूत्रवृत्तिः' इत्यादिरूपेण `अलङ्कार' इति शब्दं सगौरवं समाविशन् । अनेनापि तस्य कालस्य अलङ्कारकेन्द्रिता प्रवृत्तिः सुस्पष्टा ज्ञायते।

अस्मिन् समये भामह-दण्डि-उद्भट-वामन-रुद्रटादयः प्रतिभाशालिनः काव्यशास्त्रनिष्णाताः सञ्जाताः, यैर्नाट्यशास्त्रस्य परिधेर्बहिः काव्यशास्त्रस्य स्वतन्त्रविधारूपेण प्रतिष्ठापनं विहितम्। अस्मिन् रचनात्मककाले विश्लिष्टानां काव्यतत्त्वानां सुदृढपृष्ठभूमौ एव आगामिनि `निर्णयात्मककाले' काव्यतत्त्वप्राधान्यविषये प्रामाणिकसिद्धान्तमतस्य आविष्कारः प्रतिष्ठा च सञ्जात इति मन्यते। अतोऽस्मिन् काले सञ्जातानां काव्यशास्त्रिणाम् अवदानं विना उत्तरवर्त्तिनो काव्यशास्त्रिणः काव्यात्मतत्त्वादिविषये सीमानिर्धारणं असम्भवप्रायम् अनुमीयते। अस्मिन् समये सञ्जातानां विदुषाम् अनुपमं योगदानं साहित्यशास्त्रस्य इतिहासे अद्यापि शाश्वतम् अद्वितीयञ्चानुभूयते।

\specialmahakhanda{\textbf{चतुर्थ-प्रकाशः}}

\specialupakhanda{\textbf{निर्णयात्मकः कालः}}

\begin{shloka}[center]

\textbf{(आनन्दवर्धनात् ८००  मम्मटपर्यन्तम् १०००)}

\end{shloka}

\customparagraph{\textbf{विषयप्रवेशः}}

अस्य कालस्य प्रमुखाचार्योऽस्ति आनन्दवर्धनः । कालेऽस्मिन् नाट्यशास्त्रस्य ध्वन्यालोकलोचनस्य च टीकाकारः अभिनवगुप्तः, वक्रोक्तिसम्प्रदायस्य प्रवर्तकः कुन्तकः, ध्वनिविरोधी आचार्यः महिमभट्टश्च स्वविवेचनवैदुष्यं प्रादर्शयन् । तथा च रुद्रभट्ट-भोजराज- धनिक-धनञ्जयादयो विद्वांसोऽपि लक्षणशास्त्रीया ग्रन्थाः प्रकाशयन्ति स्म । तदा काव्यप्रणयने आवश्यकीयानां तत्त्वानां मध्ये कस्य प्राधान्यम् ? कस्याप्राधान्यम् ? को बाह्यसौन्दर्यपक्षः ? कश्चान्तरिकात्मपक्षः ? इत्यादिकासु जिज्ञासासु लक्षणशास्त्रिणोऽनिर्णीता आसन् । अस्यामवस्थायां काव्ये सामान्यार्थातिरिक्तो विशिष्टार्थोऽपि भवतीति–उद्घोषयन् आनन्दवर्धनः ध्वनितत्त्वस्य सस्वरूपं महत्त्वं प्रदर्शितवान् । तदनन्तरमेव रसवैशिष्ट्ययुतं ध्वनितत्त्वं काव्यात्मरूपेण स्थापितम् । तच्च प्रायशः काव्यशास्त्रिभिरनुमोदितञ्च । तस्मादेव समयस्यास्य नाम निर्णयात्मकोऽभवत् ।

\specialupakhanda{\textbf{आनन्दवर्धनाचार्यः}}

अस्य कालस्य प्रारम्भिकः प्रमुखश्चाचार्य आनन्दवर्धनोऽस्ति । एषो ध्वनिसम्प्रदायस्य संस्थापकरूपेण प्रसिद्धो जातः । कल्हणस्य राजतरङ्गिण्यामुद्धरणेनायं काश्मीरनिवास्यासीदिति ज्ञायते। तत्रायमाचार्यो राज्ञोऽवन्तिवर्मणः समकालिकरूपेण निर्दिश्यते यत्—

\begin{shloka}[center]

\textbf{मुक्ताकणशिवस्वामी कविरानन्दवर्धनः ।}

\textbf{प्रथां रत्नाकरश्चागात् साम्राज्येऽवन्तिवर्मणः ।।}

\end{shloka}

काश्मीराधिपतेरवन्तिवर्मणः समयः ८५० इसवीतः ८८४ इसवीयं यावत् स्वीक्रियते । ध्वन्यालोके उद्भटस्योद्धरणं प्राप्यते । उद्भटस्य कालः ८०० इसवीयशतको मन्यते । तथा च राजशेखरोऽपि काव्यमीमांसायां ध्वन्यालोकस्योद्धरणं करोति, यथा—

\begin{shloka}[center]

\textbf{ध्वनिनातिगभीरेण काव्यतत्त्वनिवेशिना ।}

\textbf{आनन्दवर्धनः कस्य नासीदानन्दवर्धनः ।।}

\end{shloka}

तस्मात् राजशेखरात् प्रागेवास्य समयो भवितुं शक्नोति । राजशेखरस्य समयश्च दशमशतकस्य पूर्वार्धः स्वीक्रियते । इत्यादिप्रमाणैरानन्दवर्धनस्य समयो नवमशताब्दी एवेत्यनुमातुं युज्यते ।

अस्य पिता नोणोपाध्याय  आसीदिति स एव स्वकीये देवीशतके उल्लिखति यत्—

\begin{shloka}[center]

\textbf{देव्याः स्वप्नोद्गमादिष्टदेवीशतकसंज्ञया ।}

\textbf{देशितानुपमामाधादतो नोणसुतो नुतिम् ।।}

\end{shloka}

अस्य पञ्च ग्रन्थाः प्राप्यन्ते । विषमवाणलीला, अर्जुनचरितम्, देवीशतकम्, तत्त्वालोकः, ध्वन्यालोकश्चेति । एतेष्वपि ध्वन्यालोकाख्येन शास्त्रीयग्रन्थेनैवानन्दवर्धनाचार्यस्य प्रसिद्धिरस्ति ।

\customparagraph{\textbf{काव्यशास्त्रीयं योगदानम्}}

आचार्योऽयं काव्यशास्त्रस्य विलक्षणप्रतिभासम्पन्नो विवेचक  आसीत् । ध्वनिसिद्धान्तस्य स्थापनयाऽयं ध्वनिप्रस्थापनपरमाचार्येति विशेषणेन विभूषितो जातः । अस्य ध्वन्यालोकः संस्कृतकाव्यशास्त्रस्य युगप्रवर्तकः, पौरस्त्यकाव्यचिन्तनपरम्परायाः प्रकाशस्तम्भश्च मन्यते । अलङ्कारशास्त्रीयेतिहासे आनन्दवर्धनं विना अलङ्कारशास्त्रं न कदापि पूर्णतामेति । शास्त्रजगति तस्य ध्वन्यालोकस्य प्रवेशः सुमहत् सौभाग्यमाधत्ते । ध्वनितत्त्वस्य व्याख्यातृत्वेन स एव सर्वान् आचार्यान् अतिशेते । एषः काव्यतत्त्वेषु ध्वनिरेव प्रधानमित्युद्घुष्य ध्वनिं सम्प्रदायकोटौ प्रापयत् । प्राग् भामहादिभिर्विद्वद्भिर्विवेचितानि काव्यतत्त्वानि गभीरपरीक्षणानन्तरं ध्वनितत्त्वस्य सहकारितया ध्वनावेवान्तर्भावितानि । अस्य प्रौढता, मौलिकता इत्यादिकं वैशिष्ट्यं पौरस्त्यकाव्यचिन्तनस्य मूलभूतसिद्धान्तरूपेण विद्वांसः कालान्तरं यावत् सम्मानितवन्तः ।

\customparagraph{\textbf{कृतयः}}

आचार्योऽयं ध्वनिप्रस्थानस्य प्रस्थापनाय ध्वन्यालोकाख्यं ग्रन्थं रचितवान् । अस्य ध्वन्यालोकातिरिक्ता विषमवाणलीला, अर्जुनचरितम्, देवीशतकम् इति त्रयः काव्यग्रन्थाः प्राप्यन्ते । काव्यशास्त्रीयग्रन्थस्तु ध्वन्यालोक इत्येक एव ।

\customparagraph{\textbf{ध्वन्यालोकः}}

ध्वन्यालोकस्य  कारिकाभागस्य वृत्तिभागस्य च रचनाया विषये मतभेदा दृश्यन्ते । तथापि प्रमाणबलात् द्वयोरेव भागयो रचयिता आनन्दवर्धन एवेति निश्चीयते ।

ध्वन्यालोके चत्वार उद्योताः सन्ति । तेषु प्रथमे ध्वन्यभाववादिनां सयुक्तिकं  खण्डनं विद्यते । आनन्दवर्धनोऽस्य प्रसङ्गस्योल्लेखः प्रथमायामेव कारिकायां करोति, यत्–

\begin{shloka}[center]

\textbf{काव्यस्यात्मा ध्वनिरिति बुधैर्यः समाम्नातपूर्व–}

\textbf{स्तस्याभावं जगदुरपरे भाक्तमाहुस्तमन्ये ।}

\textbf{केचिद्वाचां स्थितमविषये तत्त्वमूचुस्तदीयं}

\textbf{तेन ब्रूमः सहृदयमनः प्रीतये तत्स्वरूपम् ॥ (ध्वन्यालोकः)}

\end{shloka}

तत्र \textbf{प्रथमोद्योते-} ध्वनिविरोधिषु नामग्रहणप्रसङ्गेऽभाववादिनां कृते जगदुरिति परोक्षलिट्लकारः, तथा च अनिर्वचनीयवादिनां कृतेऽपि ऊचुरिति परोक्षलिट्लकार इति अदृश्यभूतकालार्थप्रकाशकस्य लकारस्यैव प्रयोगः कृतः । तस्मात् ध्वनिकारेण एते मते स्वप्रतिभाबलात् सम्भाव्योल्लिखिते इति मन्यते । भाक्तवादिनां कृते तु आहुरिति लट्लकारस्य प्रयोगादनेन भामहवामनादयोऽलङ्कारवादिनः सूच्यन्ते । एतेषामभावभाक्तानिर्वचनीयवादिनां मतानि सतर्कं विखण्ड्य अनेन समहत्त्वं ध्वनितत्त्वं स्थापितम्  । ततः परं वस्तुध्वनिः, अलङ्कारध्वनिः, रसध्वनिश्च प्रतिपादितः । \textbf{द्वितीयोद्योते–} अविवक्षितवाच्य- ध्वनेर्विवक्षितान्यपरवाच्यध्वनेश्च भेदोपभेदानां साङ्गोपाङ्गवर्णनं कृत्वा गुणानामपि विवेचनं कृतमास्ते । \textbf{तृतीयोद्योते–} पदैः, वाक्यैः, पदांशैः, रचनादिभिश्च ध्वनेः पदप्रकाश्यतायाः प्रतिपादनं कृत्वा रसानां विरोधाविरोधानां सिद्धान्ताः प्रस्तुताः । \textbf{चतुर्थोद्योते–} ध्वनेः, गुणीभूतव्यङ्ग्यस्य समावेशेन प्राक् वर्णितवस्त्वपि अपूर्वचमत्कारमावहतीति प्रतिपादितमस्ति । यथा– जीर्णेऽपि वृक्षे ऋतुराजसमागमेन वासन्ती श्रीरभिव्याप्ता भवति तेनैव प्रकारेण ध्वनिभेदेषु केनापि भेदसम्मिश्रणेन प्राचीनकविवर्णितोऽपि पदार्थः विलक्षणसौन्दर्यमभिव्यनक्ति’ इत्यादि कविप्रतिभाया आनन्त्यं समहत्त्वं वर्णितमस्ति ।

\customparagraph{\textbf{काव्यलक्षणम्}}

आनन्दवर्धनाचार्यो ध्वनिसम्प्रदायस्य संस्थापकोऽस्ति । अतो ध्वन्यात्मकमेव काव्यलक्षणं प्रस्तौति । स ध्वन्यालोके  प्रतीयमानार्थव्यञ्जनसामर्थ्यौ शब्दार्थौ एव काव्यत्वं स्वीकरोति । एतादृशे काव्ये ध्वनितत्त्वमात्मरूपेणावस्थितं भवति, तद्यथा—

सहृदयहृदयाह्लादव्यञ्जकं शब्दार्थमयत्वमेव काव्यस्य लक्षणं भवति, अर्थात् सहृदयहृदयाह्लादकौ शब्दार्थावेव काव्यमिति काव्यस्वरूपं ध्वन्यालोके प्रस्तौति, यथा-

\begin{shloka}[center]

\textbf{सोऽर्थस्तद्व्यक्तिसामर्थ्ययोगी शब्दश्च कश्चन ।}

\textbf{यत्नतः प्रत्यभिज्ञेयौ तौ शब्दार्थौ महाकवेः ।।}

\end{shloka}

क्रोञ्चद्वन्द्ववियोगदर्शनाद् हृदि आविष्कृतः शोक एव श्लोकत्वेन आदिकविवाल्मीकिमुखाद् आदिकाव्यम् अवतरितम् इति तत्र वाल्मीकिपद्यमुल्लिखति यत्-

\begin{shloka}[center]

\textbf{मा निषाद प्रतिष्ठां त्वमगमः शाश्वतीः समा ।}

\textbf{यत्क्रोञ्चमिथुनादेकमवधीः काममोहितम् ॥}

\end{shloka}

अत रसध्वनियुक्तौ शब्दार्थौ काव्यमिति फलितम् ।

\begin{shloka}[center]

\textbf{काव्यस्यात्मा स एवार्थस्तथा चादि कवेः पुरा ।}

\textbf{क्रौञ्चद्वन्द्ववियोगोत्थः शोकः श्लोकत्वमागतः ।। (ध्वन्यालोकः)}

\end{shloka}

\customparagraph{\textbf{काव्यात्मतत्त्वविचारः}}

ध्वन्यालोकस्य प्रथमपद्येन काव्यस्यात्मा ध्वनिरिति बुधैः.... इति काव्यस्यात्मस्थानीयं तत्त्वं ध्वनिरित्युद्घोषयति आनन्दवर्धनः । ध्वन्यालोककारात् पूर्ववर्त्तिनोऽऽचार्या अलङ्कार-रीति-गुणादीनि काव्याङ्गानि काव्यस्यात्मतत्त्वानीति उक्तवन्तः । तेषां मतानि विविच्य ध्यवन्यालोककारो प्रतीयमानार्थप्रतिपादकं ध्वनिमेव काव्यस्यात्मतत्त्वं स्वीकरोति । अभिधालक्षणातात्पर्याख्याभिस्तिस्रः शक्तयः शब्दार्थमेव प्रतिपादयन्ति । एषा शक्तित्रयव्यतिरिक्ता काचिद् व्यञ्जनावृत्तिर्भवति, सा ध्वन्यर्थं प्रकाशयति । ध्वनिकारो ध्वनितत्त्वप्रतिपादने बीजरूपेण वैयाकरणस्य स्फोटतत्त्वं स्वीकरोति । एषोऽलङ्कारादीनां काव्यतत्त्वानां स्थाननिर्धारणं कृत्वा ध्वनिं प्रतिपादयति । अस्य मतेऽलङ्कारादिकं केवलं वाच्यवाचकचारुत्वहेतुव्यतिरिक्तं न किमप्यस्ति । शब्दार्थौ स्वकीयं प्राधान्यतां परित्यज्य व्यञ्जनावृत्त्या वाच्यातिरिक्तं प्रतीयमानार्थं प्रतिपादयतः, स एव ध्वनिः । तदुक्तं यत्–

\begin{shloka}[center]

\textbf{यत्रार्थः शब्दो वा  तमर्थमुपसर्जनीकृतस्वार्थौ ।}

\textbf{व्यङ्क्तः काव्यविशेषः स ध्वनिरिति सूरभिः कथितः ।। (ध्वन्यालोकः)}

\end{shloka}

\customparagraph{\textbf{काव्यभेदाः}}

आनन्दवर्धनो ध्वनितत्त्वप्राधान्यं विविच्य काव्यभेदान् प्रस्तौति । यत्र ध्वनेः प्राधान्यं भवति तत्रोत्तमं ध्वनिकाव्यं स्वीक्रियते । यत्र चानुत्तमा अर्थात् गौणरूपेण ध्वनिरस्ति तत्र मध्यमं काव्यभेदं मन्यते । तथा च यत्र शब्दार्थचमत्कृतिरेवास्ति, ध्वनेः लवलेशोऽपि न दृश्यते तत्तु चित्राख्यमधमकाव्यं गण्यते । अनेन प्रकारेणायम् उत्तममध्यमाधमभेदेन काव्यस्य त्रिभेदान् स्वीकरोति, तद्यथा ध्वन्यालोके–

\begin{shloka}[center]

\textbf{गुणप्रधानभावाभ्यां व्यङ्ग्यस्यैवं व्यवस्थिते ।}

\textbf{काव्ये उभे ततोऽन्यद्यत् तच्चित्रमभिधीयते ।।}

\textbf{चित्रं शब्दार्थभेदेन द्विविधं च व्यवस्थितम् ।}

\textbf{तच्चकिञ्चिच्छब्दचित्रं वाच्यचित्रमतपरम् ।। (ध्वन्यालोकः)}

\end{shloka}

\customparagraph{\textbf{रीतिविवेचनम्}}

ध्वनिकारो रीतितत्त्वं  ध्वनेः (रसस्य) व्यञ्जकत्वेन स्वीकरोति । तत्र अलक्ष्यक्रमव्यङ्ग्यरूपस्य ध्वनिभेदस्य व्यञ्जकतत्त्वानां वर्णनावसरे ‘सङ्घटना’ नाम्ना रीतेः स्वरूपमुल्लिखति, यथा—

\begin{shloka}[center]

\textbf{असमासा समासेन मध्यमे न च भूषिता ।}

\textbf{तथा दीर्घसमासेति त्रिधा सङ्घनोदिता ।। (ध्वन्यालोकः)}

\end{shloka}

समासरहिता, मध्यमसमासयुता तथा दीर्घसमासयुता त्रिधा सङ्घटना भवति । सा च माधुर्यादीन् गुणानाश्रित्य काव्यात्मस्थानीभूतान् रसान् व्यञ्जयतीति कथयित्वा रीतितत्त्वं ध्वनेः सहायकतत्त्वरूपेण स्वीकरोति ध्वनिकारः ।

\customparagraph{\textbf{गुणविवेचनम्}}

ध्वनिकारो रीतेरपेक्षया गुणतत्त्वं रसध्वनेर्निकटसम्बन्धयुतं स्वीकरोति । स रीतिवत् गुणानामनियतावस्था नास्तीति कथयन् उल्लिखति यत्—

\textbf{गुणानां हि माधुर्यप्रसादप्रकर्षः करुणविप्रलम्भशृङ्गारविषय एव । रौद्राद्भुतादिविषयमोजः । माधुर्यप्रसादौ रसभावतदाभासविषयावेव, इति विषयनियमो व्यवस्थितः । तथा च गुणास्तु व्यङ्ग्यविशेषावभासिवाच्यप्रतिपादनसमर्थशब्दधर्मा एव} इति । \textbf{(ध्वन्यालोकः)}

\customparagraph{\textbf{अलङ्कारविवेचनम्}}

अलङ्कारविषयेऽऽनन्दवर्धनस्य मतमेवोत्तरवर्त्तिनो काव्यशास्त्रिणः सादरं स्वीकुर्वन्ति । स काव्येऽलङ्कारस्थानं काव्यात्मभूतस्य रसस्योपकारकत्वेन भवतीति स्वीकरोति । तथा च यथा मानवशरीरे कटककुण्डलादयो भवन्ति तथैव शब्दार्थमये काव्यशरीरेऽपि अलङ्कारा बाह्यसौन्दर्यवर्धका एव भवन्ति,  तद्यथा—  अङ्गाश्रितास्त्वलङ्कारा मन्तव्याः कटकादिवत् ।

तथा च ध्वनिकाव्येऽलङ्काराः काव्यप्रणयनकाले पृथक् यत्नेन विना समागता एव काव्याङ्गीभूतस्य रसस्य कृते उपयोगिनो भवन्ति । विपरीतास्तु रसप्रतीतौ–अवरोधका जायन्ते, तद्यथा ध्वन्यालोके—

\begin{shloka}[center]

\textbf{रसाक्षिप्ततया यत्र बन्धः शक्यक्रियो भवेत् ।}

\textbf{अपृथक् यत्ननिर्वर्त्य सोऽलङ्कारो ध्वनौ मतः ।। (ध्वन्यालोकः)}

\end{shloka}

`ध्वनेर्भिन्ना वाच्यालङ्काराः' इति प्रमाणीकर्तुं कथयति यत्—

\begin{shloka}[center]

\textbf{व्यङ्ग्यस्य यत्राप्राधान्यं वाच्यमात्रानुयायिनः ।}

\textbf{समासोक्त्यादयस्तत्र वाच्यालङ्कृतयः स्फुटाः ।। (ध्वन्यालोकः)}

\end{shloka}

अनेन प्रकारेण ध्वनिकारो बाह्यसौन्दर्यबर्धकतत्त्वरूपेणालङ्कारस्य स्वरूपमुल्लिखति । पश्चात् काव्यशास्त्रीयपरम्परायामनेन निर्धारितं काव्येऽलङ्कारस्य स्थानमेव सर्वस्वीकार्यरूपेण स्थिरीभूतम् ।

\customparagraph{\textbf{औचित्यसम्प्रदायस्य प्रारूपप्रदर्शकः}}

आचार्योऽयम् औचित्यसम्प्रदायस्य स्थापनार्थं स्पष्टरूपेण प्रारूपप्रदर्शनात् क्षेमेन्द्रस्य कृते उपकारकोऽस्ति । एषो ध्वन्यालोके सङ्घटनावृत्त्यलङ्कारादीनाम् औचित्यपूर्णविनिवेशनमेव काव्यात्मभूतस्य रसस्योपकारकत्वं सिद्ध्यतीति कथयति यत्—

\begin{shloka}[center]

\textbf{एतद्यथोक्तमौचित्यमेव तस्या नियामकम् ।}

\textbf{सर्वत्र गद्यबन्धेऽपि छन्दो नियमवर्जिते ।।}

\textbf{रसाद्यनुगुणत्वेन व्यवहारार्थशब्दयोः ।}

\textbf{औचित्यवान् यस्ता एता वृत्तयो द्विविधा स्थिता ।}

\textbf{रसभावादितात्पर्यमाश्रित्य विनिवेशनम् ।}

\textbf{अलङ्कृतीनां सर्वासामलङ्कारत्वसाधनम् ।।  (ध्वन्यालोकः)}

\end{shloka}

काव्यरचनायां काव्यकर्तृजनै रसस्यौचित्यानुकूलौ शब्दार्थावेव योजनीयौ इत्यस्मिन् विषये सावधानेन भाव्यम् इति ध्वन्यालोके ध्वनिकारः प्रस्तौति, यथा—

\begin{shloka}[center]

\textbf{वाच्यानां वाचकानां च यदौचित्येन योजनम् ।}

\textbf{रसादिविषयेणैतत् मुख्यं कर्म महाकवेः ।। (ध्वन्यालोकः)}

\end{shloka}

तथा च रसप्रतीतिर्विषये तु तस्य प्रतिबन्धको मुख्यरूपेण औचित्याभाव एव मन्यते, तत्र ध्वन्यालोके यथा—

\begin{shloka}[center]

\textbf{अनौचित्यादृते नान्यत् रसभङ्गस्य कारणम् ।}

\textbf{प्रसिद्धौचित्यबन्धस्तु रसस्योपनिषत्परा ।। (ध्वन्यालोकः)}

\end{shloka}

आनन्दवर्धनस्येतत्प्रकारक औचित्यविचार आचार्यक्षेमेन्द्रस्य कृते औचित्यसम्प्रदायस्थापनार्थं प्रेरकोऽभूदिति मन्तुं युज्यते ।

आनन्दवर्धनो नवमशताब्द्यां जातः काश्मीरदेशीयः संस्कृतकाव्यशास्त्रीयेतिहासस्य कश्चन युगप्रवर्तकः, ध्वनिप्रतिष्ठापकश्च आसीत्। राज्ञोऽवन्तिवर्मणः समकालिकः, नोणोपाध्यायस्य सुतोऽयं मनीषी विषमबाणलीला-अर्जुनचरित-देवीशतक-ध्वन्यालोकादिग्रन्थानां प्रणेता आसीत्, परन्तु अस्य अक्षयकीर्तेः शास्त्रीयस्तम्भः `ध्वन्यालोकः' एव वर्तते। अनेन ग्रन्थेन प्राक्तनानां ध्वनिविरोधिनां (अभाववादः, भाक्तवादः, अनिर्वचनीयवादश्चेति) मतानां तर्कसङ्गतं खण्डनं कृत्वा ``काव्यस्यात्मा ध्वनिः'' इति शाश्वतसिद्धान्तस्य सुदृढप्रतिष्ठापनं विहितम्, येन स `ध्वनिप्रस्थापनपरमाचार्यः' इति विशेषणेन विभूषितो जातः। अस्य अपरं काव्यशास्त्रीयं वैशिष्ट्यं समग्रकाव्यतत्त्वानां रसपरकनवीनव्याख्याने दृश्यते। तेन व्यङ्ग्यार्थस्य प्राधान्याप्राधान्याभ्यां काव्यस्य त्रयो भेदाः स्वीकृताः । ते उत्तमं (ध्वनिकाव्यम्), मध्यमं (गुणीभूतव्यङ्ग्यम्), अधमं (चित्रकाव्यम्) चेति। एवमेव तेन रीतिः `सङ्घटना' (असमासा, मध्यमसमासा, दीर्घसमासा) रूपेण रसाभिव्यञ्जिका स्वीकृता, गुणा रसस्य अन्तरङ्गधर्माः, अलङ्काराश्च मानवशरीरस्य कटकादिवत् काव्यशरीरस्य बाह्यसौन्दर्यविधायका उपकारकाश्चेति प्रतिपादिताः। विशेषतः \textit{``अनौचित्यादृते नान्यत् रसभङ्गस्य कारणम्''} इति उद्घोष्य तेन औचित्यस्य यत् प्रारूपं प्रदर्शितम्, तदेव कालान्तरे आचार्यक्षेमेन्द्रस्य औचित्यसम्प्रदायस्य आधारशिलारूपं सञ्जातम्। आनन्दवर्धनस्य इयं रसध्वनिगर्भा समन्वयवादिनी दृष्टिः पौरस्त्यकाव्यचिन्तनपरम्पराया अद्वितीयः प्रकाशस्तम्भोऽस्ति, योऽद्यापि संस्कृतसाहित्यशास्त्रे सर्वमान्यः सर्वोत्कृष्टश्च वर्तते।

\specialupakhanda{\textbf{भट्टनायकः}}

संस्कृतकाव्यशास्त्रस्य प्रसिद्ध आचार्योऽस्ति भट्टनायकः । अस्य एक एव ग्रन्थो हृदयदर्पणाभिधान आसीत् । किन्तु दुर्भाग्यवशात् स ग्रन्थोऽधुना नोपलभ्यते । तथापि तदुत्तरवर्त्तिभिराचार्यैः शास्त्रीयविवेचने हृदयदर्पणस्य समधिका चर्चा विहिता । अयं ध्वनिविरोधिष्वाचार्येषु अन्यतम आसीत् । अनेन स्वकीये ग्रन्थे ध्वन्यालोकस्य सिद्धान्तानां खण्डनमकार्षीत् । तस्मात् आचार्योऽयं आनन्दवर्धनस्यानन्तरं जातः । अस्य ध्वनिविरोधिनां विचाराणां खण्डनं ध्वन्यालोकलोचनकारोऽभिनवगुप्तः कृतवान् । अतः भट्टनायकोऽभिनवगुप्तस्य पूर्ववर्त्तीति ज्ञायते । एतैः प्रमाणैरस्य समयो दशमशताब्द एव निश्चीयते ।

\begin{itemize}
\item 
\end{itemize}
\customparagraph{\textit{काव्यशास्त्रीयं योगदानम्}*}

भट्टनायकः स्वोत्तरवर्त्तिनामाचार्याणां ग्रन्थेषु हेतुद्वयेन ख्यातिं प्राप्नोति यत्– ध्वनिविरोधिना आचार्यत्वेन, रसनिष्पत्तिविषयकस्य सिद्धान्तस्य प्रतिपादकत्वेन च । ध्वन्यालोकस्य ध्वनिसम्बन्धिनां  विचाराणां खण्डनं करोति भट्टनायकः । परमभिनवगुप्तो प्रबलयुक्तिसहितं ध्वनिविरोधि  मतं विखण्ड्य ध्वनिं काव्यात्मतत्त्वे स्थिरीकरोत्येव ।

रसनिष्पत्तिविषयेऽस्य स्वतन्त्रो विचारो वर्तते । एतस्य साङ्ख्यदर्शनानुकूलं रसनिष्पत्तिविषयकं मतं भुक्तिवादस्य नाम्ना काव्यशास्त्रे प्रसिद्धमस्ति । काव्यप्रकाशकारोऽपि स्वग्रन्थे भट्टनायकस्य रसनिष्पत्तिसूत्रस्य व्याख्याया उद्धरणमकरोत् । तस्यां व्याख्यायां भट्टनायकोऽभिधालक्षणाभ्यां भिन्नयोः शब्दस्य भावकत्व–भोजकत्व–व्यापारयोः कल्पनामकरोत् । तत्र भावकत्वव्यापारेण अनुकार्यभूतयोः नायकनायिकयोः साधारणीकरणं भवति, तदनु भोजकत्वव्यापारेण सामाजिकं  रसास्वादनं करोति ।

\customparagraph{\textbf{भुक्तिवादस्य वैशिष्ट्यम्}}

काव्ये अभिधया सीतारामादयो रत्यादयश्च ज्ञायन्ते । पश्चात् भावकत्वव्यापारेण सीतादिविषये अगम्या इयम्’ इत्यादिरसावरोधिज्ञानप्रतिबन्धद्वारा कान्तात्वादिरसानुकूलधर्मयुतेन सीतादयः पदार्थाः प्रतीतेर्विषयीक्रियन्ते । एवम् प्रकारेण रामसीतादेशकालवयोवस्थादिषु साधारणीकृतेषु पूर्वव्यापारविरते सति भोजकत्वाख्यस्य तृतीयव्यापारस्य महिम्ना रजस्तमसो निगीर्णयोः प्रवृद्धसत्वगुणजनितेन वेद्यान्तरस्पर्शशून्येन आत्मानन्दसाक्षात्कारेण आस्वादविषयीकृतो भावोपनीतः साधारणात्मारत्यादिस्थायीभाव एव रसः । अर्थात् अभिधयोपस्थापितेषु भावकत्वेन साधारणीकृतेषु विभावादिषु भोजकत्वेन व्यापारेण साक्षात्कारविषयतां नीतै रत्यादिस्थायिभावरूपो रसो निष्पद्यते ।

अनेन रससूत्रस्य व्याख्यायां प्रस्तुतं साधारणीकरणसिद्धान्तप्रतिपादनं पौरस्त्यकाव्यशास्त्रस्य कृते महत्त्वपूर्णमवदानं मन्यते । भावकत्वभोजकत्वाख्ययोरप्रसिद्धशब्दव्यापारयोः  परिकल्पनमेवास्य मतस्य न्यूनत्वं दृश्यते । अस्यैव सिद्धान्तस्य परिष्कारपूर्वकम् अभिनवगुप्तः सर्वमान्यं रसनिष्पत्तिविषयकं सिद्धान्तमतं प्रस्तौति ।

\customparagraph{\textbf{आचार्योऽभिनवगुप्तः}}

पौरस्त्यकाव्यशास्त्रे अभिनवगुप्तस्य नाम निर्णायकविवेचकरूपेण गृह्यते । अयं ध्वन्यालोकस्य लोचनाख्यायाष्टीकाया नाट्यशास्त्रस्य अभिनवभारतीटीकायाश्च लेखकरूपेण ख्यातिप्राप्तोऽस्ति । अस्य टीकाकृतिरपि मौलिकग्रन्थवत् सर्वैरादृता सञ्जाता । स यद्यपि काश्मीरवासी  आसीत्, किन्तु तस्य पूर्वजाः काश्मीरवासिनो नासन् । अभिनवगुप्तस्य जन्मनः प्राक् तस्य पूर्वजा उत्तरप्रदेशस्य कान्यकुब्जनगरस्य वासिनः समभवन् । तस्मिन् समये तन्नगरं विकासस्य पराकाष्ठामलभत् । तदा कान्यकुब्जसाम्राजस्याधिपतिर् यशोवर्मा अवर्तत । केनचित् कारणेन काश्मीराधिपतिर् ललितादित्योऽष्टमे शताब्दे कान्यकुब्जेश्वरराज्यं प्रति आक्रमणं चकार । तदा कान्यकुब्जनरेशो यशोवर्मा पराजितो बभूव । तस्य राज्येऽभिगुप्तनामा कश्चन विद्वान् निवसति स्म । ललितादित्योऽभिनवगुप्तं ससम्मानं स्वराजधानीमनयत् । तत्र तस्मै निवासमकारयत्, इत्यादि वृत्तान्तोऽभिनवगुप्तः स्वकीये तन्त्रालोकाभिधाने ग्रन्थे समुल्लिखति । तत्रैव स्वपितुर्नाम नृसिंहगुप्तस्तथा पितामहस्य  नाम  वराहगुप्तश्चोल्लिखितमस्ति । अभिनवगुप्तः स्वकीयानां विवृतिविमर्शिनी–क्रमस्तोत्र–भैरवस्तोत्राख्यानां ग्रन्थानां रचनाकालं निर्दिशति स्म । क्रमस्तोत्रस्य रचनाकालः ९९० ईसवीयः, भैरवस्तोत्रस्य रचनाकालः ९६२ ईसवीयः, विवृतिविमर्शिन्याश्च १०१४ ईसवीयो वर्तते । अत एवास्य समय ईसवीयदशमस्योत्तरार्धात् एकादशस्य पूर्वार्धपर्यन्तो मन्यते ।

\customparagraph{\textbf{काव्यशास्त्रीयं योगदानम्}}

अभिनवगुप्तस्य  कृतयो बह्व्यः सन्ति । परं काव्यशास्त्रसम्बद्धास्तु ध्वन्यालोकलोचनम्, अभिनवभारती (नाट्यशास्त्रस्य टीका), काव्यकौतुकाभरणञ्च तिस्रः कृतयः सन्ति । एतेषु केवलं मूलग्रन्थस्य व्याख्या एव न, अपितु काव्यशास्त्रे विवादयुतेषु प्रसङ्गेषु सतर्कं सर्वस्वीकारयोग्यं मतमुपस्थापितमस्ति । रसध्वनिगुणालङ्कारादिविषये स्थाने स्थाने समालोचका अस्य विचारं  प्रमाणरूपेणोपस्थापयन्ति । मौलिकग्रन्थान् विनैवानेन काव्यशास्त्रे सम्मानितं स्थानमाप्तम् । भरतस्य रससूत्रस्य व्याख्यायामनेन प्रस्तुतं मतम् अभिव्यक्तिवादस्य नाम्ना रसिकसमाजे समादरणीयमस्ति ।

\customparagraph{\textbf{अभिव्यक्तिवादस्य वैशिष्ट्यम्}}

लोके जनसाधारणो विभिन्नानुभूतिं करोति । स कदाचित् केनचिज्जनेन सह अनुरागं करोति तथैव द्वेषं भयं च प्रस्तौति । एतादृशानां विषयाणामनुभवं तु विनश्यति परं तेषां संस्कारस्तस्य हृदि वासनात्मतया नित्यरूपेण वसति । त एव मानवहृदयेऽवस्थिता भावाः काव्यजगति रतिशोकादिस्थायिभावाः कथ्यन्ते । यथा लोके सीतादीनि आलम्बनकारणानि, चन्द्रिकादीनि उद्दीपनकारणानि च भवन्ति । एतेषां संयोगावस्थायां रोमाञ्चादीनि, वियोगावस्थायाञ्च अश्रुपातादिकार्याणि तथा हर्षचिन्तादीनि सहायककारणानि जायन्ते । एतान्येव कार्यकारणसहकारिकारणानि सुललितकाव्येन समर्पितै रसिकहृदयं प्रविष्टै रसास्वादनचातुरीसहकृतेन काव्यस्य पुनः पुनरनुसन्धानरूपभावना–विशेषमहिम्ना    लौकिकसीतारामादिभिरसाधारणधर्मविनाशैरलौकिक–विभावानुभावव्यभिचारीशब्दाभिधेयैः सीतादिभिरालम्बनकारणैः चन्द्रिकादिभिरुद्दीपनकारणैरश्रुपातादिभिः कार्यैश्चिन्तादिभिश्च सहकारिकारणैर्मिलित्वा उत्पादितेनालौकिकव्यापारेण सच्चिदानन्दरूपस्यात्मनो दीपस्य शरावादिरावरणमिव अज्ञानावरणभङ्गकरणानन्तरं परिमितजीवधर्मेषु नष्टेषु प्रमात्रा स्वप्रकाशतया प्रत्यक्षेण निजस्य आत्मानन्देन सहानुभूयमानः प्रागेव वासनारूपेण हृदयावस्थतो रत्यादिरेव रसः ।

अस्य मते विभावादीनां रत्यादिस्थायिभावेन सह संयोगात्= व्यङ्ग्यव्यञ्जकसम्बन्धाद् रसनिष्पत्तिः=अभिव्यक्तिर्भवति । रसनिष्पत्तिविषये काव्यशास्त्रेऽयमभिनवगुप्तस्य विचारः सर्वस्वीकृतोऽस्ति ।

आचार्योऽभिनवगुप्तः काश्मीरदेशीयो भारतीयकाव्यशास्त्रीयसौन्दर्यदर्शनस्य अद्वितीयो निर्णायक आसीत्। मूलतः कान्यकुब्जदेशीयोऽयं महाभागो ध्वन्यालोकस्य `लोचन'टीकायास्तथा नाट्यशास्त्रस्य `अभिनवभारती'टीकायाः प्रणयनेन केवलं साधारणव्याख्याकर्तृत्वेनैव न, अपितु मौलिकविचारप्रदायकत्वेन काव्यशास्त्रजगति प्रतिष्ठितोऽभवत्। रस-ध्वनि-गुण-अलङ्कारादिषु अनिर्णितेषु गभीरेषु च काव्यतत्त्वेषु अस्य तर्कसङ्गतं विवेचनं सर्वैः समालोचकैरद्यापि प्रामाणिकं मार्गदर्शकञ्च स्वीकृतम्। अपरं सर्वोत्कृष्टं श्लाघनीयञ्च योगदानं भरतमुनेर्रससूत्रस्य व्याख्याप्रसङ्गे `अभिव्यक्तिवादः' इत्याख्यसम्प्रदायस्य सुदृढं प्रतिष्ठापनम् अस्ति। अस्य मते सहृदयानां हृदये वासनारूपेण नित्यस्थिता रत्यादयः स्थायिभावा एव नाट्य-काव्ययोरलौकिकविभावादिभिः सम्यक् व्यञ्जिताः सन्तो रसत्वेन परिणमन्ति। अस्य मते काव्यात्मभूतस्य रसस्य आस्वादोऽज्ञानावरणभङ्गेन समुद्भूतोऽलौकिकः स्वप्रकाशानन्दरूपो भवति, यत्र विभावादीनां स्थायिभावेन सह व्यङ्ग्यव्यञ्जकसम्बन्धस्तिष्ठति। अभिनवगुप्तस्य इयं दर्शनमयी रससमीक्षा काव्यशास्त्रस्य शिरोभूषणत्वेन राजते ।

\specialupakhanda{\textbf{आचार्यो राजशेखरः}}

अयमाचार्यः कान्यकुब्जप्रदेशस्य प्रतिहारवंशीयनरेशस्य महेन्द्रपालस्य तत्सुनोर् महीपालस्य च गुरुः, तयोः सभायां राजपण्डितश्चासीत् । एतत् तथ्यं तस्य कर्पूरमञ्जरी, बालरामायणम्, तथा बालभारतम् इत्येतेषु ग्रन्थेषु प्राप्यते । सियोदिनस्य शिलालेखेनापि महेन्द्रपालस्य शासनकालः ९०३ तमात् ९१७ ईसवीयपर्यन्तमासीत् । अतोऽस्मिन् कालेऽस्य स्थितिरुचिता ज्ञायते । अस्य सर्वप्रथमोल्लेखः सोमदेवेन स्वकीये यशस्तिलकचम्पूग्रन्थे विहितः । ततो धनपालेन स्वकीयां तिलकमञ्जर्यां राजशेखरस्य प्रशंसा कृता । एतेनापि अस्य समयो दशमशताब्द एवेति प्रमाणितो भवति । असौ मूलरूपेण महाराष्ट्रियः, कान्यकुब्जवास्तव्य आसीत् । अस्य पिता अकालजलदो दर्दुरकः, माता च शीलवत्यासीत् । असौ चौहानवंशीयां कामपि विदुषीं कन्यां परिणीतवान्, या संस्कृतप्राकृतभाषयोर् मर्मज्ञा आसीत् । तस्या एव प्रेरणया राजशेखरः प्राकृतभाषायां अवन्तीसुन्दरीनाटकस्य रचनामकरोत् ।

\customparagraph{\textbf{काव्यशास्त्रीयं योगदानम्}}

अस्य शास्त्रीयविषये काव्यमीमांसा नाम सुप्रसिद्धो ग्रन्थो विद्यते । अस्य ग्रन्थस्य अष्टादशसङ्ख्यकानि अधिकरणानि उपकल्पितानि, परम् अधुना कविरहस्याभिधं प्रथममधिकरणमेवोपलब्धमास्ते । प्रकरणेऽस्मिन्नष्टादशाध्यायाः सन्ति । तेषु कवीनाम् उपयोगिसमस्तसिद्धान्तानां वर्णनं कृतम् । एतेषु अध्यायेषु काव्यस्वरूपकाव्यप्रकारकाव्यभेद- रीतिनिरूपणकाव्यार्थयोनिशब्दहरणार्थहरणादीनां विषये विवेचनं कृतं वर्तते । केषाञ्चन विषयाणां तु अत्र नवीनतयोपन्यासं चकार राजशेखरः, यथा– काव्यपुरुषोत्पत्तिः, साहित्यविद्यावध्वा सह तस्य विवाहः, सम्बन्धश्च ।

राजशेखरो भूगोलवेत्ता, विभिन्नप्रदेशानां भाषाशैल्यादीनां सुविज्ञश्चासीत् । अत एवात्र प्राचीनभारतस्य भूगोलज्ञानस्य प्रचुरा सामग्री स्थापिता विद्यते । तथा च स्वग्रन्थे प्राचीनानां साहित्यशास्त्रिणां बहूनि नामानि उल्लिखितानि । येनात्मप्रशंसाविमुखानां तेषां समयनिर्धारणे सारल्यं भवति । काव्यशास्त्रीयविषयेऽस्य विचारा अधः प्रन्तूयन्ते ।

\customparagraph{\textbf{काव्यलक्षणम्}}

राजशेखरो गुणालङ्कारयोः सुरुचिपूर्णं वाक्यमेव काव्यमिति स्वीकरोति, तद्यथा काव्यमीमांसायाम्— ‘गुणवदलङ्कृतं वाक्यमेव काव्यम् ।’ तच्च रसप्रधानमेव भवितुं शक्नोति, तत्रैव- ‘अस्तु नाम निःसीमोऽर्थसार्थः, किन्तु रसवत एव निबन्धो युक्तो न नीरसस्य’ ।

\customparagraph{\textbf{काव्यहेतुः}}

शक्तिरेव काव्यस्य कारणम्, प्रतिभाव्युत्पत्ती तु शक्तेः कारणत्वेन मन्येते, तद्यथा काव्यमीमांसायाम्— \textbf{‘सा (शक्तिः) केवलं काव्ये हेतुः' इति यायावरीयः । शक्तिकर्तृके हि प्रतिभाव्युत्पत्तिकर्मणी । शक्तस्य प्रतिभाति, शक्तश्च व्युत्पद्यते । या शब्दग्राममर्थसार्थमलङ्कार- तन्त्रमुक्तिमार्गमन्यदपि तथाविधमधिहृदयं प्रतिभासयति सा प्रतिभा । अप्रतिभस्य पदार्थसार्थः परोक्ष इव । (काव्यमीमांसा)}

\customparagraph{\textbf{काव्यप्रयोजनम्}}

अस्य मते इहलोकपरलोकेषु सर्वत्र चिरन्तनानन्दप्रदानमेव काव्यस्य प्रयोजनम् भवतीति काव्यमीमांसायां सरलमधुरं काव्यप्रयोजनमुल्लिखति ।

\customparagraph{\textbf{रसः}}

राजशेखरः काव्यस्यात्मरूपेण रसं स्वीकरोति । नीरसस्तु काव्यमेव न भवतीति स्वकीयं मतं काव्यमीमांसायां प्रस्तौति यत्—

‘उक्तिचणं च ते वचो रस आत्मा’ ‘रसवत एव निबन्धयुक्तो न नीरसस्य’ इति । तथा च सामर्थ्ययुक्तः कविरेव रसवत् काव्यं करोति, असमर्थस्तु नीरसमेव, अतः सरसकाव्यार्थं कवेर्विलक्षणप्रतिभा एव आवश्यकी भवतीति कथयति, यत्—

\begin{shloka}[center]

\textbf{अस्तिचानुभूयमानो रसस्यानुगुणो विगुणश्चार्थः,}

\textbf{काव्ये तु कविवचनानि रसयन्ति विरसयन्ति च । (काव्यमीमांसा)}

\end{shloka}

एतादृशैः काव्यचिन्तनै राजशेखरो रसतत्त्वसमर्थक आचार्योऽस्तीति मन्तुं शक्यते ।

\customparagraph{\textbf{अलङ्कारः}}

अयमाचार्योऽलङ्कारं काव्यस्य आभूषणरूपं स्वीकरोति । तथा च वेदादीनां क्लिष्टत्वपरिहरणार्थम् अलङ्कारस्यावश्यकतां विचार्य वेदस्य सप्तममङ्गमिति घोषयति च । काव्यपुरुषमुपामादयोऽलङ्कारा अलङ्‌कुर्वन्तीति काव्यमीमांसायामुल्लिखति राजशेखरः —

\begin{shloka}[center]

\textbf{‘अनुप्रासोपमादयश्च त्वामलङ्कुर्वन्ति ।’}

\textbf{‘उपकारकत्वादलङ्कारः सप्तममङ्गम् ।’ (काव्यमीमांसा)}

\end{shloka}

\customparagraph{\textbf{रीतिः}}

राजशेखरः काव्ये वचनविन्यासक्रमरूपेण रीतितत्त्वं स्वीकरोति । अयं वामनाचार्यस्य रीतिसम्बन्धिविचाराधारे स्वमतं प्रस्तौति । स च त्रिविधा रीतीः स्वीकरोति, यथा—

‘वचनविन्यासक्रमो रीतिः ।’ सा च वैदर्भी, गौडीया, पाञ्चाली चेति तिस्रः ।’ एषु समासवदनुप्रासवद्योगवृत्तिपरम्परागर्भं जगाद सा गौडीया रीतिः । ईषदसमासमीषदनुप्रास- मुपचारगर्भञ्च जगाद सा पाञ्चाली रीतिः । स्थानानुप्रासवदसमासं योगवृत्तिगर्भं च जगाद सा वैदर्भी रीतिः ।

एतादृशैर्विवेचनैः काव्यमीमांसाया लेखको राजशेखरः संस्कृतकाव्यशास्त्रस्य कृते महत्त्वपूर्णयोगदानकर्ता अस्ति । पौरस्त्यकाव्यशास्त्रेऽस्य नाम ससम्मानेन गृह्यते ।

आचार्यराजशेखरः संस्कृतसाहित्यस्य काव्यशास्त्रस्य च क्षेत्रे देदीप्यमानो नक्षत्र आसीत् । दशमशताब्द्यां कान्यकुब्जदेशस्य प्रतिहारवंशीयनरेशस्य सभापण्डितोऽसौ महाभागो न केवलं विलक्षणः कविः, अपितु प्रकृतेर्भूमेश्च प्रकाण्डविद्वान् आसीत् । तस्य \textbf{‘काव्यमीमांसा’} इति कालजयी कृतिः काव्यशास्त्रीयपरम्पराया महत्त्वपूर्णनिधिरर्वर्तते, यत्र कवीरहस्यकाव्यस्वरूपप्रयोजनादीनां  नवीनशैल्या मौलिकं विवेचनं कृतमस्ति ।

काव्यतत्त्वानां निरूपणे राजशेखरेण प्रतिपादितं ‘गुणवदलङ्कृतं वाक्यमेव काव्यम्’ इति लक्षणम्, ‘शक्तिरेव काव्यहेतुः’ इति काव्यहेतुप्रतिपादनम्, ‘काव्यस्यात्मा रसः’ इति स्वीकारः, अलङ्कारो हि वेदस्य सप्तमाङ्गम् इति उद्घोषः, एवमेव वैदर्भी-गौडीया-पाञ्चालीति त्रिविधरीतीनां वचनविन्यासरूपेण वर्गीकरणञ्च राजशेखरस्य शास्त्रप्राविण्यं परिचाययन्ति । अत एव, रसतत्त्वसमर्थकत्वेन पौरस्त्यकाव्यशास्त्रस्य इतिहासे आचार्यराजशेखरस्य नाम सदा आदरेण गौरवेण च स्मर्यते ।

\specialupakhanda{\textbf{आचार्यो धनञ्जयः}}

काव्यशास्त्रपरम्परायां धनञ्जयस्य प्रवेशस्तत्र मौलिकचिन्तनपुनरावर्तनं द्योतयति । तस्य दशरूपकं यद्यपि नाट्यशास्त्रसंबद्धम्, तथापि रसविवेचने मौलिकं चिन्तनं प्रस्तौति । अयं काव्यशास्त्री विष्णुनाम्नः पितुः सुतः, धनिकस्याग्रजोऽस्ति । स हि मुञ्जस्य समकालिकः । तथ्यमिदं स्वयमेव दशरूपके उल्लिखति, यथा–

\begin{shloka}[center]

\textbf{विष्णोः सुतेनापि धनञ्जयेन विद्वन्मनोरागनिबन्धहेतुः ।}

\textbf{आविष्कृतं मुञ्जमहीशगोष्ठीवैदग्ध्यभाजा दशरूपमेतत्॥ (दशरूपकम्)}

\end{shloka}

मुञ्जश्च भोजस्य पितृव्यः १०३१—१०५१ वैक्रमाब्दे स्थितिमान् । अतो धनञ्जयस्यापि समयः स एव मन्यते ।

\customparagraph{\textbf{काव्यशास्त्रीयं वैशिष्ट्यम्}}

अस्य शास्त्रीयग्रन्थो दशरूपकमेव । तत्र ३०० कारिकाः सन्ति । ते च चतुर्षु प्रकाशेषु विभक्ताः सन्ति । प्रथमे प्रकाशे ग्रन्थप्रयोजनम्, नाट्यलक्षणम्, पञ्चसन्धयः, अर्थोपक्षेपकाश्च सप्रपञ्चं निरूपिताः। द्वितीये—नायिकानायकभेदवृत्त्यादीनां निरूपणम्, तृतीये—नाटकलक्षणम्, तत्स्वरूपञ्च सविस्तारं निरूपणम्, चतुर्थे प्रकाशे—रसनिरूपणं विद्यते । अस्य मते रसः सामाजिके एव उत्पद्यते, यथा–

\begin{shloka}[center]

\textbf{रस स एव स्वाद्यत्वाद्रसिकस्यैव वर्तनात् ।}

\textbf{नानुकार्यस्य वृत्तत्वात् काव्यस्यातत्परत्वतः ।। (दशरूपकम्)}

\end{shloka}

स च विभावादिभिरास्वाद्यमानीयते । रामादिगतधीरोदात्ताद्यवस्थाः परित्यक्तविशेषाः सत्यो रसहेतवो भवन्ति । काव्यर्थभेदादात्मानन्दसमुद्भवः स्वादः । स चेतसो विकासविस्तारक्षोभविक्षेपरूपेण चतुर्विधः । इत्यनेन प्रकारेण रसस्यास्वादविषये मौलिको विचारो प्रस्तुतोऽनेन । अन्ये विषयास्तु नाट्यशास्त्रस्य व्याख्यात्मकाः सन्ति ।

\specialupakhanda{\textbf{आचार्यो राजनकः कुन्तकः}}

आचार्योऽयं वक्रोक्तिसम्प्रदायस्य संस्थापकोऽस्ति । अस्य समयविषयेऽपि सङ्केतादिकं दुर्लभमेव । तथापि कुन्तक आनन्दवर्धनस्य ध्वनिप्रस्थानेन सह सुपरिचित आसीत् । असौ आनन्दवर्धनस्य ‘स्वेच्छाकेसरिण—’ इति श्लोकम्, तथा रुद्रटस्य राजशेखरस्य चाधिकानि पद्यानि स्ववक्रोक्तिजीवितग्रन्थे प्रस्तुतवान् । तस्माद्  अस्य समयो दशमशताब्दस्य उत्तरार्धो भवितुमर्हति । व्यक्तिविवेकस्य रचयिता महिमभट्टः कुन्तकस्य सिद्धान्तान् अखण्डयत् । एतस्य समय एकादशशताब्दस्य उत्तरार्धो मन्यते । ध्वन्यालोकलोचनकारोऽभिनवगुप्तः कुन्तकस्य कुत्रापि उल्लेखं नाकरोत् । अतः कुन्तकाभिनवगुप्तौ समकालिकौ आस्ताम् । अस्य समयो दशमशताब्दस्योत्तरार्धं  मन्तुं युज्यते ।

\customparagraph{\textbf{काव्यशास्त्रीयं वैशिष्ट्यम्}}

अस्योपलब्ध एक एव ग्रन्थोऽस्ति वक्रोक्तिजीवितम् । लाक्षणिकग्रन्थोऽयं चतुर्षून्मेषेषु  विभक्तोऽस्ति । अत्र कारिका वृत्तिरुदाहरणञ्चेति त्रयो भागाः सन्ति । कारिकावृत्तिभागौ कुन्तकस्य रचने, परमुदाहरणान्यन्यतः समाहृतानि सन्ति । वक्रोक्तिकाव्यजीविते चत्वार उन्मेषाः सन्ति । कुन्तकस्तत्र प्रथमे उन्मेषे काव्यप्रयोजनम्, साहित्यकल्पनायाः, वक्रोक्तेर्लक्षणस्य च प्रतिपादनम्, वक्रोक्तेः षण्णां प्रकाराणां नामनिर्देशं चकार । द्वितीये– वक्रोक्तिभेदेषु वर्णविन्यासवक्रता, पदपूर्वार्धवक्रता, प्रत्ययवक्रता इति त्रयाणां भेदानां विवेचनं कृतवान् । तृतीये– वाक्यवक्रताया सविस्तारेण वर्णनमकरोत् । वाक्यवक्रतायामेव सर्वेषामलङ्काराणामन्तर्भावं चकार । चतुर्थे च– प्रकरणवक्रतायाः प्रबन्धवक्रतायाश्च विशदं वर्णनमकरोत् ।

वक्रोक्तितत्त्वस्य कल्पना भामहेन कृता । तस्यैवोल्लेखस्य प्रेरणया असौ वक्रोक्तिं काव्यस्यात्मत्वेन स्वीकृतवान्  । अस्य वक्रोक्तिर् विशिष्टस्वरूपा विद्यते, यथा— ‘वक्रोक्तिरेव वैदग्ध्यभङ्गीभणितिरुच्यते’ इति । अस्यैव सूत्रस्य व्याख्यारूपेण अयमेव प्रस्तौति यत्— वक्रोक्तिः प्रसिद्धाभिधानव्यतिरेकिणी विचित्रैवाभिधा । कीदृशी भङ्गीभणितिः ? वैदग्ध्यं विदग्ध्यभावः कविकर्मकौशलम्, तस्य भङ्गी विच्छित्तिः, तया भणिति विचित्रैवाभिधा वक्रोक्तिरित्युच्यते । तस्मात् वक्रोक्तिरेव काव्यस्यात्मा । अनेन विना काव्ये प्राणसञ्चारोऽपि भवितुं नार्हति ।  सा च षड्विधा— वर्णविन्यासवक्रता–पदपूर्वार्द्धवक्रता–प्रत्ययाश्रयवक्रता–वाक्यवक्रता–प्रकरणवक्रता–प्रवन्धवक्रताभेदात् ।

\customparagraph{\textbf{काव्यलक्षणम्}}

कुन्तकः काव्यलक्षणमपि स्वप्रतिपादितस्य वक्रोक्तितत्त्वस्याधारे प्रस्तौति । वक्रतार्थप्रतिपादनसमर्थौ शब्दार्थौ काव्यमिति काव्यलक्षणम् उल्लिखति । तथा च काव्यस्यात्मस्थाने वक्रोक्तिरेव भवितुं शक्नोति । सामान्यार्थप्रतिपादकौ शब्दार्थौ तु काव्यरूपं न जायेते । तत्तु इतिहासादिवाङ्मयभेदे समागच्छति । तद्यथा वक्रोक्तिजीविते—

\begin{shloka}[center]

\textbf{शब्दार्थौ सहितौ वक्रकविव्यापारशालिनि ।}

\textbf{बन्धे व्यवस्थितौ तद्विदाह्लादकारिणी ।।}

\textbf{वक्रोक्तिः काव्यजीवितम्, इति च (वक्रोक्तिजीवितम्)}

\end{shloka}

\customparagraph{\textbf{काव्यहेतुः}}

काव्यस्य कारणविषये कुन्तकः साङ्केतिकरूपेण हेतुत्रयं स्वीकरोति । वक्रोक्तिभेदवर्णनप्रसङ्गे तेषां सम्पादकतत्त्वकथने उल्लिखति यत्—

\textbf{‘तच्च न प्रतिभासंरम्भमात्रसाध्यम्, किन्तु तद्विहितसमस्तसामग्रीसम्पाद्यम्’} इति ।

अपि च

\begin{shloka}[center]

\textbf{सर्वसम्पत्परिस्पन्दः सम्पाद्यं सरसात्मनाम् ।}

\textbf{अलौकिकचमत्कारकारि काव्यैकजीवितम् ।।  (वक्रोक्तिजीवितम्)}

\end{shloka}

अत्र प्रस्तुता प्रतिभा एव  न, अपितु समस्तसामग्रीसम्पाद्यमिति कथनेन व्युत्पत्तिः, अभ्यासश्च बोध्यते । तस्मात् कुन्तकोऽपि हेतुत्रयपक्षपाती एव ।

\customparagraph{\textbf{काव्यप्रयोजनम्}}

कुन्तकः काव्यस्य प्रयोजनविषये सविस्तारं वर्णयति । स सुकुमारक्रमेण धर्मादिचतुर्वर्गफलप्राप्तिः काव्यादेव भवतीति कथयन् व्यवहारज्ञानं, मानसिकानन्दप्राप्तिश्च काव्यादेव जायते, इति स्वीकरोति यद्  वक्रोक्तिजीविते—

\begin{shloka}[center]

\textbf{धर्मादिसाधनोपायः सुकुमारक्रमोदितः ।}

\textbf{काव्यबन्धोऽभिजातानां हृदयाह्लादकारकः ।।}

\textbf{व्यवहारपरिस्पन्द सौन्दर्यव्यवहारिभिः ।}

\textbf{सत्काव्याधिगमादेव नूतनौचित्यमाप्यते ।।}

\textbf{चतुर्वर्गफलास्वादमप्यतिक्रम्य तद्विदाम् ।}

\textbf{काव्यामृतरसेनान्तश्चमत्कारो वितन्यते ।।  (वक्रोक्तिजीवितम्)}

\end{shloka}

एतादृशोऽस्य काव्यप्रयोजनविचारोऽन्येषु शास्त्रीयविचारेषु समन्वयशीलो दृश्यते ।

\customparagraph{\textbf{अलङ्कारमतम्}}

अन्येषां विदुषामिव कुन्तकस्यालङ्कारविचारो नैव दृश्यते । अयं तु शब्दार्थौ अलङ्कार्यौ, एतयोरलङ्कारस्तु वक्रोक्तिरेव इति कथयित्वा एकमेव विशिष्टालङ्कारं स्वीकरोति । इतरेषामाचार्याणां मतानुरूपमलङ्कारतत्त्वं तु कुन्तको वाक्यवक्रोक्तिभेदेऽन्तर्भावयति । तद्यथा—

\begin{shloka}[center]

\textbf{उभावेतावलङ्कार्यौ तयोर्पुनरलङ्कृतिः ।}

\textbf{वक्रोक्तिरेव वैदग्ध्यभङ्गीभणितिरुच्यते ।।}

\textbf{वाक्यस्य वक्रभावोऽन्यो भिद्यते य सहस्रधा ।}

\textbf{यत्रालङ्कारवर्गोऽसौ सर्वेऽप्यन्तर्भविष्यति ।।  (वक्रोक्तिजीवितम्)}

\end{shloka}

\customparagraph{\textbf{रीतिविवेचनम्}}

आचार्यः कुन्तको रीतितत्त्वविषये प्रचलितामवधारणां प्रकारान्तरेण स्वीकरोति । स तु मार्गः’ इति नाम्ना रीतितत्त्वं वर्णयति । अयं  त्रिविधं मार्गमुल्लिखति । तद्यथा—

\begin{shloka}[center]

\textbf{सम्प्रति तत्र ये मार्गाः कविप्रस्थानहेतवः ।}

\textbf{सुकुमारो विचित्रश्च मध्यमश्चोभयात्मकः ॥  (वक्रोक्तिजीवितम्)}

\end{shloka}

\customparagraph{\textbf{वक्रोक्तिसम्प्रदायसंस्थापकः}}

‘वक्रोक्तिः’ इति शब्दस्य प्रयोगस्तु पौरस्त्यकाव्यशास्त्रे प्राचीनकालादेव प्रचलित आसीत् । भामहदण्डीवामनरुद्रटादयो विद्वांसोऽस्य विषये विचारान् प्रस्तुवन्ति, तथापि काव्यस्यानिर्वार्यतत्त्वरूपेण तु नैव प्रयुक्तवन्तः । एतेषां प्राचीनानां सङ्केतरूपान् विचारान् विमृश्य कुन्तकः पृथक् सम्प्रदायरूपेण वक्रोक्तिं प्रस्तौति । एषः प्राग्जातैराचार्यैः सामान्यरूपेण गृहीतं वक्रोक्तिं वैदग्ध्यभङ्गीभणितिर्वक्रोक्तिरिति विशिष्टस्वरूपप्रदर्शनपूर्वकं काव्यस्य जीविततत्त्वम् अर्थात् आत्मतत्त्वरूपेण प्रस्थापयति, यथा–

\textbf{शब्दार्थौ सहितौ वक्रकविव्यापारशालिनि ।}

\textbf{बन्धे व्यवस्थितौ तद्विदाह्लादकारिणी ।।}

\textbf{वक्रोक्तिः काव्यजीवितम्, इति च  (वक्रोक्तिजीवितम्)}

इदं तत्त्वं स्वकीयविचारानुरूपं काव्यस्य आत्मस्थाने स्थापयितुं कुन्तकः स्वस्य लाक्षणिकग्रन्थस्य नाम अपि ‘वक्रोक्तिकाव्यजीवितम्’ इति विदधाति । कुन्तकस्य वक्रोक्तेः कल्पना तादृशी व्यापका विद्यते, यत्र पूर्ववर्त्तिभिः काव्यशास्त्रिभिः प्रतिपादितानि सर्वाण्यलङ्कारध्वनिप्रभृतीनि काव्यतत्त्वानि समाविशन्ति । रसध्वनिविषये पूर्णतः परिचितोऽपि कुन्तकः स्वकीयप्रतिभाशक्तिबलात् तान् सर्वान् वक्रोक्तौ समाविश्य वक्रोक्तिरेव काव्यस्य जीवितमिति प्रमाणीकरोति । अस्य विश्लेषणात्मिका प्रतिभाशक्तिरतीव विलक्षणा वर्तते, तस्मादेवायं वक्रोक्तितत्त्वं काव्यस्य आत्मस्थानीयतत्त्वस्य कोटौ स्थापयितुं समर्थो जातः ।

अनेन प्रकारेण कुन्तकः स्वप्रतिभाशक्त्या प्रबलतार्किकविचारपरिपाकेन प्राग्जातैः काव्यशास्त्रिभिरलङ्कारभेदरूपेण विवेचितं वक्रोक्तितत्त्वं काव्यात्मस्थानीयमुद्घोषयति । पौरस्त्यकाव्यशास्त्रीयविकासपरम्परायां कुन्तकस्येतादृशं योगदानं सर्वैः स्मरणीयमभिनन्दनीयञ्च विद्यते ।

‘राजनकः’  इत्युपाधिभूषितः कुन्तकः संस्कृतकाव्यशास्त्रस्य इतिहासे वक्रोक्तिसम्प्रदायस्य प्रतिष्ठापकत्वेन क्रान्तिकारी विद्वान् मन्यते। दशमशताब्द्या उत्तरार्धे स्थितोऽसौ महाभागः स्वकीये ‘वक्रोक्तिजीवितम्’ इति चतुर्षून्मेषेषु विभक्ते लाक्षणिकग्रन्थे पूर्ववर्तिनां रस-ध्वनि-अलङ्कारादीनां सर्वेषां काव्यतत्त्वानां समन्वयं कृत्वा ‘वक्रोक्तिः काव्यजीवितम्’ इति नवीनं सिद्धान्तं प्रस्थापितवान् । तस्य मते कविकर्मकौशलपूर्णा विचित्रैवाभिधा अर्थात् ‘वैदग्ध्यभङ्गीभणितिः’ एव वक्रोक्तिरस्ति, या वर्णविन्यासादिषड्विधभेदेषु व्याप्ता सति काव्ये प्राणसञ्चारं करोति । कुन्तकेन काव्यलक्षणे शब्दार्थयोः वक्रकविव्यापारशालित्वं स्वीकृतम्, काव्यहेतौ प्रतिभा-व्युत्पत्ति-अभ्यासानां त्रित्वं समर्थितम्, तथा च काव्यप्रयोजने आह्लादप्रदानपूर्वकं चतुर्वर्गफलप्राप्तिरङ्गीकृता । तेन पारम्परिकम् अलङ्कारतत्त्वं वाक्यवक्रतायाम् अन्तर्भावितम् । रीतितत्त्वं ‘मार्गः’ (सुकुमारमार्गः, विचित्रमार्गः, मध्यममार्गः) नाम्ना नव्यरीत्या निरूपितम् । अतः स्वकीयया विलक्षणया विश्लेषणात्मकप्रतिभया काव्यशास्त्रं समृद्धं कुर्वाणः कुन्तकः पौरस्त्यकाव्यशास्त्रपरम्परायां सदा स्मरणीयो वर्तते ।

\specialupakhanda{\textbf{आचार्य महिमभट्टः}}

संस्कृतकाव्यशास्त्रेतिहासे महिमभट्टो ध्वनिविरोधिनि  प्रस्थाने समायाति । अस्य समयोऽपि पूर्वाचार्यवद्  अनुमानेनैव ज्ञायते । अस्य व्यक्तिविवेकाख्यो लाक्षणिकग्रन्थो विद्यते । आचार्यो महिमभट्टः स्वकीये व्यक्तिविवेके– `काव्यकाञ्चनकथाश्ममानिना कुन्तकेन निजकाव्यलक्ष्मणि'  इति कुन्तकस्य नामोल्लखं करोति । अत आचार्योऽयं कुन्तकात् परवर्त्तीति स्वीक्रियते । ईशवीयद्वादशशताब्दस्य रुय्यकोऽलङ्कारसर्वस्वेऽस्य मतस्योल्लेखमकार्षीत् । तथा चायं ध्वन्यालोकलोचनस्य सिद्धान्तानां खण्डनं चकार । अतोऽस्य समय एकादशशताब्दस्य मध्यकालं  स्वीकर्तुं शक्यते ।

\customparagraph{\textbf{काव्यशास्त्रीयं वैशिष्ट्यम्}}

अस्यैक एव व्यक्तिविवेकाभिधानो ग्रन्थः समुपलभ्यते । अत्र ‘व्यक्तिः’ शब्देन व्यञ्जनाशक्तिर्बोध्यते, ‘विवेकः’ इत्यनेन च समीक्षणम् ज्ञायते । अर्थात् व्यञ्जनायाः समीक्षणमस्य ग्रन्थस्य नामार्थो भवति । व्यञ्जनाया ध्वनेश्च खण्डनं कृत्वानुमानेऽन्तर्भावयितुं ग्रन्थोऽयं रचित इति’ स एव व्यक्तिविवेके समुल्लिखति, यत्—

\begin{shloka}[center]

\textbf{अनुमानान्तर्भावं सर्वस्यैव ध्वनेः प्रकाशयितुम् ।}

\textbf{व्यक्तिविवेकं कुरुते प्रणम्य महिमापरां वाचाम् ।। (व्यक्तिविवेकः)}

\end{shloka}

महिमभट्टो वस्तुतो नैयायिकः । न्यायशास्त्रे प्रमाणविषयेऽनुमानस्य विशदं विवेचनं कृतं विद्यते । महिमभट्टो व्यञ्जनाशक्तिं तत्रैवान्तर्भावं कल्पयति । प्रथमन्तु व्यञ्जनाशक्तिरभिधाशक्त्यनन्तरवर्त्ती विद्यते । अतो व्यञ्जनाशक्ति वाच्यार्थगामी एव, अपरम् अर्थेनार्थान्तरप्रतीतिरूपत्वाद् अनुमानप्रक्रियासादृश्यं पुनश्च सहेतुकप्रतीतित्वादस्यानुमानेऽन्तर्भावं नातीव दुष्करमिति तीव्रमालोच्य महिमभट्टो ध्वनिमार्गनिरसने प्रवृत्तो बभूव । स्वग्रन्थेऽस्य कल्पनाया विद्वत्तायाश्च विशदः प्रकाशो लभ्यते ।

व्यक्तिविवेके त्रयो विमर्शाः सन्ति । असौ प्रथमे विमर्शे— ध्वनिमनुमानेऽन्तर्भाव्य ध्वनौ दोषान् प्रादर्शयत् । अनुमानेन समेषां ध्वनिभेदानां विशदं विवरणमुपस्थाप्य स्वकीयं प्रौढं पाण्डित्यं प्रदर्शयामास ।  द्वितीये विमर्शे— अनौचित्यस्य मुख्यं कारणं काव्यदोषं प्रतिपाद्य तस्य विभिन्नानां प्रकाराणां विस्तरेण विवेचनमुपस्थापितमस्ति । तृतीये विमर्शे— महिमभट्टो पूर्वाचार्यैः काव्यात्मतत्त्वरूपेण स्थापिते ध्वनितत्त्वे तीव्रं प्रहारं करोति । चत्वारिंशत् ध्वनेरुदाहरणानि आदाय अनेन साधितं यत्— इमानि सर्वाणि ध्वनेरेव प्रकारभूतानि सन्तीति । ग्रन्थेऽस्मिन् सः प्राचीनपरम्परां विहाय नूतनविचारप्रतिपादने प्रयत्नमकरोत् । यद्यपि तत्परवर्त्तिनो विद्वांसोऽस्य ध्वनिविरोधिविचारस्य खण्डनं चक्रुः । असौ अभिधालक्षणाव्यञ्जनाशक्तिषु काव्येऽभिधाशक्तिरेव केवलं वर्तते, तयैव सर्वेषामर्थानां प्रतीतिर्जायते’ इति स्वीकरोति ।

\customparagraph{\textbf{काव्यलक्षणम्}}

\textbf{विभावादिसंयोजनात्मा रसाभिव्यक्त्यव्यभिचारिकविव्यापारकाव्यम् । (व्यक्तिविवेकः)}

महिमभट्टो विशिष्टप्रतिभासम्पन्नः काव्यशास्त्री आसीत् । अनेन स्वपूर्ववर्त्तिनामाचार्याणां ध्वनिविषयकमतानां तीव्रप्रहारपूर्वकं खण्डनमकरोत्, परं ध्वनितत्त्वस्य लोकप्रियतायाः कारणेन अस्य मतस्य मान्यता दीर्घकालिका न जाता । तथापि अस्य तार्किकविवेचनशैली पौरस्त्यकाव्यशास्त्रजगति प्रशंसिताऽस्ति ।

\specialupakhanda{\textbf{आचार्यः क्षेमेन्द्रः}}

संस्कृतकाव्यशास्त्रे प्रचलितानां सम्प्रदायानां चरमसम्प्रदायरूपेण क्षेमेन्द्रस्य औचित्यसम्प्रदायः समायाति । अनेन प्राचीनकाव्यशास्त्रिभिः सामान्यरूपेण उल्लिखितम् औचित्यतत्त्वं विशिष्टविवेचनपुरस्सरं सम्प्रदायरूपेण स्थापितम् । विविधविद्यासु पारङ्गतः चत्वरिंशदधिकग्रन्थानां च प्रणेता, काश्मीरदेशीयस्य सिन्धोः पौत्रः प्रकाशेन्द्रस्य चात्मज आसीत् । स च काश्मीरस्यानन्तराजस्य समकालिक आसीद्  इति तस्यैव औचित्यविचारचर्चाया अन्तिमश्लोके ज्ञायते यत्— ‘तस्य श्रीमदनन्तराजनृपतेः काले किलायं कृतः’ । तथा च कविकण्ठाभरणस्यान्तिमश्लोकेऽपि ‘राज्ये श्रीमदनन्तराजनृपतेः काव्योदयोऽयम्’ इति । अनन्तराजो हि १०८५—११२० मितवैक्रमाब्दवर्षेषु काश्मीरभुवं शशासेति राजतरङ्गिणीतो ज्ञायते । तस्य कलशाख्यः पुत्रः ११२० विक्रमाब्दादेव राज्येऽधिकृत आसीत् । यद्यपि अनन्तराजः ११६८ मितवैक्रमाब्दपर्यन्तं जीवित आसीत् । अतः क्षेमेन्द्रस्यापि  स्थितिः कालश्च तदा एव इति, अर्थात् अस्य समय एकादशशतकस्योत्तरार्धः स्वीक्रियते ।

\customparagraph{\textbf{काव्यशास्त्रीयं योगदानम्}}

क्षेमेन्द्राचार्येण रचिता नैका ग्रन्थाः लभ्यन्ते । तेषु औचित्यविचारचर्चा, सुवृत्ततिलकम्, कविकण्ठाभरणञ्चेति त्रयो ग्रन्थाः  काव्यशास्त्रसम्बद्धा दृश्यन्ते । औचित्यविचारचर्चा अस्य सुप्रसिद्धा कृतिरस्ति । तत्रैवानेन औचित्यतत्त्वं कव्यात्मस्थानीभूतत्त्वेन विवेचितम् । यस्मिन् १९ कारिका, तासां वृत्तिः उदाहरणानि च सन्ति । तत्र पदवाक्य–प्रबन्ध–अर्थ–गुण–अलङ्कार–रस–क्रिया–कारक–लिङ्ग– विशेषण–उपसर्ग–निपात–काल–देश–कुल–व्रत–तत्त्व–सत्त्व–अभिप्राय–स्वभाव–सार–सङ्ग्रह–प्रतिभा–अवस्था–विचार–नाम–आशीः प्रभृतिषु काव्याङ्गेषु औचित्यं निर्दिष्टमस्ति । यद्यपि औचित्यं तत्पूर्ववर्त्तिभिराचार्यैरपि निरौचित्यं काव्यं न भवतीति कथितमेव, परं तस्य काव्यजीवितत्वेन तु क्षेमेन्द्रेण निर्दिष्टमिति क्षेमेन्द्र औचित्यप्रस्थानस्य प्रवर्तक इति संस्कृतकाव्यशास्त्रे प्रसिद्धिं प्राप्नोति । कविकण्ठाभरणं कविशिक्षाविषयको ग्रन्थो वर्तते । ग्रन्थेऽस्मिन् कवित्वप्राप्ति–कविशिक्षाकथन– चमत्कारकथन–गुणदोषविभाग–परिचयप्राप्ताख्याः पञ्चसन्धयः सन्ति ।  सुवृत्ततिलके वृत्तावचय– गुणदोषकथन– वृत्तविनियोगाख्यास्त्रयो विन्यासाः सन्ति ।

\customparagraph{\textbf{औचित्यसम्प्रदायः}}

क्षेमेन्द्र औचित्यस्य महत्त्वं विवेचयन् कथयति यत्— यथा लोकव्यवहारे जन औचित्ययुक्तेन व्यवहारेण महत्प्रतिष्ठां प्राप्नोति तथैव व्यवहारस्यानौचित्येन हास्यास्पदो जायते । तद्वत् काव्येऽपि रसालङ्कारगुणरीत्यादीनां तत्त्वानामौचित्यपूर्णविनियोजनेनैव काव्यसौन्दर्यं प्रकाशते । ओैचित्याभावे तु कवि हास्यास्पदत्वं प्राप्नोति । शरीरस्यालङ्करणाय स्थानौचित्यं प्रदर्शयन् क्षेमेन्द्र काव्यौचित्यस्य महत्त्वं प्रस्तौति, यथा—

\begin{itemize}
\item 
\end{itemize}
\begin{shloka}[center]

\textit{कण्ठे मेखलया नितम्बफलके तारेण हारेण वा}*

\textbf{पाणौ नूपुरबन्धनेन चरणे केयूरपाशेन वा ।}

\textbf{शौर्येण प्रणते रिपौ करुणया नायान्ति के हास्यताम्}

*\textit{औचित्येन विना रुचिं प्रतनुते नालङ्कृतिर्नो गुणः ।। (औचित्यविचारचर्चा)}

\end{shloka}

\begin{itemize}
\item 
\end{itemize}
काव्यस्य सर्वेषु अङ्गेषु औचित्यतत्त्वस्य व्यापकत्वं प्रदर्श्य क्षेमेन्द्र औचित्यमेव काव्यात्मत्वं स्वीकरोति । औचित्यं काव्यात्मरूपेण स्वीकरणार्थम् अस्य युक्तिरस्ति यत्— रसे रसत्वम्, गुणे गुणत्वम्, अलङ्कारे अलङ्कारत्वम्, रीतौ रीतित्वं तदा एव भवति यदा एतेषां विधानं औचित्यपूर्णं स्यात्।

क्षेमेन्द्रात् प्राक् जातानां काव्यशास्त्रीणां कृतिषु अपि औचित्यस्य चर्चा विद्यते । ध्वन्यालोककार आनन्दवर्धनस्तु औचित्यविषये विस्तारपूर्विकां चर्चां करोति । परमेते विद्वांसः काव्यात्मतत्त्वरूपेण तु न स्वीकुर्वन्ति । तथापि आनन्दवर्धनस्य औचित्यविषयकसमीक्षा अस्य कृते आधारस्तम्भोऽभवदिति कथयितुं शक्यते ।

औचित्यस्य लक्षणं प्रस्तौति क्षेमेन्द्रो यथा—

\begin{shloka}[center]

\textbf{उचितं प्राहुराचार्याः सदृशं किल तस्य तत् ।}

\textbf{उचितस्य च यो भावस्तदौचित्यं प्रचक्षते ।।  (औचित्यविचारचर्चा)}

\end{shloka}

अस्य मते औचित्यं तु सर्वेषां कार्याणां कृतेऽनिवार्यविषयमस्ति । येन विना काव्यमेव न किमपि कार्यमुद्देश्यानुरूपं न भवति । अत औचित्यं काव्यस्य कृते आत्मस्थानीयम्, बाह्यम्, स्थिरम्, अस्थिरम्, इत्यादि चर्चां विनैव स्वीकरणीयं भवति ।

\customparagraph{\textbf{काव्यलक्षणम्}}

अयमाचार्योऽलङ्कारैरलङ्कृतं शब्दार्थसाहित्यमेव काव्यं स्वीकरोति । तच्च रससिद्धम्, औचित्यतत्त्वमेव आत्मारूपेणावस्थितं भवितुमर्हति । एतत् मतं क्षेमेन्द्रः प्रस्तौति स्वग्रन्थे यथा—

\begin{shloka}[center]

\textbf{काव्यं विशिष्टशब्दार्थसाहित्यसदलङ्कृतिः । (सरस्वतीकण्ठाभरणे)}

\textbf{औचित्यं रससिद्धस्य स्थिरं काव्यस्य जीवितम् । (औचित्यविचारचर्चा)}

\end{shloka}

एकादशशतकस्योत्तरार्धे काश्मीरभूमौ प्रादुर्भूतः आचार्यः क्षेमेन्द्रः संस्कृतकाव्यशास्त्रस्य औचित्यसम्प्रदायस्य प्रवर्त्तकत्वेन अद्वितीयं स्थानं भजते। बहुविद्यापारङ्गत एषश्चत्वरिंशदधिकग्रन्थानां निर्माता आसीत्, येषु मुख्यतया \textit{औचित्यविचारचर्चा}, \textit{कविकण्ठाभरणम्}, तथा \textit{सुवृत्ततिलकम्} इति त्रयो ग्रन्थाः काव्यशास्त्रीययोगदानदृष्ट्या अत्यन्तं महत्त्वपूर्णाः सन्ति। पूर्ववर्तिभिराचार्यैः सामान्यरूपेण चर्चितम् औचित्यतत्त्वं क्षेमेन्द्रेणैव काव्यस्य परमप्राणत्वेन अर्थात् `काव्यस्य स्थिरं जीवितम्' इति रूपेण संस्थाप्य विशिष्टसम्प्रदायस्य गौरवम् अदायि। तस्य मते ‘औचित्याभावे रसालङ्कारगुणादीनां काव्यतत्त्वानां विनियोजनं तथा एव निष्फलं हास्यास्पदञ्च भवति यथा शरीरे आभूषणानाम् अनुपयुक्ततया धारणम्।’ अतः काव्यस्य पद-वाक्य-प्रबन्धादिषु विविधेषु अङ्गेषु औचित्यस्य व्यापकतां प्रदर्श्य तं रससिद्धस्य काव्यस्य आत्मत्वेन प्रतिष्ठाप्य क्षेमेन्द्रः संस्कृतकाव्यशास्त्रे चरमसम्प्रदायस्य सुदृढं सोपानं निर्मितवान्, एतदेव तस्य चिरस्मरणीयं श्रेष्ठञ्च योगदानं वर्तते।

\specialupakhanda{\textbf{आचार्यो भोजराजः}}

मालवदेशस्य अधिपतिर्भोजराजः संस्कृतकविपण्डितानां केवलम् आश्रय एव नासीत्, अपितु स स्वयमपि एको धुरन्धरः समालोचक आसीत् । अस्य समयो निश्चितप्रायो वर्तते । मुञ्जराजस्यानन्तरं भोजराजस्य पिता सिन्धुराजो राजा बभूव । भोजराजस्यैकस्य दानपत्रस्य कालः १०२१ ईसवीयो वर्तते। भोजराजस्योत्तराधिकारिणो जयसिंहस्य १०५५ तमस्य ईसवीयवर्षस्य शिलालेख उपलब्धोऽस्ति । अनेन प्रमाणितो भवति यद् – भोजराजस्यान्तिमः कालः १०५४ ईसवीयो भवितुमर्हति । तस्माद्  भोजराजस्याविर्भावः काल एकादशशताब्दस्य प्रथमार्धः स्वीकरणीयो दृश्यते ।

\customparagraph{\textbf{काव्यशास्त्रीयं योगदानम्}}

भोजराजस्य सरस्वतीकण्ठाभरणम्, शृङ्गारप्रकाशश्च काव्यशास्त्रसम्बद्धौ ग्रन्थौ स्तः । सरस्वतीकण्ठाभरणे पञ्च परिच्छेदाः सन्ति । तेषु प्रथमे परिच्छेदे— दोषगुणयोर्विवेचनम्, तत्र च पदगताः षोडश, वाक्यगताः षोडश, वाक्यार्थगता षोडशेति अष्टचत्वारिंशद्  दोषाः विवेचिताः । एवमेव शब्दगताश्चतुर्विंशतिरर्थगताश्चतुर्विंशतिरिति अष्टाचत्वारिंशत् गुणा निरूपिताः सन्ति । द्वितीये परिच्छेदे— चतुर्विंशतिः शब्दालङ्काराः, तृतीये— अर्थालङ्काराः, चतुर्थे उभयालङ्काराः, पञ्चमे च रस–भाव–सन्धिवृत्तीनां निरूपणं विद्यते ।

शृङ्गारप्रकाशश्चास्यापरो ग्रन्थः । अत्र षट्त्रिंशत् प्रकाशाः सन्ति । तेषु आद्याष्टप्रकाशेषु शब्दार्थविषयकवैयाकरणमतानि निरूपितानि । नवमे— गुणाः, दशमे— दोषाः, एकादशे— महाकाव्यम्, द्वादशे च नाटकं सप्रपञ्चं विवेचितमस्ति । शेषेषु चतुर्विंशतिप्रकाशेषु साङ्गोपाङ्गं रसतत्त्वं विवेचितम् । अत्र भोजस्य रसविषये विशिष्टमतं दृश्यते यद्- अयं शृङ्गारमेव सर्वरसेश्वरं स्वीकरोति, यथा—

\begin{shloka}[center]

\textbf{शृङ्गारवीरकरुणाद्भुतरौद्रहास्यवीभत्सवत्सलभयानकशान्तनाम्नः ।}

\textbf{आम्नासिषुर्दशरसान् सुधियो वयं तु शृङ्गारमेव रसनाद्रसमामनामः ।। (शृङ्गारप्रकाशः)}

\end{shloka}

परं तन्मते यः शृङ्गारः स न केवलं यूनोः परस्पराकाङ्क्षात्मकमपितु धर्मार्थकाममोक्षाख्यचतुर्विधभेदात्मकः । अस्यैव मतस्य समर्थनाय भोजराजः शृङ्गारप्रकाशाख्यं विपुलकायं ग्रन्थमरचयत् । काव्यशास्त्रेतिहासे ग्रन्थोऽयं सर्वाधिकदीर्घकायो विद्यते ।

\customparagraph{\textbf{काव्यलक्षणम्/काव्यप्रयोजनम्}}

आचार्यो भोजराजः काव्यस्य लक्षणं प्रयोजनञ्च एकत्रैवोल्लिखति । अयं दोषरहितं, गुणसहितम्, अलङ्कारैरलङ्कृतं, रसान्वितञ्च कवेः कर्म काव्यं स्वीकरोति । तस्मात् कीर्तिः, प्रीतिश्च प्राप्येते, तद्यथा स्वग्रन्थे—

\begin{shloka}[center]

\textbf{निर्दोषं गुणवत्काव्यमलङ्कारैरलङ्कृतम् ।}

\textbf{रसान्वितं कविः कुर्वन् कीर्तिं प्रीतिं च विन्दति ।। (शृङ्गारप्रकाशः)}

\end{shloka}

\customparagraph{\textbf{गुणः}}

भोजराजोऽलङ्कारादपि महत्त्वाधिकं गुणतत्त्वं स्वीकरोति । अस्य मते गुणालङ्कारयोर्मुख्यं तु गुणतत्त्वमेवास्ति, तद्यथा स्वग्रन्थे—

\begin{shloka}[center]

\textit{\textbf{अलङ्कृतमपि श्रव्यं न काव्यं गुणवर्जितम् ।}}

\textit{\textbf{गुणयोगस्तयोर्मुख्यं गुणालङ्कारयोगयोः ।। (शृङ्गारप्रकाशः)}}

\end{shloka}

एकादशशताब्दस्य प्रथमार्धे मालवदेशस्य शासकरूपेण अवस्थितो भोजराजः संस्कृतकवीनाम् आश्रयदाता, स्वयञ्च उत्कृष्टः काव्यशास्त्री च आसीत्। तस्य ‘सरस्वतीकण्ठाभरणम्’ तथा संस्कृतकाव्यशास्त्रस्य इतिहासस्य सर्वाधिकविशालकायः ग्रन्थः ‘शृङ्गारप्रकाशः’ इति द्वौ ग्रन्थौ काव्यशास्त्रीययोगदानदृष्ट्या अत्यन्तं महत्त्वपूर्णौ स्तः। भोजराजस्य रससम्बद्धः विचारो विशिष्टोऽस्ति, यत्र तेन शृङ्गारमेव एकमात्रं सर्वरसेश्वरम् अङ्गीकृतम् । तन्मते अयं शृङ्गारः केवलं रतिभावमात्रं न, अपितु धर्म-अर्थ-काम-मोक्षात्मकश्चतुर्विधः  पुरुषार्थप्रदायकोऽस्ति। तेन निर्दोषम्, गुणयुक्तम्, अलङ्कारैरलङ्कृतम्, रसान्वितञ्च कवेः कर्म काव्यमिति काव्यस्वरूपं स्वीकृतम्, यस्य प्रयोजनं कीर्तिः प्रीतिश्चास्ति। अलङ्कारापेक्षया गुणतत्त्वस्य प्राधान्यं स्वीकुर्वन् असौ गुणहीनं काव्यं श्रव्यमपि अग्राह्यमिति उद्‌घोषयति । गुणप्राधान्यवादेन, विपुलग्रन्थरचनया च लक्षणशास्त्री भोजराजः संस्कृतकाव्यशास्त्रपरम्परायां स्मरणीयो वर्तते।

\specialupakhanda{\textbf{निर्णयात्मककाल-उपसंहारः}}

अस्मिन् काले सर्वथा कालसंज्ञानुसारमेव कार्यं सञ्जातम् । एषोऽवधिः संस्कृतकाव्यशास्त्रस्य इतिहासे तर्कवितर्कप्रवर्धकः, भ्रमनिरोधकश्च वर्तते। अस्मात् पूर्ववर्त्तिनि काले काव्यप्रणयनस्य आन्तरिकतत्त्वं प्रति आचार्याणां जिज्ञासा अनिर्णीता आसन् । यत्र अलङ्कार-रीति-गुणादिषु काव्यात्मतत्त्वं किमिति विषये महत्यो विप्रतिपत्तय आसन्। अस्यां संशयग्रस्तायामवस्थायाम् आचार्यप्रवरेण आनन्दवर्धनेन ‘काव्यस्यात्मा ध्वनिरिति’ उद्‌घोष्य ध्वनिप्रस्थानस्य युगप्रवर्तको विचार उपस्थापितः, येन काव्यशास्त्रस्य चिन्तनदिशा सर्वथा परिवर्तिता।

कालेऽस्मिन् आनन्दवर्धन-अभिनवगुप्त-राजशेखर-कुन्तक-मम्मटप्रभृतयो धुरन्धरा मूर्धन्याश्च काव्यशास्त्रिणः सञ्जाताः, यैः स्वकीयतर्कवैदुष्येण सर्वेषां काव्यतत्त्वानां सुदृढं स्थानं निर्धारितम्। तदनुसारं मानवशरीरवद् काव्यशरीरस्य एका सुव्यवस्थिता परिकल्पना स्थिरीभूता, यत्र— काव्यशरीरस्य आत्मत्वेन रसो ध्वनिर्वा अवतिष्ठते, गुणाः शौर्यादिवद् आन्तरिकधर्मा भवन्ति, अलङ्कारा वस्त्रभूषादिवद् बाह्यसौन्दर्यवर्धकाः स्वीक्रियन्ते, तथा च रीतयोऽवयवसंस्थानादिवद् अङ्गसङ्गठनात्मिकाः सञ्जायन्ते।

पूर्वपक्षिणां सिद्धान्तानां तीव्रतार्किकखण्डन-मण्डनपूर्वकं काव्यतत्त्वानां यद् वैज्ञानिकं स्वरूपम्, काव्यतत्त्वानां स्थानञ्च निर्धारितम्; तदुत्तरवर्त्तिभिराचार्यैरपि प्रायशः सादरमभिनन्दितम्। अतः पौरस्त्यकाव्यशास्त्रस्येतिहासे कालस्यास्य महत्त्वं सातिशयं विद्यते। कालोऽयं केवलं नूतनसिद्धान्तानां प्रसवभूमिरेव नासीद् अपितु संस्कृतसाहित्यशास्त्रस्य सार्वभौमम्, सुस्थिरम्, शाश्वतञ्च स्वरूपं निश्चेतुं परमनिर्णायकः सिद्धो जातः।

\specialmahakhanda{\textbf{पञ्चम-प्रकाशः}}

\specialupakhanda{\textbf{व्याख्याकालः}}

\begin{shloka}[center]

\textbf{(मम्मटात् १००० तमात् वैक्रमाब्दात् पण्डितराज-जगन्नाथं १७०० तमं वैक्रमाब्दं यावत्)}

\end{shloka}

\customparagraph{\textbf{विषयप्रवेशः}}

पौरस्त्यकाव्यशास्त्रेतिहासे मम्मटाचार्यस्योदयाद् व्याख्याकाल आरभ्यते । अत्र ध्वनिसम्प्रदायपक्षपातिनो  मम्मट-रुय्यक-हेमचन्द्र-विद्याधर-विद्यानाथप्रभृतय आचार्याः समभवन् । रससम्प्रदायस्य विश्वनाथ-शारदातनय-शिङ्गभुपाल-भानुदत्त-रूपगोस्वामिप्रभृतयः समालोचकाः सञ्जाताः । अलङ्कारसम्प्रदायस्य जगन्नाथविश्वेश्वरादयः प्रतिभाशालिनः शास्त्रकाराः समागताः । कविशिक्षापरम्पराया राजशेखर–अमरसिंह–अमरचन्द्रः–देवेश्वरादयो विद्वांसः समागताः । निर्णयात्मककाले निर्णीता विषयाः कालेऽस्मिन् स्थिरीभूताः । कालस्यास्य समयो जगन्नाथविश्वेश्वरपर्यन्तं मन्यते । तत्पश्चात्तु संस्कृतकाव्यशास्त्रस्य ह्रासकाल एवेति मन्यते । किमिति चेत् तदा काव्यशास्त्रस्य विषये प्राग्जातानाचार्यानतिरिच्य किमपि कार्यं नाभूत् । कालस्यास्य प्रमुख आचार्यो वाग्देवतावतारो मम्मट एव ।

\specialupakhanda{\textbf{मम्मटः}}

काश्मीरदेशोद्भवस्यास्य पिता जैयटः कैयट–उब्वटनामानौ–अनुजौ चाऽऽस्तामिति काव्यप्रकाशस्य सुधाकरटीकाकारो भीमसेन उल्लिखति, तद्यथा—

\begin{shloka}[center]

\textbf{शब्दब्रह्म सनातनं न विदितं शास्त्रैः क्वचित् केनचित्}

\textbf{तद्देवी हि सरस्वती स्वयमभूत् काश्मीरदेशे पुमान् ।}

\textbf{श्रीमज्जैयटगेहिनीसुजठराज्जन्माप्य युग्मानुजः}

\textbf{श्रीमन्‘मम्मट’संज्ञयाश्रिततनुं सारस्वतीं सूचयन् ।।}

\end{shloka}

भीमसेनस्य कथनमिदं मम्मटसमयात् ६०० वर्षपश्चात् १७२३ तमस्य ईशवीयवर्षस्यास्ति, उव्वटश्च स्वकीये ऋक्प्रातिशाख्ये स्वं वज्रटस्यात्मजमाह । अतः कथनमिदं सर्वस्वीकार्यं नास्ति । आचार्यो मम्मटः १०१५ ईशवीयाब्दे वर्तमानस्याभिनवगुप्तस्य, १०१० ईसवीयां नवसाहसाङ्कचरितकर्तुर्महाकवेः पद्मगुप्तस्य च स्वकीये काव्यप्रकाशे नामोल्लेखमकरोत् । उदात्तालङ्कारस्योदाहरणविषयके पद्ये विद्वज्जनान् प्रति भोजराजस्य दानशीलतापि तेन प्रशंसिता । एतेन मम्मटो भोजराजस्यानन्तरं जातः’ इति सिद्ध्यति । अलङ्कारसर्वस्वे रुय्यकश्च ११३५ ईसाब्दे मम्मटमतस्य खण्डनमकरोत् । अतः मम्मटस्य समयो भोजराजरुय्यकयोर्मध्ये एकादशशताब्दस्योत्तरार्धे स्वीक्रियते ।

\customparagraph{\textbf{काव्यशास्त्रीयं योगदानम्}}

काव्यशास्त्रीयेतिहासे मम्मटस्य योगदानं स्मरणीयमस्ति । इतः पूर्वं काव्यशास्त्रीयमार्गे ये सिद्धान्ता निर्धारिता अभवन्, तान् सर्वान् सिद्धान्तान् परिष्कुर्वन् काव्यतत्त्वविषयेषु प्रामाणिकं सर्वस्वीकार्यं च मतं प्रस्तौति मम्मटः । अस्य काव्यप्रकाशस्तथाविधं  सिद्धान्तस्रोतो वर्तते यस्मात् काव्यसिद्धान्तस्य विभिन्ना धारा निस्सृताः सन्ति । आनन्दवर्धनाचार्यकृतायाः ध्वनेरुद्भावनाया अनन्तरम् आचार्यो भट्टनायको महिमभट्टश्च ध्वनिसिद्धान्तध्वंसाय याः याः युक्तीः समुपस्थापयामासतुः तासां सर्वासां युक्तीनां सतर्कं खण्डनं कृत्वा मम्मटः ध्वनिसिद्धान्तप्रतिष्ठापनमकार्षीत् । तस्मान् मम्मटो ध्वनिप्रस्थापनपरमाचार्यस्य विशेषणेन विभूषितोऽभूत् ।

मम्मटस्य काव्यशास्त्रविषयकः सुप्रसिद्धो ग्रन्थः काव्यप्रकाशः सर्वेषां काव्यशास्त्रिणामादरणीयोऽस्ति । मम्मटस्य विद्वतां प्रमाणयितुं तस्य काव्यप्रकाश एव अलम् । ग्रन्थोऽयं गभीरकाव्यतत्त्वप्रकाशने काव्यतत्त्वस्थानस्थिरीकरणे चाद्वितीयो विद्यते । स च सूत्रात्मिकायां शैल्यां रचितोऽस्ति । अत एवास्याभिप्रायबोधे कुत्रचित् काठिन्यमनुभूयते । काव्यशास्त्रनिष्णातैर्विद्वद्भिरस्य ग्रन्थस्य शताधिकाष्टीकाग्रन्था निर्मिताः, तथापि ग्रन्थोऽयं तथैव दुर्गम एवेति कथ्यते—

\begin{shloka}[center]

\textbf{काव्यप्रकाशस्य कृता गृहे गृहे टीकास्तथाऽप्येष तथैव दुर्गमः ।}

\end{shloka}

अस्मिन् ग्रन्थे दश उल्लासाः सन्ति । एतेषु दशसु उल्लासेषु नाट्यशास्त्रसम्बन्धितान् विषयान् विहाय काव्यसम्बन्धिनां सर्वेषामङ्गानां विवेचनं करोति मम्मटः । तत्र प्रथमे उल्लासे— काव्यलक्षणस्य, काव्यप्रयोजनस्य, काव्यकारणस्य काव्यभेदस्य च निरूपणमास्ते । द्वितीये उल्लासे— शब्दशक्तीनां विस्तरेण समीक्षा कृता  वर्त्तते । तृतीये उल्लासे— शाब्दीव्यञ्जना विवेचिता । चतुर्थे— सर्वेषां ध्वनिभेदानाम्, रसानां भावानां च विस्तरेण विवेचनमास्ते । पञ्चमे— गुणीभूतव्यङ्ग्यकाव्यस्य व्याख्याया अनन्तरं व्यञ्जनाशक्तिस्थिरीकरणाय नवीना युक्तयः समुपस्थापिताः सन्ति । षष्ठे— चित्रकाव्यस्य सामान्यः परिचयो वर्तते । सप्तमे— काव्यदोषाणां विस्तारेण विवेचनमस्ति । उल्लासोऽयं काव्यलक्षणस्य ‘अदोषौ’ पदं सम्यक्तया व्याकरोति । अष्टमे— ‘सगुणौ’ इति काव्यलक्षणपदस्य व्याकृतमास्ते । तत्र मम्मटः माधुर्यौजःप्रसादाख्यान् त्रिगुणान् प्रस्तौति । नवमदशमयोरुल्लासयोः क्रमशः शब्दालङ्कारार्थालङ्कारयोर्विवेचनमस्ति । अनेन एकषष्ठ्यालङ्काराणां विवेचनं कृतम् । ग्रन्थस्यास्य रचनायां काश्मीरदेशोद्भवः कश्चिदल्लटपण्डितोऽपि सहाय्यमकरोत्’ इति काव्यप्रकाशस्य सर्वप्राचीनटीकाकारो माणिक्यचन्द्रः, रुचकपण्डितश्चालिखत—‘अथ चायं ग्रन्थोऽन्येनारब्धोऽपरेण च समापित इति द्विखण्डोऽपि सङ्घटनावशात् अखण्डायते ।’ अन्येऽपि समालोचकाः अल्लटस्य सहाय्यमासीदिति स्वीकुर्वन्ति, यत्-

\begin{shloka}[center]

\textbf{कृतोऽयं मम्मटाचार्यैर्ग्रन्थः परिकरावधिः ।}

\textbf{प्रबन्धः पूरितः शेषो विधायाल्लटसूरिणा ॥}

\end{shloka}

मम्मटाचार्यो विभिन्नेषु काव्यतत्त्वेषु कीदृशं मतमुपस्थापयतीति अत्र प्रदर्श्यते ।

\customparagraph{\textbf{काव्यलक्षणम्}}

कविकर्मकाव्यं कीदृशं भवितुमर्हतीति जिज्ञासासमाधानार्थं काव्यस्वरूपं निर्धार्यते तदेव काव्यलक्षणं मन्यते । अस्मिन् विषये प्रायः काव्यशास्त्रिभिः स्वमतमुपस्थापितं विद्यते परं तेषां सर्वेषां मतसमायोजनपूर्वकं मम्मटः स्वकीयं काव्यलक्षणं प्रस्तौति, यथा—

\begin{shloka}[center]

\textbf{‘तददोषौ शब्दार्थौ सगुणावनलङ्कृती पुनः क्वापि’ (काव्यप्रकाशः)}

\end{shloka}

अत्र अदोषौ रसाद्युद्देश्यप्रतीतिबन्धकदोषरहितौ, सगुणौ माधुर्यौजःप्रसादान्यतमगुणेन तादृशाभ्यां गुणाभ्यां वा सहितौ, अनलङ्कृती क्वचित् स्फुटरससहितस्थले तु स्फुटालङ्काररहितौ–अपि, अत्र नञोऽल्पार्थकत्वेन अनलङ्कृतीत्यस्य स्फुटालङ्काररहितौ–इत्यर्थः उपपद्यते । नञर्थाश्च—

\begin{shloka}[center]

\textbf{तत्सादृश्यमभावश्च तदन्यत्वं तदल्पता ।}

\textbf{अप्राशस्त्यं विरोधश्च नञर्थाः षट् प्रकीर्तिताः ।।}

\end{shloka}

अपि शब्दात् सर्वत्र सालङ्कारौ–इति लभ्यते । शब्दार्थौ शब्दश्च अर्थश्च । शब्दः प्रतीतिपदार्थको ध्वनिः । अर्थोऽत्र शब्दवृत्तिजन्यज्ञानविषयः । तत् काव्यमिति व्यस्तः पदार्थः । इत्यादि साभिप्रायपदार्थविवेचनेन इदृशं काव्यलक्षणं निगलितं भवति यत्— रसादिप्रतिबन्धकदोषरहितौ, प्रातिभासिकेन सत्येन वा गुणेन सहितौ, सर्वत्र सालङ्कारौ, स्फुटरससहितस्थले तु अस्फुटालङ्कारावपि शब्दार्थौ काव्यमिति ।

\customparagraph{\textbf{काव्यहेतुः}}

काव्यहेतुविषयेऽपि मम्मटमतं सर्वोत्कृष्टं मन्यते । सर्वेषां प्राग्जातानां काव्यशास्त्रिणां विचारसमायोजनपूर्वकं मम्मटः सर्वस्वीकारयोग्यं काव्यकारणं प्रस्तौति, यत्—

\begin{shloka}[center]

\textbf{शक्तिर्निपुणता लोकशास्त्रकाव्याद्यवेक्षणात् ।}

\textbf{काव्यज्ञशिक्षयाभ्यास इति हेतुस्तदुद्भवे ।। (काव्यप्रकाशः)}

\end{shloka}

\textbf{शक्तिः (प्रतिभा)-} आह्लादपरकबोधजनिका वर्णना, स चातिशयो जन्मान्तरतोऽनुवर्तमाना देवतामहापुरुषादिप्रसादजन्या ऐहिकी शक्तिः, यां विना काव्यं नोद्भवति, उद्भूतं वा उपहसनीयं जायते।

\textbf{निपुणता (व्युत्पत्तिः)}- लोकशास्त्रकाव्याद्यवेक्षणाद् **निपुणता- ‘**लोकः’ इत्यनेन लौकिकव्यवहारज्ञानम्, ‘शास्त्रम्’ इत्यनेन व्याकरणाभिधानकोषकलादिग्रन्थानामध्ययनम्, ‘काव्य’मित्यनेन वाल्मीकिकालिदासादीनां रामायणरघुवंशादीनामनुशीलनम्, ‘आदि’शब्देन स्त्रीपुरुषलक्षणग्रन्थाः, चार्वाक–बौद्ध–जैनादिधर्मदर्शनग्रन्थादयो ग्रन्थाश्चेतेषां सर्वेषां सङ्ग्रहः । इत्यादिविषयाणामवेक्षणाद्‍= विमर्शनात्, पुनः पुनः पठनाद् विचारणाच्च निपुणता= वर्णनीयविषयोपयोगिसमस्तवस्तुपौर्वापर्यपरामर्शकौशलयुता बहुज्ञता ।

\textbf{अभ्यासः—}  काव्यविषयकज्ञानयुतां विदुषां शिक्षणेन यादृशं नैपुण्यं प्राप्यते स एव अभ्यासः । अर्थात् काव्यं कर्तुं वा विचारयितुं ये जानन्ति तेषामुपदेशेन काव्यस्य  निष्पादने विचारणे यादृशं सामर्थ्यं प्राप्यते स एव अभ्यासः । एते त्रयः समुदिताः, न तु व्यस्ता एव, तत्— तस्य काव्यस्य उद्भवे निर्माणे च हेतुर्भवति, न तु हेतवः ।

त्रयो हेतवः, किमर्थं हेतुरिति एकवचनं न विवक्षितमिति चेत्— हेतुत्रयमिदं  परस्परसापेक्षं  भवितुमर्हति । यथा— दण्डचक्रचीवरादेः समुदितस्य घटत्वादि–एकधर्मावच्छिन्नं प्रति हेतुत्वं जायते तथैव शक्तिनिपुणताभ्यासानां समुदितानामेव काव्यं प्रति हेतुत्वं बोध्यम्, न तु हेतव इति तृणारणिमणिन्यायेन (तृणजवह्निं प्रति तृणस्य, अरणिजवह्निं प्रति अरणेर्मणिजवह्निं प्रति मणेश्च पृथक् पृथक् हेतुत्वम्) काव्यं प्रति शक्त्यादेर्हेतुत्वं न । अस्य सारोऽयम्— शक्तिनिपुणताभ्यासा इति त्रयो हेतव एकीभूय  हेतुत्वेनैव काव्यस्य निर्माणे कारणं भवति  ।

\customparagraph{\textbf{काव्यप्रयोजनम्}}

काव्यप्रयोजनेऽपि पौरस्त्यकाव्यशास्त्रे मम्मटस्यैव मतं प्रामाणिकं व्यावहारिकञ्च मन्यते । अत्र सर्वेषां काव्यशास्त्रिणां काव्यप्रयोजनविचाराः संयोजिताः सन्ति, तद्यथा—

\begin{shloka}[center]

\textbf{काव्यं यशसेऽर्थकृते व्यवहारविदे शिवेतरक्षतये ।}

\textbf{सद्यः परनिर्वृत्तये कान्तासम्मिततयोपदेशयुजे ।। (काव्यप्रकाशः)}

\end{shloka}

वृत्तिश्च यथा—

कालिदासादीनामिव यशः, श्रीहर्षादेः धावकादीनामिव धनम्, राजादिगतोचिताचार- परिज्ञानम्, आदित्यादेर्मयुरादीनामिव रोगनिवारणम्, सकलप्रयोजनमौलिभूतम्, समनन्तरमेव रसास्वादनसमुद्भूतम्, विगलितवेद्यान्तरमानन्दम्, प्रभुसम्मितशब्दप्रधानवेदादिशास्त्रेभ्यः, सुहृत्सम्मितार्थतात्पर्यवत्पुराणादितिहासेभ्यश्च शब्दार्थयोर्गुणभावेन रसाङ्गीभूतव्यापारप्रवणतया विलक्षणं यत्काव्यं लोकोत्तरवर्णनानिपुणं कविकर्म तत् कान्तेव सरसतापादनेनाभिमुखीकृत्य रामादिवत्प्रवर्त्तितव्यं न रावणादिवदित्युपदेशं च यथा योगं कवेः सहृदयस्य च करोतीति सर्वथा तत्र यतनीयम् ।

\textbf{‘कालिदासादीनामिव यशः’} इत्यत्र कालिदासस्येव नाम किमर्थं गृहीतमिति कालिदासस्य मातृपितृकुलं न कोऽपि जानाति । न च दानादिकं किमपि प्रसिद्धमस्ति, तथापि तस्य कीर्तिर्विश्वव्यापिपिनी  अस्ति । सा कीर्तिर्नितान्तरूपेण काव्यसिर्जिता एव । एतदेव कालिदासादीनामिव यशः’ इति कथनस्याभिप्रायोऽस्ति । आदिपदात्दण्डीभारव्यादयः । 	\textbf{श्रीहर्षादेरिव धन}मित्यस्याशयोऽस्ति यत्— धावकः कविः रत्नावलीनाटिकां कृत्वा राज्ञे श्रीहर्षाय समर्पयामास, श्रीहर्षश्च धावकाय प्रचूरं धनं दत्वा तां नाटिकां  स्वकृतत्वेन ख्यापयामासेति प्रसिद्धिरस्ति । आदिपदेन भोजप्रबन्धकारिभिर्भोजाद्  आदि ज्ञायते । एवञ्च रमणीयं काव्यं विरच्य प्रकाशनार्थं विक्रयेण तथा मूल्याङ्कने लोकप्रियत्वेन पुरस्कारयोग्यं भूत्वा तथा च तादृशं योग्यकाव्यं पर्यालोच्य शोधादिकार्येण च धनमाप्तुं शक्यते ।

\textbf{राजादिगतोचिताचारपरिज्ञान}मित्यस्य राजा, आदिपदात् राजपुत्रादीनां वर्णाश्रमाणाञ्च पृथ्वीपालनादितत्तदुचिताचारपरिज्ञानम् । \textbf{आदित्यादेर्मयुरादीनामिव रोगनिवारण}मित्यस्य— पुरा किल मयुरशर्मा कविः क्लेशमसहिष्णु–अत्युच्चतरुशाखावलम्वीशतरज्जुशिक्यमधिरुढः सन् सूर्यस्य प्रार्थनां कृतवान् । स एकैकपद्यान्ते एकैकं रज्जुखण्डं चिच्छेद । एवम् क्रियमाणः काव्यपरितुष्टो रविः सद्य एव नीरोगां रमणीयां च तत्तनुमकार्षीत्, तदेव मयुरशतकनाम्ना प्रसिद्धं स्तोत्रकाव्यम् । आदिपदात् विल्हणकवेः वधनिवृत्तिचौरपञ्चाशिकारूपकाव्यमहिम्नेत्यवधेयम् । 	\textbf{सकलप्रयोजनमौलीभूतमिति}—जगति यावन्ति वस्तुनः फलानि सन्ति तेभ्यः श्रेष्ठं काव्यफलमित्याशयः । \textbf{समनन्तरमेव} इत्यस्य झटिति अर्थात् काव्यश्रवणानन्तरमेव, न तु यागादिवत्देहान्तरोत्पादनेन, न च वृक्षरोपणादिवत् कालविलम्बनेनेत्यर्थः । 	\textbf{रसास्वादनसमुद्भूत}मित्यस्य ‘रस्यते इति रसः’ स्थायिभावस्तदास्वादनेनाऽऽस्वादकरणभूतेन विभावानुभावव्यभिचारिसंयोगेन समुद्भूतं व्यक्तमित्यर्थः । 	विगलितवेद्यान्तरमित्यस्य विगलितं तिरोहितम्, वेद्यान्तरम् अन्यसांसारिकज्ञानमिति, विषयान्तराभासरहितमित्यर्थः। 	\textbf{प्रभुसम्मितोपदेशविलक्षण}मित्यस्य यथा— प्रभोराज्ञाबोधकस्य शब्दस्य परवृत्तिमाज्ञाभङ्गं न सहते तथैव वेदादिशास्त्राण्यपि शब्दपरिवृत्त्यसहानि सर्वथा विधिपराणि च, अतस्तेभ्यो विलक्षणमित्यर्थः । 	\textbf{सुहृत्सम्मितोपदेशविलक्षण}मित्यस्य यथा— मित्राणि इदम् इत्थम्, एवम् कृते एवम् भवति, अमुकेन जनेन इत्थं कृतम्, तस्य इत्थं जातम्, तुभ्यं यथा रोचते तथा कुरु’ इति अवगमयन्ति, तथैव पुराणेतिहासाश्च केवलं सदसत् बोधयन्ति न तु नियोजयन्ति, एतद् विलक्षणमित्यर्थः । ताभ्यां प्रभुसुहृत्सम्मितोपदेशाभ्यां विलक्षणो रसप्रधानः काव्यलक्षणः । तत्र हि रसाङ्गीभूतो यो व्यापारो विभावादिसंयोजनात्मा व्यञ्जनारूपो वा, तन्निष्पाद्यरसादिव्यक्तिनिष्पादकतया शब्दार्थयोर्द्वयोरपि गुणत्वात् रसस्यैव प्राधान्यम्, एवम् विलक्षणं यत्काव्यमिति । काव्यं तु लोकोत्तरा रसोद्बोधसमर्था या वर्णना, तत्र निपुणो यः कविस्तस्य कर्म ।

\textbf{कान्तेव सरसतापादनेन} इत्यस्य— यथा कान्ता स्वप्रियं सरसतामापाद्य स्वाभिमुखीकृत्य आज्ञां वा आदेशं विनैव स्वेच्छितपथि प्रवर्तयति, तथैव काव्यमपि शृङ्गारादिरसेन श्रोतृमनो वशीकृत्य सत्पथे प्रवर्तयति ।

\textbf{रामादिवत्प्रवर्त्तितव्यं न रावणादिव}दित्यस्य काव्ये तु सदसत्पक्षयोः प्रस्तुतिर्भवति परं तत्र वर्णिताद् विषयाद् रामादिवत्सत्कर्मशीलस्येव कार्यमनुकरणीयं न तु रावणवद्  असत्कार्ये प्रवृत्तस्यापि इति । 	\textbf{यथायोग}मित्यस्य— यशोऽर्थानर्थनिवारणानि कवेरेव, व्यवहारोपदेशौ सहृदयस्यैव, परनिर्वृत्तिरुभयोरपीत्यर्थः ।

\textbf{सर्वथा तत्र यतनीय}मित्यस्य— एतावदिष्टविशिष्टं काव्यम्, अत एव तत्र कविः सहृदयश्च निर्माणेऽनुशीलने च प्रयत्नरतो भवितव्यमिति मम्मटस्य काव्यप्रयोजनस्य सारभूतोऽर्थः ।

\customparagraph{\textbf{अलङ्कारनिरूपणम्}}

काव्यशास्त्रेऽलङ्कारस्य विषयेऽपि मम्मटस्य मतमेव आधिकारिकम्, प्रामाणिकम्, सर्वस्वीकार्यं च मन्यते । अलङ्कारवादिभिः काव्यात्मस्थानित्वेन स्वीकृतमलङ्कारतत्त्वमनेन काव्यस्य बाह्यसौन्दर्यतत्त्वत्वेन स्थापितम् । एतदेव मतं पौरस्त्यकाव्यशास्त्रे सिद्धान्तमतत्वेन स्थिरीभूतम् । एतस्यालङ्कारलक्षणं यत्—

\begin{shloka}[center]

\textbf{उपकुर्वन्ति तं सन्त येऽङ्गद्वारेण जातुचित् ।}

\textbf{हारादिवदलङ्कारास्तेऽनुप्रासोपमादयः ।। (काव्यप्रकाशः)}

\end{shloka}

यथा– मानवशरीरं कटककुण्डलादयः बाह्यरूपेणैव सुन्दरं करोति । ते च अनित्या भवन्ति, अर्थात् तेषामभावेऽपि मानवस्यास्तित्वं भवत्येव तथैव काव्येऽपि उपमादयोऽलङ्कारा बाह्यसौन्दर्यवर्धका एव ।

\customparagraph{\textbf{गुणस्य स्वरूपम्}}

आचार्यो मम्मटो गुणस्य विषये यं मतमुपस्थापयति । तदेवोत्तरवर्त्तिनामाचार्याणां कृते स्वीकारयोग्यं जातम्, शास्त्रीयसन्दर्भे च प्रामाणिकं मन्यते । एतस्य गुणस्वरूपं यथा—

\begin{shloka}[center]

\textbf{ये रसस्याङ्गिनो धर्माः शौर्यादय इवात्मनः ।}

\textbf{उत्कर्षहेतवस्ते स्युरचलस्थितयो गुणाः ।। (काव्यप्रकाशः)}

\end{shloka}

यथा— मानवशरीरे शौर्यादय आत्मनो धर्मा भवन्ति । ते च स्थायिरूपेण आत्मनि अवस्थिताः सन्ति, तथैव काव्यशरीरेऽपि माधुर्यादयो गुणा आत्मरूपेणावस्थितं रसमङ्गीकृत्य स्थायित्वेनावस्थिता जायन्ते । अर्थात् गुणस्य काव्यात्मतत्त्वेन सह अलङ्कारापेक्षया निकटसम्बन्धो भवति । एतेन ध्वनिवादिनां गुणालङ्कारयोः पार्थक्यं, तयो रसेन सह सम्बन्धश्च स्पष्टो भवति, यद् — गुणा रसधर्माः, नियमितकाव्ये विद्यमानाः स्वरूपाधायकाश्च भवन्ति । अलङ्कारास्तु शब्दार्थधर्माः काव्येऽनियमितविद्यमानाश्चमत्काराधायका एव भवन्ति ।

\customparagraph{\textbf{रीतेर्विचारः}}

काव्यप्रकाशकारो मम्मटो रीतेर्विषये समहत्त्वं विवेचनं न करोति । स वृत्त्यनुप्रासाख्यशब्दालङ्कारान्तर्गतत्त्वेन रीतिं स्वीकरोति । काव्यप्रकाशस्य नवम उल्लासे वृत्त्यनुप्रासस्य भेदकथनप्रसङ्गे प्राचीनानां रीतेर्लक्षणं नामभेदेन प्रदर्शयति, यत् —

\begin{shloka}[center]

\textbf{माधुर्यव्यञ्जकैर्वर्णैरुपनागरिकोच्यते ।}

\textbf{ओजः प्रकाशकैस्तैस्तु परुषाः कोमलापरैः ।। (काव्यप्रकाशः)}

\end{shloka}

अर्थात् प्राचीनानां वैदर्भीगौडीपाञ्चालीति रीतित्रयं क्रमशः उपनागरिकापरुषाकोमलेति भेदेन स्वीकरोति, तत्र स कथयति यत्—

\begin{shloka}[center]

\textbf{केषाञ्चिदेता वैदर्भीप्रमुखा रीतयो मता । (काव्यप्रकाशः)}

\end{shloka}

\customparagraph{\textbf{काव्यभेदविवेचनम्-}}

आचार्यो मम्मटो ध्वनिवादी विचारकोऽस्ति, अतो ध्वन्याधारेणैव काव्यभेदान् विवृणोति । अस्य मते ध्वनिर्यत्र सर्वोत्कृष्टेन स्थितो भवति, तदेव उत्कृष्टं काव्यं जायते । एतदेवाधारे काव्यं चतुर्षु भेदेषु वर्गीकरोति मम्मटो यथा—

\customparagraph{\textbf{उत्तमकाव्यम्-}}

\begin{shloka}[center]

\textbf{इदमुत्तममतिशयिनि व्यङ्ग्ये वाच्यात् ध्वनिर्वुधैर्कथितः । (काव्यप्रकाशः)}

\end{shloka}

पूर्वाचार्यैर्वैयाकरणैः प्रधानभूतस्फोटरूपव्यङ्ग्यव्यञ्जकशब्दस्य ध्वनिरितिव्यवहारकृतः । ततस्तन्मतानुसारिभिरन्यैरपि न्यग्भावितवाच्यव्यङ्ग्यव्यञ्जनक्षमस्य शब्दार्थयुगलस्य ध्वनिरिति व्यवहारकृतः । अत एव वाच्यापेक्षया व्यङ्ग्यस्य प्राधान्यत्वं ध्वनित्वमिति उत्तमकाव्यलक्षणम् । अस्योदाहरणं यथा—

\begin{shloka}[center]

\textbf{निःशेषच्युतचन्दनं स्तनतटं निर्मृष्टरागोऽधरो}

\textbf{नेत्रे दूरमनञ्जने पुलकिता तन्वी तवेयं तनुः ।}

\textbf{मिथ्यावादिनि दूति बान्धवजनस्याज्ञातपीडागमे}

\textbf{वापीं स्नातुमितो गताऽसि न पुनस्तस्याधमस्यान्तिकम् ।।}

\end{shloka}

इत्यादौ सामान्यवाच्यार्थापेक्षया तदन्तिकमेव रन्तुं गतासीति प्राधान्येन व्यज्यते । अत एवोत्तमं काव्यम् ।

\customparagraph{\textbf{मध्यमकाव्यम्-}}

\begin{shloka}[center]

\textbf{अतादृशी गुणीभूतव्यङ्ग्यं व्यङ्ग्ये तु मध्यमम् । (काव्यप्रकाशः)}

\end{shloka}

यत्र व्यङ्ग्यार्थापेक्षया वाच्यार्थस्य चमत्कारित्वे सति गुणीभूतव्यङ्ग्यनामकं मध्यमकाव्यम् । तथा च व्यङ्ग्यवाच्ययोस्तुल्यत्वे सत्यपि गुणीभूतव्यङ्ग्यमेव भवति, यथा—

\begin{shloka}[center]

\textbf{ग्रामतरुणं तरुण्या नवञ्जुलमञ्जरीसनाथकरम् ।}

\textbf{पश्यन्त्या भवति मुहुर्नितरां मलिना मुखच्छाया ।।}

\end{shloka}

अत्र वञ्जुललतागृहे समागमार्थं दत्तसङ्केता नायिका नागता’ इति प्रतीयमानार्थं गुणीभूतम्, तदपेक्षया मुखकान्तेर्मलिनत्वरूपं वाच्यार्थमेव चमत्कारित्वादिति । अधमकाव्यं यथा-

\customparagraph{\textbf{अधमकाव्यम्-}}

\begin{shloka}[center]

\textbf{शब्दचित्रं वाच्यचित्रमव्यङ्ग्यं त्ववरं स्मृतम् । (काव्यप्रकाशः)}

\end{shloka}

यत्र गुणालङ्कारयुक्तं किन्तु स्फुटप्रतीयमानार्थरहितं वाक्यं चमत्कारकमस्ति तदधमकाव्यम् । स च शब्दार्थभेदेन द्विधा भवति । तेषु शब्दचित्रं यथा—

\begin{shloka}[center]

\textbf{स्वच्छन्दोच्छलदच्छकच्छकुहरच्छातेतराम्बुच्छटा–}

\textbf{मूर्च्छन्मोहमहर्षितुल्यविहितस्नानाह्निकाह्नाय वः ।}

\textbf{भिद्यादुद्यदुदारदर्दुरदरी दीर्घादरिद्रद्रुम—}

\textbf{द्रोहोद्रेकमहोर्मिमेदुरमदा मन्दाकिनी मन्दताम् ।।}

\end{shloka}

अत्र कवेर्मन्दाकिनीविषयको रतिभावो विद्यते, तथा च अन्यतीर्थापेक्षया मन्दाकिन्या  विशेषताधिक्यवर्णनाद् व्यतिरेकालङ्कारोऽपि व्यज्यते, परं कवेरत्र अनुप्राससहकारेण शब्दचमत्कारप्रदर्शनतात्पर्यत्वेन प्रतीयमानार्थस्तिरोहितत्वात् प्राधान्येन शब्दचमत्कार एव दृश्यते । अतः पद्यमिदं शब्दचित्रस्योदाहरणम् ।

\begin{shloka}[center]

\textbf{विनिर्गतं मानदमात्ममन्दिराद्भवत्युपश्रुत्य यदृच्छयापि यम् ।}

\textbf{ससभ्रमेन्द्रद्रुतपातितार्गला निमिलिताक्षीव भियाऽमरावती ।।}

\end{shloka}

अत्र हयग्रीवस्योत्साहस्थायिभावात्मकस्य वीरत्वव्यङ्ग्येऽपि कवेर्तात्पर्यमत्र उत्प्रेक्षाऽलङ्कारचमत्कारप्रदर्शनेऽस्ति । अतोऽमरावत्या नेत्रनिमीलनात्मकेन सम्भावनारूपेण अर्थगतचमत्कारजनकोत्प्रेक्षाऽलङ्कारेण प्रतीयमानार्थतिरोहितत्वाद्  अर्थाऽलङ्कारप्राधन्यमस्ति । तस्माद्  इदं काव्यमर्थचित्रस्योदाहरणम् ।

आचार्यमम्मट एकादशशताब्द्या उत्तरार्धे काश्मीरदेशे समुत्पन्नः संस्कृतकाव्यशास्त्रेतिहासस्य कश्चन युगप्रवर्तकः, निर्णायकश्च विद्वान् आसीत्। `वाग्देवतावतारः' इति उपाधिना विख्यातेन अनेन रचितः काव्यप्रकाशः सकलकाव्यशास्त्रीयसिद्धान्तानां परिष्कारकः, प्रामाणिकः, सर्वस्वीकार्यश्च आकरग्रन्थो वर्तते। ध्वनिसिद्धान्तस्य खण्डनकर्तॄणां भट्टनायक-महिमभट्टादीनां युक्तीनां सतर्कं विखण्ड्य ध्वनेः सुदृढप्रतिष्ठापनादयं **'**ध्वनिप्रस्थापनपरमाचार्यः' इत्युपाधिना समाद्रियते। मम्मटेन स्वग्रन्थे दोषरहित-गुणसहित-सालङ्कारशब्दार्थयोः काव्यस्वरूपं स्वीकृतम् । शक्ति-निपुणता-अभ्यासानां (प्रतिभा-व्युत्पत्ति-अभ्यासाः) हेतुत्रयाणां पृथक् पृथक् न अपितु समुदितरूपेणैव काव्यहेतुत्वं प्रतिपादितम्। तेन निरूपितं षड्‌विधं काव्यप्रयोजनं काव्यजगतः कृते अतीव व्यावहारिकं मार्गदर्शकञ्चास्ति। अयं गुणालङ्कारयोस्तात्त्विकं भेदं प्रदर्श्य गुणाः काव्यस्यात्मभूतस्य रसस्य स्थायिधर्माः (शरीरे शौर्यादिवत्) तथा अलङ्काराः शब्दार्थयोर्बाह्यसौन्दर्यविधायका अस्थायिधर्माः (हारादिवत्) इति निर्णायकं मतं स्थिरीकरोति। मम्मटेन ध्वनिप्राधान्याधारेणैव काव्यस्य त्रयो मुख्यभेदाः स्वीकृताः । तत्र उत्तमं (ध्वनिकाव्यम्), मध्यमं (गुणीभूतव्यङ्ग्यम्), अधमं (शब्दचित्रम् अर्थचित्रञ्च) चेति सोदाहरणं वैज्ञानिकीं वर्गीकरणपद्धतिं प्रस्तूय प्राचीनविवादाः शमिताः। यद्यपि सूत्रात्मकशैल्या संरचितत्वादयं ग्रन्थो गृहे गृहे टीका विहिता अपि दुर्गम एव मन्यते, तथापि मम्मटस्य इयं समन्वयवादिनी दृष्टिरुत्तरवर्तिनां काव्यशास्त्रिणां कृते आधारशिला सञ्जाता। काव्यस्य मूलतत्त्वपरिज्ञाने, गभीरतत्त्वाविष्करणे, जिज्ञासूनां भ्रमापहरणे च मम्मटस्य योगदानं संस्कृतसाहित्यशास्त्रे चिरस्मरणीयम्, अक्षुण्णम्, सर्वमान्यञ्च वर्तते।

\specialupakhanda{\textbf{रुय्यकः}}

आचार्योऽयमलङ्कारसम्प्रदायप्रवर्धकेषु अन्यतमो विद्यते । एषः काश्मीरदेशवास्तव्यः श्रीकण्ठचरितस्य रचयितुर्मङ्खककवेर् गुरुश्चासीत् । श्रीकण्ठचरितस्य रचनाकालः ११३५ ईसवीयतः ११४९ इसवीयं मन्यते । रुय्यकः स्वलङ्कारसर्वस्वे विक्रमाङ्कदेवचरितस्य (१०८५ ई.) कतिपयश्लोकान् उदाहरणरूपेण प्रस्तुतवान्, व्यक्तिविवेकस्य (१०५० ई.) टीकां च रचितवान् । अतोऽस्य कालो महिमभट्टस्य, विल्हणस्य, मम्मटस्य च पश्चादेव निर्धारयितुं शक्यते । तथा च अस्य शिष्यो मङ्खककवि द्वादशशताब्द्या मध्यवर्त्तिको ज्ञायते । अतोऽस्य समयो द्वादशशताब्दस्य मध्यभागः स्वीक्रियते । आचार्योऽयं ‘रुचकः’ इति अपरेणाभिधानेनापि विख्यातम्, यत्—

\begin{shloka}[center]

\textbf{ज्ञात्वा 	श्रीतिलकात् सर्वालङ्कारोपनिषद्रसम् ।}

\textbf{काव्यप्रकाशसङ्केतो रुचकेनेह लिख्यते ।।}

\end{shloka}

‘कृतिः श्रीविपश्चिद्वरराजानकतिलकात्मजः श्रीमदालङ्कारिकसमाजाग्रगण्यश्रीराजानक- रुय्यकस्य राजानकरुचकापरनाम्नोऽलङ्कारसर्वस्वकृतः ।’ इति चित्रमीमांसायाम् उल्लिखितो विद्यते । अस्य कृतिषु मध्ये साहित्यमीमांसा, अलङ्कारसर्वस्वञ्चेति द्वे कृती काव्यशास्त्रसम्बद्धे स्तः ।

\customparagraph{\textbf{काव्यशास्त्रीयं योगदानम्}}

अनेन काव्यप्रकाशस्य सङ्केतनामिका टीका लिखिता । तत्र स्थाने स्थाने मम्मटस्य मतस्यालोचनापि कृता । तथैव महिमभट्टस्य व्यक्तिविवेकोपरि व्यक्तिविवेकविचाराख्या टीका च विहिता । साहित्यमीमांसा तु अस्य प्रथमो मौलिको ग्रन्थो विद्यते । तत्र अष्टौ प्रकरणाः सन्ति । तेषु निम्नाङ्किता विषया विवेचिताः प्राप्यन्ते, प्रथमे– कविरसिकयोः प्रभेदः, द्वितीये–वृत्त्यादेर्लक्षणम्, तृतीये– दोषाणां विवेचनम्, चतुर्थे– गुणानां मीमांसा, पञ्चमे– अलङ्कारविवेचनम्, षष्ठे– रसभावयोर्समीक्षणम्, सप्तमे– कवीनां वैशिष्ट्यचतुष्टयम्, अष्टमे– आनन्दस्य स्वरूपं वर्णितं विद्यते ।

अस्य अलङ्कारसर्वस्वाख्यो लाक्षणिकग्रन्थस्तु अलङ्कारसम्प्रदायस्य कृते वरदानरूपं  दृश्यते । ग्रन्थस्यास्य वैशिष्ट्यं विद्यते यद्— अलङ्काराणां वैज्ञानिकं वर्गीकरणम् । ग्रन्थेऽस्मिन् शताधिकाष्टादशसूत्राः सन्ति । येषु एकादशसूत्रेषु षड्शब्दालङ्काराणाम्, एकसप्ततिसूत्रेषु परम्परागतानामर्थालङ्काराणाम्, त्रिषु सूत्रेषु रसवत्–प्रेय–ऊर्जस्वि–समाहितादीनाम्, शेषेषु मिश्राद्यलङ्काराणां विवेचनं कृतवान् रुय्यकः । अनेन प्राचीनानां पञ्चसप्ततिरेवालङ्काराः स्वीकृताः । सप्तसङ्ख्यकानां नूतनालङ्काराणामुद्भावेन सह समष्टौ ८२ अलङ्कारा अनेन विवेचिताः । एवमेवास्य नाटकमीमांसाश्रीकण्ठस्तवालङ्कारमञ्जरीप्रभृतयो ग्रन्थाः सम्प्रत्यनुपलब्धाः कथ्यन्ते ।

रसध्वन्यादिविषयेऽनेन पूर्वेषामाचार्याणां मतमुल्लिख्य तदनुरूपमेव स्वीकृतम् । स्पष्टतया तु नोक्तम्। तथापि वर्णनशैल्या एष ध्वनिवादी आसिदिति मन्यते । ध्वनिदृष्ट्या एव काव्यस्य वर्गीकरणं कृत्वा चित्रकाव्यवर्गेऽलङ्कारं स्थापयित्वा अलङ्कारस्य सूत्रात्मकं विवेचनं कृतम् ।

\begin{shloka}[center]

\textbf{तत्र व्यङ्ग्यस्य प्राधान्याप्राधान्याभ्यां ध्वनिगुणीभूतव्यङ्ग्याख्यौ द्वौ काव्यभेदौ ।}

\textbf{व्यङ्ग्यस्यास्फुटत्वेऽलङ्कारवत्त्वेन चित्राख्यः काव्यभेदस्तृतीयः । तत्रोत्तमो ध्वनिः ।}

\end{shloka}

\customparagraph{\textbf{(अलङ्कारसर्वस्वम्)}}

अस्य अलङ्कारशास्त्रीयपक्षे महत्त्वपूर्णं योगदानं तु अलङ्काराणां वैज्ञानिकं वर्गीकरणमेव । अलङ्कारविश्लेषणसन्दर्भेऽनेन अलङ्कारस्य वर्गीकरणमिति शीर्षकमुल्लिख्य स्पष्टरूपेण एकत्र अलङ्कारवर्गीकरणं नोक्तम् । तत्र तु ‘अथ सादृश्यमूलका अलङ्काराः’ ‘इति सादृश्यमूलका अलङ्कारा वर्णिता’ इत्यनया शैल्या विहितं साङ्केतिकवर्गीकरणं समालोचकैर्विविच्य एकत्रीकृत्य ‘रुय्यकस्यालङ्कारवर्गीकरण’मिति प्रकाशितम् । अतोऽलङ्काराध्ययने प्रवृत्ता जिज्ञासवः अलङ्कारसर्वस्वे सारल्येन अलङ्कारवर्गीकरणं कुत्र विद्यते इति ज्ञातुं न शक्नुवन्ति । अस्य वर्गीकरणमुत्तरवर्त्तिभिराचार्यैरभिनन्दितम्, अलङ्कारविवेचनेऽऽधारत्वेन स्वीकृतञ्च । रुय्यकस्य प्रसिद्धम् अलङ्कारवर्गीकरणं निम्नानुसारं वर्तते—

अलङ्कारस्य—

१ शुद्धवर्गः २ मिश्रवर्गः

शुद्धस्य—

१ अर्थालङ्कारवर्गः २ शब्दालङ्कारवर्गः

अर्थालङ्कारस्य—

१. सादृश्यविच्छित्तिः २. विशेषणविच्छित्तिः ३. गम्यतार्थविच्छित्तिः ४. विरोधविच्छित्तिः ५. शृङ्खलाविच्छित्तिः ६. न्यायविच्छित्तिः ७. गूढार्थप्रतीतिविच्छित्तिश्चेति सप्तधा सन्ति ।

\textbf{१. सादृश्यविच्छित्तिवर्गे—}  (क) भेदाभेदतुल्यतामूलके  उपमा, अनन्वयः, उपमेयोपमा, स्मरणेति चत्वारोऽलङ्काराः सङ्गृहीताः ।

(ख) अभेदप्राधान्यमूलके-  (अ) आरोपाश्रिते— रूपकम्, परिणामः, सन्देहः, भ्रान्तिमान्, उल्लेखः, अपह्नुतिश्चेति षट्, (आ) अध्यवशायाश्रिते— उत्प्रेक्षा, अतिशयोक्तिश्चेति द्वावेव ।

(ग) गम्यौपम्यमूलके-  तुल्ययोगिता, दीपकम्, प्रतिवस्तूपमा, दृष्टान्तः, निदर्शना चेति षट् ।

(घ) भेदप्राधान्यमूलके-  व्यतिरेकः, सहोक्तिः, विनोक्तिश्चेति त्रयः ।

\textbf{२. विशेषणविच्छित्तिवर्गे—} (क) केवलविशेषणविच्छित्तिभेदे-  समासोक्तिः, परिकरश्चेति द्वावेव । (आ) सविशेष्यविशेषणविच्छित्तिभेदे-  श्लेषः, अप्रस्तुतप्रशंसा, अर्थान्तरन्यासश्चेति त्रयः ।

\textbf{३. गम्यतार्थविच्छित्तिवर्गे—} पर्यायोक्तिः, व्याजस्तुतिः, आक्षेपश्चेति त्रयः ।

\textbf{४. विरोधविच्छित्तिवर्गे—} विरोधः, विभावना, विशेषोक्तिः, असङ्गतिः, विषमः, समः, विचित्रम्, व्याघातः, अधिकः, विशेषः, अन्योन्यञ्चेति एकादशः ।

\textbf{५. शृङ्खलाबन्धवर्गे—} कारणमाला, एकावली, मालादीपकम्, सारश्चेति चत्वारः ।

\textbf{६. न्यायविच्छित्तिवर्गे—}

(क) तर्कन्यायमूलके- काव्यलिङ्गम्, अनुमानञ्च द्वावेव ।

(ख) लोकन्यायमूलके- प्रत्यनीक, प्रतीप, मीलित, सामान्य, तद्गुण, अतद्गुण र उत्तरश्चेति सप्त ।

(ग) वाक्यन्यायमूलके- यथासङ्‌ख्य, परिसङ्‌ख्या, अर्थापत्ति, समुच्चय, परिवृत्ति, विकल्प, समाधि र पर्यायश्चेति अष्टौ ।

\textbf{७. गूढार्थप्रतीतिमूलकवर्गे—} सूक्ष्मः, व्याजोक्तिः, वक्रोक्तिः, स्वभावोक्तिः, भाविकः, उदात्तः, रसवत्, प्रेयः, ऊर्जस्वि, समाहितः, भावोदयः, भावशान्तिः, भावसबलता चेति त्रयोदश ।

अलङ्कारसर्वस्वकारो रुय्यको द्वादशशताब्द्या मध्यसमये काश्मीरदेशे समुत्पन्नः अलङ्कारसम्प्रदायस्य प्रवर्धकेषु अन्यतमो विद्वान् आसीत्। मङ्खककवेर्गुरुणा अनेनाचार्येण साहित्यमीमांसा-अलङ्कारसर्वस्वसदृशैः प्रौढग्रन्थैः, काव्यप्रकाशादिलक्षणशास्त्रसम्बद्धाभि टीकाभिश्च संस्कृतकाव्यशास्त्रं महता प्रकर्षेण समृद्धम्। स्पष्टतया नोक्तमपि अस्य वर्णनशैल्या एषो ध्वनिवादी आचार्यः प्रतीयते, येन काव्यस्य ध्वनिकाव्यम्, गुणीभूतव्यङ्ग्यकाव्यम्, चित्रकाव्यञ्च त्रयो भेदाः स्वीकृताः । अलङ्काराश्च चित्रकाव्यवर्गे स्थापिताः।

अस्य आचार्यस्य सर्वोपरि श्लाघनीयं योगदानम् अलङ्काराणां वैज्ञानिकं वर्गीकरणम् अस्ति। तेन परम्परागत-नूतन-अलङ्कारान् संयोज्य समष्टौ ८२ अलङ्काराणां विवेचनं प्रथमतः शुद्ध-मिश्र-भेदेन तदनन्तरं सादृश्य-विरोध-शृङ्खला-न्यायादि सप्तसु वर्गेषु कृतम्। अलङ्कारसर्वस्वे विप्रकीर्णम् अस्य साङ्केतिकं वर्गीकरणं समालोचकैरेकत्रीकृत्य `रुय्यकस्यालङ्कारवर्गीकरणम्' इति नाम्ना प्रकाशितम् । यद् उत्तरवर्तिभिराचार्यैरलङ्कराध्ययनस्य आधारत्वेन सादरं स्वीकृतम्। अतोऽस्य अलङ्कारसर्वस्वाख्यो लाक्षणिकग्रन्थोऽलङ्कारसम्प्रदायस्य कृतेऽद्वितीयं वरदानरूपं वर्तते, वर्गीकरणानेन संस्कृतसाहित्ये रुय्यकस्य स्थानम् आरादरणीयं सञ्जातम्।

\specialupakhanda{\textbf{हेमचन्द्रः}}

संस्कृतकाव्यशास्त्रे हेमचन्द्रस्य नाम रुय्यकपश्चादायाति । काव्यानुशासनस्य रचयिता आचार्योऽयं गुर्जरस्याहमदावादजनपदस्य धुन्धुकाख्ये ग्रामे १०८० ईसवीयवर्षे जन्माऽलभत  । अस्य पिता चाचः, माता च पाहिनी, शैशवाभिधानञ्च चङ्गदेव आसीत् । अस्य गुरुश्च देवचन्द्र आसीत् । सूरिपदप्राप्तोऽयं ११७२ इसवीये स्वर्गवासमलभत  ।

\customparagraph{\textbf{काव्यशास्त्रीयं योगदानम्}}

अस्य काव्यानुशासनाख्यः शास्त्रीयग्रन्थो विद्यते । सूत्रात्मिकां शैल्यां लिखितस्यास्य ग्रन्थस्य स्वयमेव विवेकाभिधानां व्याख्यामकरोत् । अष्टाध्यायेषु निबद्धस्य ग्रन्थस्यास्य प्रथमेऽध्याये— काव्यप्रयोजनस्य, काव्यहेतोः, काव्यलक्षणस्य, शब्दार्थयोः स्वरूपस्य च विवेचनमास्ते । द्वितीये— रसस्य तस्य भेदानाञ्च विशदं विवेचनं कृतं वर्तते । तृतीयेऽध्याये— दोषाणां निर्णयो वर्तते । चतुर्थेऽध्याये— ओजःप्रसादमाधुर्याणां विवेचनमास्ते । पञ्चमेऽध्याये— षट्‌प्रकाराणां शब्दालङ्काराणां तथा षष्ठेऽध्याये— एकोनविंशत्यर्थालङ्काराणाञ्च वर्णनं कृतमस्ति ।

हेमचन्द्रः पूर्वेषां मते विद्यमानस्य संसृष्टेरलङ्कारस्य सङ्करालङ्कारे एवान्तर्भावं चकार, तथा च दीपके तुल्ययोगितायाः । अनेन परिवृत्तिनामकस्य नवीनालङ्कारस्य चोद्भावनं कृतम् । अस्मिन्नेवालङ्कारे मम्मटाभिमतयोः पर्याप्तपरिवृत्त्योरलङ्कारयोरन्तर्भावं करोति । तथा च निदर्शनायामेव प्रतिवस्तूपमादृष्टान्तयोरन्तर्भावं करोति आचार्योऽयम् । अस्य काव्यानुशासनाख्यो ग्रन्थः सङ्ग्रहात्मक  आभाति । मौलिकता तु कुत्रचित् स्थाने एव दृश्यते । ग्रन्थेऽस्मिन् काव्यमीमांसायाः, काव्यप्रकाशस्य, ध्वन्यालोकलोचनस्याभिनवभारत्याश्चोदाहारणानि समुल्लिखितानि । तथापि रसालङ्कारविषये अनेन कृतं विवेचनं पौरस्त्यकाव्यशास्त्रस्य कृतेऽवदानं मन्यते ।

\customparagraph{\textbf{काव्यलक्षणम्}}

आचार्यो हेमचन्द्रः प्राग्जातानामाचार्याणामनुयायी विद्यते । अस्य काव्यलक्षणमपि भामहमम्मटमतानुरूपं शब्दार्थमूलमस्ति ।

\begin{shloka}[center]

\textbf{अदोषौ सगुणौ सालङ्कारौ च शब्दार्थौ काव्यम् ।  (काव्यानुशासनम्)}

\end{shloka}

\specialupakhanda{\textbf{आचार्यो वाग्भटः}}

हेमचन्द्रस्य समकालिको जैनधर्मावलम्बी आलङ्कारिकोऽस्ति वाग्भटः । एतस्य रचना वाग्भटालङ्काराभिधाना वर्तते । अस्यैव ग्रन्थस्य एकस्य पद्यस्य टीकया ज्ञायते यदस्य वास्तविकं नाम बाहड आसीत् । अयं सोमस्यात्मजः कस्यचन राज्ञो महामात्यपदे प्रतिष्ठित आसीत् । आचार्योऽयं संस्कृतप्राकृतभाषयोर्विशेषज्ञो ज्ञायते । अस्य ग्रन्थस्योदाहरणेषु कर्णपुत्रस्य अनहिलवाडाधिपतेश्  चालुक्यवंशीयनरेशस्य जयसिंहस्य स्तुतिः समुपलभ्यते । तस्मात् वाग्भटस्य जयसिंहेन सह घनिष्ठः सम्बन्ध आसीदिति बोध्यते । जयसिंहस्य राज्यकालः १०९३ ईसवीयतः ११४३ ईसवीयं यावत् स्वीक्रियते । अतोऽस्य कालो द्वादशशताब्दस्य पूर्वार्धः स्वीक्रियते ।

\customparagraph{\textbf{काव्यशास्त्रीयं योगदानम्}}

अस्य शास्त्रीयग्रन्थस्य नाम वाग्भटालङ्कारो वर्तते । लाक्षणिकग्रन्थोऽयं पञ्चसु परिच्छेदेषु संविभक्तोऽस्ति । तत्र काव्यशास्त्रसिद्धान्तानां सङ्क्षेपेण विवेचनं प्राप्यते । तेषु प्रथमपरिच्छेदे— काव्यस्वरूपस्य, प्रतिभाभ्यासव्युत्पत्तीनाञ्च काव्यहेतूनां वर्णनं वर्तते । द्वितीये— काव्यस्यानेकेषां भेदानां वर्णनं कृत्वा पदानाम्, वाक्यानाम्, अर्थानाञ्च दोषाणां सङ्क्षिप्तचर्चा विद्यते । तृतीये— दशगुणानां सोदाहरणं विवेचनम्, चतुर्थे— चतुर्णां शब्दालङ्काराणाम्, पञ्चत्रिंशदर्थालङ्काराणाम्, द्वयो रीत्योश्च निरूपणम् । पञ्चमे— नवानां रसानाम्, नायकनायिकयोर्भेदानाम् अन्येषां च विषयाणां विवेचनं प्राप्यते ।

संस्कृतकाव्यशास्त्रे वाग्भटस्य सामान्ययोगदानं विद्यते । तथापि अलङ्कारविवेचनप्रसङ्गे वाग्भटस्य नाम ससम्मानं  गृह्यते ।

\customparagraph{\textbf{काव्यलक्षणम्}}

आचार्यो वाग्भटः स्वकीये वाग्भटालङ्कारे प्राग्जातानामाचार्याणां काव्यलक्षणं विविच्य परिष्कारमयं सुस्पष्टं काव्यस्वरूपं प्रस्तौति । गुणालङ्काररीतियुतं रसमयं काव्यम्, काव्यप्रणेतुः कीर्तये भवतीति सप्रयोजनं काव्यलक्षणमुल्लिखति, यथा-

\begin{shloka}[center]

\textbf{साधु शब्दार्थसन्दर्भं गुणालङ्कारभूषितम् ।}

\textbf{स्फुटरीतिरसोपेतं काव्यं कुर्वीत कीर्तये ।। (वाग्भटालङ्कारः)}

\end{shloka}

\specialupakhanda{\textbf{जयदेवः}}

संस्कृतकाव्यशास्त्रस्यालङ्कारपक्षीयेतिहासे जयदेवस्यावदानं स्मरणीयं विद्यते । अस्य चन्द्रालोकाख्योऽलङ्कारविषयको ग्रन्थः शास्त्रीयज्ञानार्जनशीलान् पाठकान् कृते सरलसुबोध्यतया ग्रहणीयोऽस्ति । अस्योपनाम पीयूषवर्षोऽपि अस्तीति जयदेवः स्वयमेवोल्लिखति, यत्— \textbf{``चन्द्रालोकममुं स्वयं वितनुते पीयूषवर्षः कृती ।''} अस्य पितरौ नाम सुमित्रामहादेवौ स्तः । तद्यथा—

\begin{shloka}[center]

\textbf{महादेवः सत्रप्रमुखमखविघ्नैकचतुरः ।}

\textbf{सुमित्रा तद्भक्तिप्रणिहितमतिर्यस्य पितरौ ।।}

\end{shloka}

असौ वङ्गदेशस्य राज्ञो लक्ष्मणसेनस्य पञ्चरत्नमध्ये एकतम आसीत् । सम्प्रति वङ्गप्रदेशे ‘केंदुली’ इति नाम्ना प्रसिद्धे ग्रामे वीरभूमिमण्डले अस्य स्थानमासीत्, यस्य प्राचीनं नाम ‘विन्दुविल्लवः’ इत्यासीत् । अत्रास्य स्मृतौ अद्यापि वैष्णवमेलकं भवति । स्वकीये काव्ये प्रसन्नराघवे स आत्मानं  विदर्भदेशस्य कुण्डिनपुरनिवासिनमाख्याय कौडिन्यमकथयत्, यथा—

\begin{shloka}[center]

\textbf{विलासौ यद्वाचाममररसनिष्पन्द–मधुरः}

\textbf{कुरङ्गाक्षी विम्बाधरमधुरभावां गमयति ।}

\textbf{कवीन्द्रः कौडिन्यः स तव जयदेव श्रवणयो–}

\textbf{रयासीदातिथ्यं न किमिहमहादेवतनयः ।।}

\end{shloka}

चतुर्दशशताब्द्यां जातो विश्वनाथः स्वकीये साहित्यदर्पणे प्रसन्नराघवस्य ‘कदली कदली करभः करभे’ति श्लोकमेकमुदाहरणरूपेणोल्लिखिति । तथा च काव्यप्रकाशस्य खण्डनं कुर्वन् असौ कथयति यत्—

\begin{shloka}[center]

\textbf{अङ्गीकरोति यः काव्यं शब्दार्थावनलङ्कृती ।}

\textbf{असौ न मन्यते कस्मादनुष्णमनलङ्कृती ।। (चन्द्रालोकः)}

\end{shloka}

अतो जयदेवो मम्मटविश्वनाथयोरन्तरवर्त्ती भवितुमर्हति । आचार्योऽयं अलङ्कारसर्वस्वकारेण सह पूर्णतया परिचित आसीत् । रुय्यक एव सर्वप्रथमं विचित्रविकल्पाख्ययोरलङ्कारयोः कल्पनामकरोत् । जयदेवोऽपि चन्द्रालोके इमौ अलङ्कारौ यथा स्याताम् तथा वर्णयति । तस्मात् जयदेवस्य समय त्रयोदशशताब्दस्यावसाने मन्तुं युज्यते ।

\customparagraph{\textbf{काव्यशास्त्रीयं योगदानम्}}

अस्यालङ्कारशास्त्रीयो ग्रन्थः चन्द्रालोकाख्योऽस्ति । नाम्ना एव न अपितु कर्मणापि ग्रन्थोऽयं चन्द्रस्यालोकवत् सुशीतलानन्दवितरणपुरस्सरं पाठकानां हृदि अभूतपूर्वमानन्दमाविष्करोति । कथमिति चेत्— कठिनसूत्रात्मककाव्यशास्त्रीयाध्ययनेन शिथिलमानसिकावस्थायुताः पाठका एकस्मिन्नेव अनुष्टुप् छन्दसि निर्मिते पद्येऽलङ्कारस्य लक्षणोदाहरणयोः सरलतया संयोजनमधीत्य काव्यशास्त्रज्ञानमाप्नुवन्ति । एतादृशम् अलङ्काराणां विवेचनं काव्यशास्त्रीयेतिहासे दुर्लभमेव । उदाहरणार्थमत्रोपमालङ्कारस्य लक्षणोदाहरणमुल्लिख्यते, यत्—

\begin{shloka}[center]

\textbf{उपमा यत्र सादृश्यलक्ष्मीरुल्लसति द्वयोः ।}

\textbf{हृदये खेलतोरुच्चैस्तन्वङ्गीस्तनयोरिव ।। (चन्द्रालोकः)}

\end{shloka}

अस्यैव अलङ्कारवर्णनशैल्या लोकप्रियतामवलोक्य पश्चाद्  अप्पयदीक्षितोऽपि स्वकीये कुवलयानन्देऽस्यानुसरणं करोति, इति स एव कुवलयानन्दे स्वीकरोति यत्—

\begin{shloka}[center]

\textbf{चन्द्रालोको विजयतां शरदागमसम्भवः ।}

\textbf{हृद्यः कुवलयानन्दो यत्प्रसादादभूदयम् ।।}

\end{shloka}

एतादृशैर् गुणैर्विशिष्टे चन्द्रालोके दशमयूखाः सन्ति । तेषु त्रिशतोपरिपञ्चाशच्छ्लोका  अनुष्टुप्छन्दसि सन्ति । अत्र सरला, सुगमा च भाषाशैली विद्यते । अस्य प्रथमे मयूखे— काव्यलक्षणस्य, काव्यहेतोः, शब्दभेदानाञ्च विवेचनमस्ति । द्वितीये— दोषाणां समीक्षणम्, तृतीये— लक्षणाख्यस्य काव्याङ्गस्य विवेचनं विद्यते । चतुर्थे मयूखे— दशानां गुणानां सविस्तारविवेचनं कृतं वर्तते । पञ्चमे मयूखे— पञ्चानां शब्दालङ्काराणाम्, शतसङ्ख्यकानाञ्चार्थालङ्काराणां विवेचनं कृतमनेन आचार्येण । षष्ठे मयूखे— रसानाम्, भावानाम्, रीतीनाम्, पञ्चवृत्तीनाञ्च समीक्षणं प्राप्यते । सप्तमे— व्यञ्जनायाः, ध्वनिकाव्यस्य च सभेदं निरूपणमस्ति । अष्टमे मयूखे— गुणीभूतव्यङ्ग्यकाव्यस्य प्रकारसहितं विवेचनं लभ्यते । अवसाने स्थितयोद्वयोर्मयूखयोर्लक्षणायाः, अभिधायाश्च वर्णनमास्ते ।

\customparagraph{\textbf{अलङ्कारविचारः}}

चन्द्रालोके वर्णितेषु विषयेषु जयदेवस्यालङ्कारपक्षे विशेषा रुचिर् दृश्यते । तस्मादयं धुरन्धर  आलङ्कारिको मन्यते । आचार्योऽयं कथयति यद् — यथा औष्ण्येन विना अग्नेः कल्पना अपि न भवति, तथैव अलङ्कारेण विना काव्यस्य निर्माणमसम्भवमेव । यद्यपि जयदेवः प्राग्वर्णितेषु काव्यतत्त्वेषु सानुकूला धारणा प्रस्तौति तथापि प्राधान्येनायं काव्यस्यानिवार्यतत्त्वरूपेण अलङ्कारतत्त्वमेव स्वीकरोति, यत्—

\begin{shloka}[center]

\textbf{अङ्गीकरोति यः काव्यं शब्दार्थावनलङ्कृती ।}

\textbf{असौ न मन्यते कस्मादनुष्णमनलङ्कृती ॥ (चन्द्रालोकः)}

\end{shloka}

\customparagraph{\textbf{काव्यलक्षणम्}}

जयदेवस्य काव्यलक्षणं वाग्भटवत् सर्वतत्त्वसमन्वितं विद्यते । अयमपि प्राग्जातानां काव्यशास्त्रिणां काव्यलक्षणं परिष्कृत्य सुस्पष्टं काव्यलक्षणं प्रस्तौति, यथा-

\begin{shloka}[center]

\textbf{निर्दोषा लक्षणवती सरीतिर्गुणभूषणा ।}

\textbf{सालङ्काररसाऽनेकवृत्तिर्वाक् काव्यनामभाक् ॥ (चन्द्रालोकः)}

\end{shloka}

त्रयोदशशताब्दस्य अवसाने आविर्भूतः आचार्यो जयदेवः (पीयूषवर्षः) संस्कृतकाव्यशास्त्रस्य अलङ्कारप्रधाने इतिहासे देदीप्यमानं स्मरणीयं च स्थानं बिभर्ति। महादेव-सुमित्रयोरात्मजेन अनेन विरचितः \textbf{`चन्द्रालोकः'} इत्याख्यो ग्रन्थः सरलसुबोधशैलीमवलम्ब्य कठिनकाव्यशास्त्रीयसिद्धान्तान् सरलीकृत्य पाठकानां हृदये अलङ्कारज्ञानसौकर्यं जनयति । अस्य अभूतपूर्वां लोकप्रियतां दृष्ट्वा परवर्त्ती अप्पयदीक्षितः `कुवलयानन्दः' इत्याख्यस्य अलङ्कारग्रन्थस्य रचनां कृतवान्। जयदेवः काव्ये अलङ्कारतत्त्वस्य परमसमर्थकः आसीत् । मम्मतस्य मतं खण्डयन् सः दृढतापूर्वकं प्रतिपादयति यत्- अलङ्कारहीनां वाचमेव काव्यं मन्यमानो जनोऽग्निम् उष्णतारहितं किमर्थं न स्वीकरोति ? अनेन प्रतिपादितं काव्यलक्षणं दोषराहित्य-गुण-रीति-रस-अलङ्कारादि-सर्वतत्त्वसमन्वितं वर्तते। अतः सिद्धान्तानां सरलीकरणेन, अलङ्कारस्य अनिवार्यत्वप्रतिपादनेन, अनुपमविवेचनशैल्या च जयदेवः संस्कृतकाव्यशास्त्रपरम्परायां सर्वैरादृतः काव्यशास्त्री वर्तते।

\specialupakhanda{\textbf{विद्याधरः}}

संस्कृतकाव्यशास्त्रे ‘एकावली’ इति ग्रन्थस्य प्रणेतुर्विद्याधरस्य नाम स्मरणीयं विद्यते । अयं स्वकीये ग्रन्थे सर्वाण्येवोदाहरणानि स्वरचितानि प्रस्तौति । ग्रन्थोऽयं स्वाश्रयदातुरुत्कलाधिपस्य राज्ञो नरसिंहस्य तोषणार्थं लिखितो वर्तते । अस्य समयो निश्चेतुं सारल्यं भवति । आचार्योऽयं स्वग्रन्थे रुय्यकस्योल्लेखं चकार । अत एवास्य कालस्योत्तरावधिः द्वादशशताब्दो वर्तते ।

\customparagraph{\textbf{काव्यशास्त्रीयं योगदानम्}}

अस्य एकावलीति एक एव लाक्षणग्रन्थो लभ्यते । तत्राष्टौ उन्मेषाः सन्ति । तेषु क्रमशः काव्यस्य स्वरूपम्, प्रवृत्तेर्विचारः, ध्वनेर्भेदः, गुणीभूतव्यङ्ग्यम्, गुणाः, रीतयः, दोषाः, शब्दालङ्काराः, अर्थालङ्काराश्च वर्णिताः सन्ति । ध्वनिसम्प्रदायस्याचार्योऽयम् । अस्य ग्रन्थस्योपजीव्यो ग्रन्थः काव्यप्रकाशोऽस्ति । अयं पूर्वेषामालङ्कारिकाणां शास्त्रीयमतानि स्वतर्कशक्त्या सुदृढतया समर्थयति ।

\customparagraph{\textbf{काव्यलक्षणम्}}

अस्य काव्यलक्षणं भामहानुरूपं शब्दप्रधानं विद्यते, यथा-

\begin{shloka}[center]

\textbf{शब्दार्थवपुस्तावत् काव्यमिति । (एकावली)}

\end{shloka}

\specialupakhanda{\textbf{विद्यानाथः}}

आचार्योऽयं प्रतापरुद्रयशोभूषणस्य रचयिता अस्ति । अनेन प्रतापरुद्रप्रशस्तये ग्रन्थोऽयं रचितः । तद्यथा—

\begin{shloka}[center]

\textbf{प्रतापरुद्रदेवस्य गुणानाश्रित्य निर्मितः ।}

\textbf{अलङ्कारप्रबन्धोऽयमन्तकरणोत्सवाय वः ।।}

\end{shloka}

प्रतापरुद्रः काकतीयवंशस्य अतीव साहसी राजा आसीत् । नृपोऽयं यादववंशीयं सेवणाख्यं नरेशं पराजितवान् । प्रतापरुद्रो विद्यानाथस्याश्रयदाता आसीत् । प्रतापरुद्रस्य समयः १२९८ ईसवीतः १३१७ ईसवीयं यावदास्ते । अतो विद्यानाथस्यापि अयमेव समयोऽस्ति ।

\customparagraph{\textbf{काव्यशास्त्रीयं योगदानम्}}

अस्य काव्यशास्त्रीयो ग्रन्थः प्रतापरुद्रयशोभूषणाख्योऽस्ति । प्रतापरुद्रस्य संस्तुतौ रचितस्यास्य ग्रन्थस्योदाहरणानि एतेन स्वयमेव प्रतापरुद्रस्य यशोवर्णने रचितानि सन्ति । ग्रन्थस्यास्य तृतीयेऽध्यायेऽलङ्काराणामुदाहरणे तु स्वयमेव प्रतापकल्याणनामकं नाटकं विरचय्य  तस्यैव श्लोकास्तत्र प्रस्तुताः सन्ति । ग्रन्थेऽस्मिन् नव प्रकरणानि सन्ति । तेषु क्रमशः नायकस्य, काव्यस्य, नाटकस्य, रसस्य, दोषाणाम्, गुणानाम्, शब्दालङ्काराणाम्, अर्थालङ्काराणाञ्च विवेचनं प्राप्यते । विद्यानाथेन अलङ्कारेतरविषये मम्मटस्य काव्यप्रकाशस्यानुकरणं कृतम्, अलङ्कारविषये त्वनेनालङ्कारसर्वस्वस्य मतमङ्गीकृतम् । अतोऽत्र रुय्यकस्य ग्रन्थानुसारेणायं परिणामोल्लेखविचित्रविकल्पानामलङ्काराणां विवेचनं करोति । काव्यप्रकाशे एतेषां चतुर्णामलङ्काराणां वर्णनं नास्ति ।

आचार्यस्यास्य काव्यशास्त्रीयं योगदानं पूर्वेषामाचार्याणां विचारसुदृढीकरणार्थमेव उन्मुखमस्ति । अनेनापि काव्यशास्त्रीयेतिहासप्रवर्धने किञ्चिद्  अवदानं तु कृतमेव ।

\customparagraph{\textbf{काव्यलक्षणम्}}

विद्यानाथोऽपि प्राग्जाताचार्यसरण्यामेव स्वकीयं काव्यलक्षणं प्रस्तौति । अस्य मते दोषरहितौ गुणालङ्कारसमन्वितौ शब्दार्थावेव काव्यं भवति । तच्च गद्यपद्यशैल्यां विरच्यते, तद्यथा-

\begin{shloka}[center]

\textbf{गुणालङ्कारसहितौ शब्दाथौ दोषवर्जितौ ।}

\textbf{गद्योपद्योभयमयं काव्यं काव्यविदो विदुः ॥ (प्रतापरुद्रयशोभूषणम्)}

\end{shloka}

\specialupakhanda{\textbf{कविराजो विश्वनाथः}}

काव्यशास्त्रीयेतिहासस्य उत्तरकाले सञ्जातेषु लक्षणशास्त्रिषु कविराजो विश्वनाथः काव्यशास्त्रमर्मज्ञः, मूर्धन्यविद्वान्, रसवादी च मन्यते। अनेन प्रथमतो मम्मटकृत-काव्यप्रकाशस्य ‘दर्पणः’ आख्या टीका लिखिता। टीकालेखनकारणाद् अनेन काव्यशास्त्रीयविषयेषु गभीरं ज्ञानं लब्धम्। टीकालेखनसमये काव्यप्रकाशकारस्य मम्मटस्य कतिपयविचारेषु अस्य मतवैभिन्न्यं सञ्जातम्। अतोऽनेन स्वकीयस्वतन्त्रप्रतिभया ‘साहित्यदर्पणः’ नामकं लाक्षणिकग्रन्थं विरचय्य मम्मटमतमतिरिच्य स्वकीयं मौलिकं मतं प्रदर्शितम्।

\customparagraph{\textbf{वंशपरिचयः स्थितिकालश्च}}

‘कविराजः’ इत्युपाधिनाऽलङ्कृतोऽयम् आचार्य उत्कलप्रदेशस्थिते समृद्धे पण्डितकुले समुद्भूत आसीत्। अस्य पूर्वजा अपि काव्यशास्त्रज्ञा आसन्। अस्य प्रपितामहो नारायणः, पितामहानुजः चण्डीदासः, पिता च महाकविः चन्द्रशेखर आसीत्। एतत् स्वयमेव साहित्यदर्पणे सादरं प्रस्तौति, यत्—

\begin{shloka}[center]

\textit{\textbf{श्रीचन्द्रशेखरमहाकविचन्द्रसूनु-}}

\textit{\textbf{श्रीविश्वनाथकविराजकृतं प्रबन्धम्।}}

\textit{\textbf{साहित्यदर्पणममुं सुधियो विलोक्य}}

\textit{\textbf{साहित्यतत्त्वमखिलं सुखमेव वित्तः।।}} (साहित्यदर्पणः १.२)

\end{shloka}

विश्वनाथः पिता चन्द्रशेखरश्च कलिङ्गनरेशस्य राजसभायां ‘सन्धिविग्रहिकः’ इति गरिमामये पदे प्रतिष्ठितौ आस्ताम्। अयं महानुभावो वैष्णवमतानुयायी आसीत्। अस्य पाण्डित्येन प्रभाविताः समकालिका राजानो विद्वांसश्च तम् \textbf{महापात्र-कविसूक्तिसुधाकर-आलङ्कारिकचक्रवर्त्ति}-प्रभृतिभिः समुत्कृष्टोपाधिभिर्विभूषितवन्तः।

अस्य स्थितिकालविषये ग्रन्थान्ते स्वोल्लेखेन ज्ञायते यत्- अयं स्वग्रन्थे जयदेवस्य \textit{गीतगोविन्दस्य}, श्रीहर्षस्य \textit{नैषधीयचरितस्य}, रुय्यकस्य \textit{अलङ्कारसर्वस्वस्य} च श्लोकान् उदाहरणरूपेण उद्धरति। एते सर्वे द्वादश-त्रयोदशशताब्द्या आचार्याः सन्ति। अपरतः, पञ्चदशशताब्द्यां सञ्जातस्य मल्लिनाथस्य पुत्रः कुमारस्वामी स्वकीये ‘रत्ननारायण’ टीकायां विश्वनाथस्य उल्लेखं करोति। अतः एतयोर्मध्यवर्ती कालः अर्थात् चतुर्दशशताब्दी (१३००-१३८० ई.) एव विश्वनाथस्य प्रामाणिकः समयो मन्यते।

\customparagraph{\textbf{काव्यशास्त्रीयं साहित्यिकञ्च योगदानम्}}

आचार्यविश्वनाथः केवलं लक्षणग्रन्थकार एव नासीत्, अपितु एष उच्चकोटिकः कविरपि आसीत्। तेन राघवविलासाख्यं महाकाव्यम्, कुवलयाश्वचरितं काव्यम्, प्रभावतीपरिणाख्या नाटिका, चन्द्रकला नाटिका, प्रशस्तिरत्नावली, काव्यप्रकाशदर्पणम् (काव्यप्रकाशटीका), नरसिंहविजयाख्यं काव्यम्, साहित्यदर्पणाख्यो लाक्षणिकग्रन्थश्चेति विविधविषयसम्बद्धा ग्रन्था लिखिताः । एतेषु काव्यशास्त्रसम्बन्धिग्रन्थः साहित्यदर्पणोऽस्ति । अत्र संस्कृतसाहित्ये प्रचलितानां सर्वेषां विधानां सुबोध्यशैल्या शास्त्रीयसमीक्षा विद्यते । तस्मात् सम्पूर्णकाव्यशास्त्रस्य प्रतिनिधिभूतोऽयं ग्रन्थः समग्रे विद्वत्समाजे समादरणीयोऽस्ति ।

साहित्यदर्पणे दश परिच्छेदाः सन्ति, येषु काव्यशास्त्रस्य साङ्गोपाङ्गं विवेचनं प्राप्यते-

\begin{itemize}
\item प्रथमपरिच्छेदे काव्यलक्षणकाव्यप्रयोजनादिनां काव्याड्गानां विवेचनमास्ते। अत्रैव मम्मट-भामह-आनन्दवर्धनादीनां काव्यलक्षणानां तीव्रं खण्डनं विद्यते।


\item द्वितीयपरिच्छेदे वाक्य-पदयोः स्वरूपं निरूप्य शब्दशक्तीनाम् अभिधा-लक्षणा-व्यञ्जनानां सविस्तारं समीक्षणं विहितम्।


\item तृतियपरिच्छेदे काव्यात्मभूतस्य रसस्य निष्पत्तिः, नव रसाः, एकोनपञ्चाशद्भावाः, नायक-नायिकानां भेदाश्च सविस्तरेण वर्णिताः। अयं परिच्छेदो ग्रन्थस्य मुख्यो विशालतमश्चास्ति ।


\item चतुर्थपरिच्छेदे ध्वनिवादानुसारं ध्वनिकाव्यस्य (उत्तमकाव्यस्य) गुणीभूतव्यङ्ग्यकाव्यस्य (मध्यमकाव्यस्य) च भेदप्रभेदा निरूपिताः।


\item पञ्चमपरिच्छेदे व्यञ्जनावृत्तेरनिवार्यतां प्रदर्श्य, मीमांसकादीनां विरोधिमतानां सतर्कं खण्डनमण्डनं कृत्वा व्यञ्जनायाः स्थापना विहिता।


\item षष्ठपरिच्छेदे सर्गबन्धादीनां श्रव्यकाव्यभेदानाम्, दशरूपक-उपरूपकाणा| दृश्यकाव्यभेदानाम्, नाटकस्याङ्गभूतानां सन्धिसन्ध्यङ्गादीनाञ्च विस्तृतं विवेचनं कृतम्। दृश्यकाव्यदृष्ट्या अयं परिच्छेदोऽतीव महत्त्वपूर्णोऽस्ति।


\item सप्तमपरिच्छेदे रसस्यापकर्षकाणां काव्यदोषाणां सारल्येन समीक्षणं विहितम्।


\item अष्टमपरिच्छेदे माधुर्य-ओज-प्रसादानां काव्यगुणानां सोदाहरणं स्वरूपं प्रदर्शितम्।


\item नवमपरिच्छेदे रीतिसिद्धान्तानुसारं वैदर्भी-गौडी-पाञ्चाली-लाटीति रीतिभेदानां लक्षणोदाहरणसहितं विश्लेषणं कृतम्।


\item दशमे परिच्छेदे अनुप्रास-उपमा-रूपकादीनां शब्दार्थालङ्काराणां सभेदं सोदाहरणञ्च विशदं विवेचनं विहितम्।


\end{itemize}
\customparagraph{\textbf{काव्यलक्षणम्}}

कविराजो विश्वनाथो मम्मटस्य ‘तददोषौ शब्दार्थौ...’ इति लक्षणं दोषपूर्णं प्रदर्श्य स्वकीयं काव्यलक्षणं प्रस्तौति— \textbf{‘वाक्यं रसात्मकं काव्यम्’} इति। आचार्योऽयं रसवादी अस्ति। अस्य आशयोऽस्ति यत् रसास्वादयुक्तं वाक्यमेव काव्यपदवाच्यं भवति, रसविहीनं वाक्यं तु केवलं वार्तामात्रमेव।

\customparagraph{\textbf{काव्यप्रयोजनम्}}

काव्यप्रवृत्तिं प्रति पाठकान् प्रेरयितुं विश्वनाथः सरलं प्रयोजनं वर्णयति, यत्-

\begin{shloka}[center]

\textit{\textbf{चतुर्वर्गफलप्राप्तिः सुखादल्पधियामपि।}}

\textit{\textbf{काव्यादेव यतस्तेन तत्स्वरूपं निरूप्यते।। (साहित्यदर्पणः १.२)}}

\end{shloka}

अर्थात् धर्म-अर्थ-काम-मोक्ष इति चतुर्वर्गफलस्य प्राप्तिः शास्त्रैस्तु कठिना भवति, परन्तु काव्यात् सुकुमारमतीनां राजकुमारादीनां मन्दबुद्धीनाञ्च सुखपूर्वकं जायते।

\customparagraph{\textbf{काव्यभेदाः}}

विश्वनाथः पूर्ववर्तिनाम् उत्तम-मध्यम-अधमख्यं त्रिविधकाव्यभेदं तिरस्कृत्य, अधमं (चित्रकाव्यम्) काव्यमेव न मनुते। यतो यत्र प्रतीयमानो रसो नास्ति, तत् काव्यं भवितुं नार्हति। अतो विश्वनाथः काव्यस्य द्वावेव भेदौ स्वीकरोति ।

\begin{shloka}[center]

\textit{\textbf{काव्यं ध्वनिर्गुणीभूतव्यङ्ग्यं चेति द्विधा मतम्।}} (साहित्यदर्पणः ४.१)

\begin{itemize}
\item ध्वनिकाव्यम् (उत्तमम्): यत्र वाच्यार्थापेक्षया व्यङ्ग्यार्थः प्रधानो भवति।


\item गुणीभूतव्यङ्ग्यम् (मध्यमम्): यत्र व्यङ्ग्यार्थो वाच्यार्थापेक्षया गौणो भवति।


\end{itemize}
\end{shloka}

\customparagraph{\textbf{रसस्वरूपम्}}

विश्वनाथोऽभिनवगुप्तस्य रसस्वरूपं दृढतया स्वीकृत्य ब्रूते यत्- सहृदयानां हृदये स्थितो रत्यादिस्थायीभाव एव विभावादिभिरभिव्यक्तो भूत्वा रसरूपतां प्राप्नोति ।

\begin{shloka}[center]

\textit{\textbf{विभावेनानुभावेन व्यक्तः सञ्चारिणा तथा।}}

\textit{\textbf{रसतामेति रत्यादिः स्थायीभावः सचेतसाम्।।}} (साहित्यदर्पणः ३.१)

\end{shloka}

\customparagraph{\textbf{अलङ्कार-गुण-रीतयः}}

विश्वनाथः अलङ्कार-गुण-रीतीनां सम्बन्धं रसस्य आधारेण मानवशरीरेण सह तोलयित्वा स्पष्टं करोति-

\begin{itemize}
\item अलङ्कारः - शब्दार्थयोः अस्थिराः धर्माः सन्ति, ये शरीरस्थ कटक-कुण्डलादिवत् रसस्य शोभां वर्धयन्ति। \textit{\textbf{(अलङ्कारास्तेऽङ्गदादिवत्)}}


\item गुणः - रसस्य आत्मभूतस्य स्थायिधर्माः सन्ति, यथा मानवस्य शौर्यादयः गुणाः।             \textit{\textbf{(रसस्याङ्गित्वमाप्तस्य धर्माः शौर्यादयो यथा गुणाः...)}}


\item रीतिः - पदसङ्घटनारूपा रीतिर्मानवशरीरस्य अङ्गसंस्थानवद् रसस्य उपकारं करोति। \textit{\textbf{(पदसङ्घटना रीतिरङ्गसंस्थाविशेषवत्...)}}


\end{itemize}
कविराजो विश्वनाथः संस्कृतकाव्यशास्त्रपरम्परायां सर्वैः काव्यशास्त्रजिज्ञासुभिरादृतो विद्यते। तस्य ‘साहित्यदर्पणः’ न केवलं प्राचीनसिद्धान्तानां सङ्ग्रहग्रन्थः, अपितु सरल-सुबोध-प्राञ्जलशैल्यां लिखितः समन्वयवादी ग्रन्थोऽस्ति। मम्मटस्य जटिलतां परिहृत्य, नाट्यशास्त्रस्य काव्यशास्त्रस्य च एकस्मिन्नेव स्थाने सुन्दरं समन्वयम् उपस्थाप्य विश्वनाथेन परवर्तिपाठकानां कृते महान् उपकारः कृतः। अतः काव्यशास्त्रीयपरम्परायाम् अस्य स्थानं समादरणीयं वर्तते।

\specialupakhanda{\textbf{केशवमिश्रः}}

केशवमित्रस्य लाक्षणिकग्रन्थस्य नाम \textbf{अलङ्कारशेखरो}ऽस्ति । ग्रन्थस्यास्य प्रारम्भे अन्त्ये च अलङ्कारशेखराख्यो ग्रन्थो धर्मचन्द्रपुत्रस्य माणिक्यचन्द्रस्यानुरोधेन लिखितः इत्युल्लेखो लभ्यते । धर्मचन्द्रो रामचन्द्राख्यस्य राज्ञः पुत्र आसीत् । स काबुलदेशस्य मुगलनृपमजयत् । कनिङ्‌घममतानुसारेण काङ्‌गडाप्रदेशस्य नृपो माणिक्यचन्द्रो धर्मचन्द्रस्यानन्तरं १५६३ ईशवीये राज्यं प्राप्य दशवर्षाणि यावद् राज्यमकरोत् । तस्मात् केशवमिश्रस्य समयः षोडसशताब्दस्योत्तरार्धः इत्यनुमीयते ।

अस्यालङ्कारशेखराख्ये लाक्षणिकग्रन्थे कारिका, वृत्तिः, उदाहरणश्चेति त्रयो भागाः सन्ति । असौ काव्यादर्शकाव्यप्रकाशकाव्यमीमांशाध्वन्यालोकप्रभृतिभ्यः प्रसिद्धेभ्यो लाक्षणिकग्रन्थेभ्यः सामग्रीं सङ्‌कलय्य काव्यशास्त्रीयतत्त्वानि सारल्येन प्रस्तौति । अलङ्कारशेखरे अष्टौ रत्नानि विंशतिर्मरीचयश्च सन्ति । तत्र काव्यलक्षणम्, रीतिः, शब्दशक्तिः, दोषः, गुणः, अलङ्कारः, रूपकभेदाश्च वर्णिताः सन्ति ।

\specialupakhanda{\textbf{शारदातनयः}}

शारदातनयस्य काव्यशास्त्रीयग्रन्थो भावप्रकाशनम् अस्ति । स्वकीये ग्रन्थे भोजराजस्य शृङ्गारप्रकाशस्य मम्मटस्य काव्यप्रकाशस्य चोल्लेखं करोति । सिंहभूपालोऽस्य मतस्योल्लेखं करोति । सिंहभूपालस्य समयः १३२० ख्रिष्टाब्दः इति निर्धारितो विद्यते । तस्मादस्य समयः द्वादशशताब्दस्योत्तरार्द्धम् अनुमीयते । भावप्रकाशने दश अधिकारा विद्यन्ते । तेषु प्रथमद्वितीययो भावरसयोः स्वरूपवर्णनम् तृतीये रसभेदानां वर्णनम् चतुर्थे पञ्चमे च नायकनायिकयोर्भेदवर्णनम् षष्ठे शब्दार्थसम्बन्धवर्णनम् सप्तमे नाट्येतिहासवर्णनम् अष्टमे दशरूपकवर्णनम् नवमे नृत्यभेदवर्णनम् दशमे नाट्यभेदवर्णनं विद्यते ।

\specialupakhanda{\textbf{शिङ्गभूपालः}}

शिङ्गभूपालो नाट्यशास्त्रस्य सङ्गीतशास्त्रस्य च मर्मज्ञ आसीत् । नाट्यशास्त्रानन्तरं शार्ङ्गदेवस्य सङ्गीतरत्नाकरः सङ्गीतशास्त्रस्य बृहद् ग्रन्थो विद्यते । अस्य ग्रन्थस्य सङ्गीतसुधाकराख्याया टीकायाः प्रणेता शिङ्गभूपालो विद्यते । अनन्तरमनेन रसार्णवसुधाकराख्य नाट्यशास्त्रविषयिणी कृती रचिता ।

शिङ्गभूपालस्य कालविषये मतभेदाः सन्ति । केचित् त्रयोदशशताब्दस्य मध्यभागः, केचित् षोडशशताब्द एवास्य समयो वदन्ति । पूर्वापरप्रसङ्गमधीत्य अस्य समय त्रयोदशशताब्दस्योत्तरार्धः इति अनुमीयते ।

रसार्णवसुधाकराख्यो ग्रन्थः त्रिषु विलासेषु सङ्‌घटितो विद्यते । रञ्जकोल्लासाभिधाने प्रथमे विलासे नायकनायिकयोः स्वरूपम् चतसृणां वृत्तीनां च स्वरूपं वर्णितं विद्यते । रसिकोल्लासाभिधाने द्वितीये विलासे रसस्य साङ्गोपाङ्गवर्णनं विद्यते । भावविलासाख्ये तृतीये विलासे रूपकाणां साङ्गोपाङ्गवर्णनं विद्यते । दशरूपकापेक्षया ग्रन्थोऽयं स्पष्टं विशदञ्च मन्यते ।

\specialupakhanda{\textbf{भानुदत्तः}}

रससम्प्रदायस्य सम्वर्धकोऽयं काव्यशास्त्री रसमञ्जर्या अन्तिमे श्लोके स्वकीयं परिचयं प्रस्तोति यत्

\begin{shloka}[center]

\textbf{तातो यस्य गणेश्वरः कविकुलालङ्कारचूडामणि-}

\textbf{र्देशो यस्य विदेहभूः सुरसरित् कल्लोलकीर्त्या युतः ॥}

\end{shloka}

इत्यनुसारमयं मैथिलः, अस्य पितुर्नाम च गणेश्वर आसीत् । सूचिग्रन्थेष्वपि भानुदत्तो मैथिल इति गणितो विद्यते । अस्य पितु गणेश्वरस्य पुत्र चण्डेश्वरो विवादरत्नाकराख्यां कृतिं प्रणैषित् । १३५३ तमे ईशवीये चण्डेश्वरस्य सुवर्णमयं तुलादानमभूदिति सङ्केतो लभ्यते । असौ स्वकीये ग्रन्थे शृङ्गारतिलकदशरूपकयोर्निर्देशमकरोत् । गोपालाचार्यश्च १४२८ ख्रीष्टाब्दे रसमञ्जर्या विकासाभिधानां टीकामलिखत् । इत्यादि सङ्केतैर्ज्ञायते यत्- त्रयोदशशताब्दस्यस्यान्ते चतुर्दशशताब्दस्यारम्भेऽस्य समयः इति ।

भानुदत्तस्य रसमञ्जरीरसतरङ्गिण्याख्यौ द्वौ ग्रन्थौ विद्येते । रसमञ्जर्यां नायकनायिकयोर्भेदादिकं विस्तरेण वर्णितं विद्यते । रसमञ्जर्याः समधिका टीका लभ्यते । तेषु नागेशस्य प्रकाशाख्या टीका प्रसिद्धा अस्ति । अस्य रसतरङ्गिण्यां रसस्य विस्तृतं वर्णनं विद्यते । तस्मिन् अष्टौ तरङ्गाः सन्ति । तेषु रसानां रसस्याङ्गभूतानां तत्त्वानां विस्तृतं वर्णनं लभ्यते । अस्य च नवसङ्‌ख्यका टीका प्राप्यते ।

\specialupakhanda{\textbf{रूपगोस्वामी}}

पौरस्त्यकाव्यशास्त्रीयेतिहासे कैश्चिद् विद्वद्भिर्वैष्णवभक्तिधारायां भक्तिरसस्य शास्त्रीया चर्चा विहिता । तेषु गौडीयवैष्णवसम्प्रदाये प्रसिद्धेषु षड् गोस्वामिषु अन्यतमस्य श्रीमद्रूपगोस्वामिनः स्थानं सर्वोत्कृष्टं सर्वमान्यञ्च विद्यते । असौ भक्तिसाहित्यप्रवर्तकस्य श्रीकृष्णचैतन्यमहाप्रभोः परमप्रियः शिष्यः, मुकुन्दस्य पौत्रः, कुमारस्य पुत्रः, सनातनगोस्वामिनो भ्राता आसीत् । अस्य समयः पञ्चदशशताब्दस्यान्तः षोडशशताब्दस्य प्रारम्भ (१४८९–१५६४ ख्रिष्टाब्दः) मन्यते । रूपगोस्वामिना नाट्यशास्त्र-काव्यशास्त्र-भक्तिशास्त्राणां अपूर्वं समन्वयं विधाय रससिद्धान्ते एका युगान्तकारिणी क्रान्तिः समानीता । अस्य नाटकचन्द्रिका, भक्तिरसामृतसिन्धुः, उज्ज्वलनीलमणिश्चेति त्रयो रसशास्त्रविषयकाः प्रधानग्रन्थाः प्रकाशिताः समुपलभ्यन्ते ।

\customparagraph{\textbf{नाटकचन्द्रिका}}

ग्रन्थेऽस्मिन् नाटकस्य नाट्यशास्त्रस्य च साङ्गोपाङ्गं मौलिकं वर्णनमास्ते । आचार्यभरतस्य नाट्यशास्त्रम्, शिङ्गभूपालस्य रसार्णवसुधाकरश्चेति ग्रन्थद्वयस्य गहनमध्ययनं कृत्वा वैष्णवनाट्यपरम्परानुकूलो ग्रन्थोऽयं प्रणीतः’ इति स्वयं रूपगोस्वामी नाटकचन्द्रिकायाः प्रारम्भे सादरं समुल्लिखति । विशेषतो विश्वनाथस्य साहित्यदर्पणोक्तनाट्यलक्षणविरुद्धं भरतसम्मतं मतमत्र समर्थितम् । अष्टसु प्रकाशेषु विभक्तेऽस्मिन् ग्रन्थे नाटकलक्षणम्, नायकनायिकाभेदवर्णनम्, रूपकविभागः, नाट्याङ्गानां वैज्ञानिकं विवेचनम्, नाटकस्य अङ्कदृश्यविभाजनम्, भाषाविधानम्, वृत्तिविचारः, वृत्तीनां रसानुकूलप्रयोगप्रभृतयो नाट्यशास्त्रोपयोगिनः सर्वेऽपि विषयाः सूक्ष्मतया वर्णिता सन्ति । अत्र स्वरचितविदग्धमाधवललितमाधवादिवैष्णवग्रन्थेभ्यः प्रायशो भगवल्लीलामयानि उदाहरणानि समुद्धृतानि विद्यन्ते ।

\customparagraph{\textbf{भक्तिरसामृतसिन्धुः}}

नाम्ना एव ग्रन्थोऽयं भक्तिरसस्य अपारसागररूपोऽस्तीति ज्ञायते । मम्मटादिप्राचीनकाव्य- शास्त्रिणां ``भक्तिर्देवादिविषयिका रतिर्भाव एव न तु रसः'' इति मतं सप्रमाणं खण्डयित्वा भक्तिरसस्य अलौकिकप्रधानरसत्वं प्रतिष्ठापयितुं ग्रन्थोऽयं महदुपकरोति । अस्य रचनासमयः १५४१ ख्रिष्टाब्दो मन्यते । ग्रन्थोऽयमाध्यात्मिकसाहित्यिकसौन्दर्यशास्त्रीयविषययोरनुपमेय आकरग्रन्थो मन्यते । भक्तिरसामृतसिन्धुः पूर्वविभागः, दक्षिणविभागः, पश्चिमविभागः, उत्तरविभागश्चेति चतुर्षु प्रधानेषु विभागेषु विभक्तोऽस्ति । एते विभागाः पुनरनेकासु लहरिकासु विभक्ताः सन्ति ।

\begin{itemize}
\item \textbf{प्रथमविभाग:} – अत्र भक्तेः सामान्यलक्षणम्, तस्याः षड्विधं वैशिष्ट्यम्; साधनभक्तिः, भावभक्तिः, प्रेमाभक्तिश्चेति भक्तेः क्रमिकविकासात्मका भेदाः सोदाहरणं वर्णिताः सन्ति ।


\item \textbf{दक्षिणविभाग:} – अत्र रससामग्रीणाम् अर्थाद् विभाव-अनुभाव-सात्त्विकभाव- व्यभिचारिभावानां स्थायिभावस्य च शास्त्रीयवर्णनानन्तरं मुख्यभक्तिरसस्य सामान्यविवरणं प्रदत्तमास्ते ।


\item \textbf{पश्चिमविभाग:} – अत्र मुख्यभक्तिरसस्य विशिष्टं प्रपञ्चनम् अस्ति । तत्र विविधासु लहरिकासु क्रमशः शान्तभक्तिरसः, प्रीतभक्तिरसः (दास्यम्), प्रेयोभक्तिरसः (सख्यम्), वत्सलभक्तिरसः, मधुरभक्तिरसः (शृङ्गारः) प्रभृतीनां पञ्चानां मुख्यभक्तिरसानां साङ्गोपाङ्गसहितं निरूपणमास्ते । रसनिरूपणप्रसङ्गे रूपगोस्वामी दृढतया प्रस्तौति यत्- मूलरसस्तु श्रीकृष्णविषयिणी भक्तिरेव, अन्ये लौकिका रसास्तु तस्यैव विकृतयः सन्तीति ।


\item \textbf{उत्तरविभाग:} –अत्र भक्तिरसविकृतिभूतानां गौणभूतानाञ्च सप्त रसानां वर्णनं विद्यते । ते च हास्य-अद्भुत-वीर-करुण-रौद्र-भयानक-बीभत्सप्रभृतयः सन्ति । ग्रन्थान्ते प्राग्वर्णितानां रसानां परस्परं मैत्रिविरोधयोः सूक्ष्मं शास्त्रीयविवेचनं वर्तते ।


\end{itemize}
\customparagraph{\textbf{उज्ज्वलनीलमणिः}}

ग्रन्थोऽयं भक्तिरसामृतसिन्धोः परिशिष्टत्वेन पूरकत्वेन रचितोऽस्ति, यत्र भक्तिरसामृतसिन्धोः पश्चिमविभागे वर्णितस्य सर्वोत्कृष्टस्य `मधुरभक्तिरसस्य' सर्वाङ्गविवेचनं विद्यते । ग्रन्थेऽस्मिन् लौकिकशृङ्गाराद् विलक्षणः `उज्ज्वलः' इति शब्दं प्रयुज्य तस्यार्थोऽलौकिक-दिव्य-मधुर-भक्तिरस इत्युल्लिख्य तस्यैव परा अवस्था प्रदर्शिता विद्यते । अत्र नायकशिरोमणेः श्रीकृष्णस्य, नायिकाशिरोमण्याः श्रीराधायाः स्वरूपं संवर्ण्य, राधायाः कायव्यूहरूपाणां सखीनां वर्णनम्, कृष्णसखानाञ्च सविस्तरं वर्णनमास्ते । अनन्तरं मधुररसस्य विभाव-अनुभाव-सात्त्विक- व्यभिचारि-स्थायिभावानां गहनं विश्लेषणं विधाय, विप्रलम्भसम्भोगयोरलौकिकप्रकारद्वयोः साङ्गोपाङ्गवर्णनं प्राप्यते । एतादृशो विषयवस्तुसंवलितोऽयं ग्रन्थो भक्तिरसवर्णनदृष्ट्या अद्वितीयो विशालकायात्मकश्चास्ति । ग्रन्थोऽयं भक्तेरसस्य, अलङ्कारशास्त्रस्य, नाट्यसाहित्यस्य च दृष्ट्या सर्वत्र श्लाघनीयो ग्रन्थो मन्यते । अस्य ग्रन्थस्य जीवगोस्वामिकृता `लोचनरोचनी' तथा विश्वनाथचक्रवर्त्तिकृता `आनन्दचन्द्रिका' चेति द्वे टीके सुप्रसिद्धे स्तः ।

श्रीमद्रूपगोस्वामिविरचिता इमे ग्रन्थाः केवलं सम्प्रदायविशेषस्य धार्मिकसाहित्यमेव न, अपितु संस्कृतकाव्यशास्त्रस्य नाट्यशास्त्रस्य च इतिहासे समादरणीयाः सन्ति । तेन प्राचीनालङ्कारिकाणां शृङ्गारादिरससीमामतिरिच्य `भक्तिरसः' एव ब्रह्मानन्दसहोदरः, सर्वानन्दचमत्कारकारी, परमपुरुषार्थमयश्चेति सुदृढत्वेन संस्थापितः । अतो रससिद्धान्तस्य व्यापकदृष्ट्या, नाट्यलक्षणस्य परिष्कारदृष्ट्या, आध्यात्मिकसौन्दर्यशास्त्रस्य प्रतिष्ठापनदृष्ट्या च रूपगोस्वामिनोऽवदानं संस्कृतवाङ्मयेऽविस्मरणीयम्, आदरणीयञ्च विद्यते ।

\specialupakhanda{\textbf{अप्पयदीक्षितः}}

आचार्योऽयं दक्षिणभारतस्य ग्रन्थकारेषु विशिष्टो वर्तते । संस्कृतकाव्यशास्त्रस्योत्तरवर्त्तिसमयेऽऽलङ्कारिकसमाजेऽयं समादरणीयोऽस्ति । अयं वेङ्कटपतिनाम्नो नृपाश्रये वसति स्म । तद्यथा कुवलयानन्दे —

\begin{shloka}[center]

\textbf{अमुं कुवलयानन्दमकरोदप्पयदीक्षितः ।}

\textbf{नियोगाद् वेङ्कटपतेर्निरुपाधिकृपानिधेः ।। (कुवलयानन्दः)}

\end{shloka}

दक्षिणभारते वेङ्कटपतिर्नामकयोर्द्वयोर्नृपयोः  स्थितिर्लभ्यते । एगेलिनमहोदयस्य मते विजयनगरस्य वेङ्कटपतिरेवास्याश्रयदाताऽऽसीत् । अस्य समयः १५३५ ईसवीयो मन्यते परं हुल्त्समहोदयस्य मते तु अयं पेन्नकोण्डाया नृपतेर्वेङ्कटपतिद्वितीयस्याश्रयेऽयमासीदिति । अस्य शासनकालश्च १५८५ ईसवीयतः १६१३ ईसवीयपर्यन्तं मन्यते । एतैः प्रमाणैर् दीक्षितस्य समयः षष्ठदशशताब्दस्योत्तरार्धारभ्य सप्तमशतकस्य प्रारम्भपर्यन्तं स्वीकर्तुं युज्यते । अप्पयदीक्षितानुजस्याच्चनस्य पौत्रः नीलकण्ठदीक्षितः शिवलीलार्णवे लिखति यत् ‘दीक्षितोऽयं द्विसप्ततिवर्षान्तमजीव्यत्, शतसङ्ख्यकाश्च ग्रन्था अरचयत् ।’ दीक्षितः स्वस्यां शिवादित्यमणिदीपिकायां स्वपितुर्नाम रङ्गराजः, पितामहश्च आचार्यः दीक्षितः, इति अलिखत् । अस्योत्तरवर्त्तिना जगन्नाथेन सह साम्प्रदायिको विवाद आसीत् । यवनकन्यया सह सम्बन्धस्य कारणात् दीक्षितो जगन्नाथं सतिरस्कारं जातिच्युतमकरोत् । अत एव पण्डितराजो जगन्नाथोऽस्य तीव्रमालोचनं करोति । रसगङ्गाधरस्य स्थाने स्थाने तथा चित्रमीमांसाखण्डनाख्यं ग्रन्थं विरचय्य  स्वकीयं प्रतिशोधं पूरयति जगन्नाथः । पश्चात् स्वानुजस्य पौत्रो नीलकण्ठदीक्षित चित्रमीमांसादोषधिक्काराख्यं ग्रन्थं विरचय्य  जगन्नाथमतखण्डनं करोति ।

अनेन विविधविधासम्बद्धाः शताधिका  ग्रन्था विरचिता इति सङ्केतः प्राप्यते, परन्तु अधुना तेषां सर्वेषामुपलब्धिर्नास्ति । तथापि अद्वैतवेदान्तविषयस्य, पूर्वमीमांसाविषयस्य, व्याकरणविषयस्य, भक्तिविषयस्य, काव्यशास्त्रविषयस्य च ग्रन्था लभ्यन्ते । तेषु विषयप्रसङ्गेनात्र काव्यशास्त्रविषयस्यैव चर्चा क्रियते ।

\customparagraph{\textbf{काव्यशास्त्रीयं योगदानम्}}

काव्यशास्त्रविषयकग्रन्था वृत्तिवार्तिकम्, चित्रमीमांसा, कुवलयानन्दश्चेति त्रयः सन्ति ।

\textbf{वृत्तिवार्तिकम्-} ग्रन्थोऽयं शब्दवृत्तीनां विवेचने लिखितोऽतीव  लघ्वाकारो  विद्यते । द्विपरिच्छेदयुते लघुग्रन्थेऽस्मिन् क्रमशोऽभिधालक्षणयोर्विवेचनं विद्यते ।

\textbf{चित्रमीमांसा}- काव्यालङ्कारसौन्दर्यवर्णने समर्पितोऽयं ग्रन्थ अतिशयोक्त्यलङ्कारपर्यन्तमेव लभ्यते । ग्रन्थस्यान्तेऽवस्थिता इयं कारिका ग्रन्थस्य अपूर्णत्वं सूचयति, यत्—

\begin{shloka}[center]

\textbf{अप्यर्धचित्रमीमांसा न मुदे कस्य मांसला ।}

\textbf{अनूरुरिव घर्मांशोरर्धेन्दुरिव धूर्जटेः ।। (चित्रमीमांसा)}

\end{shloka}

अस्य ग्रन्थस्यालङ्कारविवेचनं प्रति पण्डितराजः सपूर्वाग्रहमिव तीव्रं खण्डनं प्रस्तौति ।

\customparagraph{\textbf{कुवलयानन्दः}}

ग्रन्थोऽयमलङ्कारवर्णने समर्पितोऽस्ति । अस्योपजीव्यं तु जयदेवस्य चन्द्रालोको वर्तते’ इति दीक्षित एवोल्लिखति, यत्—

\begin{shloka}[center]

\textbf{चन्द्रालोको विजयतां शरदागमसम्भवः ।}

\textbf{हृद्यः कुवलयानन्दो यत्प्रसादादभूदयम् ।। (कुवलयानन्दः)}

\end{shloka}

चन्द्रलोकपथमनुसरन् असौ स्वकीये कुवलयानन्दे अनुष्टुप् छन्दसि रचितेषु श्लोकेषु एकपादमलङ्कारस्य लक्षणार्थं तथा अपरपादं तस्यैवोदाहरणार्थं योजयति । तस्मात् स्वल्पाध्ययनेनैव कस्यापि अलङ्कारस्य परिचयं ग्रहीतुं सारल्येन शक्यते । अनेनात्र सर्वेषां पूर्वजानामालङ्कारिकाणां मतसंयोजनेन अलङ्कारसङ्ख्यापि शतोपरिपञ्चविंशतिरधिका  प्रस्तुताः । प्रायः पूर्वजानां केषाञ्चित् स्वकीयप्रतिभोत्पादितानामलङ्काराणां विवेचनयुतोऽयं ग्रन्थो पूर्णतया नवीनस्तु नैव परम् अलङ्काराणां सामान्यज्ञानाय अतीवोपयोगीति निश्चप्रचम् । इदमेवास्य लोकप्रियतायाः कारणमपि ।

विविधविधासु पारङ्गतोऽप्पयदीक्षितः काव्यशास्त्रेऽपि सम्मानितोऽस्ति । यद्यपि अनेन चित्रमीमांसाकुवलयानन्दाख्ययोर्लाक्षणिकग्रन्थयोरलङ्काराणामेव विवेचनं कृतम्, तथापि तेषां सारल्येन समीक्षणमस्य स्मरणीयं वैशिष्ट्यमस्ति ।

\specialupaupakhanda{\textbf{पण्डितराजो जगन्नाथः}}

आचार्योऽयं काव्यशास्त्रीयेतिहासस्योत्कर्षकाले सर्वाधिकपाण्डित्यपूर्णं रसगङ्गाधराख्यं ग्रन्थं विरचय्य संस्कृतवाङ्मयस्य साहित्यशास्त्रीयपक्षं सर्वोत्कृष्टस्थानं प्रापयति स्म । स मुगलशासकस्य शाहजहाँनृपस्य कृपापात्र आसीत् । स एव दिल्लीश्वर इमं पण्डितराजोपाधिना विभूषयति स्म । तत्रैव एतेन सातिशयसुन्दररूपवत्या यवनकन्यया सह पाणिग्रहणमपि कृतवान्, तस्मादेव कारणाद् अयं पण्डितसमाजाद् बहिष्कृतो जातः, इति किंवदन्ती  श्रूयते । अस्याश्रयदातुः शाहजहाँनृपस्य समयः १६२८ ईसवीयतः १६५८ ईसवीयं यावद् मन्यते । तथा च अयमाचार्यो वादविवादखण्डनादीनामाधारे दीक्षितस्य समकालिक आभाति । अतोऽस्यापि समयः स एवेति मन्तुं युज्यते । पण्डितराजो हि तैलङ्गब्राह्मणः । अस्य पिता पेरुभट्टः, माता च लक्ष्मीदेवी आसीत् । अस्य जीवनविषये कर्णपरम्परया श्रूयते यत्— असौ जयपुरे संस्कृतपाठशालामस्थापयत् । स तत्रैकं यवनविद्वांसं शास्त्रार्थे पराजितवान् । पराजयं प्राप्योऽयं काजीति नामधेयो विद्वान् देहलीमेत्य शाहजहाँसमक्षं पण्डितराजस्य विषये प्रशंसापूर्णं प्रतिभावैशिष्ट्यवर्णनं कृतवान् । तेन राज्ञा मौलवीयानामाक्षेपनिराकरणार्थं पण्डितराजं स्वराजसभायामनयत् । तत्र शास्त्रार्थे मौलवीयानां मानमर्दनञ्चाकारयत् । एतेन प्रसन्नो भूत्वा मुगलसम्राट् शाहजहाँ जगन्नाथं पण्डितराजोपाधिना विभूषितवान् । एकदा शाहजहाँनृपस्यान्तःपुरनिवासिन्यां यवनसुन्दर्यां गभीरानुरागेण मुग्धो जगन्नाथो नृपसमक्षं यवनीकन्यामयाचत्—

\begin{shloka}[center]

\textbf{न याचे गजालिं न वा वाजिराजिं न वित्तेषु चित्तं मदीयं कदाचित् ।}

\textbf{इयं सुस्तनी मस्तकन्यस्तहस्ता लवङ्गीकुरङ्गी मदङ्गीकरोतु ।।}

\end{shloka}

एतादृशी प्रार्थना पण्डितराजस्य यवनकन्यया सह स्नेहस्यासामान्यामवस्थां सूचयति । अस्य यवनकन्यां प्रति कीदृशमनुरागमासीदिति स्वयमेव वर्णयति, यत्—

\begin{shloka}[center]

\textbf{यवनीनवनीतकोमलाङ्गी शयनीये यदि नीयते कदापि ।}

\textbf{अवनीतलमेव साधु मन्ये नवनीमाधवनीविनोदहेतुः ।।}

\end{shloka}

तथा च—

\begin{shloka}[center]

\textbf{यवनी रमणी विपदः शमनी कमनीयतमा नवनीतसमा ।}

\textbf{उहि उहि वचोऽमृतपूर्णमुखी स सुखी जगतीह यदङ्कगता ।।}

\end{shloka}

यवनीकन्यया सह गभीरस्नेहस्य पराकाष्ठा विवाहरूपे परिणता । एतस्यैव कारणेन काशीनगर्या विद्वत्सभा एनं सतिरष्कारं बहिष्कारञ्चाकरोत् । बहिष्कारिकायां विद्वद्सभायामप्पयदीक्षितः प्रमुख  आसीत् । अतः पूर्वाग्रहपीडितः पण्डितराजो दीक्षितस्यालङ्कारमतं तीव्रतया खण्डनमकरोदित्यनुमीयते । यवनकन्यया सह विवाहकारणात् काशीनगरे स्थितैः पण्डितैरपमानितो जगन्नाथो मानसिकपीडया विक्षिप्तः सन् स्वरचिताया गङ्गालहर्याः पाठं कुर्वन् गङ्गायामेव लीनो जात इति श्रूयते।

\customparagraph{\textbf{काव्यशास्त्रीयं योगदानम्}}

काव्यशास्त्रस्य विषयेऽस्य वैचारिकदृष्ट्या वर्णनशैल्या चातीव प्रौढाः कृतिर्विद्यते रसगङ्गाधरः । ग्रन्थोऽयं ध्वन्यालोककाव्यप्रकाशादिग्रन्थवत् सम्मानितोऽस्ति । तत्र पण्डितराजो नव्यन्यायशैल्या काव्यसिद्धान्तान् व्याख्याति तथा विभिन्नविषयेषु उदाहरणानि स्वयमेवोद्भावितानि सन्तीति रसगङ्गाधरे स्वयमेव , यत्—

\begin{shloka}[center]

\textbf{निर्माय नूतनमुदाहरणानुरूपं}

\textbf{कामं मयाऽत्र विहितं न परस्य किञ्चित् ।}

\textbf{किं सेव्यतां सुमनसां मनसाऽपि गन्धः}

\textbf{कस्तूरिकाजननशक्तिभृता मृगेण ।। (रसगङ्गाधरः)}

\end{shloka}

रसगङ्गाधरो यद्यपि अपूर्णो विद्यते तथापि अस्य प्रौढप्रबन्धनं संस्कृतकाव्यशास्त्रस्य कृते महत्त्वपूर्णा उपलब्धिर्मन्यते । आननद्वये विभक्तेऽस्मिन् ग्रन्थे शास्त्रीयसमीक्षणम् उत्तरालङ्कारपर्यन्तमेव प्राप्यते । तत्र \textbf{प्रथमे आनने}— प्राग्जातानां लक्षणशास्त्रिणां काव्यलक्षणस्य खण्डनं कृत्वा ‘रमणीयार्थप्रतिपादकः शब्दः काव्यमिति प्रस्तौति । काव्यनिर्माणे प्रतिभा एव कारणत्वेन भवतीति चोल्लिखति । तथा च काव्यस्य चत्वारो भेदाः, रसस्य साङ्गोपाङ्गविवेचनमत्रैव लभ्यते । \textbf{द्वितीये आनने—} ध्वनेर्विवेचनम्, अभिधालक्षणयोः समीक्षणम्, अलङ्काराणाञ्च नवीनशैल्या सविस्तारं वर्णनं विहितमस्ति । रसगङ्गाधस्य विविधासु टीकासु नागेशभट्टेन कृता गुरुमर्मप्रकाशिनी इत्याख्यां टीका प्रसद्धा जाता ।  जगन्नाथस्य अन्या कृतयः-

\textit{\textbf{१. गङ्गालहरी (द्विपञ्चाशत्पद्यबद्धा गङ्गास्तुतिः)}}

\textit{\textbf{२. सुधालहरी (द्वात्रिंशत्पद्यात्मिका सूर्यस्तुतिः)}}

\textit{\textbf{३. अमृतलहरी (एकादशश्लोकयुता यमुनाप्रशस्तिः)}}

\textit{\textbf{४. करुणालहरी (चतुषष्ठीश्लोकयुता विष्णुस्तुतिः)}}

\textit{\textbf{५. लक्ष्मीलहरी (एकचत्वारिंशत्पद्यात्मिका लक्ष्मीस्तुतिः)}}

\textit{\textbf{६. यमुनावर्णनचम्पूः (चम्पूकाव्यम्)}}

\textit{\textbf{७. आसफविलासः (आसफखाँमहोदयस्य प्रशस्तिः)}}

\textit{\textbf{८. प्राणाभरणम् (कामरूपाधीशस्य प्राणनारायणस्य प्रशस्तिः)}}

\textit{\textbf{९.जगदाभरणम् (उदयपुराधीशराणाकर्णसिंहसूनोर्जगत्सिंहस्य प्रशस्तिः)}}

\textit{\textbf{१०. चित्रमीमांसाखण्डनम् (अप्पयदीक्षितस्य चित्रमीमांसाखण्डनम्)}}

\textit{\textbf{११. मनोरमाकुचमर्दनम् (मनोरमाख्यायाः सिद्धान्तकौमुदीटीकायाः खण्डनम्)}}

\textit{\textbf{१२. भामिनीविलासः (कवितासङ्‌ग्रहः)}}

\customparagraph{\textbf{काव्यलक्षणम्}}

आचार्योऽयं प्राचीनविदुषां काव्यलक्षणान्नधीत्य परिष्कृतं काव्यस्वरूपं प्रस्तौति । पूर्वेषां शब्दार्थमयं काव्यमिति काव्यस्वरूपे मताधिक्यमपि जगन्नाथः शब्द एव काव्यं कथयति, परमस्य रमणीयार्थप्रतिपादकत्वेन विशेषणेन ध्वनिवादिनो मतमनुसरणमस्तीति सूच्यते । अस्य काव्यलक्षणं रसगङ्गाधरे यथा—

\textit{\textbf{रमणीयार्थप्रतिपादकः शब्दः काव्यम्’ इति । रमणीयता च लोकोत्तरह्लादजनकज्ञानगोचरता । लोकोत्तरं चाह्लादगतश्चमत्कारत्वापरपर्यायोऽनुभवसाक्षिको जातिविशेषः । कारणं च तदवच्छिन्ने भावनाविशेषः पुनः पुनरनुसन्धात्मा । `पुत्रस्ते जातः' `धनं ते दास्यामि' इति वाक्यार्थधीजन्याह्लादस्य न लोकोत्तरत्तरत्वम्, अतो न तस्मिन् वाक्ये काव्यतत्वप्रसक्तिः ।}}

\textit{\textbf{>इत्थं चमत्कारजनकभावनाविषयार्थप्रतिपादकशब्दत्वम्}}

\textit{\textbf{>यत्प्रतिपादितार्थविषयकभावनात्वं चमत्कारजनकतावच्छेदकं तत्त्वम्}}

\textit{\textbf{>स्वविशिष्टजनकतावच्छेदकार्थप्रतिपादकतासंसर्गेण चमत्कारत्ववत्वमेव वा काव्यत्वमिति फलितम् । इति रसगङ्गाधरे काव्यलक्षणमूलम् । (रसगङ्गाधरः)}}

शब्दस्यार्थास्तु निरसाः, निष्प्रयोजनाऽपि भवितुं शक्नुवन्ति, तस्मादत्र रमणीयार्थ- सुन्दरप्रतीयमानार्थः, प्रतिपादकः- अभिव्यञ्जनक्षमः, इति विशेषणं न्यस्तम् । एतादृशो विलक्षणार्थव्यञ्जकः शब्दः काव्यं भवतीति जगन्नाथस्याशयः ।

\textbf{काव्यलक्षणे सन्निविष्टानां शब्दानां वैशिष्ट्यम्}

रमणीयार्थप्रतिपादक इत्यस्य स्थाने केवलम् अर्थप्रतिपादकः शब्दः काव्यमित्युच्यते चेत् तदपि अतिव्याप्तिर्लक्षणदोषयुतं  यत्- `घटमानय' इत्यादीनां वाक्यानामपि अर्थप्रतिपादकत्वं वर्तत एव । इत्यादिषु वाक्येषु काव्यस्यातिव्याप्तिवारणाय \textbf{`रमणीय'} इति विशेषणपदस्य सन्निवेशो विहितः ।

तथा च \textbf{`अर्थ'} इति शब्दप्रयोगेऽपि विशेषो विद्यते यत्- व्याकरणे शब्दप्रतिपादकत्वं विद्यते । व्याकरणं तु काव्यं न भवति । तस्माद्  व्याकरणात् काव्यस्य भेदं प्रदर्शयितुं रमणीयार्थप्रतिपादक इत्यत्रार्थशब्दस्य सन्निवेशो विहितः ।

तथैव \textbf{`प्रतिपादकः'} शब्दस्य निवेशने कारणं यद् - वाचकाः, लक्षकाः, व्यञ्जकाश्च शब्दा भवन्ति । एतेषां स्थाने प्रतिपादक इति सामान्यपदं न्यस्तम् । तस्मात् प्रतिपादकशब्दप्रयोगेण न केवलं रमणीयार्थवाचकोऽपितु रमणीयार्थलक्षकः रमणीयार्थव्यञ्जकश्च शब्दः काव्यत्वेन व्यपदेशो जातः ।

रमणीयार्थप्रतिपादकः शब्द इत्यत्र \textbf{`शब्दः'} इति पदं विशेष्यं भवति । तत्र विशेषणं तु रमणीयार्थप्रतिपादकः । अत्र `शब्दः' इति विशेष्यपदस्य सन्निवेषाभावात् रमणीयानुरागस्य प्रतिपादकः कान्तायाः कटाक्षभुजाक्षेपादिः काव्यं स्यादिति अतिव्याप्तिर्लक्षणदोषो भवेत् ।

‘\textbf{रमणीय’} इति पदं घटमानय इत्यादि वाक्यव्यावृत्तये। ‘\textbf{अर्थ’} इति पदं रमणीय शब्दप्रतिपादके व्याकरणादावतिव्याप्तिनिरासाय । ‘\textbf{प्रतिपादकम्’} इति  पदं व्यङ्ग्यादिसङ्ग्रहाय । ‘\textbf{शब्द’} इति पदं कटाक्षादिवारणाय   ।

तदयं सङ्क्षेपः - \textbf{रमणीया - लोकोत्तराह्लादजनकभावनाविषया, येऽर्थाः - वाच्यलक्ष्यव्यङ्ग्यरूपाः। तत्प्रतिपादक - तद्विषय-बोधकः,लक्ष्यकः,व्यञ्जको वा, शब्दः - सुप्तिङन्तरूपः काव्यमिति फलितोऽर्थः।}

एवम् रमणीयार्थप्रतिपादकत्वे सति शब्दत्वं काव्यस्य लक्षणं निर्दुष्टमिति प्रस्तौति जगन्नाथः ।

पण्डितराजो जगन्नाथः स्वकृते काव्यलक्षणे एव सम्भावितान् दोषान् आशङ्क्य त्रीन् विकल्पान् प्रदर्श्य परिष्कृतं काव्यलक्षणं प्रस्तौति  ।

\customparagraph{\textbf{काव्यलक्षणपरिष्कारः}}

रमणीयार्थप्रतिपादकः शब्दः काव्यम् । रमणीयार्थप्रतिपादकत्वे सति शब्दत्वं काव्यत्वम्  । इयं च रमणीयता लौकिक्या रमणीयताया विशिष्टा । रमणीयता च \textbf{लोकोत्तराह्लादजनकज्ञानगोचरता} । लोकोत्तरत्वं च चमत्कारजनकमेव । इदमनुभवैकवेद्यं भवति । न तु प्रमाणान्तरै वेदितुं शक्यते । लोकोत्तराह्लादे कारणं पुनः पुनरनुसन्धानात्मा भावनाविशेषः वर्तते । पुनः पुनः समानविषयकः स्मृतिविशेषो  भावना भवति । पुत्रस्ते जातः धनं ते दास्यामि इत्यादिवाक्यार्थज्ञानेनापि आनन्दो जायते परमसौ न काव्यानन्दः । यतो हि अत्र लोकोत्तरत्वं नास्ति । काव्यानन्दः लोकोत्तरोऽस्ति । अतः पूर्वोक्तानां काव्यानां न काव्यत्वप्रसक्तिः । इति जगन्नाथस्याशयः ।

रमणीयार्थप्रतिपादकः शब्दः काव्यमिति लक्षणे रमणीयताशब्दस्य परिष्कारं प्रदर्श्यानेन काव्यलक्षणमपि त्रिधा परिष्कृतम् । इमे परिष्कारा जगन्नाथेन नव्यन्यायसरणिमाश्रित्य कृता इति ज्ञायते ।

काव्यलक्षणे संघटितस्य रमणीयपदस्य विवृत्तौ \textbf{`रमणीयता च लोकोत्तरज्ञानगोचरता'} इत्यत्र ज्ञानशब्दस्य सन्निवेशे `ज्ञानं समूहालम्बनात्मकं भवति'  तस्माद्  घटमानय इत्यत्रापि काव्यत्वप्रसक्तिः स्यात् । अतः चमत्कारजनकज्ञानविषयार्थप्रतिपादकशब्दत्वं काव्यत्वमिति ज्ञानशब्दप्रयोगे समूहालम्बनरूपस्य उक्तज्ञानस्य स्थले घटमानय इत्यस्य काव्यत्वं न भवेत् इत्यतिव्याप्तेर्वारणाय ज्ञानपदस्य स्थाने `स्मृतिप्रवाहरूपा (ज्ञानधारा) भावना' इत्यर्थकं भावनापदं प्रयुज्य प्रथमः परिष्कारो विहितः ।

(शब्दार्थः> लोकोत्तरस्य‍= अलौकिकस्य आह्लादस्य= आनन्दस्य जनकम्= उत्पादकं यज् ज्ञानम् तद् गोचरता= तन्निरूपितविषयतारूपा अर्थनिष्ठा रमणीयता इति । ‘\textbf{लोकोत्तरत्वम्’} चाह्लादगतश्चमत्कारत्वापरपर्यायोऽनुभवसाक्षिको जातिविशेषः । ‘\textbf{लोकोत्तराह्लादे कारणम्’} पुनः पुनरनुसन्धानात्मा भावना विशेषः)

\customparagraph{\textbf{प्रथमपरिष्कारः}}

\textbf{चमत्कारजनकभावनाविषयार्थप्रतिपादकशब्दत्वं} काव्यत्वमिति प्रथमः काव्यलक्षण- परिष्कारः । चमत्कारस्य जनिका या भावना (ज्ञानधारा) तस्या विषयो योऽर्थः तस्य प्रतिपादकत्वे सति शब्दत्वं काव्यत्वम् । अस्मिन् प्रथमे परिष्कारे ज्ञानपदस्य स्थाने भावनापदप्रयोगस्य गभीरार्थोऽस्ति यत्- ज्ञानस्य क्वचित् समूहालम्बनरूपत्वेन `घटम् आनय' इति वाक्ये काव्यत्वं प्रसक्तं भवेत् । परं क्वचिद् भावना अपि समूहालम्बनरूपा स्वीकृता । तदा उक्त अतिव्याप्तिदोषः काव्यलक्षणे भावनापदप्रयोगेऽपि समागत एव । अत एतद्दोषपरिहाराय द्वितीयं लक्षणपरिष्कारं प्रतिपादयति जगन्नाथः ।

\customparagraph{\textbf{द्वितीयपरिष्कारः}}

यत्प्रतिपादितार्थविषयकभावनात्वं चमत्कारजनकतावच्छेदकं तत्त्वमिति द्वितीयो लक्षणपरिष्कारः । परमत्र यत्तदोः प्रयोगेण अर्थस्य अननुगतत्वं नाम दोषः । यत्तदोः शब्दयोः कारणेन काव्यलक्षणस्य गुरुत्वमपि दृश्यते ।  एतद्दोषपरिहाराय तृतीयं लक्षणपरिष्कारं प्रदर्शयति जगन्नाथः ।

\customparagraph{\textbf{तृतीयपरिष्कारः}}

स्वविशिष्टजनकतावच्छेदकार्थप्रतिपादकतासंसर्गेण चमत्कारत्ववत्वमेव वा काव्यमिति तृतीयः परिष्कारः । अनेन च चमत्कारत्ववत्त्वमेव काव्यत्वमित्युक्त्वा पूर्वोक्तदोषपरिहारः । चमत्कारत्वविशिष्टत्वे सति शब्दत्वं काव्यत्वमिति लक्षणं फलितम् ।

\customparagraph{\textbf{काव्यप्रयोजनम्}}

\textbf{कीर्तिपरमाह्लादगुरुराजदेवताप्रसादाद्यनेकप्रयोजनकस्य काव्यस्य ’} इति । \textbf{(रसगङ्गाधरः)}

\customparagraph{\textbf{काव्यहेतुः}}

\textbf{तस्य काव्यस्य कारणं केवला कविगतप्रतिभा, सा च काव्यघटनानुकूलशब्दार्थोपस्थितिः, तद्गतं च प्रतिभात्वं काव्यकारणताऽवच्छेदकतया सिद्धो जातिविशेषः उपाधिरूपं वा खण्डम् । (रसगङ्गाधरः)}

प्रज्ञानवनवोन्मेषशालिनी प्रतिभोच्यते’ इति च प्रतिभायाः लक्षणम् । प्रतिभायाश्च हेतुः क्वचिद्देवतामहापुरुषप्रसादादिजन्यमदृष्टम्, क्वचिच्च विलक्षणव्युत्पत्तिकाव्यकारणाभ्यासौ’ इति च । अर्थात् काव्यस्य कारणं प्रतिभा, प्रतिभायाश्च कारणं क्वचिद्देवतामहापुरुषप्रसादादिजन्यमदृष्टं प्रतिभा, क्वचिच्च विलक्षणव्युत्पत्तिकाव्यकारणाभ्यासजन्यं दृष्टं प्रतिभा एव कारणमिति आशयः ।

\customparagraph{\textbf{काव्यभेदाः}}

काव्यभेदविषये जगन्नाथोऽपि पूर्वेषां ध्वनिवादिनां मतमनुसरति, परन्तु अस्य मते काव्यस्य चत्वारो भेदा भवन्ति । अत्रास्य विचारो विद्यते यत्— शब्दचित्रकाव्यापेक्षया अर्थचित्रकाव्यस्य किञ्चित् सौन्दर्यं भवति । अत एव अनेन उत्तमोत्तम–उत्तम–मध्यम–अधमश्चेति काव्यस्य चत्वारो भेदा स्वीकृताः । तद्यथा—

\textbf{उत्तमोत्तमकाव्यस्वरूपम्}

शब्दार्थौ यत्र गुणीभावितात्मानौ कमप्यर्थमभिव्यङ्क्तस्तदाद्यम् ।

\textbf{उत्तमकाव्यस्वरूपम्-}

यत्र व्यङ्ग्यमप्रधानमेव सच्चमत्कारकारणं तद्वितीयम् ।

\textbf{मध्यमकाव्यस्वरूपम्-}

यत्र व्यङ्ग्यचमत्कारसमानाधिकरणो वाच्यचमत्कारस्तत्तृतीयम् ।

\textbf{अधमकाव्यस्वरूपम्-}

यत्रार्थचमत्कृत्युपस्कृता शब्दचमत्कृतिर्प्रधानं तदधमं चतुर्थम् । \textbf{(रसगङ्गाधरः)}

एतादृशैः काव्यशास्त्रीयविवेचनैः पण्डितराजः पौरस्त्यकाव्यशास्त्रे समादरणीयं स्थानं भजते । काव्यशास्त्रस्योत्तरसमये जातानामलङ्कारशास्त्रिणां समाजेऽस्य विलक्षणप्रतिभावैशिष्ट्यम् अद्वितीयं दृश्यते । काव्यलक्षणकाव्यभेदयोरस्य नवीनदृष्टिकोणो विद्यते । रसविषयेऽनेन उत्पत्ति-अनुमिति-भुक्ति–अभिव्यक्तिरूपाणां निष्पत्तिर्विषयकसिद्धान्तानामतिरिक्ता अन्ये सिद्धान्ताः प्रस्तुताः सन्ति, परं ते सर्वे स्वकीयप्रतिभाबलात् स्वयमेवाविष्कृता मन्यन्ते । अभिधा–लक्षणाविषययोरस्य पूर्ववत् सैद्धान्तिकविचारो वर्तते, परं विवेचनशैल्या तु विलक्षणं दृश्यते। अलङ्काराणां समीक्षणपद्धतिरस्य नवीनाऽस्ति ।  तत्र पण्डितराजो वैयाकरणनैयायिकयोः शैल्या विवेचनं करोति । वैयाकरणे प्रसिद्धः शाब्दबोधोऽलङ्कारप्रकरणेऽपि समाविशति । समीक्षणप्रकारोऽयं सामान्यपाठकानां कृते कठिनत्वमाभाति, परं काव्यशास्त्ररसिकानां कृते तु गीष्मर्तौ सुस्वस्थस्वच्छगभीरजलयुते समुद्रे निमज्जनवत्  परमानन्दमुत्पादयति ।

\specialupakhanda{\textbf{व्याख्याकाल-उपसंहारः}}

पौरस्त्यकाव्यशास्त्रस्य दीर्घतरे इतिहासे मम्मटाचार्याद् आरभ्य पण्डितराजजगन्नाथं यावद् व्याप्तोऽयं `व्याख्याकालः' काव्यशास्त्रीयसिद्धान्तानां स्थिरीकरणाय, परिष्करणाय, समन्वयाय च पूर्णतया विशिष्टसमय आसीत् । निर्णयात्मककाले उद्भाविता विभिन्नसिद्धान्ता अस्मिन् काले प्रौढताङ्गताः स्थिरीभूताश्च । यद्यपि कालस्यास्य पश्चात् संस्कृतकाव्यशास्त्रस्य ह्रासकालो मन्यते, तथापि अस्मिन् कालखण्डे जाता काव्यशास्त्रिणो न केवलं प्राचीनानां व्याख्याकाराः, अपितु स्वकीयप्रतिभाबलेन विविधानां काव्यशास्त्रीयतत्त्वानां मार्गदर्शका अभवन् । अस्मिन् युगे ध्वनिसम्प्रदायपक्षपातिना वाग्देवतावतारभूतेन मम्मटाचार्येण स्वकीये काव्यप्रकाशे विरोधिमतानां सतर्कं तिरस्कृत्य ध्वनिसिद्धान्तस्य स्थिरता प्रतिष्ठापिता । तदनन्तरं काश्मीरदेशोद्भवेन आचार्यरुय्यकेन `अलङ्कारसर्वस्वे' अलङ्काराणां यद् वैज्ञानिकं वर्गीकरणं कृतम्, तत् समग्रकाव्यशास्त्रस्य कृते वरदानमेव मन्यते । गुर्जरदेशीयो हेमचन्द्रः, महामात्यो वाग्भटश्च सङ्ग्रहात्मकशैल्या काव्यतत्त्वानि विवेचयन्ति । पीयूषवर्षेण जयदेवेन `चन्द्रालोके' अनुष्टुप्छन्दसा एकस्मिन्नेव श्लोके लक्षणोदाहरणयोः संयोजनं कृत्वा कठिनत्वेनानुभूतं शास्त्रं काव्यशास्त्रप्रवेष्टुमुत्सुकानां बालानां कृतेऽपि सुबोधं सरलञ्च कृतम्, यस्य अनुकरणम् अग्रे अप्पयदीक्षितेन `कुवलयानन्दे' शतोपरिपञ्चविंशत्यलङ्काराणां निरूपणाय कृतम् । अत्रैव एकावलीकारो विद्याधरः, प्रतापरुद्रयशोभूषणस्य रचयिता विद्यानाथश्च स्वरचितराजप्रशस्तिपरकैर्नैकैर्ग्नन्थैः काव्यशास्त्रीयेतिहासं  समवर्धयताम् ।

अस्मिन्नेव काले रससम्प्रदायस्य भक्तिधारायाश्च अभूतपूर्वो विकासः सञ्जातः । उत्कलदेशीयः कविराजो विश्वनाथो मम्मटमतं खण्डयित्वा `वाक्यं रसात्मकं काव्यम्' इति समुद्‌घोषयन् काव्यशास्त्रं रसमयमकरोत्, दृश्यश्रव्यकाव्ययोरद्वितीयं शास्त्रीयसमीक्षणञ्च प्रस्तुतवान् । शारदातनयशिङ्गभूपालभानुदत्तादयो लक्षणशास्त्रिणो नैकेषु ग्रन्थेषु रस-भाव-नायकनायिका- भेदादीनां गम्भीरं विवेचनं चक्रुः । वैष्णवभक्तिधाराया रूपगोस्वामिना भक्तिरसामृतसिन्धु- उज्ज्वलनीलमणि-नाटकचन्द्रिकादिभिर्ग्रन्थैर्लौकिक-अलौकिकयोः समन्वयम् विधाय, अन्यरसान् भक्तिरस्यैव विकृतिं मन्वानः साहित्यिकदृष्ट्या `भक्तिरस' एव मूलरसत्वेन प्रतिष्ठापितः । कालस्यास्य उत्तरार्धे केशवमिश्रोऽलङ्कारशेखराख्ये ग्रन्थे विंशतिमरीचिषु काव्यशास्त्रं सङ्‌क्षेपेण व्याख्यातवान् ।

अस्य व्याख्याकालस्य चरमोत्कर्षः पण्डितराजजगन्नाथस्य रसगङ्गाधराख्ये प्रौढग्रन्थे दृश्यते । तेन प्राग्जातानां सर्वेषां काव्यलक्षणं सतर्कं खण्डयित्वा `रमणीयार्थप्रतिपादकः शब्दः काव्यम्' इति स्वकीयं स्वतन्त्रं मतं स्थापितम् । जगन्नाथेन नव्यन्यायस्य परिष्कृतशैल्या काव्यलक्षणस्य तिसृभिः कोटिभिर्यद् विवेचनं विहितम्, यच्च स्वरचितनूतनोदाहरणानां साहाय्येन अलङ्कारप्रकरणे वैयाकरण-नैयायिकानां शाब्दबोधः समाविष्टः, तेन विलक्षणं काव्यशास्त्रीयं वैशिष्ट्यमाविष्कृतम् । मम्मटस्य सूत्रात्मकशैल्या आरभ्य रूपगोस्वामिनो भक्तिरसामृतधारां स्पृष्ट्वा पण्डितराजस्य नव्यन्यायात्मकवैचारिकप्रौढत्वं प्राप्तोऽयं व्याख्याकालः काव्यशास्त्रचिन्तकैः स्मरणीयः प्रकाशस्तम्भ इवानुभूयते इति ।

\specialmahakhanda{\textbf{षष्ठ-प्रकाशः}}

\specialupakhanda{\textbf{ह्रासकालः (अर्वाचीनकालः)}}

\begin{shloka}[center]

\textbf{(जगन्नाथपश्चात्)}

\end{shloka}

\specialupakhanda{\textbf{आचार्यो विश्वेश्वरः}}

पण्डितराजजगन्नाथस्यानन्तरं संस्कृतसमालोचनाक्षेत्रे संस्कृतविदुषां तर्कशक्तिः क्षीणा इव प्रतीयते । तादृशे समये दीप इव आचार्यो विश्वेश्वरोऽजायत । अस्य पूर्वजा अल्मोडाजनपदस्य पाटियाग्रामवास्तव्या भारतद्वाजगोत्रियाः कान्यकुब्जब्राह्मणा आसन् । अस्य पिता लक्ष्मीदत्तपाण्डेयो जीवनस्योत्तरार्द्धे पाटियग्रामात् काशीमेत्य गङ्गातटे मणिकर्णिकाघट्टे वासमकरोत् । काश्यामेवाष्टादशशताब्द्याः प्रारम्भे विश्वेश्वरस्य जन्माभवत् । अस्य विद्यागुरुरपि पिता एवेति मन्दारमञ्जर्यायां स्वपितुर्वन्दना कृताऽनेन, तद्यथा-

\begin{shloka}[center]

\textbf{जयति यथाजातानां वाग्जातसुजातपारिजातश्रीः ।}

\textbf{श्रीलक्ष्मीधरविबुधावतंशचरणाब्जरेणुनिकरः ।।}

\end{shloka}

अयं पौरस्त्यकाव्यशास्त्रस्य तीव्रविकासस्यान्तिमप्रायो विवेचकोऽस्ति । पण्डितवरस्य विश्वेश्वरस्य रचनाबाहुल्येऽपि प्रसिद्धा अलङ्कारशास्त्रसम्बद्धाः पञ्चग्रन्थाः सन्ति, यत्— अलङ्कारकौस्तुभः, अलङ्कारमुक्तावली, अलङ्कारप्रदीपः, रसचद्रिका, कवीन्द्रकण्ठाभरणञ्चेति । तेषु अलङ्कारकौस्तुभे— उपमादय एकषष्टिरर्थालङ्काराः सप्रपञ्चं विवेचिताः । तत्र अन्यैर्वर्णिता अलङ्काराः काव्यप्रकाशवर्णितेष्वलङ्कारेष्वन्तर्भावात् तत्रस्था अलङ्कारा एव विवेचिता इति विश्वेश्वरः स्वयमेव अलङ्कारकौस्तुभस्यान्तिमपद्ये समुल्लिखति । अत एवास्य अलङ्कारविकासे नूतनं प्रणयनं नास्ति । तस्मादयं काव्यशास्त्रीयेतिहासस्य ह्रासकाले परिगणितोऽस्ति ।

\begin{shloka}[center]

\textbf{अन्यैरुदीरितमलङ्करणान्तरं यत् काव्यप्रकाशकथितं तदनुप्रवेशात् ।}

\textbf{सङ्क्षेपतो बहुनिबन्धविभावनेनालङ्कारजातमिह चारु मया न्यरूपि ॥}

\textbf{(अलङ्कारकौस्तुभ पृ. ४१९)}

\end{shloka}

अस्यानेनैव कृता उपज्ञाख्या व्याख्या लभ्यते । ग्रन्थेऽस्मिन् अलङ्कारविवेचनं विशदं तर्कपूर्णञ्च विद्यते। अलङ्कारमुक्तावली तु अलङ्कारकौस्तुभस्यैव सङ्क्षिप्तरूपाऽस्ति । अस्य विषये तेनैवोक्तं यत्—

\begin{shloka}[center]

\textbf{नानापक्षविभावनकुतुकमलङ्कारकौस्तुभं कृत्वा ।}

\textbf{सुखबोधाय शिशूनां क्रियते मुक्तावली तेषाम् ।।  (अलङ्कारमुक्तावली)}

\end{shloka}

सरलपरिपाट्या अलङ्कारान्नवबोधाय रचितोऽलङ्कारप्रदीपोऽस्यापरो ग्रन्थः । तत्र कुवलयानन्दशैल्याम् अलङ्काराणां सुबोध्यवर्णनं विद्यते । ग्रन्थेऽस्मिन् शतोपरित्रिंशद्  अलङ्कारास्तद्भेदाश्च प्रस्तुता आचार्येण । रसचन्द्रिका तु रसविवेचनसम्बद्धा । अत्र रसविषये सविस्तारेण विवेचनं लभ्यते । कवीन्द्रकण्ठाभरणं तु शब्दचित्रकाव्यसम्बद्धो ग्रन्थः । तस्मिन् चत्वारः परिच्छेदाः सन्ति । तेषु अष्टपञ्चाशत् वृत्तजातयः वर्णिताः । रसचन्द्रिकायां रसस्य निरूपणं कृतम् ।

इत्यादीनां ग्रन्थानां रचयिता विश्वेश्वरः संस्कृतकाव्यशास्त्रेतिहासस्य ह्रासोन्मुखकाले सञ्जातो विशिष्टो विद्वान् अस्ति । अस्य कृतयः काव्यशास्त्रानुशीलनसम्बद्धानां पाठकानां कृते उपयोगिनः सन्ति ।

पौरस्त्यकाव्यशास्त्रीयेतिहासे पण्डितराजजगन्नाथपश्चात् विश्वेश्वरानन्तरं नवीनशास्त्रीयविचारप्रकाशकाः काव्यशास्त्रिणः प्रायशः नोपलभ्यन्ते । कतिचिद् विद्यन्ते ते सर्वे प्राचीनानां काव्यतत्त्वानां पक्षपोषका एवाभिलक्ष्यते । केवलं तेषां नामानि प्रकाशिता ग्रन्थाश्चात्र उल्लिख्यन्ते ।

\textbf{ह्रासकाले सञ्जातानां  काव्यशास्त्रिणां नामनिर्देशः}

१. \textbf{नञ्जराजयशोभूषणम्} नरसिंहः१७३९ १७५९

२ \textbf{साहित्योद्देशः}, साहित्यसिद्धान्त सीताराम शास्त्री १९२३

३ \textbf{व्यञ्जनावादः} यदुनाथ मिश्र १९३७

४ \textbf{साहित्यमञ्जरी} श्रीपाद शास्त्री हसूरकर १८८८ १९४७

५ \textbf{भक्तिरसार्णवः} स्वामीकरपात्रीजी १९०७ १९८२

६ \textbf{रसचन्द्रिका} लेखनाथ झा १९४०

७ \textbf{लोकमान्यालङ्कारः} श्री गजानन शास्त्री करमलकर १९४१

८ \textbf{साहित्यविमर्शः} श्री कौत्स अप्पल सोमेश्वर शर्मा १९५१

९ \textbf{साहित्यसारः} श्री सर्वेश्वर कवि १९९२ तिरुपति विद्यापीठ

१० \textbf{साहित्यबिन्दुः} छज्जुराम शास्त्री १९०५ १९७९

११ \textbf{कारिकाग्रन्थ} सिद्धान्तशिरोमणि विश्वेश्वर १०६० १९८०

१२ \textbf{अभिनवकाव्यप्रकाशः} गिरिधरलाल व्यास शास्त्री १९६६

१३ \textbf{काव्यालङ्कारकारिका} रेवाप्रसाद द्विवेदी

१४ \textbf{ध्वनिकल्लोलिनी} आनन्द झा १९७८ कामेश्वर संस्कृत विश्वविद्यालय दरभङ्गा

१५ \textbf{काव्यसत्यालोकः} ब्रह्मानन्द झा १९२३ २०००

१६ \textbf{रसचन्द्रिका, व्यञ्जनावृत्तिविचारः} रामावतारमिश्र

१७ \textbf{साहित्यसन्दर्भः} प्रो शिवजी उपाध्याय

१८ \textbf{सौन्दर्यकारिका} प्रो जगन्नाथ पाठक

१९ \textbf{मैथिलीकाव्यविवेकः} कविशेखर पं बदरीनाथ झा १९९४ बदरीनाथ ग्रन्थावली

२० \textbf{सौन्दर्यदर्शनविमर्शः} गोविन्द्रचन्द्र पाण्डेय

२१ \textbf{काव्यतत्त्वविवेकः} डा रमाशङ्कर तिवारी

२२ \textbf{काव्यात्मनिर्णयः} डा हरिशचन्द्र दीक्षित

२३ \textbf{अभिनवकाव्यशास्त्रम्} डा शङ्करदेव अवतरे

२४ \textbf{काव्यसिद्धान्तकारिका} प्रो अमरनाथ पाण्डेय २०२१ अखिलभारतीय संस्कृत परिषद्

२५ \textbf{रसवसुमूर्ति} प्रो चन्द्रमौलि द्विवेदी

२६ \textbf{चमत्कारविचारचर्चा} प्रो रामप्रताप वेदालङ्कार

२७ \textbf{अभिनवकाव्यालङ्कारसूत्रम्} प्रो राधावल्लभ त्रिपाठी

२८ \textbf{अभिराजयशोभूषणम्} प्रो राजेन्द्र मिश्र

२९ \textbf{लघुच्छन्दोऽलङ्कारदर्पण} पं नित्यानन्दशास्त्री

३० \textbf{नव्यकाव्यतत्त्वमीमांसा} प्रो रहसबिहारी द्विवेदी २००५

३१ \textbf{अलङ्कारविद्योतनम्} कृष्णमाधव झा २००८

\textbf{स्रोतः} \textit{संस्कृत काव्यशास्त्र की अर्वाचीन परम्परा- राजमङ्गल यादव (प्रतिभा प्रकाशन-२०११)}

समालोचकैः सङ्केतिता अन्येऽपि ग्रन्थकारा दृश्यन्ते । तथाप्यत्र प्रकाशिताः काव्यशास्त्रीयग्रन्था एव समुल्लिखिताः । एतेषां भूमिका केवलं पूर्वाचार्यै स्थापितानां शास्त्रीयमतानां परिपोषिका एव दृश्यते । एते ग्रन्थाः पृथक् नवीनसैद्धान्तिकतत्त्वान्वेषणे समर्था न विद्यन्ते । अत एव जगन्नाथहेमचन्द्रानन्तरं संस्कृतकाव्यशास्त्रस्य अवसानकालनाम्ना समालोचकैरुद्‌घोषितमिति मन्ये ।

\specialupaupakhanda{\textbf{ह्रासकाल-उपसंहारः}}

पौरस्त्यकाव्यशास्त्रस्य तीव्रतरो विकासः पण्डितराजजगन्नाथस्य प्रौढविवेचनानन्तरं शिथिलोऽनुभूयते। जगन्नाथपश्चात् प्रारब्धः कालखण्डः समालोचकैः `ह्रासकालः' इति नाम्ना अनुमन्यते। अस्मिन् काले काव्यसिद्धान्तानां सम्प्रदायानां वा नूतना उद्भावना क्षीणा अभवत् । तस्मिन् मन्दप्राये शास्त्रीयविकासे अष्टादशशताब्द्याः प्रारम्भे काश्यां प्रादुर्भूतो भारद्वाजगोत्रीयः कान्यकुब्जब्राह्मणो विश्वेश्वरः सञ्जातः । यद्यपि विश्वेश्वरेण अलङ्कारकौस्तुभः, अलङ्कारमुक्तावली, अलङ्कारप्रदीपः, रसचन्द्रिका, कवीन्द्रकण्ठाभरणम् इति पञ्च प्रमुखा ग्रन्था रचिताः, येषु तर्कपूर्णं विवेचनमपि कृतम्; तथापि तेनैव अलङ्कारकौस्तुभस्यान्ते स्वयमेव स्वीकृतं यद् ‘अन्यैरुदीरितानां नवीनालङ्काराणामन्तर्भावो मम्मटभट्टस्य काव्यप्रकाश एव भवतीति ।’ अतस्तस्य कार्यं प्राचीनानां सिद्धान्तानां व्यवस्थापनं परिपोषणञ्चासीत्, न तु सर्वथा नूतनसिद्धान्तनिर्माणम्, यत् अस्य कालस्य ह्रासोन्मुखतां प्रमाणीकरोति ।

विश्वेश्वरानन्तरं नरसिंह-अमरनाथपाण्डेय-  सीतारामशास्त्रि-यदुनाथमिश्र-श्रीपादशास्त्रि- करपात्री-लेखनाथझा-गजाननशास्त्रि-रेवाप्रसादद्विवेदि-आनन्दझा-ब्रह्मानन्दझा -राधावल्लभत्रिपाठि-राजेन्द्रमिश्र-रहसबिहारीद्विवेदिप्रभृतयः प्रकाण्डविद्वांसः काव्यशास्त्रीयपरम्पराया व्याख्या- विवेचनादिमाध्यमेन प्रवर्धनं चक्रुः । एतेषां सर्वेषाम् आचार्याणां मुख्यं प्रतिपाद्यं नवीनप्रतिभाप्रदर्शनेन सह पूर्वाचार्यैः स्थापितानां सिद्धान्तानां व्याख्यानम्, सरलीकरणम्, पाठ्यग्रन्थरूपेण सङ्कलञ्च आसीत् ।

पण्डितराजजगन्नाथोत्तरकालः सैद्धान्तिकदृष्ट्या ह्रासकालः प्रतीयते, तथापि ग्रन्थरचना- प्रकाशनदृष्ट्या संरक्षणदृष्ट्या च अयं कालोऽतीव उर्वरः समृद्धश्च ‍दृश्यते । ये पूर्ववर्तिनः काव्यशास्त्रिणः काव्यशास्त्रीयमार्गान् स्थापितवन्तः, ह्रासकालस्य अर्वाचीना विद्वांसस्तेषामेव मार्गाणां संरक्षकाः, परिष्कर्तारः, प्रचारकाश्च सञ्जाताः । अतः पृथक् सैद्धान्तिकतत्त्वान्वेषणाभावेऽपि, पौरस्त्यसाहित्यशास्त्रस्य अविरलप्रवाहे अस्य ह्रासकालस्य तत्र सञ्जातानां विदुषाञ्च अवदानं शास्त्रीयपरम्पराया गौरवरक्षणार्थमतीव महत्त्वपूर्णम्, आदरणीयं चास्तीति मन्यते ।

\specialkhanda{\textbf{नेपालदेशोद्भवानां काव्यशास्त्रिणां परिचयः}}

ऐतिहासिकतथ्यानुसारं भारतभूमिः संस्कृतभाषायाः विकासभूमिरासीत् । तस्मिन् समये अधुना पृथक् नाम्ना विकसिता सार्वभौमदेशा अपि भारतभूमिरिति नाम्ना ज्ञायन्ते स्म । पौराणिकेतिहासानुसारं नेपालभूमिर्वशिष्ठविश्वामित्रादीनां मुनीनां साधनाभूमिरासीदिति सङ्केतादयः प्राप्यन्ते । अतोऽधुना सीमाङ्कितायां नेपालभूमावपि संस्कृतस्य विकासोऽजायत इत्यनुमितुं शक्यते । आधुनिकनैपालदेशे तु पृथक् सैद्धान्तिकप्रणयनसहितः काव्यशास्त्रीयविकासो नोपलभ्यते । नेपालीभाषावाङ्मयस्य विकासारभ्भ एव संस्कृतग्रन्थानामनुवादकार्यात् सञ्जात इति ऐतिहासिकप्रमाणेन ज्ञायते । अतोऽधुनापि प्राचीनकाले भारतवर्षे विकसिता संस्कृतभाषासम्बद्धा ये सिद्धान्ताः सन्ति तेषामेव पक्षपोषणमत्र सञ्जात इति अधुना प्रकाशितास्तदनुकूलग्नन्थानवलोक्य ज्ञायते । तथापि पौरस्त्यकाव्यशास्त्रस्य विकासे आधुनिकनेपालदेशस्यापि किञ्चद् योगदानमस्तीति स्वीकर्तुं शक्यते । पौरस्त्यवाङ्मयस्य काव्याशास्त्रीयविषयमवलम्ब्य स्वविचारप्रकाशका नैपालकविद्वांसः तेषां ग्नन्थाश्च अत्र उल्लिख्यन्ते ।

\customparagraph{\textbf{पं. कुलचन्द्रगौतमः}}

पं कुलचन्द्रगौतमः प्राचीननैपालदेशीयः काव्यशास्त्रीयो ग्रन्थकार आसीत् । अनेन राघवालङ्कारः, अलङ्कारचन्द्रोदय इति द्वौ शास्त्रीयग्नन्थौ रचितौ । अनेन द्वयोरेव ग्रन्थयोरेकस्मिन्नेव पद्ये पूर्वार्धे लक्षणम्, अपरार्धे चोदाहरणमिति चन्द्रालोकस्य अलङ्कारविश्लेषणशैली स्वीकृता दृश्यते । अलङ्कारचन्द्रोदये (१९७८) पौरस्त्यकाव्यशास्त्रानुसारम् अलङ्कारलक्षणं समुल्लिख्य उदाहरणे नेपालदेशीयतात्कालिकशासकस्य चन्द्रसमसेरस्य प्रशस्तिमुल्लिख्य मौलिकाः कविता रचिता विद्यन्ते । राघवालङ्कारे (२०१६) पौरस्यकाव्यशास्त्रानुसारमलङ्कारलक्षणमुल्लिख्य उदाहरणे दशरथपुत्रस्य रामचन्द्रस्य प्रशस्तिवर्णिताः श्लोकाः समुल्लिखिताः सन्ति । ग्रन्थेऽस्मिन् ९ शब्दालङ्काराः ९३ अर्थालङ्काराश्च वर्णिताः सन्ति ।

\customparagraph{\textbf{सोमनाथसिग्द्यालः}}

पौरस्त्यसमालोचासिद्धान्तस्य नेपालीभाषायां प्रचारकः सोमनाथसिग्द्यालोऽस्ति । अनेन साहित्यप्रदीप इति नाम्ना सैद्धान्तिकं ग्रन्थं विरच्य कठिनमित्यनुभूतः पौरस्त्यकाव्यशास्त्रीयसिद्धान्तः नेपालीभाषायां सुबोध्यरूपेण अवतारितः । \textbf{साहित्यप्रदीपः} इत्याख्यो ग्रन्थः २०१५ तमे वैक्रमाब्दे नेपाल एकेडेमी इति संस्थया प्रकाशित आसीत् । सिग्द्यालेन तत्र सर्वेषां काव्याङ्गानां सुबोध्यरूपेण चर्चा विहिता । अधुनापि ग्रन्थोऽयं नैपालकानां संस्कृतकाव्यशास्त्रजिज्ञासुजनानां कृते समादरणीयोऽस्ति ।

\customparagraph{\textbf{ऋषभदेवः}}

ऋषभदेवः शर्मा रेग्मी इति नाम्ना प्रसिद्धेन विदुषा \textbf{काव्यमञ्जरी}ति नाम्ना काव्यशास्त्रीयग्रन्थः रचितः । २०१५ तमाद् वैक्रमाब्दाद् आरभ्य २०२८ तमे वैक्रमाब्दे षष्ठसंस्करणत्वेनापि प्रकाशितोऽयं ग्रन्थो लघुकायात्मको विद्यते परं काव्यस्य कृते प्राधान्येन आवश्यकानि तत्त्वानि वर्णितानि सन्ति । तानि सर्वाणि सैद्धान्तिकतत्त्वानि प्राचीनानां पौरस्त्यकाव्यशास्त्रिणां सिद्धान्तानुकूलानि विद्यन्ते ।

उपर्युल्लिखितान् ग्रन्थान्नतिरिच्य नेपालीभाषासम्बद्धा पौरस्त्यकाव्यशास्त्रीयविषयानुकूला ग्रन्थकाराः पूर्वसिद्धान्तपोषका एव विद्यन्ते । अत्र नेपालीसमालोचनाक्षेत्रे प्रसिद्धानां काव्यशास्त्रीयग्रन्थकाराणां  नामानि ग्नन्थाश्चोल्लिख्यन्ते ।

\textbf{१. भरतराजपन्तः} नेपाली अलङ्कार परिचय

\textbf{२. लक्ष्मीप्रसादाचार्यः} अलङ्कार सागर सोपान

\textbf{३. हेमाङ्गराजाधिकारी} पूर्वीय समालोचना सिद्धान्त

\textbf{४. हिमांशुथापा} साहित्यपरिचय

\textbf{५. ईश्वरवरालः} एवम् अन्य नेपाली सहित्यकोश

\textbf{६. केशवप्रसादोपाध्यायः} साहित्य प्रकाश

पौराणिकेतिहासाद् आरभ्य ऋषिमुनीनां पवित्रतपोभूमित्वेन संस्तुतायाम् अस्यां नेपालभूमौ पौरस्त्यकाव्यशास्त्रस्य संरक्षणे, संवर्धने च नैपालकविदुषां प्रयासोऽतीव श्लाघनीयः प्रतीयते । यद्यपि आधुनिकनेपालदेशे कश्चन सर्वथा मौलिको नवीनः सैद्धान्तिकः सम्प्रदायो नोद्भावितः, तथापि प्राचीनभारतवर्षे विरचितानां काव्यशास्त्रीयसिद्धान्तानां दार्ढ्याय, समन्वयाय च अत्रत्यैर्मनीषिभिः सातिशयं योगदानं दत्तम् । अस्मिन् क्रमे अग्रणीकाव्यशास्त्री कुलचन्द्रगौतमो जयदेवकृतस्य चन्द्रालोकस्य शैल्याम् एकस्मिन्नेव श्लोके लक्षणोदाहरणयोरद्भुतं सामञ्जस्यं कृत्वा `अलङ्कारचन्द्रोदयः' तथा च नव-शब्दालङ्काराणां त्रिनवत्यर्थालङ्काराणाञ्च महता प्रपञ्चेन युक्तो राघवालङ्कारश्चेति ग्रन्थद्वयं प्रणीतवान्, यत्र समसामयिकशासकस्य चन्द्रशमशेरस्य मर्यादापुरुषोत्तमरामस्य च उदात्ताः प्रशस्तयः सन्ति । तथैव `साहित्यप्रदीपः' इत्याख्यस्य शास्त्रीयग्रन्थस्य रचयिता सोमनाथसिग्द्यालवर्यः कठिनतमान् पौरस्त्यकाव्यशास्त्रीयसिद्धान्तान् नेपालस्य राष्ट्रभाषायां सारल्येन उपस्थाप्य नैपालदेशीयसमालोचनापरम्परायां पौरस्त्यकाव्यशास्त्रस्य पूर्णं व्याख्यानं कृतवान् । ऋषभदेव-भरतराजपन्त-लक्ष्मीप्रसादाचार्य-हेमाङ्गराजाधिकारि-हिमांशुथापा-ईश्वरबराल-केशवप्रसादोपाध्यायाश्चेति विद्वांसः पौरस्त्यकाव्यतत्त्वानां विस्तृतां चर्चाम् अकुर्वन् । अनेन प्रकारेण नेपालीसाहित्यवाङ्मयस्य शास्त्रीयभित्तिं सुदृढां कर्तुम्, पौरस्त्यसाहित्यशास्त्रस्य ज्ञानगङ्गां नेपालभूमावपि प्रवाहयितुं च नेपालदेशोद्भवानां काव्यशास्त्रिणाम् इयं व्याख्यानपरम्परा फलीभूता दृश्यते ।

\specialupakhanda{\textbf{काव्यशास्त्रप्रवेशिका- पुस्तकस्य निष्कर्षः/उपसंहारः}}

पौरस्त्यकाव्यशास्त्रपरम्परायां लक्षणशास्त्रिभिर्विलक्षणप्रतिभया प्रणीता ज्ञाताज्ञाताश्च शतशो ग्रन्था विद्यन्ते । तेषु मुख्यतः काव्यशास्त्रशरीरस्य सर्वाङ्गसमीक्षकाः, अलङ्कारमात्रसमीक्षकाः, ध्वनिसमीक्षकाः, रसमात्रसमीक्षकाः, नाट्यसमीक्षकाश्च विभिन्नविधासम्बद्धा ग्रन्थाः सन्ति । एतेषु भामहस्य काव्यालङ्कारः, दण्डिनः काव्यादर्शः, वामनस्य काव्यालङ्कारसूत्रवृत्तिः, रुद्रटस्य काव्यालङ्कारः, सरस्वतीकण्ठाभरणञ्च, कुन्तकस्य वक्रोक्तिजीवितम्, मम्मटस्य काव्यप्रकाशः, हेमचन्द्रस्य काव्यानुशासनम्, वाग्भटस्य वाग्भटालङ्कारः, जयदेवस्य चन्द्रालोकः, विश्वनाथस्य साहित्यदर्पणः, जगन्नाथस्य रसगङ्गाधरश्च ग्रन्थाः काव्यशास्त्रस्य सर्वाङ्गविवेचकाः सन्ति । अलङ्कारमात्रसमीक्षणयुतेषु ग्रन्थेषु—  उद्भटस्य काव्यालङ्कारसारसङ्ग्रहः, शोभाकरस्यालङ्काररत्नाकरः,  अप्पयदीक्षितस्य कुवलयानन्दश्चित्रमीमांसा च विश्वेश्वरस्य अलङ्कारकौस्तुभः, अलङ्कारप्रदीपः, अलङ्कारमुक्तावली च, नरसिंहाचार्यस्य अलङ्कारेन्दुशेखरः, केशवाचार्यस्य अलङ्कारशेखरश्च वर्तन्ते । तथा च ध्वनिमात्रसमीक्षितेषु— आनन्दवर्धनस्य ध्वन्यालोकः, मम्मटस्य शब्दव्यापारविचारः, अप्पयस्य वृत्तिवार्तिकश्च ग्रन्था लभ्यन्ते । अत्रैव ध्वनिविरोधिनो  ग्रन्थाः— भट्टनायकस्य हृदयदर्पणः, महिमभट्टस्य व्यक्तिविवेकः, मुकुलभट्टस्याभिधावृत्तिमातृका च सन्ति । तथैव रुद्रभट्टस्य शृङ्गारतिलकम्, भानुदत्तस्य रसमञ्जरी रसतरङ्गिणी च, रसमात्रसमीक्षायुता ग्रन्थाः सन्ति । दृश्यकाव्यतत्त्वमात्रविवेचकेषु— भरतस्य नाट्यशास्त्रम्, धनञ्जयस्य दशरूपकम्, सागरनन्दिनो नाटकलक्षणरत्नकोशः, रामचन्द्रगुणचन्दयोर्नाट्यदर्पणः, विद्यानाथस्य प्रतापरुद्रयशोभूषणम्, शारदातनयस्य भावप्रकाशनम्, रूपगोस्वामिनो नाटकचन्द्रिका च समागच्छन्ति । अस्मिन्नेव मार्गे विचारप्रधाना ग्रन्था अपि सन्ति । तेषु रुद्रटस्य हृदयलीला, क्षेमेन्द्रस्य कविकण्ठाभरणम्, सुवृत्ततिलकमौचित्यविचारचर्चा च, राजशेखरस्य काव्यमीमांसा, विश्वेश्वरस्य कवीन्द्रकण्ठाभरणम्, रामानन्दनस्य रसिकजीवनञ्च समाविशन्ति ।

अनेन प्रकारेण विलक्षणप्रतिभायुता विद्वांसः पृथक् पृथक् काव्यतत्त्वमङ्गीकृत्य मौलिकविचारानाविष्कुर्वन्तो विविधान् ग्रन्थान् रचितवन्तः। तत्र मुख्यवैचारिकसिद्धान्तेषु रस–अलङ्कार–रीति–ध्वनि–वक्रोक्ति–औचित्येति षट् काव्यशास्त्रीयसम्प्रदायाः सन्ति । एते एव पौरस्त्यकाव्यशास्त्रस्य शिरोभूषणभूताः, संस्कृतभाषाया गौरवविषयकविचाराश्च मन्यन्ते । एतेषां सैद्धान्तिकविचाराणां संस्थापकस्य श्रेयःप्राप्ता विद्वांसः सन्ति यत्— रसस्य भरतमुनिः, अलङ्कारस्य भामहः, रीतेर्वामनः, ध्वनेरानन्दवर्धनः, वक्रोक्तेः कुन्तकः, औचित्यस्य च क्षेमेन्द्रमहोदयाः । एतेषां विदुषां केन्द्रीयविचारे अन्य समालोचकेषु केचित् समर्थयन्ति, केचित् परिष्कारयन्ति, केचित्तु पूर्णरूपेण विपरीतविचाराः प्रस्तुवन्ति । इत्यनेन  समर्थनखण्डनमण्डनादिकारणेन संस्कृतकाव्यशास्त्रस्य समृद्धिर्विश्ववाङ्मयस्य शिरोभूषणभूता दृश्यते ।

संस्कृतलक्षणशास्त्रस्येतिहासो नैकैर्विचारैरलङ्कृतत्वात् काव्यप्रणयने कानि तत्त्वानि प्राधान्येनोपयोज्यानि, कानि च सामान्येन तथा कानि प्रायेणानुपयोज्यानीति विषयं  निर्णेतुं काठिन्यं  प्रतिभाति, परं सूक्ष्मविवेचनद्वारा मम्मटविश्वनाथजगन्नाथानां मतानि सर्वाधिकमान्यानि दृश्यन्ते । काव्यलक्षणहेतुप्रयोजनेषु मम्मटस्य मतं स्वीकरणीयम्, तथा काव्यभेदविषये जगन्नाथस्य मतं मान्यं मन्यते । नाट्यविषये भरतधनञ्जययोर्विचारपरिष्कारेण प्रस्तुतं विश्वनाथस्य मतं श्रेष्ठमनुभूयते । महाकाव्यविषयेऽपि विश्वनाथस्यैव मतं सुस्पष्टं दृश्यते ।  गुणस्यावस्थितिर्विषये आनन्दवर्धनस्य, रसप्रतीतिविषयेऽभिनवगुप्तस्य, अलङ्कारस्थानविषये मम्मटस्य, औचित्यरीतिवक्रोक्ति–आदिषु विषयेषु ध्वनिवादिनो मतं श्रेयष्करं दृश्यते ।

एतेषां सर्वेषां विचारमन्थनेन नवनीतरूपं सर्वाङ्गपूर्णं काव्यशरीरं निम्नानुसारम् उपयुक्तमिति उल्लिख्यते यत्—

ब्रह्मानन्दमयो रस एव काव्यशरीरस्यात्मा भवतीति विचारे केवलनामान्तरमेवास्ति, अर्थात् प्रकारान्तरेणापि तथ्यमिदं सर्वे काव्यसमालोचकाः स्वीकुर्वन्ति एव । गुणालङ्काररीतयस्तु काव्यशरीरस्योत्कर्षहेतव एव । अलङ्कारा रसमेवालङ्‌कुर्वन्ति, किञ्च–आत्मैवालङ्कार्यो भवति । रीतिवादिनां मते रीतयो गुणात्मका भवन्ति । माधुर्यादयो गुणाश्च रसस्यैव धर्माः । तस्माद्  रसप्राधान्ये रचितं काव्यशरीरं काव्यनामभाग् भवति । वक्रोक्तेः साध्योऽपि रस एव । ध्वनिरपि तद्वादिनां मते रस एव । एतेन सिद्ध्यति यत् केवलं विभिन्नकाव्यतत्त्वपरिपुष्टो रस एव काव्यात्मतयाऽवस्थितो भवति । अन्ये तु तस्य साधकाः परिपोषकाश्च जायन्ते ।

पौरस्त्यवाङ्मयस्य कमनीयमाभूषणं शास्त्रीयपक्षं विद्यते । अस्यैव पक्षस्य योगदानेन विश्ववाङ्मये संस्कृतवाङ्मयस्य सम्मानिता स्थितिरस्ति । अधुनापि पाश्चात्त्यसमालोचकाः काव्यतत्त्वानुसन्धाने पौरस्त्यशास्त्रीयपक्षमङ्गीकुर्वन्ति । अष्टादशशताब्दादारभ्यास्य विकासे तादृशी तीव्रता न दृश्यते तथापि भारतीयः, नेपालदेशीयश्च विद्वांसः पौरस्त्यकाव्यशास्त्रस्य व्याख्याविस्तारे सावधानेन प्रयासरताः सन्ति । तस्मादस्य परम्परा सुदीर्घकालं यावत् सञ्चलन् गमिष्यतीति निःशङ्कं सम्भाव्यते ।

\begin{visheshshlokatwo}[center]{84}

\textbf{इति काव्यशास्त्रजिज्ञासुजनानभिलक्ष्य कृता काव्यशास्त्रप्रवेशिका समाप्ता}

\end{visheshshlokatwo}

\specialupakhanda{\textbf{सन्दर्भसामग्रीसूची}}

\begin{flushleft}

अग्रवालः, हंसराजः, (सन् १९८५)  \textit{\textbf{संस्कृतसाहित्येतिहासः,}} चौखम्बा सुरभारती प्रकाशन ।

आनन्दवर्धनः,(सन् २००४)\textit{\textbf{ध्वन्यालोकः लोचनसहित,}}(व्या. जगन्नाथ),चौखम्बा विद्याभवन ।

.............(२००३) \textit{\textbf{ध्वन्यालोकः}},  (व्या. बदरीनाथ शर्मा), चौखम्बा संस्कृत सिरिज आफिस।

आप्टे, शिवराम,(सन् १९९९) \textit{\textbf{संस्कृत हिन्दी शब्दकोशः,}} दिल्ली : न्यु भारतीय बुक कर्पोरेसन।

उपाध्याय, बलदेव,(सन् १९९०) \textit{\textbf{संस्कृत साहित्य का सङ्क्षिप्त इतिहास,}} शारदा निकेतन।

..............(सन् १९६३)\textit{\textbf{भारतीय साहित्यशास्त्र,}} (प्रथम खण्ड),नन्दकिशोर एण्ड सन्स।

ऋषभदेवः (वि. सं. २०१५) \textit{\textbf{काव्यमञ्जरी,}} नेपाल भारत मैत्री संघ ।

काणे, पी. वी., (सन् २००२)	\textit{\textbf{संस्कृत काव्यशास्त्र का इतिहास,}} ,मोतीलाल बनारसी दास।

कुन्तकः, (सन् १९९८) \textit{\textbf{वक्रोक्तिजीवितम्}}, (व्या. राधेश्याममिश्रः), चौखम्बा संस्कृत विद्याभवन।

कुमार विमल, (सन् १९९८) \textit{\textbf{सौन्दर्यशास्त्र के तत्त्व,}} दिल्ली : राजकमल प्रकाशन।

क्षेमेन्द्रः,(सन् २०१०) \textit{\textbf{औचित्यविचारचर्चा,}} (तृ. सं.), चौखम्बा कृष्णदास अकादमी ।

खिस्ते, श्रीनारायणशास्त्री, (सन् २००३) \textit{\textbf{अलङ्कारसारमञ्जरी,}} चौखम्बा संस्कृत सिरिज आफिस।

गुप्तः, अभिनवः, (सन् १९६४)   \textit{\textbf{अभिनवभारती,}} गायकवाड ओरियन्टल सिरिज।

गैरोला, वाचस्पति, (सन् १९९७) \textit{\textbf{संस्कृतसाहित्य का इतिहास,}} (पं. सं.), चौखम्बा विद्याभवन ।

जगन्नाथः,(सन् २०१२) \textit{\textbf{रसगङ्गाधरः,}} (व्या. बदरीनाथ झा), चौखम्बा विद्याभवन।

जयदेवः, (सन् १९९९) \textit{\textbf{चन्द्रालोकः,}} (व्या. कृष्णमणि त्रिपाठी), चौखम्बा सुरभारती प्रकाशन।

दण्डी, (सन् १९९७)\textit{\textbf{काव्यादर्शः,}} (व्या. रामचन्द्रमिश्रः), चौखम्बा विद्याभवन ।

दीक्षितः,अप्पयः, (सन् १९९७) \textit{\textbf{कुवलयानन्दः,}} (व्या.भोलाशङ्करव्यासः),चौखम्बा विद्याभवन।

धनञ्जयः,(सन् १९८९) \textit{\textbf{दशरूपकम्,}} (सावलोकम्), (चन्द्रकलाव्याख्योपेतम्),चौखम्बा विद्याभवन, ।

नागेन्द्रः,(सन् १९६६)  \textit{\textbf{भारतीय काव्यशास्त्र की भूमिका,}} (भाग–२), नेसनल पब्लिसिङ् हाउस।

पाण्डेय, कान्तिचन्द्र,(सन् १९७८)  \textit{\textbf{स्वतन्त्र कलाशास्त्र,}} (भाग–२), चौखम्बा विद्याभवन।

.............(सन् १९६७) \textit{\textbf{स्वतन्त्र कलाशास्त्र,}} (भाग–१), चौखम्बा संस्कृत सिरिज अफिस।

पोखरेल, बालकृष्ण,(वि सं २०५२) \textit{\textbf{नेपाली बृहत् शब्दकोश,}} नेपाल राजकीय प्रज्ञा प्रतिष्ठान।

भरतमुनिः, (सन् १९८०) \textit{\textbf{नाट्यशास्त्रम्,}}  ओरियन्टल इन्टिच्युट।

भट्टोद्भटः, (सन् १९२५) \textit{\textbf{काव्यालङ्कारसारसङ्ग्रहः,}} प्राच्य विद्या संशोधन मन्दिर ।

..........(सन् २००४) \textit{\textbf{नाट्यशास्त्रम्}}, च. भा., (व्या. पारसनाथ द्विवेदी), चौखम्बा विद्याभवन।

भामहः, (सन् १०८५) \textit{\textbf{काव्यालङ्कारः,}} बिहार राष्ट्रभाषा परिषद्।

भोजराजः, (सन् १९७६) \textit{\textbf{सरस्वतीकण्ठाभरणम्}}, चौखम्बा ओरियन्टालिया ।

मम्मटः, (सन् १९९०) \textit{\textbf{काव्यप्रकाशः,}} सम्पूर्णानन्दसंस्कृतविश्वविद्यालयः  ।

..........(सन् १९८३) \textit{\textbf{काव्यप्रकाशः}}, (व्या. श्री निवास शास्त्री), अ. सं., साहित्य भण्डार ।

महिमभट्टः, (संवत् २०३९) \textit{\textbf{व्यक्तिविवेकः,}} चौखम्बा संस्कृत संस्थान ।

मिश्रः, केशवः,  (सन् १०८४)\textit{\textbf{अलङ्कारशेखरः,}} चौखम्बा संस्कृत सिरिज अफिस।

मिश्रः, रामचन्द्रः,(सन् ११९०) \textit{\textbf{संस्कृतसाहित्येतिहासः,}} चौखम्बा विद्याभवन ।

यादव राजमङ्गल, (सन् २०११) \textit{\textbf{संस्कृत काव्यशास्त्र की अर्वाचीन परम्परा,}} प्रतिभा प्रकाशन ।

राजशेखरः, (वि सं २०३०) \textit{\textbf{काव्यमीमांसा,}} (अनु.डिल्लीराम तिमिल्सिना), ने. रा. प्र. प्र. ।

रुद्रटः, (सन् १९६६) \textit{\textbf{काव्यालङ्कारः}}, चौखम्बा विद्या भवन ।

रुय्यकः, (सन् २०३६) \textit{\textbf{अलङ्कारसर्वस्वम्,}} चौखम्बा संस्कृत संस्थान ।

वामनः, (सन् २०१०)\textit{\textbf{काव्यालङ्कारसूत्रवृत्तिः,}} चौखम्बा कृष्णदास अकादमी।

वाग्भटः, (सन् १९५७), \textit{\textbf{वाग्भटालङ्कार,}} चौखम्बा विद्याभवन ।

विश्वनाथः, (सन् १९९५) \textit{\textbf{साहित्यदर्पणः,}} (व्या.लोकमणिदहाल), चौखम्बा सुरभारती प्रकाशन ।

................(सन् १९९९) \textit{\textbf{साहित्यदर्पणः,}} (व्या.शेषराजरेग्मी),कृष्णदास अकादमी।

विद्याधरः (सन् १९०३) \textit{\textbf{एकावली,}} राजकीय ग्रन्थशाला ।

विद्यानाथः, (सन् १९०९) \textit{\textbf{प्रतापरुद्रयशोभूषणम्,}} राजकीय ग्रन्थमाला ।

विश्वेश्वरः (सन् १८९८) \textit{\textbf{अलङ्कारकौस्तुभम्}}, निर्णय सागर मुद्रणालय ।

व्यास, (वि सं २०६७) \textit{\textbf{अग्निपुराणम्}}, (अनु.तिलकप्र. लुइँटेल), विद्यार्थी पुस्तक भण्डार  ।

व्यास, भोलाशङ्कर,(सन् २०२२)\textit{\textbf{भारतीय साहित्यशास्त्र और काव्यालङ्कार}}, चौखम्बा विद्याभवन।

सहाय, राजवंश,(सन् २०१४) \textit{\textbf{संस्कृतसाहित्यकोशः,}} चौखम्बा विद्याभवन।

……………(सन् १९९७)\textit{\textbf{अलङ्कारानुशीलन,}} चौखम्बा विद्याभवन।

सिग्द्याल सोमनाथः, (वि सं २०५०) \textit{\textbf{आदर्शराघव}} (पञ्चम सं) साझाप्रकाशन ।

सिंहः, अमरः,(सन् १९९७) \textit{\textbf{अमरकोशः,}} (रामाश्रमी टीका), चौखम्बा संस्कृत प्रतिष्ठान ।

हेमचन्द्रः (सन् १९३४) \textit{\textbf{काव्यानुशासनम्}}, द्वि. सं., पाण्डुरङ्ग जाबजी।

\end{flushleft}



\end{document}